% concatenated from Modal_definability_in_classes_of_Euclidean_frames_61.tex
%
%
%
%
\documentclass[draft]{article}
%
%
%
%
\hyphenation{a-bo-ve ad-mis-si-bi-li-ty as-su-me be-low bet-ween Boo-le-an clas-ses clus-ter clus-ter-fi-ni-te com-ple-te com-pu-ta-bi-li-ty con-di-tion con-di-tions con-se-quent-ly Con-se-quent-ly con-si-der cons-tant cons-tants cor-res-ponds cor-res-pon-den-ce de-ci-da-ble de-ci-ding de-fi-na-bi-li-ty de-fi-na-ble ele-men-ta-ry equi-va-len-ce Eu-cli-dean fa-mi-ly flo-wers fol-lo-wing fol-lows for-ma-li-ze for-mu-la fra-me fra-mes ga-me ge-ne-ra-ted Hen-ce he-re-di-ta-ri-ly ima-ge in-duc-ti-ve-ly ins-tan-ce in-tui-tio-nis-tic Lem-ma Lem-mas lo-cal-ly lo-gics Lo-gic mi-ni-mal mi-ni-ma-li-ty mo-dal mo-del mo-dels more-over Mo-re-over ne-ces-sa-ry Ne-ver-the-less non-empty non-ne-ga-ti-ve no-ti-ce nul-la-ri-ness ob-ser-ved or-der po-si-ti-ve pro-blem pro-blems pro-po-si-tion Pro-po-si-tion pro-po-si-tion Pro-po-si-tions quan-ti-fier ran-ges re-la-tion-ship re-la-tion-ships res-pect re-sult Sahlq-vist Scien-ce se-cond sen-ten-ce sen-ten-ces sin-ce stan-dard stra-te-gy substi-tu-ti-on theo-ry the-re-fo-re un-de-ci-da-bi-li-ty un-de-ci-da-ble uni-fi-cation va-lua-tion wi-thin}
%
%
%
%
\usepackage{amsfonts,amsmath,amssymb,color,graphicx,latexsym}
%
%
%
%
\newtheorem{lemma}{Lemma}
%
%
\newtheorem{theorem}{Theorem}
%
%
\newtheorem{proposition}{Proposition}
%
%
\newtheorem{corollary}{Corollary}
%
%
%
%
\newenvironment{proof}{{\bf Proof:}}{~$\dashv$\\}
%
%
\newenvironment{example}{{\bf Example:}}{\\}
%
%
\newenvironment{claim}{{\em Claim:}\/}{}
%
%
\newenvironment{proofclaim}{{\em Proof:}\/}{}
%
%
%
%
\def\sf{\mathtt{sf}}
%
%
\def\R{\mathbb{R}}
%
%
\def\Q{\mathbb{Q}}
%
%
\def\Z{\mathbb{Z}}
%
%
\def\N{\mathbb{N}}
%
%
\def\st{\mathtt{st}}
%
%
\def\mea{\mathtt{mea}}
%
%
\def\degre{\mathtt{deg}}
%
%
\def\cl{\mathtt{cl}}
%
%
\def\fcl{\mathtt{fcl}}
%
%
\def\len{\mathtt{len}}
%
%
\def\lsem{{\lbrack\mid}}
%
%
\def\rsem{{\mid\rbrack}}
%
%
\def\qd{\mathtt{qd}}
%
%
\def\qdd{\mathtt{qdd}}
%
%
\def\red{\mathtt{red}}
%
%
\def\Kframe{\mathtt{fr}}
%
%
\def\SQEMA{\mathtt{SQEMA}}
%
%
\def\dust{\mathtt{dust}}
%
%
\def\lcl{\mathtt{lcl}}
%
%
\def\nonrooted{\mathtt{nr}}
%
%
\def\nonended{\mathtt{ne}}
%
%
\def\Log{\mathtt{Log}}
%
%
\def\Th{\mathtt{Th}}
%
%
\def\Fr{\mathtt{Fr}}
%
%
\def\fod{\mathtt{fod}}
%
%
\def\mdef{\mathtt{md}}
%
%
\def\size{\mathtt{len}}
%
%
\def\id{\mathtt{id}}
%
%
\def\hd{\mathtt{hd}}
%
%
\def\tl{\mathtt{tl}}
%
%
\def\kernel{\mathtt{kernel}}
%
%
\def\root{\mathtt{root}}
%
%
\def\succ{\mathtt{succ}}
%
%
\def\translation{\mathtt{tr}}
%
%
\def\llbrack{\lbrack\lbrack}
%
%
\def\rrbrack{\rbrack\rbrack}
%
%
\def\llangle{\langle\langle}
%
%
\def\rrangle{\rangle\rangle}
%
%
\def\all{\mathtt{all}}
%
%
\def\bd{\mathtt{bd}}
%
%
\def\cnt{\mathtt{cnt}}
%
%
\def\fin{\mathtt{fin}}
%
%
\def\infini{\mathtt{inf}}
%
%
\def\dom{\mathtt{dom}}
%
%
\def\mp{\mathtt{mp}}
%
%
\def\MPRE{\mathbf{MPRE}}
%
%
\def\LIN{\mathbf{L}}
%
%
\def\tra{\mathtt{tr}}
%
%
\def\card{\mathtt{Card}}
%
%
\def\QBF{\mathbf{QBF}}
%
%
\def\coNP{\mathbf{coNP}}
%
%
\def\NP{\mathbf{NP}}
%
%
\def\NEXPTIME{\mathbf{NEXPTIME}}
%
%
\def\PSPACE{\mathbf{PSPACE}}
%
%
\def\EXPSPACE{\mathbf{EXPSPACE}}
%
%
\def\EXP{\mathbf{EXP}}
%
%
\def\P{\mathbf{P}}
%
%
\def\vvar{\mathtt{var}}
%
%
\def\ppar{\mathtt{par}}
%
%
\def\L{\mathbf{L}}
%
%
\def\BL{\mathbf{BL}}
%
%
\def\PAL{\mathbf{PAL}}
%
%
\def\S{\mathbf{S}}
%
%
\def\K{\mathbf{K}}
%
%
\def\KD{\mathbf{KD}}
%
%
\def\KB{\mathbf{KB}}
%
%
\def\A{\mathbf{A}}
%
%
\def\GL{\mathbf{GL}}
%
%
\def\MC{\mathbf{MC}}
%
%
\def\KT{\mathbf{KT}}
%
%
\def\KDB{\mathbf{KDB}}
%
%
\def\KTB{\mathbf{KTB}}
%
%
\def\KG{\mathbf{KG}}
%
%
\def\KDG{\mathbf{KDG}}
%
%
\def\KTG{\mathbf{KTG}}
%
%
\def\moddel{\mathbf{M}}
%
%
\def\Alt{\mathbf{Alt}}
%
%
\def\DAlt{\mathbf{DAlt}}
%
%
\def\alt{\mathbf{alt}}
%
%
\def\MONA{\mathbf{MONA}}
%
%
\def\PVAR{\mathbf{PVAR}}
%
%
\def\IVAR{\mathbf{IVAR}}
%
%
\def\fiv{\mathtt{fiv}}
%
%
\def\qd{\mathtt{qd}}
%
%
\def\procedure{\mathtt{procedure}}
%
%
\def\val{\mathtt{val}}
%
%
\def\LTL{\mathbf{LTL}}
%
%
\def\T{\mathbf{T}}
%
%
\def\D{\mathbf{D}}
%
%
\def\VAR{\mathbf{VAR}}
%
%
\def\BIT{\mathbf{BIT}}
%
%
\def\FOR{\mathbf{FOR}}
%
%
\def\MOD{\mathbf{MOD}}
%
%
\def\TRE{\mathbf{TRE}}
%
%
\def\SUB{\mathbf{SUB}}
%
%
\def\CHA{\mathbf{CHA}}
%
%
\def\for{\mathbf{for}}
%
%
%
%
\begin{document}
%
%
%
%
\title{Modal definability in Euclidean modal logics}
%
%
%
%
\author{Philippe Balbiani$^{a}$\footnote{Email: philippe.balbiani@irit.fr.}
%
%
\hspace{0.24cm}
%
%
Tinko Tinchev$^{b}$\footnote{Email: tinko@fmi.uni-sofia.bg.}}
%
%
%
%
\date{$^{a}$Toulouse Institute of Computer Science Research
%
%
\\
%
%
CNRS~---~INPT~---~UT, Toulouse, France
%
%
\\
%
%
$^{b}$Faculty of Mathematics and Informatics
%
%
\\
%
%
Sofia University St.~Kliment Ohridski, Sofia, Bulgaria}
%
%
%
%
\maketitle
%
%
%
%
\begin{abstract}
%
%
This paper is about the computability of the modal definability problem in classes of frames determined by Euclidean modal logics.
We characterize those Euclidean modal logics such that the classes of frames they determine give rise to an undecidable modal definability problem.
%
%
\end{abstract}
%
%
%
%
{\bf Keywords:}
%
%
Propositional Modal Logic.
Model Theory of Classical First-Order Logic.
Correspondence Theory.
Modal definability.
Stable classes of frames.
%
%
%
%
\section{Introduction}\label{section:introduction}
%
%
A modal formula is first-order definable in a class of frames if there exists a sentence that is true exactly in the frames of that class validating the modal formula.
%The first-order definability problem in a class of frames is the following decision problem: determine whether a given modal formula is first-order definable in that class.
A sentence is modally definable in a class of frames if there exists a modal formula that is valid exactly in the frames of that class satisfying the sentence.
%The modal definability problem in a class of frames is the following decision problem: determine whether a given sentence is modally definable in that class.
A modal formula and a sentence correspond in a class of frames if they are respectively valid and true in the same frames of that class.
%The correspondence problem in a class of frames is the following decision problem: determine whether a given modal formula and a given sentence correspond in that class.
%
%
\\
\\
%
%
The question of the correspondence between modal formulas and sentences is concomitant with the creation of the relational semantics of modal logic.
Kripke~\cite{Kripke:1963} already observed that some sentences are modally definable: transitivity vs $\Box p\rightarrow\Box\Box p$, symmetry vs $p\rightarrow\Box\Diamond p$, etc.
Less than $20$ years have elapsed between Kripke's observation and the development of Correspondence Theory culminating in the publication of the book ``Modal Logic and Classical Logic''~\cite{vanBenthem:1983}: in 1975, Sahlqvist~\cite{Sahlqvist:1975} isolated a large set of modal formulas which guarantee completeness with respect to first-order definable classes of frames, whereas van Benthem~\cite{vanBenthem:1975} and Goldblatt~\cite{Goldblatt:1975} independently noticed that McKinsey formula $\Box\Diamond p\rightarrow\Diamond\Box p$ has no first-order correspondent.
%
%
\\
\\
%
%
%Since the first-order conditions corresponding to Sahlqvist formulas are effectively computable~\cite{Sahlqvist:1975}, then it is natural to ask whether Sahlqvist fragment contains all modal formulas possessing first-order correspondents.
Since the first-order conditions corresponding to Sahlqvist formulas are effectively computable~\cite{Sahlqvist:1975}, then it is natural to ask whether Sahlqvist fragment modally defines all first-order correspondents of modal formulas.
This question has received a negative answer, the conjunction $(\Box\Diamond p\rightarrow\Diamond\Box p)\wedge(\Box p\rightarrow\Box\Box p)$ possessing a first-order correspondent not modally definable by a Sahlqvist formula~\cite[Chapter~$3$]{Blackburn:deRijke:Venema:2001}.
However, notice that this statement should be carefully considered, seeing that the ``Sahlqvist formulas'' considered by Blackburn {\it et al.}\/ constitute a strict subset of the formulas considered by Sahlqvist.
Indeed, we do not know whether the first-order correspondent of $(\Box\Diamond p\rightarrow\Diamond\Box p)\wedge(\Box p\rightarrow\Box\Box p)$ is modally definable by a formula considered by Sahlqvist.
%
%
\\
\\
%
%
In any case, it is natural to ask whether the first-order definability problem, the modal definability problem and the correspondence problem are decidable.
%
%
%\\
%\\
%
%
The answers to these questions have been firstly obtained by Chagrova in her doctoral thesis~\cite{Chagrova:1989} and then further developed in~\cite{Chagrov:Chagrova:1995,Chagrov:Chagrova:2006,Chagrov:Chagrova:2007,Chagrova:1991}:
%
%
%\begin{itemize}
%
%
%\item there is no algorithm which when presented with an arbitrary modal formula as input, will terminate after finitely many steps returning either ``yes'' if there is a first-order correspondent, or ``no'' if there is not,
%
%
%\item there is no algorithm which when presented with an arbitrary sentence as input, will terminate after finitely many steps returning either ``yes'' if there is a modal correspondent, or ``no'' if there is not,
%
%
%\item there is no algorithm which when presented with an arbitrary modal formula and an arbitrary sentence as input, will terminate after finitely many steps returning either ``yes'' if they correspond, or ``no'' if they do not.
%
%
%\end{itemize}
the first-order definability problem, the modal definability problem and the correspondence problem are undecidable.
%
%
%\\
%\\
%
%
Chagrova's results concern the class of all frames.
They have been obtained by reductions from accessibility problems in Minsky machines~\cite{Chagrov:Chagrova:2006}.
As a result, it is not obvious how similar results can be obtained with respect to restricted classes of frames (the class of all transitive frames, the class of all symmetric frames, etc).
%However, notice that Chagrov and Chagrova~\cite{Chagrov:Chagrova:1995} have also considered the first-order definability problem with respect to the class of all irreflexive transitive frames.
%
%
\\
\\
%
%
With respect to the class of frames determined by $\S5$, it follows from $\S5$~Reduction Theorem in~\cite[Page~$51$]{Hughes:Cresswell:1968} and Lemma~$9.7$ in~\cite[Page~$99$]{vanBenthem:1983} that every modal formula is first-order definable.
Moreover, as proved in~\cite{Balbiani:Tinchev:2005,Balbiani:Tinchev:2006}, modal definability with respect to the class of frames determined by $\S5$ is $\PSPACE$-complete.
%
%
%\\
%\\
%
%
With respect to the classes of frames determined by $\KD45$ and $\K45$, similar results hold as well~\cite{Georgiev:2017a,Georgiev:2017b}.
%
%
\\
\\
%
%
Therefore, a first question comes: is there an extension of $\K5$ such that the class of frames it determines gives rise both to a trivial first-order definability problem~---~i.e. every modal formula is first-order definable in that class~---~and an undecidable modal definability problem?
This question has been positively answered: with respect to the class of frames determined by $\K5$, first-order definability is trivial, whereas modal definability is undecidable~\cite{Balbiani:Georgiev:Tinchev:2018}.
%
%
\\
\\
%
%
The truth is that the classes of frames determined by all extensions of $\K5$ give rise to a trivial first-order definability problem.
%
%
%\\
%\\
%
%
Therefore, a second question comes: characterize those extensions of $\K5$ such that the classes of frames they determine give rise to an undecidable modal definability problem.
%We answer this question in Section~\ref{section:about:computability:of:modal:definability} where we give a necessary and sufficient condition (Proposition~\ref{proposition:beta:beta:nsc:for:decidability:theory}) characterizing those Euclidean modal logics for which the theories of the classes of frames they define are decidable in exponential space, a necessary and sufficient condition (Theorem~\ref{proposition:finale}) characterizing those Euclidean modal logics for which the modal definability problem with respect to the classes of frames they define and the correspondence problem with respect to the classes of frames they define are decidable in exponential space.
We answer this question in Section~\ref{section:about:computability:of:modal:definability} where we
%We also prove that one can determine in non-deterministic exponential time whether the problem of deciding the modal definability of sentences with respect to a given Euclidean modal logic is decidable (Theorem~\ref{corollary:modal:definability:complexity}) and whether the problem of deciding the correspondence of modal formulas and sentences with respect to a given Euclidean modal logic is decidable (Theorem~\ref{corollary:correspondence:complexity}).
also prove that one can solve in non-deterministic exponential time the following decision problem: determine whether the problem of deciding the modal definability of sentences with respect to a given Euclidean modal logic is decidable.
%
%
\\
\\
%
%
As for the rest of the paper, it is organized as follows.
Although we assume the reader is at home with basic tools and techniques in Propositional Modal Logic and Model Theory of Classical First-Order Logic, we include in Sections~\ref{section:about:preliminaries}--\ref{section:first:order:syntax:and:semantics} the needed material about these topics.
The definability problems are introduced in Section~\ref{section:definability:problems}.
Standard theorems about relative interpretations and relativizations are adapted to our context in Sections~\ref{section:about:relative:interpretations} and~\ref{section:about:relativization}.
Theorems~\ref{proposition:finale} and~\ref{proposition:finale:correspondence:problem} state in Section~\ref{section:about:computability:of:modal:definability} the main results of this paper.
%An Appendix includes the proofs of some of our results.
This paper includes the proofs of some of our results.
Some of these proofs are relatively simple and we have included them here just for the sake of the completeness.
%
%
%
%
\section{Preliminaries}\label{section:about:preliminaries}
%
%
For all $k,l{\in}\Z$, let $\lsem k,l\rsem{=}\{m{\in}\Z:\ k{\leq}m{\leq}l\}$.
%
%
\\
\\
%
%
For all sets $S$, $\wp(S)$ denotes the set of all subsets of $S$.
%
%
\\
\\
%
%
For all sets $S$, ${\parallel}S{\parallel}$ denotes the cardinality of $S$.
%For all sets $S$, $\id_{S}$ denotes the identity function on $S$.
%
%
\\
\\
%
%
For all sets $A,B$, for all functions $\rho:\ A{\longrightarrow}\wp(B)$ and for all $D{\in}\wp(B)$, let $\rho^{{-}1}(D){=}\{s{\in}A:\ \rho(s){=}D\}$.
%For all sets $A,B$, for all functions $\rho:\ A{\longrightarrow}\wp(B)$, for all $s{\in}A$ and for all $t{\in}B$, $s$ is an {\it ancestor of $t$}\/ if either $\rho(s){=}\{t\}$, or $\rho(s){=}B{\setminus}\{t\}$.
%
%
\\
\\
%
%
For all sets $W$, for all binary relations $R$ on $W$ and for all $s,t{\in}W$, let $R(s){=}\{w{\in}
%
%
%
%
%
%
%
%
%
%
%
%
$\linebreak$
%
%
%
%
%
%
%
%
%
%
%
%
W:\ s{R}w\}$ and $R^{{-}1}(t){=}\{w{\in}W:\ w{R}t\}$.
%For all sets $W$, for all binary relations $R$ on $W$ and for all $X,Y{\subseteq}W$, let $R(X){=}\bigcup\{R(s):\ s{\in}X\}$ and $R^{{-}1}(Y){=}
%
%
%
%
%
%
%
%
%
%
%
%
%$\linebreak$
%
%
%
%
%
%
%
%
%
%
%
%
%\bigcup\{R^{{-}1}(t):\ t{\in}Y\}$.
%
%
\\
\\
%
%
Let $\mathbf{N}^{+}{=}\N{\setminus}\{0\}$ and $\mathbf{N}^{-}{=}\N{\cup}\{{-}1\}$.
%
%
%The proofs of the following results are relatively simple.
%We have included the proof of Lemma~\ref{lemma:if:S:L:setminus:some:pairs:is:infinite:then:two:consequences} in the Appendix just for the sake of the completeness.
%
%
%\begin{lemma}
%
%
%For all $m{\in}\mathbf{N}^{+}{\cup}\{0\}$ and for all $n{\in}\mathbf{N}^{-}$, $(m,n){\in}\{(0,0)\}{\cup}(\mathbf{N}^{+}{\times}\mathbf{N}^{-})$ if and only if either $m{\not=}0$, or $n{=}0$.
\\
\\
%
%
Obviously, for all $m{\in}\mathbf{N}^{+}{\cup}\{0\}$ and for all $n{\in}\mathbf{N}^{-}$, $(m,n){\in}\{(0,0)\}{\cup}(\mathbf{N}^{+}{\times}\mathbf{N}^{-})$ if and only if either $m{\not=}0$, or $n{=}0$.
%
%
\\
\\
%
%
Let $\ll$ be the well-founded partial order on $\mathbf{N}^{+}{\times}\mathbf{N}^{-}$ such that for all $m,m^{\prime}{\in}\mathbf{N}^{+}$ and for all $n,n^{\prime}{\in}\mathbf{N}^{-}$,
%
%
%\begin{itemize}
%
%
%\item $(m,n){\ll}(m^{\prime},n^{\prime})$ if and only if at least one of the following conditions holds:
%
%
%\begin{itemize}
%
%
%\item $n{\in}\N$, $n^{\prime}{\in}\N$, $m{\leq}m^{\prime}$ and $n{\leq}n^{\prime}$,
%
%
%\item $n{=}{-}1$, $n^{\prime}{\in}\N$ and $m{\leq}m^{\prime}{+}n^{\prime}$,
%\item $n{=}{-}1$, $n^{\prime}{\in}\N$ and $m{\leq}m^{\prime}$,
%
%
%\item $n{=}{-}1$, $n^{\prime}{=}{-}1$ and $m{\leq}m^{\prime}$.
%
%
%\end{itemize}
%
%
%\end{itemize}
$(m,n){\ll}(m^{\prime},n^{\prime})$ if and only if $m{\leq}m^{\prime}$ and $n{\leq}n^{\prime}$.
%
%
%\end{lemma}
%
%
%\begin{lemma}\label{lemma:about:m:mprime:n:nprime:q:TWOq}
%
%
%Let $q{\geq}3$.
%For all $m,m^{\prime}{\in}\mathbf{N}^{+}$ and for all $n,n^{\prime}{\in}\mathbf{N}^{-}$, if $(m,n){\ll}(m^{\prime},
%
%
%
%
%
%
%
%
%
%
%
%
%$\linebreak$
%
%
%
%
%
%
%
%
%
%
%
%
%n^{\prime})$, $m^{\prime}{\leq}q$ and $n^{\prime}{\leq}q$ then $m{\leq}2{\times}q$ and $n{\leq}q$.
%
%
%\end{lemma}
%
%
\begin{lemma}\label{lemma:property:of:the:well:founded:order:between:pairs}
%
%
For all $m{\in}\mathbf{N}^{+}$ and for all $n{\in}\mathbf{N}^{-}$, if $m{\geq}2$ and $n{\geq}2$ then for all $m^{\prime}{\in}\mathbf{N}^{+}$ and for all $n^{\prime}{\in}\mathbf{N}^{-}$, if $m^{\prime}{<}m$ then $(m^{\prime},2){\ll}(m,n)$ and if $n^{\prime}{<}n$ then $(2,n^{\prime}){\ll}(m,n)$.
%Moreover, for all $m{\in}\mathbf{N}^{+}$ and for all $n{\in}\mathbf{N}^{-}$, if $m{\geq}2$ and $n{\geq}2$ then for all $m^{\prime}{\in}\mathbf{N}^{+}$ and for all $n^{\prime}{\in}\mathbf{N}^{-}$, if $m^{\prime}{<}m$ then $(m^{\prime},2){\ll}(m,n)$ and if $n^{\prime}{<}n$ then $(2,n^{\prime}){\ll}(m,n)$.
%
%
\end{lemma}
%
%
%\\
%\\
For all $m{\in}\mathbf{N}^{+}$ and for all $n{\in}\mathbf{N}^{-}$, let
%
%
\begin{itemize}
%
%
\item $\pi(m,n){=}\{(m^{\prime},n^{\prime}):\ m^{\prime}{\in}\mathbf{N}^{+}$, $n^{\prime}{\in}\mathbf{N}^{-}$ and $(m^{\prime},n^{\prime}){\ll}(m,n)\}$.
%
%
\end{itemize}
%
%
A subset $\mathtt{S}$ of $\mathbf{N}^{+}{\times}\mathbf{N}^{-}$ is {\it closed}\/ if for all $m{\in}\mathbf{N}^{+}$ and for all $n{\in}\mathbf{N}^{-}$, if $(m,n){\in}\mathtt{S}$ then $\pi(m,n){\subseteq}\mathtt{S}$.
%
%
\\
\\
%
%
Obviously, for all $m{\in}\mathbf{N}^{+}$ and for all $n{\in}\mathbf{N}^{-}$, $\pi(m,n)$ is a closed subset of $\mathbf{N}^{+}{\times}\mathbf{N}^{-}$.
%
%
\begin{lemma}\label{lemma:ll:finitely:many:smaller:pairs}
%
%
For all $m{\in}\mathbf{N}^{+}$ and for all $n{\in}\mathbf{N}^{-}$, $\pi(m,n)$ is finite.
%
%
\end{lemma}
%
%
%\begin{lemma}
%
%
%For all $m{\in}\mathbf{N}^{+}$ and for all $n{\in}\mathbf{N}^{-}$, $\pi(m,n)$ is a closed subset of $\mathbf{N}^{+}{\times}\mathbf{N}^{-}$.
%
%
%\end{lemma}
%
%
\begin{lemma}\label{lemma:if:S:L:setminus:some:pairs:is:infinite:then:two:consequences}
%
%
For all closed subsets $\mathtt{S}$ of $\mathbf{N}^{+}{\times}\mathbf{N}^{-}$, if $\mathtt{S}{\setminus}((\{1\}{\times}\mathbf{N}^{-}){\cup}(\mathbf{N}^{+}{\times}\{{-}1,0,
%
%
%
%
%
%
%
%
%
%
%
%
$\linebreak$
%
%
%
%
%
%
%
%
%
%
%
%
1\}))$ is infinite then either $\{2\}{\times}\mathbf{N}^{-}{\subseteq}\mathtt{S}$, or $\mathbf{N}^{+}{\times}\{2\}{\subseteq}\mathtt{S}$.
%
%
\end{lemma}
%
%
\begin{proof}
%
%
Let $\mathtt{S}$ be a closed subset of $\mathbf{N}^{+}{\times}\mathbf{N}^{-}$.
Suppose $\mathtt{S}{\setminus}((\{1\}{\times}\mathbf{N}^{-}){\cup}(\mathbf{N}^{+}{\times}
%
%
%
%
%
%
%
%
%
%
%
%
$\linebreak$
%
%
%
%
%
%
%
%
%
%
%
%
\{{-}1,0,1\}))$ is infinite.
For the sake of the contradiction, suppose neither $\{2\}{\times}
%
%
%
%
%
%
%
%
%
%
%
%
$\linebreak$
%
%
%
%
%
%
%
%
%
%
%
%
\mathbf{N}^{-}{\subseteq}\mathtt{S}$, nor $\mathbf{N}^{+}{\times}\{2\}{\subseteq}\mathtt{S}$.
Hence, let $m{\in}\mathbf{N}^{+}$ and $n{\in}\mathbf{N}^{-}$ be such that $(m,2){\not\in}\mathtt{S}$ and $(2,n){\not\in}\mathtt{S}$.
Since $\mathtt{S}{\setminus}((\{1\}{\times}\mathbf{N}^{-}){\cup}(\mathbf{N}^{+}{\times}\{{-}1,0,1\}))$ is infinite, then by
%
%
%
%
%
%
%
%
%
%
%
%
\linebreak
%
%
%
%
%
%
%
%
%
%
%
%
Lemma~\ref{lemma:ll:finitely:many:smaller:pairs}, $\mathtt{S}{\setminus}((\{1\}{\times}\mathbf{N}^{-}){\cup}(\mathbf{N}^{+}{\times}\{{-}1,0,1\})){\not\subseteq}\pi(m,n)$.
Thus, let $m^{\prime}{\in}\mathbf{N}^{+}$ and $n^{\prime}{\in}\mathbf{N}^{-}$ be such that $(m^{\prime},n^{\prime}){\in}\mathtt{S}{\setminus}((\{1\}{\times}\mathbf{N}^{-}){\cup}(\mathbf{N}^{+}{\times}\{{-}1,0,1\}))$ and $(m^{\prime},n^{\prime}){\not\ll}
%
%
%
%
%
%
%
%
%
%
%
%
$\linebreak$
%
%
%
%
%
%
%
%
%
%
%
%
(m,n)$.
Consequently, $(m^{\prime},n^{\prime}){\in}\mathtt{S}$, $m^{\prime}{>}1$ and $n^{\prime}{>}1$.
Hence, $m^{\prime}{\geq}2$ and $n^{\prime}{\geq}2$.
Since $(m^{\prime},n^{\prime}){\not\ll}(m,n)$, then either $m{<}m^{\prime}$, or $n{<}n^{\prime}$.
Since $m^{\prime}{\geq}2$ and $n^{\prime}{\geq}2$, then by Lemma~\ref{lemma:property:of:the:well:founded:order:between:pairs}, either $(m,2){\ll}(m^{\prime},n^{\prime})$, or $(2,n){\ll}(m^{\prime},n^{\prime})$.
Since $\mathtt{S}$ is a closed subset of $\mathbf{N}^{+}{\times}\mathbf{N}^{-}$ and $(m^{\prime},n^{\prime}){\in}\mathtt{S}$, then either $(m,2){\in}\mathtt{S}$, or $(2,n){\in}\mathtt{S}$: a contradiction.
Thus, either $\{2\}{\times}\mathbf{N}^{-}{\subseteq}\mathtt{S}$, or $\mathbf{N}^{+}{\times}\{2\}{\subseteq}\mathtt{S}$.
%
%
\medskip
%
%
\end{proof}
%
%
%
%
\section{Frames, galaxies and flowers}\label{section:preliminary:definitions}
%
%
In this section, we introduce $3$~types of relational structures: frames, galaxies and flowers.
%
%
\begin{figure}[ht]
%
%
\begin{center}
%
%
\begin{picture}(-50,50)(-50,50)
%
%
\put(-120.00,50.00){\circle*{3}}
\put(-115.00,50.00){$f$}
\put(-160.00,50.00){\circle*{3}}
\put(-155.00,50.00){$e$}
\put(-200.00,50.00){\circle*{3}}
\put(-195.00,50.00){$d$}
\put(-120.00,90.00){\circle*{3}}
\put(-115.00,90.00){$c$}
\put(-160.00,90.00){\circle*{3}}
\put(-155.00,90.00){$b$}
\put(-200.00,90.00){\circle*{3}}
\put(-195.00,90.00){$a$}
\put(-160.00,90.00){\vector(+1,-1){40.00}}
\put(-160.00,90.00){\vector(-1,-1){40.00}}
\put(-120.00,90.00){\vector(-1,-1){40.00}}
\put(-120.00,90.00){\vector(0,-1){40.00}}
\put(-220.00,40.00){\line(+1,0){120.00}}
\put(-220.00,60.00){\line(+1,0){120.00}}
\put(-220.00,40.00){\line(0,+1){20.00}}
\put(-100.00,40.00){\line(0,+1){20.00}}
%
%
\put(+40.00,50.00){\circle*{3}}
\put(45.00,50.00){$l$}
\put(00.00,50.00){\circle*{3}}
\put(05.00,50.00){$k$}
\put(-40.00,50.00){\circle*{3}}
\put(-35.00,50.00){$j$}
\put(+40.00,90.00){\circle*{3}}
\put(+45.00,90.00){$i$}
\put(00.00,90.00){\circle*{3}}
\put(05.00,90.00){$h$}
\put(-40.00,90.00){\circle*{3}}
\put(-35.00,90.00){$g$}
\put(-40.00,90.00){\vector(0,-1){40.00}}
\put(00.00,90.00){\vector(+1,-1){40.00}}
\put(00.00,90.00){\vector(-1,-1){40.00}}
\put(+40.00,90.00){\vector(-1,-1){40.00}}
\put(00.00,90.00){\vector(0,-1){40.00}}
\put(-60.00,40.00){\line(+1,0){120.00}}
\put(-60.00,60.00){\line(+1,0){120.00}}
\put(-60.00,40.00){\line(0,+1){20.00}}
\put(+60.00,40.00){\line(0,+1){20.00}}
%
%
\end{picture}
%
%
\caption{Example of an Euclidean frame where, on one hand, $d$, $e$ and $f$ are pairwise connected and, on the other hand, $j$, $k$ and $l$ are pairwise connected.}
%
%
\end{center}
%
%
\end{figure}
%
%
%
%
\paragraph{Frames}
%
%
A {\it frame}\/ is a tuple of the form $(W,R)$ where $W$ is a non-empty set and $R$ is a binary relation on $W$.
%
%
\\
\\
%
%
A frame $(W,R)$ is {\it irreflexive}\/ if for all $s{\in}W$, not $s{R}s$.
A frame $(W,R)$ is {\it symmetric}\/ if for all $s,t{\in}W$, if $s{R}t$ then $t{R}s$.
%Let ${\mathcal C}_{is}$ be the class of all irreflexive symmetric frames.
A frame $(W,R)$ is {\it transitive}\/ if for all $s,t,u{\in}W$, if $s{R}t$ and $t{R}u$ then $s{R}u$.
A frame $(W,R)$ is {\it Euclidean}\/ if for all $s,t,u{\in}W$, if $s{R}t$ and $s{R}u$ then $t{R}u$ and $u{R}t$.
%
%
%\begin{lemma}
%
%
%For all Euclidean frames $(W,R)$ and for all $w{\in}W$, if $R(w){=}\emptyset$ then $R^{{-}1}(w){=}\emptyset$.
\\
\\
Obviously, for all Euclidean frames $(W,R)$ and for all $w{\in}W$, if $R(w){=}\emptyset$ then $R^{{-}1}(w){=}\emptyset$.
%
%
%\end{lemma}
%
%
%
%
\paragraph{Dusts, roots and kernels}
%
%
%A frame $(W,R)$ is {\it finite}\/ if $W$ is finite.
%Obviously, for all $k{\in}\N$, the set ${\mathcal N}_{k}$ of all pairwise non-isomorphic finite frames $(W,R)$ such that ${\parallel}W{\parallel}{\leq}k$ can be effectively computed.
%As a result, for all $k{\in}\N$, let $\mathbf{N}_{k}{=}{\parallel}{\mathcal N}_{k}{\parallel}$.
%A frame $(W,R)$ is {\it simple}\/ if for all $s{\in}W$, either ${\parallel}R(s){\parallel}{\leq}1$, or ${\parallel}R(R(s)){\setminus}R(s){\parallel}{\leq}1$.
%Obviously, for all $k{\in}\N$, the set ${\mathcal N}^{s}_{k}$ of all pairwise non-isomorphic finite simple frames $(W,R)$ such that ${\parallel}W{\parallel}{\leq}k$ can be effectively computed.
%As a result, for all $k{\in}\N$, let $\mathbf{N}^{s}_{k}{=}{\parallel}{\mathcal N}^{s}_{k}{\parallel}$.
Let the {\it dust of the Euclidean frame $(W,R)$}\/ be the set (denoted $\dust(W,R)$) of all $w$ in $W$ such that $R(w){=}\emptyset$.
Let the {\it root of the Euclidean frame $(W,R)$}\/ be the set (denoted $\root(W,R)$) of all $w$ in $W$ such that $R(w){\not=}\emptyset$ and $R^{{-}1}(w){=}\emptyset$.
Let the {\it kernel of the Euclidean frame $(W,R)$}\/ be the set (denoted $\kernel(W,R)$) of all $w$ in $W$ such that $R^{{-}1}(w){\not=}\emptyset$.
%
%
\begin{figure}[ht]
%
%
\begin{center}
%
%
\begin{picture}(-50,50)(-50,50)
%
%
\put(-40.00,50.00){\circle*{3}}
\put(-35.00,50.00){$f$}
\put(-80.00,50.00){\circle*{3}}
\put(-75.00,50.00){$e$}
\put(-120.00,50.00){\circle*{3}}
\put(-115.00,50.00){$d$}
\put(-40.00,90.00){\circle*{3}}
\put(-35.00,90.00){$c$}
\put(-80.00,90.00){\circle*{3}}
\put(-75.00,90.00){$b$}
\put(-120.00,90.00){\circle*{3}}
\put(-115.00,90.00){$a$}
\put(-80.00,90.00){\vector(+1,-1){40.00}}
\put(-80.00,90.00){\vector(-1,-1){40.00}}
\put(-40.00,90.00){\vector(-1,-1){40.00}}
\put(-40.00,90.00){\vector(0,-1){40.00}}
\put(-140.00,40.00){\line(+1,0){120.00}}
\put(-140.00,60.00){\line(+1,0){120.00}}
\put(-140.00,40.00){\line(0,+1){20.00}}
\put(-20.00,40.00){\line(0,+1){20.00}}
%
%
\end{picture}
%
%
\caption{Example of an Euclidean frame with dust $\{a\}$, root $\{b,c\}$ and kernel $\{d,e,f\}$.}
%
%
\end{center}
%
%
\end{figure}
%
%
\begin{figure}[ht]
%
%
\begin{center}
%
%
\begin{picture}(-50,50)(-50,50)
%
%
\put(-40.00,50.00){\circle*{3}}
\put(-35.00,50.00){$f$}
\put(-80.00,50.00){\circle*{3}}
\put(-75.00,50.00){$e$}
\put(-40.00,50.00){\circle*{3}}
\put(-115.00,50.00){$d$}
\put(-40.00,90.00){\circle*{3}}
\put(-35.00,90.00){$c$}
\put(-80.00,90.00){\circle*{3}}
\put(-75.00,90.00){$b$}
\put(-120.00,90.00){\circle*{3}}
\put(-115.00,90.00){$a$}
\put(-120.00,90.00){\vector(0,-1){40.00}}
%\put(-80.00,90.00){\vector(+1,-1){40.00}}
\put(-80.00,90.00){\vector(-1,-1){40.00}}
\put(-40.00,90.00){\vector(-1,-1){40.00}}
\put(-80.00,90.00){\vector(0,-1){40.00}}
\put(-140.00,40.00){\line(+1,0){120.00}}
\put(-140.00,60.00){\line(+1,0){120.00}}
\put(-140.00,40.00){\line(0,+1){20.00}}
\put(-20.00,40.00){\line(0,+1){20.00}}
%
%
\end{picture}
%
%
\caption{Example of an Euclidean frame with dust $\emptyset$, root $\{a,b,c\}$ and kernel $\{d,e,f\}$.}
%
%
\end{center}
%
%
\end{figure}
%
%
%\begin{lemma}
%
%
%For all Euclidean frames $(W,R)$, $\dust(W,R)$, $\root(W,R)$ and
%
%
%
%
%
%
%
%
%
%
%
%
%\linebreak$
%
%
%
%
%
%
%
%
%
%
%
%
%\kernel(W,R)$ constitute a partition of $W$.
%
%
%\end{lemma}
%
%
\\
\\
%
%
Obviously, for all Euclidean frames $(W,R)$, $\dust(W,R)$, $\root(W,R)$ and
%
%
%
%
%
%
%
%
%
%
%
%
\linebreak$
%
%
%
%
%
%
%
%
%
%
%
%
\kernel(W,R)$ constitute a partition of $W$.
%
%
%\begin{lemma}
%
%
%For all Euclidean frames $(W,R)$, $\root(W,R){=}\bigcup\{R^{{-}1}(w){\setminus}
%
%
%
%
%
%
%
%
%
%
%
%
%$\linebreak$
%
%
%
%
%
%
%
%
%
%
%
%
%\kernel(W,R):\ w{\in}W\}$.
%
%
%\end{lemma}
%
%
%
%
\paragraph{Indexed families of Euclidean frames}
%
%
The {\it size of an indexed family $\{(W_{i},
%
%
%
%
%
%
%
%
%
%
%
%
$\linebreak$
%
%
%
%
%
%
%
%
%
%
%
%
R_{i}):\ i{\in}I\}$ of Euclidean frames}\/ is the cardinality of $I$.
An indexed family $\{(W_{i},R_{i}):\ i{\in}I\}$ of Euclidean frames is {\it non-empty}\/ if its size is positive.
%
%
\\
\\
%
%
A non-empty indexed family $\{(W_{i},R_{i}):\ i{\in}I\}$ of Euclidean frames is {\it disjoint}\/ if for all $i,j{\in}I$, if $i{\not=}j$ then $W_{i}{\cap}W_{j}{=}\emptyset$.
%
%
\\
\\
%
%
The {\it union of a disjoint indexed family $\{(W_{i},R_{i}):\ i{\in}I\}$ of Euclidean frames}\/ is the Euclidean frame $(W,R)$ such that $W{=}\bigcup\{W_{i}:\ i{\in}I\}$ and $R{=}\bigcup\{R_{i}:\ i{\in}I\}$.
%
%
%\begin{lemma}
%
%
%The union of a disjoint indexed family of Euclidean frames is Euclidean.
%
%
%\end{lemma}
%
%
%
%
\paragraph{Dust-finite indexed families of Euclidean frames}
%
%
A disjoint indexed family $\{(W_{i},R_{i}):\ i{\in}I\}$ of Euclidean frames is {\it dust-finite}\/ if for all $i{\in}I$, $\dust(W_{i},
%
%
%
%
%
%
%
%
%
%
%
%
$\linebreak$
%
%
%
%
%
%
%
%
%
%
%
%
R_{i})$ is finite.
For all $k{\in}\N$, a dust-finite indexed family $\{(W_{i},R_{i}):\ i{\in}I\}$ of Euclidean frames is {\it $k$-dust-bounded}\/ if for all $i{\in}I$, ${\parallel}\dust(W_{i},R_{i}){\parallel}{\leq}k$.
A dust-finite indexed family $\{(W_{i},R_{i}):\ i{\in}I\}$ of Euclidean frames is {\it dust-bounded}\/ if there exists $k{\in}\N$ such that $\{(W_{i},R_{i}):\ i{\in}I\}$ is $k$-dust-bounded.
%
%
%
%
%\paragraph{Quasi-root-finite indexed families of frames}
%
%
%A disjoint indexed family $\{(W_{i},R_{i}):\ i{\in}I\}$ of frames is {\it quasi-root-finite}\/ if for all $i{\in}I$ and for all $w{\in}\kernel(W_{i},R_{i})$, $R_{i}^{{-}1}(w){\setminus}\kernel(W_{i},R_{i})$ is finite.
%For all $k{\in}\N$, a quasi-root-finite indexed family $\{(W_{i},R_{i}):\ i{\in}I\}$ of frames is {\it $k$-quasi-root-bounded}\/ if for all $i{\in}I$ and for all $w{\in}\kernel(W_{i},R_{i})$, ${\parallel}R_{i}^{{-}1}(w){\setminus}\kernel(W_{i},R_{i}){\parallel}{\leq}k$.
%A quasi-root-finite indexed family $\{(W_{i},R_{i}):\ i{\in}I\}$ of frames is {\it quasi-root-bounded}\/ if there exists $k{\in}\N$ such that $\{(W_{i},R_{i}):\ i{\in}I\}$ is $k$-quasi-root-bounded.
%
%
%
%
\paragraph{Root-finite indexed families of Euclidean frames}
%
%
A disjoint indexed family $\{(W_{i},R_{i}):\ i{\in}I\}$ of Euclidean frames is {\it root-finite}\/ if for all $i{\in}I$, $\root(W_{i},
%
%
%
%
%
%
%
%
%
%
%
%
$\linebreak$
%
%
%
%
%
%
%
%
%
%
%
%
R_{i})$ is finite.
For all $k{\in}\N$, a root-finite indexed family $\{(W_{i},R_{i}):\ i{\in}I\}$ of Euclidean frames is {\it $k$-root-bounded}\/ if for all $i{\in}I$, ${\parallel}\root(W_{i},R_{i}){\parallel}{\leq}k$.
A root-finite indexed family $\{(W_{i},R_{i}):\ i{\in}I\}$ of Euclidean frames is {\it root-bounded}\/ if there exists $k{\in}\N$ such that $\{(W_{i},R_{i}):\ i{\in}I\}$ is $k$-root-bounded.
%
%
%
%
\paragraph{Kernel-finite indexed families of Euclidean frames}
%
%
A disjoint indexed family $\{(W_{i},R_{i}):\ i{\in}I\}$ of Euclidean frames is {\it kernel-finite}\/ if for all $i{\in}I$, $\kernel(W_{i},R_{i})$ is finite.
For all $k{\in}\N$, a kernel-finite indexed family $\{(W_{i},R_{i}):\ i{\in}I\}$ of Euclidean frames is {\it $k$-kernel-bounded}\/ if for all $i{\in}I$, ${\parallel}\kernel(W_{i},R_{i}){\parallel}{\leq}k$.
A kernel-finite indexed family $\{(W_{i},R_{i}):\ i{\in}I\}$ of Euclidean frames is {\it kernel-bounded}\/ if there exists $k{\in}\N$ such that $\{(W_{i},R_{i}):\ i{\in}I\}$ is $k$-kernel-bounded.
%
%
%
%
%\paragraph{Prefinite indexed families of frames}
%
%
%A disjoint indexed family $\{(W_{i},R_{i}):\ i{\in}I\}$ of frames is {\it prefinite}\/ if for all $i{\in}I$, $W_{i}$ is finite.
%For all $k{\in}\N$, a prefinite indexed family $\{(W_{i},R_{i}):\ i{\in}I\}$ of frames is {\it $k$-prebounded}\/ if for all $i{\in}I$, ${\parallel}W_{i}{\parallel}{\leq}k$.
%A prefinite indexed family $\{(W_{i},R_{i}):\ i{\in}I\}$ of frames is {\it prebounded}\/ if there exists $k{\in}\N$ such that $\{(W_{i},R_{i}):\ i{\in}I\}$ is $k$-prebounded.
%
%
%
%
\paragraph{Galaxies}
%
%
For all sets $A,B$ and for all functions $\rho:\ A{\longrightarrow}\wp(B)$, if $A{\cap}B{=}\emptyset$ and $A{\cup}B{\not=}\emptyset$ then let ${\mathcal F}_{A,B}^{\rho}{=}(W_{A,B}^{\rho},R_{A,B}^{\rho})$ be the Euclidean frame (called {\it galaxy}) such that
%
%
\begin{itemize}
%
%
\item $W_{A,B}^{\rho}{=}A{\cup}B$,
%
%
\item $R_{A,B}^{\rho}{=}\bigcup\{\{s\}{\times}\rho(s):\ s{\in}A\}{\cup}(B{\times}B)$.
%
%
\end{itemize}
%
%
Obviously, for all galaxies ${\mathcal F}_{A,B}^{\rho}$, $\dust({\mathcal F}_{A,B}^{\rho}){=}\rho^{{-}1}(\emptyset)$, $\root({\mathcal F}_{A,B}^{\rho}){=}A{\setminus}\rho^{{-}1}(\emptyset)$ and $\kernel({\mathcal F}_{A,B}^{\rho}){=}B$.
%
%
\\
\\
%
%
The {\it upper part of the galaxy ${\mathcal F}_{A,B}^{\rho}$}\/ is the set $A$.
The {\it lower part of the galaxy ${\mathcal F}_{A,B}^{\rho}$}\/ is the set $B$.
%
%
\\
\\
%
%
A galaxy ${\mathcal F}_{A,B}^{\rho}$ is {\it finite}\/ if $A$ and $B$ are finite.
%
%
\\
\\
%
%
A galaxy ${\mathcal F}_{A,B}^{\rho}$ is {\it headed}\/ if there exists $s{\in}A$ such that $\rho(s){\not=}\emptyset$.
%
%
\\
\\
%
%
A galaxy ${\mathcal F}_{A,B}^{\rho}$ is {\it simple}\/ if $B{\not=}\emptyset$ and for all $s{\in}A$, either ${\parallel}\rho(s){\parallel}{\leq}1$, or ${\parallel}B{\setminus}\rho(s){\parallel}{\leq}
%
%
%
%
%
%
%
%
%
%
%
%
$\linebreak$
%
%
%
%
%
%
%
%
%
%
%
%
1$.
%
%
\\
\\
%
%
Obviously, for all galaxies ${\mathcal F}_{A,B}^{\rho}$, ${\mathcal F}_{A,B}^{\rho}$ is simple if and only if $B{\not=}\emptyset$ and for all $s{\in}A$, either $\rho(s){=}\emptyset$, or $\rho(s){=}B$, or there exists $t{\in}B$ such that $\rho(s){=}\{t\}$, or there exists $t{\in}B$ such that $\rho(s){=}B{\setminus}\{t\}$.
%
%
\begin{figure}[ht]
%
%
\begin{center}
%
%
\begin{picture}(-50,50)(-50,50)
%
%
\put(-40.00,50.00){\circle*{3}}
\put(-35.00,50.00){$f$}
\put(-80.00,50.00){\circle*{3}}
\put(-75.00,50.00){$e$}
\put(-120.00,50.00){\circle*{3}}
\put(-115.00,50.00){$d$}
\put(-40.00,90.00){\circle*{3}}
\put(-35.00,90.00){$c$}
\put(-80.00,90.00){\circle*{3}}
\put(-75.00,90.00){$b$}
\put(-120.00,90.00){\circle*{3}}
\put(-115.00,90.00){$a$}
\put(-80.00,90.00){\vector(-1,-1){40.00}}
\put(-80.00,90.00){\vector(0,-1){40.00}}
\put(-80.00,90.00){\vector(+1,-1){40.00}}
\put(-40.00,90.00){\vector(0,-1){40.00}}
\put(-40.00,90.00){\vector(0,-1){40.00}}
\put(-140.00,40.00){\line(+1,0){120.00}}
\put(-140.00,60.00){\line(+1,0){120.00}}
\put(-140.00,40.00){\line(0,+1){20.00}}
\put(-20.00,40.00){\line(0,+1){20.00}}
%
%
\end{picture}
%
%
\caption{Example of a simple galaxy with dust $\{a\}$, root $\{b,c\}$ and kernel $\{d,e,f\}$ and where $A{=}\{a,b,c\}$, $B{=}\{d,e,f\}$, $\rho(a){=}\emptyset$, $\rho(b){=}\{d,e,f\}$ and $\rho(c){=}\{f\}$.}
%
%
\end{center}
%
%
\end{figure}
%
%
%\begin{lemma}
%
%
%For all galaxies ${\mathcal F}_{A,B}^{\rho}$, $\dust({\mathcal F}_{A,B}^{\rho}){=}\rho^{{-}1}(\emptyset)$, $\root({\mathcal F}_{A,B}^{\rho}){=}A{\setminus}\rho^{{-}1}(\emptyset)$ and $\kernel({\mathcal F}_{A,B}^{\rho}){=}B$.
%
%
%\end{lemma}
%
%
\begin{lemma}\label{lemma:every:frame:disjoint:union:galaxies}
%
%
Every Euclidean frame is the union of a disjoint indexed family of galaxies.
%
%
\end{lemma}
%
%
%\begin{lemma}\label{lemma:simple:galaxy:property:about:their:rho:function}
%
%
%For all galaxies ${\mathcal F}_{A,B}^{\rho}$, ${\mathcal F}_{A,B}^{\rho}$ is simple if and only if $B{\not=}\emptyset$ and for all $s{\in}A$, either $\rho(s){=}\emptyset$, or $\rho(s){=}B$, or there exists $t{\in}B$ such that $\rho(s){=}\{t\}$, or there exists $t{\in}B$ such that $\rho(s){=}B{\setminus}\{t\}$.
%
%
%\end{lemma}
%
%
%Moreover, for all sets $A,B$ and for all functions $\rho:\ A{\longrightarrow}\wp(B)$, if $A{\cap}B{=}\emptyset$ and $A{\cup}B{\not=}\emptyset$ then for all $C,D{\in}\wp(B)$, if $\rho^{{-}1}(C){\cap}\rho^{{-}1}(D){\not=}\emptyset$ then $C{=}D$.
%
%
Let ${\mathcal K}_{2}$ be the class of all galaxies ${\mathcal F}^{\rho}_{A,B}$ such that ${\parallel}A{\parallel}{\geq}4$, ${\parallel}B{\parallel}{\geq}4$ and for all $s{\in}A$, ${\parallel}\rho(s){\parallel}{=}2$.
Let ${\mathcal L}_{2}$ be the class of all galaxies ${\mathcal F}^{\rho}_{A,B}$ such that ${\parallel}A{\parallel}{\geq}4$, ${\parallel}B{\parallel}{\geq}4$ and for all $s{\in}A$, ${\parallel}B{\setminus}\rho(s){\parallel}{=}2$.
%
%
\begin{figure}[ht]
%
%
\begin{center}
%
%
\begin{picture}(-50,50)(-50,50)
%
%
\put(-20.00,50.00){\circle*{3}}
\put(-15.00,50.00){$h$}
\put(-60.00,50.00){\circle*{3}}
\put(-55.00,50.00){$g$}
\put(-100.00,50.00){\circle*{3}}
\put(-95.00,50.00){$f$}
\put(-140.00,50.00){\circle*{3}}
\put(-135.00,50.00){$e$}
\put(-20.00,90.00){\circle*{3}}
\put(-15.00,90.00){$d$}
\put(-60.00,90.00){\circle*{3}}
\put(-55.00,90.00){$c$}
\put(-100.00,90.00){\circle*{3}}
\put(-95.00,90.00){$b$}
\put(-140.00,90.00){\circle*{3}}
\put(-135.00,90.00){$a$}
\put(-60.00,90.00){\vector(+1,-1){40.00}}
\put(-60.00,90.00){\vector(-1,-1){40.00}}
\put(-20.00,90.00){\vector(-1,-1){40.00}}
\put(-20.00,90.00){\vector(0,-1){40.00}}
\put(-140.00,90.00){\vector(0,-1){40.00}}
\put(-140.00,90.00){\vector(+1,-1){40.00}}
\put(-100.00,90.00){\vector(-1,-1){40.00}}
\put(-100.00,90.00){\vector(+1,-1){40.00}}
\put(-160.00,40.00){\line(+1,0){160.00}}
\put(-160.00,60.00){\line(+1,0){160.00}}
\put(-160.00,40.00){\line(0,+1){20.00}}
\put(-00.00,40.00){\line(0,+1){20.00}}
%
%
\end{picture}
%
%
\caption{Example of a galaxy in ${\mathcal K}_{2}$ and ${\mathcal L}_{2}$, with dust $\emptyset$, root $\{a,b,c,d\}$ and kernel $\{e,f,g,h\}$ and where $A{=}\{a,b,c,d\}$, $B{=}\{e,f,g,h\}$, $\rho(a){=}\{e,f\}$, $\rho(b){=}\{e,g\}$, $\rho(c){=}\{f,h\}$ and $\rho(d){=}\{g,h\}$.}
%
%
\end{center}
%
%
\end{figure}
%
%
%
%
\paragraph{Flowers}
%
%
Let ${\mathcal F}_{0}^{0}{=}(W_{0}^{0},R_{0}^{0})$ be the Euclidean frame (called {\it isolated point}) such that
%
%
\begin{itemize}
%
%
\item $W_{0}^{0}{=}\{0\}$,
%
%
\item $R_{0}^{0}{=}\emptyset$.
%
%
\end{itemize}
%
%
%\begin{figure}[ht]
%
%
%\begin{center}
%
%
%\begin{picture}(-50,50)(-50,50)
%
%
%\put(-80.00,50.00){\circle*{3}}
%
%
%\end{picture}
%
%
%\caption{Isolated point.}\label{figure:lemma:about:isolated:point}
%
%
%\end{center}
%
%
%\end{figure}
%
%
%\begin{lemma}
%
%
%$\dust({\mathcal F}_{0}^{0}){=}\{0\}$, $\root({\mathcal F}_{0}^{0}){=}\emptyset$ and $\kernel({\mathcal F}_{0}^{0}){=}\emptyset$.
Obviously, $\dust({\mathcal F}_{0}^{0}){=}\{0\}$, $\root({\mathcal F}_{0}^{0}){=}\emptyset$ and $\kernel({\mathcal F}_{0}^{0}){=}\emptyset$.
%
%
%\end{lemma}
%
%
%\begin{lemma}
%
%
%The isolated point is a finite galaxy.
Moreover, the isolated point is a finite galaxy.
%
%
%\end{lemma}
%
%
%\begin{lemma}\label{lemma:about:flowers:simple:galaxies:F00}
%
%
%${\mathcal F}_{0}^{0}$ is not simple.
In other respect, ${\mathcal F}_{0}^{0}$ is not simple.
%
%
%\end{lemma}
\\
\\
%
%
%For all sets $A,B$ and for all functions $\rho:\ A{\longrightarrow}\wp(B)$, if $A{\cap}B{=}\emptyset$ and $A{\cup}B{\not=}\emptyset$ then the $(A,B,\rho)$-galaxy is {\it simple}\/ if for all $s{\in}A$, ${\parallel}\rho(s){\parallel}{\leq}1$.
%For all sets $A,B$ and for all functions $\rho:\ A{\longrightarrow}\wp(B)$, if $A{\cap}B{=}\emptyset$ and $A{\cup}B{\not=}\emptyset$ then for all $C{\in}\wp(B)$, let $\rho^{{-}1}(C){=}\{t{\in}A:\ \rho(t){=}C\}$.
For all $m{\in}\mathbf{N}^{+}$ and for all $n{\in}\mathbf{N}^{-}$, let ${\mathcal F}_{m}^{n}{=}(W_{m}^{n},R_{m}^{n})$ be the Euclidean frame (called {\it flower}) such that if $n{\in}\N$ then
%
%
\begin{itemize}
%
%
\item $W_{m}^{n}{=}\lsem0,m{+}n\rsem$,
%
%
\item $R_{m}^{n}{=}(\{0\}{\times}\lsem1,m\rsem){\cup}(\lsem1,m{+}n\rsem{\times}\lsem1,m{+}n\rsem)$,
%
%
\end{itemize}
%
%
and if $n{=}{-}1$ then
%
%
\begin{itemize}
%
%
\item $W_{m}^{n}{=}\lsem1,m\rsem$,
%
%
\item $R_{m}^{n}{=}\lsem1,m\rsem{\times}\lsem1,m\rsem$.
%
%
\end{itemize}
%
%
Obviously, for all flowers ${\mathcal F}_{m}^{n}$, if $n{\in}\N$ then $\dust({\mathcal F}_{m}^{n}){=}\emptyset$, $\root({\mathcal F}_{m}^{n}){=}\{0\}$ and $\kernel({\mathcal F}_{m}^{n}){=}\lsem1,m{+}n\rsem$ and if $n{=}{-}1$ then $\dust({\mathcal F}_{m}^{n}){=}\emptyset$, $\root({\mathcal F}_{m}^{n}){=}\emptyset$ and
%
%
%
%
%
%
%
%
%
%
%
%
\linebreak$
%
%
%
%
%
%
%
%
%
%
%
%
\kernel({\mathcal F}_{m}^{n}){=}\lsem1,m\rsem$.
Moreover, every flower is a finite galaxy.
%
%
\begin{figure}[ht]
%
%
\begin{center}
%
%
\begin{picture}(-50,50)(-50,50)
%
%
\put(-120.00,50.00){\circle*{3}}
\put(-115.00,50.00){$3$}
\put(-160.00,50.00){\circle*{3}}
\put(-155.00,50.00){$2$}
\put(-200.00,50.00){\circle*{3}}
\put(-195.00,50.00){$1$}
\put(-160.00,90.00){\circle*{3}}
\put(-155.00,90.00){$0$}
\put(-160.00,90.00){\vector(-1,-1){40.00}}
\put(-160.00,90.00){\vector(0,-1){40.00}}
\put(-220.00,40.00){\line(+1,0){120.00}}
\put(-220.00,60.00){\line(+1,0){120.00}}
\put(-220.00,40.00){\line(0,+1){20.00}}
\put(-100.00,40.00){\line(0,+1){20.00}}
%
%
\put(40.00,50.00){\circle*{3}}
\put(45.00,50.00){$3$}
\put(00.00,50.00){\circle*{3}}
\put(+05.00,50.00){$2$}
\put(-40.00,50.00){\circle*{3}}
\put(-35.00,50.00){$1$}
\put(-60.00,40.00){\line(+1,0){120.00}}
\put(-60.00,60.00){\line(+1,0){120.00}}
\put(-60.00,40.00){\line(0,+1){20.00}}
\put(+60.00,40.00){\line(0,+1){20.00}}
%
%
\end{picture}
%
%
\caption{Flowers ${\mathcal F}_{2}^{1}$ and ${\mathcal F}_{3}^{-1}$.}
%
%
\end{center}
%
%
\end{figure}
%
%
%\begin{lemma}
%
%
%For all flowers ${\mathcal F}_{m}^{n}$, if $n{\in}\N$ then $\dust({\mathcal F}_{m}^{n}){=}\emptyset$, $\root({\mathcal F}_{m}^{n}){=}\{0\}$ and $\kernel({\mathcal F}_{m}^{n}){=}\lsem1,m{+}n\rsem$ and if $n{=}{-}1$ then $\dust({\mathcal F}_{m}^{n}){=}\emptyset$, $\root({\mathcal F}_{m}^{n}){=}\emptyset$ and $\kernel({\mathcal F}_{m}^{n}){=}\lsem1,m\rsem$.
%
%
%\end{lemma}
%
%
%\begin{lemma}
%
%
%Every flower is a finite galaxy.
%
%
%\end{lemma}
%
%
\begin{lemma}\label{lemma:about:flowers:simple:galaxies}
%
%
For all flowers ${\mathcal F}_{m}^{n}$, ${\mathcal F}_{m}^{n}$ is simple if and only if either $m{=}1$, or $n{=}{-}1$, or $n{=}0$, or $n=1$.
%
%
\end{lemma}
%
%
Notice that for all $m{\in}\mathbf{N}^{+}{\cup}\{0\}$ and for all $n{\in}\mathbf{N}^{-}$, if either $m{\not=}0$, or $n{=}0$ then either ${\mathcal F}_{m}^{n}$ is the isolated point, or ${\mathcal F}_{m}^{n}$ is a flower.
%
%
%
%
\section{Generated subframes, bounded morphic images and Ehrenfeucht-Fra\"\i ss\'e games}
%
%
In this section, we introduce the standard concepts of generated subframes, bounded morphic images and Ehrenfeucht-Fra\"\i ss\'e games.
%
%
%
%
\paragraph{Generated subframes}
%
%
Let $(W,R),(W^{\prime},R^{\prime})$ be frames.
%
%
\\
\\
%
%
$(W,R)$ is a {\it generated subframe of $(W^{\prime},R^{\prime})$}\/ if the following conditions are satisfied:
%
%
\begin{itemize}
%
%
\item $W{\subseteq}W^{\prime}$,
%
%
\item $R{\subseteq}R^{\prime}$,
%
%
\item for all $s{\in}W$ and for all $t^{\prime}{\in}W^{\prime}$, if $s{R^{\prime}}t^{\prime}$ then $t^{\prime}{\in}W$ and $s{R}t^{\prime}$.
%
%
\end{itemize}
%
%
For all $s{\in}W$, let $(W_{s},R_{s})$ be the least generated subframe of $(W,R)$ containing $s$.
For all $s{\in}W$, {\it $s$ generates $(W,R)$}\/ if the following conditions are satisfied:
%
%
\begin{itemize}
%
%
\item $W_{s}{=}W$,
%
%
\item $R_{s}{=}R$.
%
%
\end{itemize}
%
%
$(W,R)$ is {\it rooted}\/ if there exists $s{\in}W$ such that $s$ generates $(W,R)$.
See~\cite[Chapter~$3$]{Blackburn:deRijke:Venema:2001}.
%
%
\\
\\
%
%
Obviously, if $(W,R)$ is finite, rooted and Euclidean then either $(W,R)$ is isomorphic with the isolated point, or $(W,R)$ is isomorphic with a flower.
%
%
\begin{lemma}\label{lemma:generated:subframe:from:a:single:point:has:specific:shapes}
%
%
If $(W,R)$ is Euclidean then for all $s{\in}W$, either $R_{s}{=}W_{s}{\times}W_{s}$, or $W_{s}{=}\{s\}{\cup}C$ and $R_{s}{=}(\{s\}{\times}B){\cup}(C{\times}C)$ for some sets $B,C$ such that either $B{\not=}\emptyset$, or $C{=}\emptyset$, $B{\subseteq}C$ and $s{\not\in}C$.
%
%
\end{lemma}
%
%
\begin{lemma}\label{lemma:about:generated:subframe:and:disjoint:unions}
%
%
Let $\{(W_{i},R_{i}):\ i{\in}I\}$ be a disjoint indexed family of Euclidean frames.
If $(W,R)$ is the union of $\{(W_{i},R_{i}):\ i{\in}I\}$ then for all $i{\in}I$, $(W_{i},R_{i})$ is a generated subframe of $(W,R)$.
%
%
\end{lemma}
%
%
%\begin{lemma}
%
%
%For all galaxies ${\mathcal F}_{A,B}^{\rho},{\mathcal F}_{A^{\prime},B^{\prime}}^{\rho^{\prime}}$, the following conditions are equivalent:
%
%
%\begin{enumerate}
%
%
%\item ${\mathcal F}_{A,B}^{\rho}$ is a generated subframe of ${\mathcal F}_{A^{\prime},B^{\prime}}^{\rho^{\prime}}$,
%
%
%\item $A{\subseteq}A^{\prime}$, either $B{=}\emptyset$, or $B{=}B^{\prime}$ and for all $s{\in}A$, $\rho(s){=}\rho^{\prime}(s)$.
%
%
%\end{enumerate}
%
%
%\end{lemma}
%
%
%\begin{lemma}
%
%
%For all flowers ${\mathcal F}_{m}^{n},{\mathcal F}_{m^{\prime}}^{n^{\prime}}$, the following conditions are equivalent:
%
%
%\begin{enumerate}
%
%
%\item ${\mathcal F}_{m}^{n}$ is a generated subframe of ${\mathcal F}_{m^{\prime}}^{n^{\prime}}$,
%
%
%\item either $n{\in}\N$, $n^{\prime}{\in}\N$, $m{=}m^{\prime}$ and $n{=}n^{\prime}$, or $n{=}{-}1$, $n^{\prime}{\in}\N$ and $m{=}m^{\prime}{+}n^{\prime}$, or $n{=}{-}1$, $n^{\prime}{=}{-}1$ and $m{=}m^{\prime}$.
%
%
%\end{enumerate}
%
%
%\end{lemma}
%
%
%\begin{lemma}\label{lemma:about:finite:rooted:Euclidean:frames}
%
%
%If $(W,R)$ is finite, rooted and Euclidean then either $(W,R)$ is isomorphic with the isolated point, or $(W,R)$ is isomorphic with a flower.
%
%
%\end{lemma}
%
%
%
%
\paragraph{Bounded morphic images}
%
%
Let $q{\geq}3$.
%
%
\\
\\
%
%
Let $(W,R),(W^{\prime},R^{\prime})$ be frames.
%
%
\\
\\
%
%
A function $f:\ W^{\prime}{\longrightarrow}W$ is a {\it bounded morphism from $(W^{\prime},R^{\prime})$ to $(W,R)$}\/ if the following conditions are satisfied:
%
%
\begin{itemize}
%
%
\item for all $s^{\prime},t^{\prime}{\in}W^{\prime}$, if $s^{\prime}{R^{\prime}}t^{\prime}$ then $f(s^{\prime}){R}f(t^{\prime})$,
%
%
\item for all $s^{\prime}{\in}W^{\prime}$ and for all $t{\in}W$, if $f(s^{\prime}){R}t$ then there exists $t^{\prime}{\in}W^{\prime}$ such that $s^{\prime}{R^{\prime}}t^{\prime}$ and $f(t^{\prime})=t$.
%
%
\end{itemize}
%
%
$(W,R)$ is a {\it bounded morphic image of $(W^{\prime},R^{\prime})$}\/ if there exists a surjective bounded morphism from $(W^{\prime},R^{\prime})$ to $(W,R)$.
See~\cite[Chapter~$3$]{Blackburn:deRijke:Venema:2001}.
%
%
%\begin{lemma}
%
%
%For all galaxies ${\mathcal F}_{A,B}^{\rho},{\mathcal F}_{A^{\prime},B^{\prime}}^{\rho^{\prime}}$, the following conditions are equivalent:
%
%
%\begin{enumerate}
%
%
%\item ${\mathcal F}_{A,B}^{\rho}$ is a bounded morphic image of ${\mathcal F}_{A^{\prime},B^{\prime}}^{\rho^{\prime}}$,
%
%
%\item zzzzz
%
%
%\end{enumerate}
%
%
%\end{lemma}
%
%
\begin{lemma}\label{lemma:about:bounded:morphisms:in:the:situation:of F:m:n:frames}
%
%
For all flowers ${\mathcal F}_{m}^{n},{\mathcal F}_{m^{\prime}}^{n^{\prime}}$, ${\mathcal F}_{m}^{n}$ is a bounded morphic image of ${\mathcal F}_{m^{\prime}}^{n^{\prime}}$ if and only if $(m,n){\ll}(m^{\prime},n^{\prime})$.
%
%
\end{lemma}
%
%
\begin{lemma}\label{lemma:about:k:s:new:element:bmi:of:W:R}
%
%
Let $k$ be an integer such that $k{\geq}4$.
Let $B,C$ be sets such that $B{\subseteq}C$, ${\parallel}C{\parallel}{>}k$, ${\parallel}B{\parallel}{>}1$ and ${\parallel}C{\setminus}B{\parallel}{>}1$.
Let $m{=}\min\{{\parallel}B{\parallel},k\}$ and $n{=}\min\{{\parallel}C{\setminus}B{\parallel},k\}$.
Let $s$ be a new element.
If $W{=}\{s\}{\cup}C$ and $R{=}(\{s\}{\times}B){\cup}(C{\times}C)$ then ${\mathcal F}_{m}^{n}$ is a bounded morphic image of $(W,R)$.
%
%
\end{lemma}
%
%
\begin{lemma}\label{lemma:about:partitions:q:greater:smaller}
%
%
If $R{=}W{\times}W$, $R^{\prime}{=}W^{\prime}{\times}W^{\prime}$ and either ${\parallel}W^{\prime}{\parallel}{\leq}q$ and ${\parallel}W^{\prime}{\parallel}{=}{\parallel}W{\parallel}$, or ${\parallel}W^{\prime}{\parallel}{>}q$ and ${\parallel}W{\parallel}{=}q$ then $(W,R)$ is a bounded morphic image of $(W^{\prime},R^{\prime})$.
%
%
\end{lemma}
%
%
\begin{lemma}\label{lemma:inclusion:Bprime:B:for:generated:subframes:is:bmi}
%
%
If $W{=}\{s\}{\cup}B$, $R{=}(\{s\}{\times}B){\cup}(B{\times}B)$, $W^{\prime}{=}\{s^{\prime}\}{\cup}B^{\prime}$ and $R^{\prime}{=}(\{s^{\prime}\}{\times}
%
%
%
%
%
%
%
%
%
%
%
%
$\linebreak$
%
%
%
%
%
%
%
%
%
%
%
%
B^{\prime}){\cup}(B^{\prime}{\times}B^{\prime})$ for some sets $B,B^{\prime}$ and for some elements $s,s^{\prime}$ such that $s{\not\in}B$, $B{\not=}\emptyset$, $s^{\prime}{\not\in}B^{\prime}$, $B^{\prime}{\not=}\emptyset$ and ${\parallel}B{\parallel}{\leq}{\parallel}B^{\prime}{\parallel}$ then $(W,R)$ is a bounded morphic image of $(W^{\prime},R^{\prime})$.
%
%
\end{lemma}
%
%
\begin{lemma}\label{lemma:inclusion:Bprime:B:for:generated:subframes:is:bmi:bis}
%
%
If $W{=}\{s\}{\cup}B$, $R{=}\{(s,u)\}{\cup}(B{\times}B)$, $W^{\prime}{=}\{s^{\prime}\}{\cup}B^{\prime}$ and $R^{\prime}{=}\{(s^{\prime},
%
%
%
%
%
%
%
%
%
%
%
%
$\linebreak$
%
%
%
%
%
%
%
%
%
%
%
%
u^{\prime})\}{\cup}(B^{\prime}{\times}B^{\prime})$ for some sets $B,B^{\prime}$ and for some elements $s,s^{\prime},u,u^{\prime}$ such that $s{\not\in}B$, $u{\in}B$, $B{\not=}\emptyset$, $s^{\prime}{\not\in}B^{\prime}$, $u^{\prime}{\in}B^{\prime}$, $B^{\prime}{\not=}\emptyset$, ${\parallel}B{\parallel}{\leq}{\parallel}B^{\prime}{\parallel}$ and if ${\parallel}B{\parallel}{\leq}{\parallel}2$ then ${\parallel}B{\parallel}{=}
%
%
%
%
%
%
%
%
%
%
%
%
$\linebreak$
%
%
%
%
%
%
%
%
%
%
%
%
{\parallel}B^{\prime}{\parallel}$ then $(W,R)$ is a bounded morphic image of $(W^{\prime},R^{\prime})$.
%
%
\end{lemma}
%
%
\begin{lemma}\label{lemma:inclusion:Bprime:B:for:generated:subframes:is:bmi:ter}
%
%
If $W{=}\{s\}{\cup}B$, $R{=}(\{s\}{\times}(B{\setminus}\{u\})){\cup}(B{\times}B)$, $W^{\prime}{=}\{s^{\prime}\}{\cup}B^{\prime}$ and $R^{\prime}{=}
%
%
%
%
%
%
%
%
%
%
%
%
$\linebreak$
%
%
%
%
%
%
%
%
%
%
%
%
(\{s^{\prime}\}{\times}(B^{\prime}{\setminus}\{u^{\prime}\})){\cup}(B^{\prime}{\times}B^{\prime})$ for some sets $B,B^{\prime}$ and for some elements $s,s^{\prime},u,u^{\prime}$ such that $s{\not\in}B$, $u{\in}B$, $B{\not=}\emptyset$, $s^{\prime}{\not\in}B^{\prime}$, $u^{\prime}{\in}B^{\prime}$, $B^{\prime}{\not=}\emptyset$, ${\parallel}B{\parallel}{\leq}{\parallel}B^{\prime}{\parallel}$ and if ${\parallel}B{\parallel}{\leq}{\parallel}2$ then ${\parallel}B{\parallel}{=}{\parallel}B^{\prime}{\parallel}$ then $(W,R)$ is a bounded morphic image of $(W^{\prime},R^{\prime})$.
%
%
\end{lemma}
%
%
\begin{lemma}\label{lemma:inclusion:Bprime:B:for:generated:subframes:is:bmi:quater}
%
%
If $W{=}B$, $R{=}B{\times}B$, $W^{\prime}{=}B^{\prime}$ and $R^{\prime}{=}B^{\prime}{\times}B^{\prime}$ for some sets $B,B^{\prime}$ such that $B{\not=}\emptyset$, $B^{\prime}{\not=}\emptyset$, ${\parallel}B{\parallel}{\leq}{\parallel}B^{\prime}{\parallel}$ and if ${\parallel}B{\parallel}{\leq}{\parallel}2$ then ${\parallel}B{\parallel}{=}{\parallel}B^{\prime}{\parallel}$ then $(W,R)$ is a bounded morphic image of $(W^{\prime},R^{\prime})$.
%
%
\end{lemma}
%
%
\begin{lemma}\label{lemma:about:frames:C:B:smaller:greater}
%
%
If $W{=}\{s\}{\cup}C$, $R{=}(\{s\}{\times}B){\cup}(C{\times}C)$, $W^{\prime}{=}\{s^{\prime}\}{\cup}C^{\prime}$ and $R^{\prime}{=}(\{s^{\prime}\}{\times}
%
%
%
%
%
%
%
%
%
%
%
%
$\linebreak$
%
%
%
%
%
%
%
%
%
%
%
%
B^{\prime}){\cup}(C^{\prime}{\times}C^{\prime})$ for some sets $B,C,B^{\prime},C^{\prime}$ and for some elements $s,s^{\prime}$ such that either $B{\not=}\emptyset$, or $C{=}\emptyset$, $B{\subseteq}C$, $s{\not\in}C$, either $B^{\prime}{\not=}\emptyset$, or $C^{\prime}{=}\emptyset$, $B^{\prime}{\subseteq}C^{\prime}$, $s^{\prime}{\not\in}C^{\prime}$, either ${\parallel}B^{\prime}{\parallel}{\leq}q$ and ${\parallel}B^{\prime}{\parallel}{=}{\parallel}B{\parallel}$, or ${\parallel}B^{\prime}{\parallel}{>}q$ and ${\parallel}B{\parallel}{=}q$ and either ${\parallel}C^{\prime}{\setminus}B^{\prime}{\parallel}{\leq}q$ and ${\parallel}C^{\prime}{\setminus}B^{\prime}{\parallel}{=}{\parallel}C{\setminus}B{\parallel}$, or ${\parallel}C^{\prime}{\setminus}B^{\prime}{\parallel}{>}q$ and ${\parallel}C{\setminus}B{\parallel}{=}q$ then $(W,R)$ is a bounded morphic image of $(W^{\prime},R^{\prime})$.
%
%
\end{lemma}
%
%
\begin{lemma}\label{lemma:about:what:happens:when:it:exists:bmi:flower:integer}
%
%
If $(W,R)$ is Euclidean then for all $s{\in}W$ and for all flowers ${\mathcal F}_{m}^{n}$, if $n{\in}\N$ and there exists a surjective bounded morphism $f$ from $(W_{s},R_{s})$ to ${\mathcal F}_{m}^{n}$ such that $f(s)$ generates ${\mathcal F}_{m}^{n}$ then $W_{s}{=}\{s\}{\cup}C$ and $R_{s}{=}(\{s\}{\times}B){\cup}(C{\times}C)$ for some sets $B,C$ such that $B{\subseteq}C$, $s{\not\in}C$, ${\parallel}B{\parallel}{\geq}m$ and ${\parallel}C{\setminus}B{\parallel}{\geq}n$.
%
%
\end{lemma}
%
%
\begin{lemma}\label{lemma:about:what:happens:when:it:exists:bmi:flower:minus:one}
%
%
If $(W,R)$ is Euclidean then for all $s{\in}W$ and for all flowers ${\mathcal F}_{m}^{n}$, if $n{=}{-}1$ and there exists a surjective bounded morphism $f$ from $(W_{s},R_{s})$ to ${\mathcal F}_{m}^{n}$ such that $f(s)$ generates ${\mathcal F}_{m}^{n}$ then either $R_{s}{=}W_{s}{\times}W_{s}$ and ${\parallel}W_{s}{\parallel}{\geq}m$, or $W_{s}{=}\{s\}{\cup}C$ and $R_{s}{=}(\{s\}{\times}B){\cup}(C{\times}C)$ for some sets $B,C$ such that $B{\subseteq}C$, $s{\not\in}C$ and ${\parallel}B{\parallel}{\geq}m$.
%
%
\end{lemma}
%
%
%
%
\paragraph{Ehrenfeucht-Fra\"\i ss\'e games}
%
%
Let $q{\geq}3$.
%
%
\\
\\
%
%
Let $(W,R)$ and $(W^{\prime},R^{\prime})$ be Euclidean frames.
%
%
\\
\\
%
%
The {\it Ehrenfeucht-Fra\"\i ss\'e game ${\mathcal G}_{q}((W,R),(W^{\prime},R^{\prime}))$}\/ is a game played by $2$~{\it players}\/ making $q$~{\it moves.}
For all $q^{\prime}{\in}\lsem1,q\rsem$, in its $q^{\prime}$-th move, the first player picks either a member $s_{q^{\prime}}$ of $(W,R)$, or a member $s_{q^{\prime}}^{\prime}$ of $(W^{\prime},R^{\prime})$.
For all $q^{\prime}{\in}\lsem1,q\rsem$, if the first player picked a member $s_{q^{\prime}}$ of $(W,R)$ in its $q^{\prime}$-th move then the second player picks a member $s_{q^{\prime}}^{\prime}$ of $(W^{\prime},R^{\prime})$ in its $q^{\prime}$-th move, whereas if the first player picked a member $s_{q^{\prime}}^{\prime}$ of $(W^{\prime},R^{\prime})$ in its $q^{\prime}$-th move then the second player picks a member $s_{q^{\prime}}$ of $(W,R)$ in its $q^{\prime}$-th move.
%At the end of each {\it play,} elements $s_{1},\ldots,s_{q}$ in $W$ and elements $s^{\prime}_{1},\ldots,s^{\prime}_{q}$ in $W^{\prime}$ have been chosen.
%
%
\\
\\
%
%
For all $q^{\prime}{\in}\lsem0,q\rsem$, for all $s_{1},\ldots,s_{q^{\prime}}{\in}W$ and for all $s^{\prime}_{1},\ldots,s^{\prime}_{q^{\prime}}{\in}W^{\prime}$, the couple $((s_{1},\ldots,s_{q^{\prime}}),(s^{\prime}_{1},\ldots,s^{\prime}_{q^{\prime}}))$ is a {\it position of the game.}
When $q^{\prime}{=}0$, this position is called the {\it empty position of the game.}
%
%
\\
\\
%
%
For all $q^{\prime}{\in}\lsem0,q\rsem$, for all $s_{1},\ldots,s_{q^{\prime}}{\in}W$ and for all $s^{\prime}_{1},\ldots,s^{\prime}_{q^{\prime}}{\in}W^{\prime}$, the second player {\it wins the position $((s_{1},\ldots,s_{q^{\prime}}),(s^{\prime}_{1},\ldots,s^{\prime}_{q^{\prime}}))$ of the game}\/ if for all $q^{\prime\prime},q^{\prime\prime\prime}{\in}
%
%
%
%
%
%
%
%
%
%
%
%
$\linebreak$
%
%
%
%
%
%
%
%
%
%
%
%
\lsem1,q^{\prime}\rsem$,
%
%
\begin{itemize}
%
%
\item $s_{q^{\prime\prime}}{=}s_{q^{\prime\prime\prime}}$ if and only if $s^{\prime}_{q^{\prime\prime}}{=}s^{\prime}_{q^{\prime\prime\prime}}$,
%
%
\item $s_{q^{\prime\prime}}{R}s_{q^{\prime\prime\prime}}$ if and only if $s^{\prime}_{q^{\prime\prime}}{R^{\prime}}s^{\prime}_{q^{\prime\prime\prime}}$.
%
%
\end{itemize}
%
%
%As a result, the second player wins the empty position of the game.
A {\it strategy of the second player in the game}\/ is a deterministic procedure telling how the second player chooses its moves during the game.
%
%
%\\
%\\
%
%
A strategy of the second player in the game is {\it winning}\/ if the second player wins each position of the game that is reachable by applying the strategy.
See~\cite[Chapter~$2$]{Ebbinghaus:Flum:1999}.
%
%
\begin{lemma}\label{lemma:about:partitions:q:greater:smaller:games}
%
%
If $R{=}W{\times}W$, $R^{\prime}{=}W^{\prime}{\times}W^{\prime}$ and either ${\parallel}W^{\prime}{\parallel}{\leq}q$ and ${\parallel}W^{\prime}{\parallel}{=}{\parallel}W{\parallel}$, or ${\parallel}W^{\prime}{\parallel}{>}q$ and ${\parallel}W{\parallel}{=}q$ then the second player has a winning strategy in the Ehrenfeucht-Fra\"\i ss\'e game ${\mathcal G}_{q}((W,R),(W^{\prime},R^{\prime}))$.
%
%
\end{lemma}
%
%
\begin{lemma}\label{lemma:about:frames:C:B:smaller:greater:games}
%
%
If $W{=}\{s\}{\cup}C$, $R{=}(\{s\}{\times}B){\cup}(C{\times}C)$, $W^{\prime}{=}\{s^{\prime}\}{\cup}C^{\prime}$ and $R^{\prime}{=}(\{s^{\prime}\}{\times}
%
%
%
%
%
%
%
%
%
%
%
%
$\linebreak$
%
%
%
%
%
%
%
%
%
%
%
%
B^{\prime}){\cup}(C^{\prime}{\times}C^{\prime})$ for some sets $B,C,B^{\prime},C^{\prime}$ and for some elements $s,s^{\prime}$ such that either $B{\not=}\emptyset$, or $C{=}\emptyset$, $B{\subseteq}C$, $s{\not\in}C$, either $B^{\prime}{\not=}\emptyset$, or $C^{\prime}{=}\emptyset$, $B^{\prime}{\subseteq}C^{\prime}$, $s^{\prime}{\not\in}C^{\prime}$, either ${\parallel}B^{\prime}{\parallel}{\leq}q$ and ${\parallel}B^{\prime}{\parallel}{=}{\parallel}B{\parallel}$, or ${\parallel}B^{\prime}{\parallel}{>}q$ and ${\parallel}B{\parallel}{=}q$ and either ${\parallel}C^{\prime}{\setminus}B^{\prime}{\parallel}{\leq}q$ and ${\parallel}C^{\prime}{\setminus}B^{\prime}{\parallel}{=}{\parallel}C{\setminus}B{\parallel}$, or ${\parallel}C^{\prime}{\setminus}B^{\prime}{\parallel}{>}q$ and ${\parallel}C{\setminus}B{\parallel}{=}q$ then the second player has a winning strategy in the Ehrenfeucht-Fra\"\i ss\'e game ${\mathcal G}_{q}((W,R),(W^{\prime},R^{\prime}))$.
%
%
\end{lemma}
%
%
%
%
\section{Reductions}
%
%
In this section, we introduce $3$~types of reductions: the $\alpha$-reduction of a galaxy, the $\gamma$-reduction of a simple galaxy and the $\delta$-reduction of a disjoint indexed family of galaxies.
%
%
\\
\\
%
%
Let $q{\geq}3$.
%
%
%
%
\paragraph{$\alpha$-reductions}
%
%
The main effect of an $\alpha$-reduction on a galaxy is to reduce the cardinality of its upper part.
%Moreover, the $\alpha$-reduction of a disjoint indexed family of galaxies will be dust-bounded.
%
%
\\
\\
%
%
Let ${\mathcal F}_{A,B}^{\rho}$ be a galaxy.
%
%
\\
\\
%
%
Let $\sim$ be the equivalence relation on $A$ defined by $u{\sim}v$ if and only if $\rho(u){=}\rho(v)$.
For all $u{\in}A$, let $\lbrack u\rbrack_{\sim}$ be the equivalence class of $u$ modulo $\sim$.
%
%
\\
\\
%
%
For all $u{\in}A$, if ${\parallel}\lbrack u\rbrack_{\sim}{\parallel}{\leq}q$ then let $C_{u}$ be $\lbrack u\rbrack_{\sim}$ else let $C_{u}$ be an arbitrary subset of $\lbrack u\rbrack_{\sim}$ of cardinality $q$.
Let $A^{\prime}$ be $\bigcup\{C_{u}:\ u{\in}A\}$.
Obviously, from each equivalence class modulo $\sim$, $A^{\prime}$ contains at most $q$ elements.
Therefore, if $B$ is finite then $A^{\prime}$ contains at most $q{\times}2^{{\parallel}B{\parallel}}$ elements.
%
%
%\\
%\\
%
%
%Let $g:\ A{\longrightarrow}A^{\prime}$ be a function such that for all $s{\in}A$, either $s{\in}A^{\prime}$ and $g(s){=}s$, or $s{\not\in}A^{\prime}$ and $g(s){\sim}s$.
%Notice that the existence of such a function is a consequence of the definition of $A^{\prime}$.
%Notice also that $g$ is surjective.
%Moreover, for all $s{\in}A$, $g(s){\sim}s$.
%
%
\\
\\
%
%
Let $B^{\prime}$ be $B$.
%
%
\\
\\
%
%
Let ${\mathcal F}_{A^{\prime},B^{\prime}}^{\rho^{\prime}}$ be the galaxy such that for all $u^{\prime}{\in}A^{\prime}$, $\rho^{\prime}(u^{\prime}){=}\rho(u^{\prime})$.
${\mathcal F}_{A^{\prime},B^{\prime}}^{\rho^{\prime}}$ is a {\it $q$-$\alpha$-reduction of ${\mathcal F}_{A,B}^{\rho}$.}
%
%
%\\
%\\
%
%
%Let $f:\ A{\cup}B{\longrightarrow}A^{\prime}{\cup}B^{\prime}$ be the function such that for all $s{\in}A$, $f(s){=}g(s)$ and for all $s{\in}B$, $f(s){=}s$.
%Obviously, $(W_{A^{\prime},B^{\prime}}^{\rho^{\prime}},R_{A^{\prime},B^{\prime}}^{\rho^{\prime}})$ is simple if and only if $(W_{A,B}^{\rho},R_{A,B}^{\rho})$ is simple.
%Notice that $f$ is surjective.
%
%
%\\
%\\
%
%
%Let us define a strategy $\pi$ of the second player in the Ehrenfeucht-Fra\"\i ss\'e game ${\mathcal G}_{q}({\mathcal F}_{A,B}^{\rho},{\mathcal F}_{A^{\prime},B^{\prime}}^{\rho^{\prime}})$.
%For all $q^{\prime}{\in}\lsem0,q{-}1\rsem$, we shall say that a member $s$ of $A$ is {\it left-fresh in a position $((s_{1},\ldots,s_{q^{\prime}}),(s^{\prime}_{1},\ldots,s^{\prime}_{q^{\prime}}))$ of the Ehrenfeucht-Fra\"\i ss\'e game ${\mathcal G}_{q}({\mathcal F}_{A,B}^{\rho},{\mathcal F}_{A^{\prime},B^{\prime}}^{\rho^{\prime}})$}\/ if for all $q^{\prime\prime}{\in}\lsem1,q^{\prime}\rsem$, $s{\not=}s_{q^{\prime\prime}}$.
%For all $q^{\prime}{\in}\lsem0,q{-}1\rsem$, we shall say that a member $s^{\prime}$ of $A^{\prime}$ is {\it right-fresh in a position $((s_{1},\ldots,s_{q^{\prime}}),(s^{\prime}_{1},\ldots,s^{\prime}_{q^{\prime}}))$ of the Ehrenfeucht-Fra\"\i ss\'e game ${\mathcal G}_{q}({\mathcal F}_{A,B}^{\rho},{\mathcal F}_{A^{\prime},B^{\prime}}^{\rho^{\prime}})$}\/ if for all $q^{\prime\prime}{\in}\lsem1,q^{\prime}\rsem$, $s^{\prime}{\not=}s^{\prime}_{q^{\prime\prime}}$.
%Let $q^{\prime}{\in}\lsem1,q\rsem$.
%In its $q^{\prime}$-th move, the second player will apply either $\mathbf{(\pi_{1})}$, or $\mathbf{(\pi_{2})}$, or $\mathbf{(\pi_{3})}$, or $\mathbf{(\pi_{4})}$, depending on the $q^{\prime}$-th move of the first player.
%
%
%\\
%\\
%
%
%$\mathbf{(\pi_{1})}$~If the first player picked in its $q^{\prime}$-th move a member $s$ of ${\mathcal F}_{A,B}^{\rho}$ such that $s{\in}A$ then let the second player consider the following $2$~cases.
%
%
%\\
%\\
%
%
%{\bf Case where $s$ is left-fresh in the position $((s_{1},\ldots,s_{q^{\prime}{-}1}),(s^{\prime}_{1},\ldots,s^{\prime}_{q^{\prime}{-}1}))$:}
%As a consequence of the construction of ${\mathcal F}_{A^{\prime},B^{\prime}}^{\rho^{\prime}}$, there exists $s^{\prime}{\in}A^{\prime}$ such that $s{\sim}s^{\prime}$ and $s^{\prime}$ is right-fresh in the position $((s_{1},\ldots,s_{q^{\prime}{-}1}),(s^{\prime}_{1},\ldots,s^{\prime}_{q^{\prime}{-}1}))$.
%Thus, let the second player pick in its $q^{\prime}$-th move the member $s^{\prime}$ of ${\mathcal F}_{A^{\prime},B^{\prime}}^{\rho^{\prime}}$.
%
%
%\\
%\\
%
%
%{\bf Case where $s$ is not left-fresh in the position $((s_{1},\ldots,s_{q^{\prime}{-}1}),(s^{\prime}_{1},
%
%
%
%
%
%
%
%
%
%
%
%
%$\linebreak$
%
%
%
%
%
%
%
%
%
%
%
%
%\ldots,s^{\prime}_{q^{\prime}{-}1}))$:}
%Consequently, let $q^{\prime\prime}{\in}\lsem1,q^{\prime}{-}1\rsem$ be such that $s{=}s_{q^{\prime\prime}}$.
%Hence, let the second player pick in its $q^{\prime}$-th move the member $s^{\prime}_{q^{\prime\prime}}$ of ${\mathcal F}_{A^{\prime},B^{\prime}}^{\rho^{\prime}}$.
%
%
%\\
%\\
%
%
%$\mathbf{(\pi_{2})}$~If the first player picked in its $q^{\prime}$-th move a member $s^{\prime}$ of ${\mathcal F}_{A^{\prime},B^{\prime}}^{\rho^{\prime}}$ such that $s^{\prime}{\in}A^{\prime}$ then let the second player consider the following $2$~cases.
%
%
%\\
%\\
%
%
%{\bf Case where $s^{\prime}$ is right-fresh in the position $((s_{1},\ldots,s_{q^{\prime}{-}1}),(s^{\prime}_{1},\ldots,
%
%
%
%
%
%
%
%
%
%
%
%
%$\linebreak$
%
%
%
%
%
%
%
%
%
%
%
%
%s^{\prime}_{q^{\prime}{-}1}))$:}
%As a consequence of the construction of ${\mathcal F}_{A^{\prime},B^{\prime}}^{\rho^{\prime}}$, there exists $s{\in}A$ such that $s{\sim}s^{\prime}$ and $s$ is left-fresh in the position $((s_{1},\ldots,s_{q^{\prime}{-}1}),(s^{\prime}_{1},\ldots,s^{\prime}_{q^{\prime}{-}1}))$.
%Thus, let the second player pick in its $q^{\prime}$-th move the member $s$ of ${\mathcal F}_{A,B}^{\rho}$.
%
%
%\\
%\\
%
%
%{\bf Case where $s^{\prime}$ is not right-fresh in the position $((s_{1},\ldots,s_{q^{\prime}{-}1}),(s^{\prime}_{1},\ldots,
%
%
%
%
%
%
%
%
%
%
%
%
%$\linebreak$
%
%
%
%
%
%
%
%
%
%
%
%
%s^{\prime}_{q^{\prime}{-}1}))$:}
%Consequently, let $q^{\prime\prime}{\in}\lsem1,q^{\prime}{-}1\rsem$ be such that $s^{\prime}{=}s^{\prime}_{q^{\prime\prime}}$.
%Hence, let the second player pick in its $q^{\prime}$-th move the member $s_{q^{\prime\prime}}$ of ${\mathcal F}_{A,B}^{\rho}$.
%
%
%\\
%\\
%
%
%$\mathbf{(\pi_{3})}$~If the first player picked in its $q^{\prime}$-th move a member $s$ of ${\mathcal F}_{A,B}^{\rho}$ such that $s{\in}B$ then let the second player pick in its $q^{\prime}$-th move the same member $s$ of ${\mathcal F}_{A^{\prime},B^{\prime}}^{\rho^{\prime}}$.
%
%
%\\
%\\
%
%
%$\mathbf{(\pi_{4})}$~If the first player picked in its $q^{\prime}$-th move a member $s^{\prime}$ of ${\mathcal F}_{A^{\prime},B^{\prime}}^{\rho^{\prime}}$ such that $s^{\prime}{\in}B^{\prime}$ then let the second player pick in its $q^{\prime}$-th move the same member $s^{\prime}$ of ${\mathcal F}_{A,B}^{\rho}$.
%
%
%\begin{lemma}\label{lemma:about:alpha:reductions:and:pi:reachable:positions}
%
%
%Let $q^{\prime}{\in}\lsem0,q\rsem$ and $((s_{1},\ldots,s_{q^{\prime}}),(s^{\prime}_{1},\ldots,s^{\prime}_{q^{\prime}}))$ be a $\pi$-reachable position in the Ehrenfeucht-Fra\"\i ss\'e game ${\mathcal G}_{q}({\mathcal F}_{A,B}^{\rho},{\mathcal F}_{A^{\prime},B^{\prime}}^{\rho^{\prime}})$.
%For all $q^{\prime\prime}{\in}\lsem1,q^{\prime}\rsem$, $s_{q^{\prime\prime}}{\in}A$ if and only if $s^{\prime}_{q^{\prime\prime}}{\in}A^{\prime}$ and $s_{q^{\prime\prime}}{\in}B$ if and only if $s^{\prime}_{q^{\prime\prime}}{\in}B^{\prime}$.
%Moreover, if $s_{q^{\prime\prime}}{\in}A$ and $s^{\prime}_{q^{\prime\prime}}{\in}A^{\prime}$ then $s_{q^{\prime\prime}}{\sim}s^{\prime}_{q^{\prime\prime}}$ and if $s_{q^{\prime\prime}}{\in}B$ and $s^{\prime}_{q^{\prime\prime}}{\in}B^{\prime}$ then $s_{q^{\prime\prime}}{=}s^{\prime}_{q^{\prime\prime}}$.
%
%
%\end{lemma}
%
%
%For all $q^{\prime}{\in}\lsem0,q\rsem$, we shall say that a position $((s_{1},\ldots,s_{q^{\prime}}),(s^{\prime}_{1},\ldots,s^{\prime}_{q^{\prime}}))$ of the Ehrenfeucht-Fra\"\i ss\'e game ${\mathcal G}_{q}({\mathcal F}_{A,B}^{\rho},{\mathcal F}_{A^{\prime},B^{\prime}}^{\rho^{\prime}})$ is {\it good}\/ if for all $q^{\prime\prime}{\in}\lsem1,q^{\prime}\rsem$, at least one of the following conditions holds:
%for all $s{\in}A{\cup}B$ and for all $s^{\prime}{\in}A^{\prime}{\cup}B^{\prime}$, we shall say that the pair $(s,s^{\prime})$ is {\it appropriate for a position $((s_{1},\ldots,s_{q^{\prime}}),(s^{\prime}_{1},\ldots,s^{\prime}_{q^{\prime}}))$ of the Ehrenfeucht-Fra\"\i ss\'e game ${\mathcal G}_{q}({\mathcal F}_{A,B}^{\rho},{\mathcal F}_{A^{\prime},B^{\prime}}^{\rho^{\prime}})$}\/ if at least one of the following conditions holds:
%
%
%\begin{itemize}
%
%
%\item $s_{q^{\prime\prime}}{\in}A$, $s^{\prime}_{q^{\prime\prime}}{\in}A^{\prime}$, $s_{q^{\prime\prime}}$ is left-fresh in $((s_{1},\ldots,s_{q^{\prime\prime}{-}1}),(s^{\prime}_{1},\ldots,s^{\prime}_{q^{\prime\prime}{-}1}))$, $s^{\prime}_{q^{\prime\prime}}$ is right-fresh in $((s_{1},\ldots,s_{q^{\prime\prime}{-}1}),(s^{\prime}_{1},\ldots,s^{\prime}_{q^{\prime\prime}{-}1}))$ and $s_{q^{\prime\prime}}{\sim}s^{\prime}_{q^{\prime\prime}}$,
%
%
%\item $s_{q^{\prime\prime}}{\in}A$, $s^{\prime}_{q^{\prime\prime}}{\in}A^{\prime}$, $s_{q^{\prime\prime}}$ is not left-fresh in $((s_{1},\ldots,s_{q^{\prime\prime}{-}1}),(s^{\prime}_{1},\ldots,s^{\prime}_{q^{\prime\prime}{-}1}))$, $s^{\prime}_{q^{\prime\prime}}$ is not right-fresh in $((s_{1},\ldots,s_{q^{\prime\prime}{-}1}),(s^{\prime}_{1},\ldots,s^{\prime}_{q^{\prime\prime}{-}1}))$ and there exists $q^{\prime\prime\prime}{\in}\lsem1,q^{\prime\prime}{-}1\rsem$ such that $s_{q^{\prime\prime\prime}}{\in}A$, $s_{q^{\prime\prime\prime}}{=}s_{q^{\prime\prime}}$, $s^{\prime}_{q^{\prime\prime\prime}}{\in}A^{\prime}$ and $s^{\prime}_{q^{\prime\prime\prime}}{=}s^{\prime}_{q^{\prime\prime}}$,
%
%
%\item $s_{q^{\prime\prime}}{\in}B$, $s^{\prime}_{q^{\prime\prime}}{\in}B^{\prime}$ and $s_{q^{\prime\prime}}{=}s^{\prime}_{q^{\prime\prime}}$.
%
%
%\end{itemize}
%
%
%\begin{lemma}\label{lemma:about:good:positions:and:left:right:fresh:elements}
%
%
%For all $q^{\prime}{\in}\lsem0,q\rsem$ and for all positions $((s_{1},\ldots,s_{q^{\prime}}),(s^{\prime}_{1},\ldots,s^{\prime}_{q^{\prime}}))$ of the Ehrenfeucht-Fra\"\i ss\'e game ${\mathcal G}_{q}({\mathcal F}_{A,B}^{\rho},{\mathcal F}_{A^{\prime},B^{\prime}}^{\rho^{\prime}})$, if $((s_{1},\ldots,s_{q^{\prime}}),(s^{\prime}_{1},\ldots,s^{\prime}_{q^{\prime}}))$ is good then for all $q^{\prime\prime}{\in}\lsem1,q^{\prime}\rsem$, if $s_{q^{\prime\prime}}{\in}A$ and $s^{\prime}_{q^{\prime\prime}}{\in}A^{\prime}$then $s_{q^{\prime\prime}}{\sim}s^{\prime}_{q^{\prime\prime}}$ and if $s_{q^{\prime\prime}}{\in}B$ and $s^{\prime}_{q^{\prime\prime}}{\in}B^{\prime}$then $s_{q^{\prime\prime}}{=}s^{\prime}_{q^{\prime\prime}}$.
%
%
%\end{lemma}
%
%
%\begin{lemma}\label{lemma:about:good:positions:and:sim:relation}
%
%
%For all $q^{\prime}{\in}\lsem0,q{-}1\rsem$ and for all positions $((s_{1},\ldots,s_{q^{\prime}}),(s^{\prime}_{1},\ldots,s^{\prime}_{q^{\prime}}))$ of the Ehrenfeucht-Fra\"\i ss\'e game ${\mathcal G}_{q}({\mathcal F}_{A,B}^{\rho},{\mathcal F}_{A^{\prime},B^{\prime}}^{\rho^{\prime}})$, if $((s_{1},\ldots,s_{q^{\prime}}),(s^{\prime}_{1},\ldots,s^{\prime}_{q^{\prime}}))$ is good then
%
%
%\begin{itemize}
%
%
%\item for all $s{\in}A$, if $s$ is left-fresh in $((s_{1},\ldots,s_{q^{\prime}}),(s^{\prime}_{1},\ldots,s^{\prime}_{q^{\prime}}))$ then there exists $s^{\prime}{\in}A^{\prime}$ such that $s^{\prime}$ is right-fresh in $((s_{1},\ldots,s_{q^{\prime}}),(s^{\prime}_{1},\ldots,s^{\prime}_{q^{\prime}}))$ and $s{\sim}s^{\prime}$,
%
%
%\item for all $s{\in}B$, if $s$ is left-fresh in $((s_{1},\ldots,s_{q^{\prime}}),(s^{\prime}_{1},\ldots,s^{\prime}_{q^{\prime}}))$ then there exists $s^{\prime}{\in}B^{\prime}$ such that $s^{\prime}$ is right-fresh in $((s_{1},\ldots,s_{q^{\prime}}),(s^{\prime}_{1},\ldots,s^{\prime}_{q^{\prime}}))$ and $s{=}s^{\prime}$,
%
%
%\item for all $s^{\prime}{\in}A^{\prime}$, if $s^{\prime}$ is right-fresh in $((s_{1},\ldots,s_{q^{\prime}}),(s^{\prime}_{1},\ldots,s^{\prime}_{q^{\prime}}))$ then there exists $s{\in}A$ such that $s$ is left-fresh in $((s_{1},\ldots,s_{q^{\prime}}),(s^{\prime}_{1},\ldots,s^{\prime}_{q^{\prime}}))$ and $s^{\prime}{\sim}s$,
%
%
%\item for all $s^{\prime}{\in}B^{\prime}$, if $s^{\prime}$ is right-fresh in $((s_{1},\ldots,s_{q^{\prime}}),(s^{\prime}_{1},\ldots,s^{\prime}_{q^{\prime}}))$ then there exists $s{\in}B$ such that $s$ is left-fresh in $((s_{1},\ldots,s_{q^{\prime}}),(s^{\prime}_{1},\ldots,s^{\prime}_{q^{\prime}}))$ and $s^{\prime}{=}s$.
%
%
%\end{itemize}
%
%
%\end{lemma}
%
%
%\begin{lemma}\label{lemma:good:positions:are:winning:alpha}
%
%
%For all $q^{\prime}{\in}\lsem0,q\rsem$ and for all positions $((s_{1},\ldots,s_{q^{\prime}}),(s^{\prime}_{1},\ldots,s^{\prime}_{q^{\prime}}))$ of the Ehrenfeucht-Fra\"\i ss\'e game ${\mathcal G}_{q}({\mathcal F}_{A,B}^{\rho},{\mathcal F}_{A^{\prime},B^{\prime}}^{\rho^{\prime}})$, if $((s_{1},\ldots,s_{q^{\prime}}),(s^{\prime}_{1},\ldots,s^{\prime}_{q^{\prime}}))$ is good then the second player wins the position $((s_{1},\ldots,s_{q^{\prime}}),(s^{\prime}_{1},\ldots,s^{\prime}_{q^{\prime}}))$ of the game.
%
%
%\end{lemma}
%
%
%\begin{lemma}\label{lemma:empty:position:is:good:alpha}
%
%
%The empty position of the Ehrenfeucht-Fra\"\i ss\'e game ${\mathcal G}_{q}({\mathcal F}_{A,B}^{\rho},
%
%
%
%
%
%
%
%
%
%
%
%
%$\linebreak$
%
%
%
%
%
%
%
%
%
%
%
%
%{\mathcal F}_{A^{\prime},B^{\prime}}^{\rho^{\prime}})$ is good.
%
%
%\end{lemma}
%
%
%\begin{lemma}\label{lemma:about:extending:winning:positions:a}
%
%
%For all $q^{\prime}{\in}\lsem0,q{-}1\rsem$, for all winning positions $((s_{1},\ldots,s_{q^{\prime}}),(s^{\prime}_{1},\ldots,
%
%
%
%
%
%
%
%
%
%
%
%
%$\linebreak$
%
%
%
%
%
%
%
%
%
%
%
%
%s^{\prime}_{q^{\prime}}))$ of the Ehrenfeucht-Fra\"\i ss\'e game ${\mathcal G}_{q}({\mathcal F}_{A,B}^{\rho},{\mathcal F}_{A^{\prime},B^{\prime}}^{\rho^{\prime}})$, for all $s{\in}A{\cup}B$ and for all $s^{\prime}{\in}A^{\prime}{\cup}B^{\prime}$, if $(s,s^{\prime})$ is appropriate for $((s_{1},\ldots,s_{q^{\prime}}),(s^{\prime}_{1},\ldots,s^{\prime}_{q^{\prime}}))$ then $((s_{1},\ldots,s_{q^{\prime}},s),(s^{\prime}_{1},\ldots,s^{\prime}_{q^{\prime}},s^{\prime}))$ is a winning position.
%
%
%\end{lemma}
%
%
%\begin{lemma}\label{lemma:about:extending:winning:positions:b}
%
%
%For all $q^{\prime}{\in}\lsem0,q{-}1\rsem$ and for all positions $((s_{1},\ldots,s_{q^{\prime}}),(s^{\prime}_{1},\ldots,s^{\prime}_{q^{\prime}}))$ of the Ehrenfeucht-Fra\"\i ss\'e game ${\mathcal G}_{q}({\mathcal F}_{A,B}^{\rho},{\mathcal F}_{A^{\prime},B^{\prime}}^{\rho^{\prime}})$, if $((s_{1},\ldots,s_{q^{\prime}}),(s^{\prime}_{1},\ldots,s^{\prime}_{q^{\prime}}))$ is good then
%
%
%\begin{itemize}
%
%
%\item for all $s{\in}A{\cup}B$, there exists $s^{\prime}{\in}A^{\prime}{\cup}B^{\prime}$ such that $((s_{1},\ldots,s_{q^{\prime}},s),(s^{\prime}_{1},\ldots,
%
%
%
%
%
%
%
%
%
%
%
%
%$\linebreak$
%
%
%
%
%
%
%
%
%
%
%
%
%s^{\prime}_{q^{\prime}},s^{\prime}))$ is good,
%
%
%\item for all $s^{\prime}{\in}A^{\prime}{\cup}B^{\prime}$, there exists $s{\in}A{\cup}B$ such that $((s_{1},\ldots,s_{q^{\prime}},s),(s^{\prime}_{1},\ldots,
%
%
%
%
%
%
%
%
%
%
%
%
%$\linebreak$
%
%
%
%
%
%
%
%
%
%
%
%
%s^{\prime}_{q^{\prime}},s^{\prime}))$ is good.
%
%
%\end{itemize}
%
%
%\end{lemma}
%
%
\begin{lemma}\label{lemma:the:alpha:reduction:is:always:dust:finite}
%
%
\begin{enumerate}
%
%
\item $A^{\prime}{\subseteq}A$ and if ${\parallel}A^{\prime}{\parallel}{\leq}2$ then $A^{\prime}{=}A$,
%
%
\item $B^{\prime}{=}B$,
%
%
\item for all $s^{\prime}{\in}A^{\prime}$, $\rho^{\prime}(s^{\prime}){=}\rho(s^{\prime})$,
%
%
\item ${\parallel}{\rho^{\prime}}^{{-}1}(\emptyset){\parallel}{\leq}q$ and ${\parallel}{\rho^{\prime}}^{{-}1}(B^{\prime}){\parallel}{\leq}q$,
%
%
\item for all $t^{\prime}{\in}B^{\prime}$, ${\parallel}{\rho^{\prime}}^{{-}1}(\{t^{\prime}\}){\parallel}{\leq}q$ and for all $t^{\prime}{\in}B^{\prime}$, ${\parallel}{\rho^{\prime}}^{{-}1}(B^{\prime}{\setminus}\{t^{\prime}\}){\parallel}{\leq}q$,
%
%
\item if ${\mathcal F}_{A,B}^{\rho}$ is simple then ${\mathcal F}_{A^{\prime},B^{\prime}}^{\rho^{\prime}}$ is simple, whereas if ${\mathcal F}_{A,B}^{\rho}$ is non-simple and $B$ is finite then ${\mathcal F}_{A^{\prime},B^{\prime}}^{\rho^{\prime}}$ is non-simple, $B^{\prime}$ is finite and ${\parallel}A^{\prime}{\parallel}{\leq}q{\times}2^{{\parallel}B{\parallel}}$,
%
%
\item ${\mathcal F}_{A^{\prime},B^{\prime}}^{\rho^{\prime}}$ is a bounded morphic image of ${\mathcal F}_{A,B}^{\rho}$,
%
%
\item the second player has a winning strategy in the Ehrenfeucht-Fra\"\i ss\'e game ${\mathcal G}_{q}({\mathcal F}_{A,B}^{\rho},{\mathcal F}_{A^{\prime},B^{\prime}}^{\rho^{\prime}})$.
%
%
\end{enumerate}
%
%
\end{lemma}
%
%
\begin{proof}
%
%
$\mathbf{(1)}$--$\mathbf{(6)}$~By construction of ${\mathcal F}_{A^{\prime},B^{\prime}}^{\rho^{\prime}}$.
%
%
\\
\\
%
%
$\mathbf{(7)}$~Let $g:\ A{\longrightarrow}A^{\prime}$ be a function such that for all $s{\in}A$, either $s{\in}A^{\prime}$ and $g(s){=}s$, or $s{\not\in}A^{\prime}$ and $g(s){\sim}s$.
Notice that the existence of such a function is a consequence of the definition of $A^{\prime}$.
Notice also that $g$ is surjective.
Moreover, for all $s{\in}A$, $g(s){\sim}s$.
Let $f:\ A{\cup}B{\longrightarrow}A^{\prime}{\cup}B^{\prime}$ be the function such that for all $s{\in}A$, $f(s){=}g(s)$ and for all $s{\in}B$, $f(s){=}s$.
Notice that $f$ is a bounded morphism from ${\mathcal F}_{A,B}^{\rho}$ to ${\mathcal F}_{A^{\prime},B^{\prime}}^{\rho^{\prime}}$.
Notice also that $f$ is surjective.
%
%
\\
\\
%
%
$\mathbf{(8)}$~Let $q^{\prime}{\in}\lsem1,q\rsem$ and $((s_{1},\ldots,s_{q^{\prime}{-}1}),(s^{\prime}_{1},\ldots,s^{\prime}_{q^{\prime}{-}1}))$ be a winning position of the Ehrenfeucht-Fra\"\i ss\'e game ${\mathcal G}_{q}({\mathcal F}_{A,B}^{\rho},{\mathcal F}_{A^{\prime},B^{\prime}}^{\rho^{\prime}})$.
Suppose the first player picked a member $s$ of ${\mathcal F}_{A,B}^{\rho}$ in its $q^{\prime}$th move.
Then
%
%
\begin{itemize}
%
%
\item if $s$ is in $A$ then
%
%
\begin{itemize}
%
%
\item if $s{=}s_{q^{\prime\prime}}$ for some $q^{\prime\prime}{\in}\lsem1,q^{\prime}{-}1\rsem$ then the second player picks $s^{\prime}_{q^{\prime\prime}}$ in $A^{\prime}$ in its $q^{\prime}$th move,
%
%
\item otherwise, the second player picks an element in $\lbrack s\rbrack_{\sim}{\cap}(A^{\prime}{\setminus}\{s^{\prime}_{1},\ldots,
%
%
%
%
%
%
%
%
%
%
%
%
$\linebreak$
%
%
%
%
%
%
%
%
%
%
%
%
s^{\prime}_{q^{\prime}{-}1}\})$ in its $q^{\prime}$th move,
%
%
\end{itemize}
%
%
\item if $s$ is in $B$ then the second player picks $s$ in $B^{\prime}$ in its $q^{\prime}$th move.
%
%
\end{itemize}
%
%
Suppose the first player picked a member $s^{\prime}$ of ${\mathcal F}_{A^{\prime},B^{\prime}}^{\rho^{\prime}}$ in its $q^{\prime}$th move.
Then
%
%
\begin{itemize}
%
%
\item if $s^{\prime}$ is in $A^{\prime}$ then
%
%
\begin{itemize}
%
%
\item if $s^{\prime}{=}s^{\prime}_{q^{\prime\prime}}$ for some $q^{\prime\prime}{\in}\lsem1,q^{\prime}{-}1\rsem$ then the second player picks $s_{q^{\prime\prime}}$ in $A$ in its $q^{\prime}$th move,
%
%
\item otherwise, the second player picks an element in $\lbrack s^{\prime}\rbrack_{\sim}{\cap}(A{\setminus}\{s_{1},\ldots,
%
%
%
%
%
%
%
%
%
%
%
%
$\linebreak$
%
%
%
%
%
%
%
%
%
%
%
%
s_{q^{\prime}{-}1}\})$ in its $q^{\prime}$th move,
%
%
\end{itemize}
%
%
\item if $s^{\prime}$ is in $B^{\prime}$ then the second player picks $s^{\prime}$ in $B$ in its $q^{\prime}$th move.
%
%
\end{itemize}
%
%
Obviously, $((s_{1},\ldots,s_{q^{\prime}{-}1},s),(s^{\prime}_{1},\ldots,s^{\prime}_{q^{\prime}{-}1},s^{\prime}))$ is a winning position of the Eh\-renfeucht-Fra\"\i ss\'e game ${\mathcal G}_{q}({\mathcal F}_{A,B}^{\rho},{\mathcal F}_{A^{\prime},B^{\prime}}^{\rho^{\prime}})$.
%
%
\medskip
%
%
\end{proof}
%
%
%Let $q{\geq}3$.
Let $\{{\mathcal F}_{A_{i},B_{i}}^{\rho_{i}}:\ i{\in}I\}$ be a disjoint indexed family of galaxies.
%We shall say that $\{(W_{A_{i}^{\prime},B_{i}^{\prime}}^{\rho_{i}^{\prime}},R_{A_{i}^{\prime},B_{i}^{\prime}}^{\rho_{i}^{\prime}}):\ i{\in}I\}$ is {\it $q$-$\alpha$-reduced from $\{(W_{A_{i},B_{i}}^{\rho_{i}},R_{A_{i},B_{i}}^{\rho_{i}}):\ i{\in}I\}$}\/ if for all $i{\in}I$, either $(W_{A_{i}^{\prime},B_{i}^{\prime}}^{\rho_{i}^{\prime}},R_{A_{i}^{\prime},B_{i}^{\prime}}^{\rho_{i}^{\prime}})$ is equal to $(W_{A_{i},B_{i}}^{\rho_{i}},R_{A_{i},B_{i}}^{\rho_{i}})$, or $(W_{A_{i}^{\prime},B_{i}^{\prime}}^{\rho_{i}^{\prime}},R_{A_{i}^{\prime},B_{i}^{\prime}}^{\rho_{i}^{\prime}})$ is a $q$-$\alpha$-reduction of $(W_{A_{i},B_{i}}^{\rho_{i}},R_{A_{i},B_{i}}^{\rho_{i}})$.
$\{{\mathcal F}_{A_{i}^{\prime},B_{i}^{\prime}}^{\rho_{i}^{\prime}}:\ i{\in}I\}$ is a {\it $q$-$\alpha$-reduction of $\{{\mathcal F}_{A_{i},B_{i}}^{\rho_{i}}:\ i{\in}I\}$}\/ if for all $i{\in}I$, ${\mathcal F}_{A_{i}^{\prime},B_{i}^{\prime}}^{\rho_{i}^{\prime}}$ is a $q$-$\alpha$-reduction of ${\mathcal F}_{A_{i},B_{i}}^{\rho_{i}}$.
%
%
\begin{proposition}\label{lemma:good:properties:of:alpha:reductions:indexed}
%
%
%Let $q{\geq}3$.
%Let $\{{\mathcal F}_{A_{i},B_{i}}^{\rho_{i}}:\ i{\in}I\}$ be a disjoint indexed family of galaxies.
Let $I^{ns}$ be the set of all $i{\in}I$ such that ${\mathcal F}_{A_{i},B_{i}}^{\rho_{i}}$ is non-simple.
%We shall say that $\{(W_{A_{i}^{\prime},B_{i}^{\prime}}^{\rho_{i}^{\prime}},R_{A_{i}^{\prime},B_{i}^{\prime}}^{\rho_{i}^{\prime}}):\ i{\in}I\}$ is {\it $q$-$\alpha$-reduced from $\{(W_{A_{i},B_{i}}^{\rho_{i}},R_{A_{i},B_{i}}^{\rho_{i}}):\ i{\in}I\}$}\/ if for all $i{\in}I$, either $(W_{A_{i}^{\prime},B_{i}^{\prime}}^{\rho_{i}^{\prime}},R_{A_{i}^{\prime},B_{i}^{\prime}}^{\rho_{i}^{\prime}})$ is equal to $(W_{A_{i},B_{i}}^{\rho_{i}},R_{A_{i},B_{i}}^{\rho_{i}})$, or $(W_{A_{i}^{\prime},B_{i}^{\prime}}^{\rho_{i}^{\prime}},R_{A_{i}^{\prime},B_{i}^{\prime}}^{\rho_{i}^{\prime}})$ is a $q$-$\alpha$-reduction of $(W_{A_{i},B_{i}}^{\rho_{i}},R_{A_{i},B_{i}}^{\rho_{i}})$.
If $\{{\mathcal F}_{A_{i}^{\prime},B_{i}^{\prime}}^{\rho_{i}^{\prime}}:\ i{\in}I\}$ is a $q$-$\alpha$-reduction of $\{{\mathcal F}_{A_{i},B_{i}}^{\rho_{i}}:\ i{\in}I\}$, $(W,R)$ is the union of $\{{\mathcal F}_{A_{i},B_{i}}^{\rho_{i}}:\ i{\in}I\}$ and $(W^{\prime},R^{\prime})$ is the union of $\{{\mathcal F}_{A_{i}^{\prime},B_{i}^{\prime}}^{\rho_{i}^{\prime}}:\ i{\in}I\}$ then
%
%
\begin{enumerate}
%
%
\item $\{{\mathcal F}_{A_{i}^{\prime},B_{i}^{\prime}}^{\rho_{i}^{\prime}}:\ i{\in}I\}$ is dust-bounded,
%
%
\item if $\{{\mathcal F}_{A_{i},B_{i}}^{\rho_{i}}:\ i{\in}I^{ns}\}$ is kernel-bounded then $\{{\mathcal F}_{A_{i}^{\prime},B_{i}^{\prime}}^{\rho_{i}^{\prime}}:\ i{\in}I^{ns}\}$ is root-bounded and kernel-bounded,
%
%
\item $(W^{\prime},R^{\prime})$ is a bounded morphic image of $(W,R)$,
%
%
\item the second player has a winning strategy in the Ehrenfeucht-Fra\"\i ss\'e game ${\mathcal G}_{q}((W,R),(W^{\prime},R^{\prime}))$.
%
%
\end{enumerate}
%
%
\end{proposition}
%
%
\begin{proof}
%
%
By Lemma~\ref{lemma:the:alpha:reduction:is:always:dust:finite}.
%
%
\medskip
%
%
\end{proof}
%
%
%
%
%\paragraph{$\beta$-reductions}
%
%
%Let $(W_{A,B}^{\rho},R_{A,B}^{\rho})$ be a simple galaxy.
%Let $q{\in}\N$.
%Let $(W_{A^{\prime},B^{\prime}}^{\rho^{\prime}},
%
%
%
%
%
%
%
%$\linebreak$
%
%
%
%
%
%
%
%R_{A^{\prime},B^{\prime}}^{\rho^{\prime}})$ be a simple galaxy obtained from $(W_{A,B}^{\rho},R_{A,B}^{\rho})$ by keeping for all $s{\in}B$, at most $\max\{q,1\}$ elements in $\rho^{{-}1}(s)$ and by deleting all non-chosen $\rho^{{-}1}(s)$'s elements.
%We shall say that $(W_{A^{\prime},B^{\prime}}^{\rho^{\prime}},R_{A^{\prime},B^{\prime}}^{\rho^{\prime}})$ is a {\it $q$-$\beta$-reduction of $(W_{A,B}^{\rho},R_{A,B}^{\rho})$.}
%
%
%\begin{lemma}\label{lemma:beta:reduction:quasi:root:bounded:dust:OK}
%
%
%Let $(W_{A,B}^{\rho},R_{A,B}^{\rho})$ be a simple galaxy, $q{\in}\N$ and $(W_{A^{\prime},B^{\prime}}^{\rho^{\prime}},R_{A^{\prime},B^{\prime}}^{\rho^{\prime}})$ be a $q$-$\beta$-reduction of $(W_{A,B}^{\rho},R_{A,B}^{\rho})$.
%
%
%Then,
%
%
%\begin{enumerate}
%
%
%\item $\{s^{\prime}{\in}A^{\prime}:\ \rho^{\prime}(s^{\prime}){=}\emptyset\}{=}\{s{\in}A:\ \rho(s){=}\emptyset\}$,
%
%
%\item $\{s^{\prime}{\in}A^{\prime}:\ \rho^{\prime}(s^{\prime}){=}B^{\prime}\}{\subseteq}\{s{\in}A:\ \rho(s){=}B\}$,
%
%
%\item for all $s^{\prime}{\in}B^{\prime}$, ${\parallel}{\rho^{\prime}}^{{-}1}(s^{\prime}){\parallel}{\leq}\max\{q,1\}$,
%
%
%\item $B^{\prime}{=}B$.
%
%
%\end{enumerate}
%
%
%\end{lemma}
%
%
%Let $\{(W_{A_{i},B_{i}}^{\rho_{i}},R_{A_{i},B_{i}}^{\rho_{i}}):\ i{\in}I\}$ and $\{(W_{A_{i}^{\prime},B_{i}^{\prime}}^{\rho_{i}^{\prime}},R_{A_{i}^{\prime},B_{i}^{\prime}}^{\rho_{i}^{\prime}}):\ i{\in}I\}$ be disjoint families of simple galaxies indexed by the same set $I$ and $q{\in}\N$.
%We shall say that $\{(W_{A_{i}^{\prime},B_{i}^{\prime}}^{\rho_{i}^{\prime}},R_{A_{i}^{\prime},B_{i}^{\prime}}^{\rho_{i}^{\prime}}):\ i{\in}I\}$ is {\it $q$-$\beta$-reduced from $\{(W_{A_{i},B_{i}}^{\rho_{i}},R_{A_{i},B_{i}}^{\rho_{i}}):\ i{\in}I\}$}\/ if for all $i{\in}I$, either $(W_{A_{i}^{\prime},B_{i}^{\prime}}^{\rho_{i}^{\prime}},R_{A_{i}^{\prime},B_{i}^{\prime}}^{\rho_{i}^{\prime}})$ is equal to $(W_{A_{i},B_{i}}^{\rho_{i}},R_{A_{i},B_{i}}^{\rho_{i}})$, or $(W_{A_{i}^{\prime},B_{i}^{\prime}}^{\rho_{i}^{\prime}},R_{A_{i}^{\prime},B_{i}^{\prime}}^{\rho_{i}^{\prime}})$ is a $q$-$\beta$-reduction of $(W_{A_{i},B_{i}}^{\rho_{i}},R_{A_{i},B_{i}}^{\rho_{i}})$.
%
%
%
%
\paragraph{$\gamma$-reductions}
%
%
%In this section, we introduce a second type of reduction: the $\gamma$-reduction of a galaxy.
The main effect of a $\gamma$-reduction on a simple galaxy is to reduce both the cardinality of its upper part and the cardinality of its lower part.
%
%
\\
\\
%
%
Let ${\mathcal F}_{A,B}^{\rho}$ be a simple galaxy satisfying the conditions $(4)$ and $(5)$ of Lemma~\ref{lemma:the:alpha:reduction:is:always:dust:finite}.
%
%
\\
\\
%
%
%A {\it component of ${\mathcal F}_{A,B}^{\rho}$}\/ is a triple $(u,X,Y)$ such that $u{\in}B$, $X{=}\rho^{{-}1}(\{u\})$ and $Y{=}\rho^{{-}1}(B{\setminus}\{u\})$.
%Let $\chi_{A,B}^{\rho}$ be the set of all components of ${\mathcal F}_{A,B}^{\rho}$.
%Let $\approx$ be the equivalence relation on $\chi_{A,B}^{\rho}$ defined by $(u,X,Y){\approx}(v,Z,T)$ if and only if ${\parallel}X{\parallel}{=}{\parallel}Z{\parallel}$ and ${\parallel}Y{\parallel}{=}{\parallel}T{\parallel}$.
%By extension, for all $u,v{\in}$, we shall write $u{\approx}v$ when $(u,\rho^{{-}1}(\{u\}),\rho^{{-}1}(B{\setminus}\{u\})){\approx}(v,\rho^{{-}1}(\{v\}),\rho^{{-}1}(B{\setminus}\{v\}))$.
%Let $h:\ \chi_{A,B}^{\rho}{/}{\approx}
%Let $\approx$ be the equivalence relation on $B$ defined by $s{\approx}t$ if and only if ${\parallel}\rho^{{-}1}(\{s\}){\parallel}{=}{\parallel}\rho^{{-}1}(\{t\}){\parallel}$.
Let $\approx$ be the equivalence relation on $B$ defined by $u{\approx}v$ if and only if ${\parallel}\rho^{{-}1}(\{u\}){\parallel}
%
%
%
%
%
%
%
%
%
%
%
%
$\linebreak$
%
%
%
%
%
%
%
%
%
%
%
%
{=}{\parallel}\rho^{{-}1}(\{v\}){\parallel}$ and ${\parallel}\rho^{{-}1}(B{\setminus}\{u\}){\parallel}{=}{\parallel}\rho^{{-}1}(B{\setminus}\{v\}){\parallel}$.
For all $u{\in}B$, let $\lbrack u\rbrack_{\approx}$ be the equivalence class of $u$ modulo $\approx$.
%
%
\\
\\
%
%
For all $u{\in}B$, if ${\parallel}\lbrack u\rbrack_{\approx}{\parallel}{\leq}q$ then let $D_{u}$ be $\lbrack u\rbrack_{\approx}$ else let $D_{u}$ be an arbitrary subset of $\lbrack u\rbrack_{\approx}$ of cardinality $q$.
Let $B^{\prime}$ be $\bigcup\{D_{u}:\ u{\in}B\}$.
Obviously, from each equivalence class modulo $\approx$, $B^{\prime}$ contains at most $q$ elements.
Therefore, $B^{\prime}$ contains at most $q{\times}(q{+}1)^{2}$ elements.
%
%
\\
\\
%
%
%Let $h:\ B{\longrightarrow}B^{\prime}$ be a function such that for all $t{\in}B$, either $t{\in}B^{\prime}$ and $h(t){=}t$, or $t{\not\in}B^{\prime}$ and $h(t){\approx}t$.
%Notice that the existence of such a function is a consequence of the definition of $B^{\prime}$.
%Notice also that $h$ is surjective.
%Moreover, for all $t{\in}B$, $h(t){\approx}t$.
%
%
Let $A^{\prime}$ be $\rho^{{-}1}(\emptyset){\cup}\rho^{{-}1}(B){\cup}\bigcup\{\rho^{{-}1}(\{u\}):\ u{\in}B^{\prime}\}{\cup}\bigcup\{\rho^{{-}1}(B{\setminus}\{u\}):\ u{\in}B^{\prime}\}$.
%Notice that for all $s{\in}A$, if $s{\not\in}A^{\prime}$ then either there exists $t{\in}B$ such that $\rho(s){=}\{t\}$, or there exists $t{\in}B$ such that $\rho(s){=}B{\setminus}\{t\}$.
Since $B^{\prime}$ contains at most $q{\times}(q{+}1)^{2}$ elements, then $A^{\prime}$ contains at most $2{\times}q{\times}(q
%
%
%
%
%
%
%
%
%
%
%
%
$\linebreak$
%
%
%
%
%
%
%
%
%
%
%
%
{\times}(q{+}1)^{2}+1)$ elements.
%
%
%\\
%\\
%
%
%For all $t,u{\in}B$, if $t{\approx}u$ then let $a_{t}^{u}:\ \rho^{{-}1}(\{t\}){\longrightarrow}\rho^{{-}1}(\{u\})$ and $b_{t}^{u}:\ \rho^{{-}1}(B{\setminus}\{t\}){\longrightarrow}\rho^{{-}1}(B{\setminus}\{u\})$ be bijective functions.
%For all $t{\in}B$, let $a_{t}:\ \rho^{{-}1}(\{t\}){\longrightarrow}\rho^{{-}1}(\{h(t)\})$ and $b_{t}:\ \rho^{{-}1}(B{\setminus}\{t\}){\longrightarrow}\rho^{{-}1}(B{\setminus}\{h(t)\})$ be bijective functions.
%such that if $t{\in}B^{\prime}$ then $a_{t}{=}\id_{\rho^{{-}1}(\{t\})}$ and $b_{t}{=}\id_{\rho^{{-}1}(B{\setminus}\{t\})}$.
%Notice that the existence of such bijective functions is a consequence of the definition of $\approx$.
%Let $g:\ A{\longrightarrow}A^{\prime}$ be the function such that for all $s{\in}A$, either $\rho(s){=}\emptyset$ and $g(s){=}s$, or $\rho(s){=}B$ and $g(s){=}s$, or there exists $t{\in}B$ such that $\rho(s){=}\{t\}$ and $g(s){=}a_{t}(s)$, or there exists $t{\in}B$ such that $\rho(s){=}B{\setminus}\{t\}$ and $g(s){=}b_{t}(s)$.
%Notice that $g$ is surjective.
%Let ${\mathcal F}_{A^{\prime},B^{\prime}}^{\rho^{\prime}}$ be the galaxy such that for all $u^{\prime}{\in}A^{\prime}$, $\rho^{\prime}(u^{\prime}){=}\rho(u^{\prime}){\cap}B^{\prime}$.
%We shall say that ${\mathcal F}_{A^{\prime},B^{\prime}}^{\rho^{\prime}}$ is a {\it $q$-$\gamma$-reduction of ${\mathcal F}_{A,B}^{\rho}$.}
%
%
\\
\\
%
%
Let ${\mathcal F}_{A^{\prime},B^{\prime}}^{\rho^{\prime}}$ be the galaxy such that for all $u^{\prime}{\in}A^{\prime}$, $\rho^{\prime}(u^{\prime}){=}\rho(u^{\prime}){\cap}B^{\prime}$.
${\mathcal F}_{A^{\prime},B^{\prime}}^{\rho^{\prime}}$ is a {\it $q$-$\gamma$-reduction of ${\mathcal F}_{A,B}^{\rho}$.}
%
%
%\\
%\\
%
%
%Let $f:\ A{\cup}B{\longrightarrow}A^{\prime}{\cup}B^{\prime}$ be the function such that for all $s{\in}A$, $f(s){=}g(s)$ and for all $s{\in}B$, $f(s){=}h(s)$.
%Notice that $f$ is surjective.
%Let $\chi:\ B{/}{\approx}{\longrightarrow}B$ be such that for all $u{\in}B$, $\chi(\lbrack u\rbrack){\approx}u$.
%
%
%\\
%\\
%
%
%Let us define a strategy $\pi$ of the second player in the Ehrenfeucht-Fra\"\i ss\'e game ${\mathcal G}_{q}({\mathcal F}_{A,B}^{\rho},{\mathcal F}_{A^{\prime},B^{\prime}}^{\rho^{\prime}})$.
%For all $q^{\prime}{\in}\lsem0,q{-}1\rsem$, we shall say that a member $s$ of $A$ is {\it left-fresh in a position $((s_{1},\ldots,s_{q^{\prime}}),(s^{\prime}_{1},\ldots,s^{\prime}_{q^{\prime}}))$ of the Ehrenfeucht-Fra\"\i ss\'e game ${\mathcal G}_{q}({\mathcal F}_{A,B}^{\rho},{\mathcal F}_{A^{\prime},B^{\prime}}^{\rho^{\prime}})$}\/ if for all $q^{\prime\prime}{\in}\lsem1,q^{\prime}\rsem$, $s{\not=}s_{q^{\prime\prime}}$.
%For all $q^{\prime}{\in}\lsem0,q{-}1\rsem$, we shall say that a member $s^{\prime}$ of $A^{\prime}$ is {\it right-fresh in a position $((s_{1},\ldots,s_{q^{\prime}}),(s^{\prime}_{1},\ldots,s^{\prime}_{q^{\prime}}))$ of the Ehrenfeucht-Fra\"\i ss\'e game ${\mathcal G}_{q}({\mathcal F}_{A,B}^{\rho},{\mathcal F}_{A^{\prime},B^{\prime}}^{\rho^{\prime}})$}\/ if for all $q^{\prime\prime}{\in}\lsem1,q^{\prime}\rsem$, $s^{\prime}{\not=}s^{\prime}_{q^{\prime\prime}}$.
%For all $q^{\prime}{\in}\lsem0,q{-}1\rsem$, we shall say that a member $s$ of $B$ is {\it left-fresh in a position $((s_{1},\ldots,s_{q^{\prime}}),(s^{\prime}_{1},\ldots,s^{\prime}_{q^{\prime}}))$ of the Ehrenfeucht-Fra\"\i ss\'e game ${\mathcal G}_{q}({\mathcal F}_{A,B}^{\rho},{\mathcal F}_{A^{\prime},B^{\prime}}^{\rho^{\prime}})$}\/ if for all $q^{\prime\prime}{\in}\lsem1,q^{\prime}\rsem$, $s{\not=}s_{q^{\prime\prime}}$.
%For all $q^{\prime}{\in}\lsem0,q{-}1\rsem$, we shall say that a member $s^{\prime}$ of $B^{\prime}$ is {\it right-fresh in a position $((s_{1},\ldots,s_{q^{\prime}}),(s^{\prime}_{1},\ldots,s^{\prime}_{q^{\prime}}))$ of the Ehrenfeucht-Fra\"\i ss\'e game ${\mathcal G}_{q}({\mathcal F}_{A,B}^{\rho},{\mathcal F}_{A^{\prime},B^{\prime}}^{\rho^{\prime}})$}\/ if for all $q^{\prime\prime}{\in}\lsem1,q^{\prime}\rsem$, $s^{\prime}{\not=}s^{\prime}_{q^{\prime\prime}}$.
%For all $q^{\prime}{\in}\lsem0,q{-}1\rsem$, we shall say that a member $s$ of $A$ is {\it left-fresh in a position $((s_{1},\ldots,s_{q^{\prime}}),(s^{\prime}_{1},\ldots,s^{\prime}_{q^{\prime}}))$ of the Ehrenfeucht-Fra\"\i ss\'e game ${\mathcal G}_{q}({\mathcal F}_{A,B}^{\rho},{\mathcal F}_{A^{\prime},B^{\prime}}^{\rho^{\prime}})$}\/ if zzzzz
%For all $q^{\prime}{\in}\lsem0,q{-}1\rsem$, we shall say that a member $s$ of $B$ is {\it left-fresh in a position $((s_{1},\ldots,s_{q^{\prime}}),(s^{\prime}_{1},\ldots,s^{\prime}_{q^{\prime}}))$ of the Ehrenfeucht-Fra\"\i ss\'e game ${\mathcal G}_{q}({\mathcal F}_{A,B}^{\rho},{\mathcal F}_{A^{\prime},B^{\prime}}^{\rho^{\prime}})$}\/ if zzzzz
%For all $q^{\prime}{\in}\lsem0,q{-}1\rsem$, we shall say that a member $s^{\prime}$ of $A^{\prime}$ is {\it right-fresh in a position $((s_{1},\ldots,s_{q^{\prime}}),(s^{\prime}_{1},\ldots,s^{\prime}_{q^{\prime}}))$ of the Ehrenfeucht-Fra\"\i ss\'e game ${\mathcal G}_{q}({\mathcal F}_{A,B}^{\rho},{\mathcal F}_{A^{\prime},B^{\prime}}^{\rho^{\prime}})$}\/ if zzzzz
%For all $q^{\prime}{\in}\lsem0,q{-}1\rsem$, we shall say that a member $s^{\prime}$ of $B^{\prime}$ is {\it right-fresh in a position $((s_{1},\ldots,s_{q^{\prime}}),(s^{\prime}_{1},\ldots,s^{\prime}_{q^{\prime}}))$ of the Ehrenfeucht-Fra\"\i ss\'e game ${\mathcal G}_{q}({\mathcal F}_{A,B}^{\rho},{\mathcal F}_{A^{\prime},B^{\prime}}^{\rho^{\prime}})$}\/ if zzzzz
%For all $q^{\prime}{\in}\lsem0,q\rsem$, we shall say that a position $((s_{1},\ldots,s_{q^{\prime}}),(s^{\prime}_{1},\ldots,s^{\prime}_{q^{\prime}}))$ of the Ehrenfeucht-Fra\"\i ss\'e game ${\mathcal G}_{q}({\mathcal F}_{A,B}^{\rho},{\mathcal F}_{A^{\prime},B^{\prime}}^{\rho^{\prime}})$ is {\it good}\/ if for all $q^{\prime\prime}{\in}\lsem1,q^{\prime}\rsem$, at least one of the following conditions holds:
%for all $s{\in}A{\cup}B$ and for all $s^{\prime}{\in}A^{\prime}{\cup}B^{\prime}$, we shall say that the pair $(s,s^{\prime})$ is {\it appropriate for a position $((s_{1},\ldots,s_{q^{\prime}}),(s^{\prime}_{1},\ldots,s^{\prime}_{q^{\prime}}))$ of the Ehrenfeucht-Fra\"\i ss\'e game ${\mathcal G}_{q}({\mathcal F}_{A,B}^{\rho},{\mathcal F}_{A^{\prime},B^{\prime}}^{\rho^{\prime}})$}\/ if at least one of the following conditions holds:
%
%
%\begin{itemize}
%
%
%\item $s_{q^{\prime\prime}}{\in}A$, $s^{\prime}_{q^{\prime\prime}}{\in}A^{\prime}$, $s_{q^{\prime\prime}}$ is left-fresh in $((s_{1},\ldots,s_{q^{\prime\prime}{-}1}),(s^{\prime}_{1},\ldots,s^{\prime}_{q^{\prime\prime}{-}1}))$, $s^{\prime}_{q^{\prime\prime}}$ is right-fresh in $((s_{1},\ldots,s_{q^{\prime\prime}{-}1}),(s^{\prime}_{1},\ldots,s^{\prime}_{q^{\prime\prime}{-}1}))$ and $s_{q^{\prime\prime}}{\sim}s^{\prime}_{q^{\prime\prime}}$,
%
%
%\item $s_{q^{\prime\prime}}{\in}A$, $s^{\prime}_{q^{\prime\prime}}{\in}A^{\prime}$, $s_{q^{\prime\prime}}$ is not left-fresh in $((s_{1},\ldots,s_{q^{\prime\prime}{-}1}),(s^{\prime}_{1},\ldots,s^{\prime}_{q^{\prime\prime}{-}1}))$, $s^{\prime}_{q^{\prime\prime}}$ is not right-fresh in $((s_{1},\ldots,s_{q^{\prime\prime}{-}1}),(s^{\prime}_{1},\ldots,s^{\prime}_{q^{\prime\prime}{-}1}))$ and there exists $q^{\prime\prime\prime}{\in}\lsem1,q^{\prime\prime}{-}1\rsem$ such that $s_{q^{\prime\prime\prime}}{\in}A$, $s_{q^{\prime\prime\prime}}{=}s_{q^{\prime\prime}}$, $s^{\prime}_{q^{\prime\prime\prime}}{\in}A^{\prime}$ and $s^{\prime}_{q^{\prime\prime\prime}}{=}s^{\prime}_{q^{\prime\prime}}$,
%
%
%\item $s_{q^{\prime\prime}}{\in}B$, $s^{\prime}_{q^{\prime\prime}}{\in}B^{\prime}$, $s_{q^{\prime\prime}}$ is left-fresh in $((s_{1},\ldots,s_{q^{\prime\prime}{-}1}),(s^{\prime}_{1},\ldots,s^{\prime}_{q^{\prime\prime}{-}1}))$, $s^{\prime}_{q^{\prime\prime}}$ is right-fresh in $((s_{1},\ldots,s_{q^{\prime\prime}{-}1}),(s^{\prime}_{1},\ldots,s^{\prime}_{q^{\prime\prime}{-}1}))$ and $s_{q^{\prime\prime}}{\approx}s^{\prime}_{q^{\prime\prime}}$,
%
%
%\item $s_{q^{\prime\prime}}{\in}B$, $s^{\prime}_{q^{\prime\prime}}{\in}B^{\prime}$, $s_{q^{\prime\prime}}$ is not left-fresh in $((s_{1},\ldots,s_{q^{\prime\prime}{-}1}),(s^{\prime}_{1},\ldots,s^{\prime}_{q^{\prime\prime}{-}1}))$, $s^{\prime}_{q^{\prime\prime}}$ is not right-fresh in $((s_{1},\ldots,s_{q^{\prime\prime}{-}1}),(s^{\prime}_{1},\ldots,s^{\prime}_{q^{\prime\prime}{-}1}))$ and there exists $q^{\prime\prime\prime}{\in}\lsem1,q^{\prime\prime}{-}1\rsem$ such that $s_{q^{\prime\prime\prime}}{\in}B$, $s_{q^{\prime\prime\prime}}{=}s_{q^{\prime\prime}}$, $s^{\prime}_{q^{\prime\prime\prime}}{\in}B^{\prime}$ and $s^{\prime}_{q^{\prime\prime\prime}}{=}s^{\prime}_{q^{\prime\prime}}$.
%
%
%\item $s_{q^{\prime\prime}}{\in}B$, $s^{\prime}_{q^{\prime\prime}}{\in}B^{\prime}$ and $s_{q^{\prime\prime}}{=}s^{\prime}_{q^{\prime\prime}}$.
%
%
%\end{itemize}
%the following conditions hold:
%
%
%\begin{itemize}
%
%
%\item $s_{q^{\prime\prime}}{\in}A$ if and only if $s^{\prime}_{q^{\prime\prime}}{\in}A^{\prime}$,
%
%
%\item $s_{q^{\prime\prime}}{\in}B$ if and only if $s^{\prime}_{q^{\prime\prime}}{\in}B^{\prime}$,
%
%
%\item if $s_{q^{\prime\prime}}{\in}A$ and $s^{\prime}_{q^{\prime\prime}}{\in}A^{\prime}$ then $\rho(s_{q^{\prime\prime}}){=}\emptyset$ if and only if $\rho^{\prime}(s^{\prime}_{q^{\prime\prime}}){=}\emptyset$,
%
%
%\item if $s_{q^{\prime\prime}}{\in}A$ and $s^{\prime}_{q^{\prime\prime}}{\in}A^{\prime}$ then $\rho(s_{q^{\prime\prime}}){=}B$ if and only if $\rho^{\prime}(s^{\prime}_{q^{\prime\prime}}){=}B^{\prime}$,
%
%
%\item if $s_{q^{\prime\prime}}{\in}A$ and $s^{\prime}_{q^{\prime\prime}}{\in}A^{\prime}$ then there exists $t{\in}B$ such that $\rho(s_{q^{\prime\prime}}){=}\{t\}$ if and only if there exists $t^{\prime}{\in}B^{\prime}$ such that $\rho^{\prime}(s^{\prime}_{q^{\prime\prime}}){=}\{t^{\prime}\}$,
%
%
%\item if $s_{q^{\prime\prime}}{\in}A$ and $s^{\prime}_{q^{\prime\prime}}{\in}A^{\prime}$ then there exists $t{\in}B$ such that $\rho(s_{q^{\prime\prime}}){=}B{\setminus}\{t\}$ if and only if there exists $t^{\prime}{\in}B^{\prime}$ such that $\rho^{\prime}(s^{\prime}_{q^{\prime\prime}}){=}B^{\prime}{\setminus}\{t^{\prime}\}$,
%
%
%\item if $s_{q^{\prime\prime}}{\in}B$ and $s^{\prime}_{q^{\prime\prime}}{\in}B^{\prime}$ then $s_{q^{\prime\prime}}{\approx}s^{\prime}_{q^{\prime\prime}}$,
%
%
%\item if $s_{q^{\prime\prime}}{\in}A$, $s^{\prime}_{q^{\prime\prime}}{\in}A^{\prime}$, there exists $t{\in}B$ such that $\rho(s_{q^{\prime\prime}}){=}\{t\}$ and there exists $t^{\prime}{\in}B^{\prime}$ such that $\rho^{\prime}(s^{\prime}_{q^{\prime\prime}}){=}\{t^{\prime}\}$ then $t{\approx}t^{\prime}$,
%
%
%\item if $s_{q^{\prime\prime}}{\in}A$, $s^{\prime}_{q^{\prime\prime}}{\in}A^{\prime}$, there exists $t{\in}B$ such that $\rho(s_{q^{\prime\prime}}){=}B{\setminus}\{t\}$ and there exists $t^{\prime}{\in}B^{\prime}$ such that $\rho^{\prime}(s^{\prime}_{q^{\prime\prime}}){=}B^{\prime}{\setminus}\{t^{\prime}\}$ then $t{\approx}t^{\prime}$,
%
%
%\item for all $q^{\prime\prime\prime}{\in}\lsem1,q^{\prime}\rsem$, $s_{q^{\prime\prime}}{=}s_{q^{\prime\prime\prime}}$ if and only if $s^{\prime}_{q^{\prime\prime}}{=}s^{\prime}_{q^{\prime\prime\prime}}$,
%
%
%\item if $s_{q^{\prime\prime}}{\in}A$ and $s^{\prime}_{q^{\prime\prime}}{\in}A^{\prime}$ then for all $q^{\prime\prime\prime}{\in}\lsem1,q^{\prime}\rsem$, if $s_{q^{\prime\prime\prime}}{\in}B$ and $s^{\prime}_{q^{\prime\prime\prime}}{\in}B^{\prime}$ then $s_{q^{\prime\prime\prime}}{\in}\rho(s_{q^{\prime\prime}})$ if and only if $s^{\prime}_{q^{\prime\prime\prime}}{\in}\rho^{\prime}(s^{\prime}_{q^{\prime\prime}})$.
%
%
%\end{itemize}
%
%
%\begin{lemma}\label{lemma:about:good:positions:and:left:right:fresh:elements:bis}
%
%
%For all $q^{\prime}{\in}\lsem0,q\rsem$ and for all positions $((s_{1},\ldots,s_{q^{\prime}}),(s^{\prime}_{1},\ldots,s^{\prime}_{q^{\prime}}))$ of the Ehrenfeucht-Fra\"\i ss\'e game ${\mathcal G}_{q}({\mathcal F}_{A,B}^{\rho},{\mathcal F}_{A^{\prime},B^{\prime}}^{\rho^{\prime}})$, if $((s_{1},\ldots,s_{q^{\prime}}),(s^{\prime}_{1},\ldots,s^{\prime}_{q^{\prime}}))$ is good then for all $q^{\prime\prime}{\in}\lsem1,q^{\prime}\rsem$, if $s_{q^{\prime\prime}}{\in}A$ and $s^{\prime}_{q^{\prime\prime}}{\in}A^{\prime}$ then $s_{q^{\prime\prime}}{\sim}s^{\prime}_{q^{\prime\prime}}$ and if $s_{q^{\prime\prime}}{\in}B$ and $s^{\prime}_{q^{\prime\prime}}{\in}B^{\prime}$ then $s_{q^{\prime\prime}}{\approx}s^{\prime}_{q^{\prime\prime}}$.
%
%
%\end{lemma}
%
%
%\begin{lemma}\label{lemma:about:good:positions:and:sim:relation:bis}
%
%
%For all $q^{\prime}{\in}\lsem0,q{-}1\rsem$ and for all positions $((s_{1},\ldots,s_{q^{\prime}}),(s^{\prime}_{1},\ldots,s^{\prime}_{q^{\prime}}))$ of the Ehrenfeucht-Fra\"\i ss\'e game ${\mathcal G}_{q}({\mathcal F}_{A,B}^{\rho},{\mathcal F}_{A^{\prime},B^{\prime}}^{\rho^{\prime}})$, if $((s_{1},\ldots,s_{q^{\prime}}),(s^{\prime}_{1},\ldots,s^{\prime}_{q^{\prime}}))$ is good then
%
%
%\begin{itemize}
%
%
%\item for all $s{\in}A$, if $s$ is left-fresh in $((s_{1},\ldots,s_{q^{\prime}}),(s^{\prime}_{1},\ldots,s^{\prime}_{q^{\prime}}))$ then there exists $s^{\prime}{\in}A^{\prime}$ such that $s^{\prime}$ is right-fresh in $((s_{1},\ldots,s_{q^{\prime}}),(s^{\prime}_{1},\ldots,s^{\prime}_{q^{\prime}}))$ and $s{\sim}s^{\prime}$,
%
%
%\item for all $s{\in}B$, if $s$ is left-fresh in $((s_{1},\ldots,s_{q^{\prime}}),(s^{\prime}_{1},\ldots,s^{\prime}_{q^{\prime}}))$ then there exists $s^{\prime}{\in}B^{\prime}$ such that $s^{\prime}$ is right-fresh in $((s_{1},\ldots,s_{q^{\prime}}),(s^{\prime}_{1},\ldots,s^{\prime}_{q^{\prime}}))$ and $s{\approx}s^{\prime}$,
%
%
%\item for all $s^{\prime}{\in}A^{\prime}$, if $s^{\prime}$ is right-fresh in $((s_{1},\ldots,s_{q^{\prime}}),(s^{\prime}_{1},\ldots,s^{\prime}_{q^{\prime}}))$ then there exists $s{\in}A$ such that $s$ is left-fresh in $((s_{1},\ldots,s_{q^{\prime}}),(s^{\prime}_{1},\ldots,s^{\prime}_{q^{\prime}}))$ and $s^{\prime}{\sim}s$,
%
%
%\item for all $s^{\prime}{\in}B^{\prime}$, if $s^{\prime}$ is right-fresh in $((s_{1},\ldots,s_{q^{\prime}}),(s^{\prime}_{1},\ldots,s^{\prime}_{q^{\prime}}))$ then there exists $s{\in}B$ such that $s$ is left-fresh in $((s_{1},\ldots,s_{q^{\prime}}),(s^{\prime}_{1},\ldots,s^{\prime}_{q^{\prime}}))$ and $s^{\prime}{\approx}s$.
%
%
%\end{itemize}
%
%
%\end{lemma}
%
%
%\begin{lemma}\label{lemma:good:positions:are:winning:gamma}
%
%
%For all $q^{\prime}{\in}\lsem0,q\rsem$ and for all positions $((s_{1},\ldots,s_{q^{\prime}}),(s^{\prime}_{1},\ldots,s^{\prime}_{q^{\prime}}))$ of the Ehrenfeucht-Fra\"\i ss\'e game ${\mathcal G}_{q}({\mathcal F}_{A,B}^{\rho},{\mathcal F}_{A^{\prime},B^{\prime}}^{\rho^{\prime}})$, if $((s_{1},\ldots,s_{q^{\prime}}),(s^{\prime}_{1},\ldots,s^{\prime}_{q^{\prime}}))$ is good then the second player wins the position $((s_{1},\ldots,s_{q^{\prime}}),(s^{\prime}_{1},\ldots,s^{\prime}_{q^{\prime}}))$ of the game.
%
%
%\end{lemma}
%
%
%\begin{lemma}\label{lemma:empty:position:is:good:gamma}
%
%
%The empty position of the Ehrenfeucht-Fra\"\i ss\'e game ${\mathcal G}_{q}({\mathcal F}_{A,B}^{\rho},
%
%
%
%
%
%
%
%
%
%
%
%
%$\linebreak$
%
%
%
%
%
%
%
%
%
%
%
%
%{\mathcal F}_{A^{\prime},B^{\prime}}^{\rho^{\prime}})$ is good.
%
%
%\end{lemma}
%
%
%\begin{lemma}\label{lemma:about:extending:winning:positions:b:gamma}
%
%
%For all $q^{\prime}{\in}\lsem0,q{-}1\rsem$ and for all positions $((s_{1},\ldots,s_{q^{\prime}}),(s^{\prime}_{1},\ldots,s^{\prime}_{q^{\prime}}))$ of the Ehrenfeucht-Fra\"\i ss\'e game ${\mathcal G}_{q}({\mathcal F}_{A,B}^{\rho},{\mathcal F}_{A^{\prime},B^{\prime}}^{\rho^{\prime}})$, if $((s_{1},\ldots,s_{q^{\prime}}),(s^{\prime}_{1},\ldots,s^{\prime}_{q^{\prime}}))$ is good then
%
%
%\begin{itemize}
%
%
%\item for all $s{\in}A{\cup}B$, there exists $s^{\prime}{\in}A^{\prime}{\cup}B^{\prime}$ such that $((s_{1},\ldots,s_{q^{\prime}},s),(s^{\prime}_{1},\ldots,
%
%
%
%
%
%
%
%
%
%
%
%
%$\linebreak$
%
%
%
%
%
%
%
%
%
%
%
%
%s^{\prime}_{q^{\prime}},s^{\prime}))$ is good,
%
%
%\item for all $s^{\prime}{\in}A^{\prime}{\cup}B^{\prime}$, there exists $s{\in}A{\cup}B$ such that $((s_{1},\ldots,s_{q^{\prime}},s),(s^{\prime}_{1},\ldots,
%
%
%
%
%
%
%
%
%
%
%
%
%$\linebreak$
%
%
%
%
%
%
%
%
%
%
%
%
%s^{\prime}_{q^{\prime}},s^{\prime}))$ is good.
%
%
%\end{itemize}
%
%
%\end{lemma}
%
%
%\begin{lemma}\label{lemma:B:is:not:equal:B:prime:when:cardinalities:are:different:first}
%
%
%For all $t^{\prime}{\in}B^{\prime}$, if ${\parallel}\rho^{{-}1}(\{t^{\prime}\}){\parallel}{\leq}q$ and ${\parallel}{\rho^{\prime}}^{{-}1}(\{t^{\prime}\}){\parallel}{>}q$ then $B^{\prime}{\not=}B$.
%
%
%\end{lemma}
%
%
%\begin{lemma}\label{lemma:B:is:not:equal:B:prime:when:cardinalities:are:different:second}
%
%
%For all $t^{\prime}{\in}B^{\prime}$, if ${\parallel}\rho^{{-}1}(B{\setminus}\{t^{\prime}\}){\parallel}{\leq}q$ and ${\parallel}{\rho^{\prime}}^{{-}1}(B^{\prime}{\setminus}\{t^{\prime}\}){\parallel}{>}q$ then $B^{\prime}{\not=}B$.
%
%
%\end{lemma}
%
%
\begin{lemma}\label{lemma:gamma:reduction:prebounded:dust:quasi:root:bounded:OK}
%
%
\begin{enumerate}
%
%
\item $A^{\prime}{\subseteq}A$,
%
%
\item $B^{\prime}{\subseteq}B$ and if ${\parallel}B^{\prime}{\parallel}{\leq}2$ then $B^{\prime}{=}B$,
%
%
\item for all $s^{\prime}{\in}A^{\prime}$, $\rho^{\prime}(s^{\prime}){\subseteq}\rho(s^{\prime})$,
%
%
\item ${\parallel}{\rho^{\prime}}^{{-}1}(\emptyset){\parallel}{\leq}q$ and ${\parallel}{\rho^{\prime}}^{{-}1}(B^{\prime}){\parallel}{\leq}q$,
%
%
\item for all $t^{\prime}{\in}B^{\prime}$, ${\parallel}{\rho^{\prime}}^{{-}1}(\{t^{\prime}\}){\parallel}{\leq}q$ and for all $t^{\prime}{\in}B^{\prime}$, ${\parallel}{\rho^{\prime}}^{{-}1}(B^{\prime}{\setminus}\{t^{\prime}\}){\parallel}{\leq}q$,
%
%
\item ${\mathcal F}_{A^{\prime},B^{\prime}}^{\rho^{\prime}}$ is simple, ${\parallel}B^{\prime}{\parallel}{\leq}q{\times}(q{+}1)^{2}$ and ${\parallel}A^{\prime}{\parallel}{\leq}2{\times}q{\times}(q{\times}(q{+}1)^{2}+1)$,
%
%
\item for all $s^{\prime}{\in}A^{\prime}{\cup}B^{\prime}$, there exists $s{\in}A{\cup}B$ such that the least generated subframe of ${\mathcal F}_{A^{\prime},B^{\prime}}^{\rho^{\prime}}$ containing $s^{\prime}$ is a bounded morphic image of the least generated subframe of ${\mathcal F}_{A,B}^{\rho}$ containing $s$,
%$f$ is a surjective bounded morphism from ${\mathcal F}_{A,B}^{\rho}$ to ${\mathcal F}_{A^{\prime},B^{\prime}}^{\rho^{\prime}}$,
%
%
\item the second player has a winning strategy in the Ehrenfeucht-Fra\"\i ss\'e game ${\mathcal G}_{q}({\mathcal F}_{A,B}^{\rho},{\mathcal F}_{A^{\prime},B^{\prime}}^{\rho^{\prime}})$.
%
%
\end{enumerate}
%
%
\end{lemma}
%
%
\begin{proof}
%
%
$\mathbf{(1)}$--$\mathbf{(6)}$~By construction of ${\mathcal F}_{A^{\prime},B^{\prime}}^{\rho^{\prime}}$.
%
%
\\
\\
%
%
$\mathbf{(7)}$~Let $s^{\prime}{\in}A^{\prime}{\cup}B^{\prime}$.
Hence, either $\mathbf{(i)}$~$s^{\prime}{\in}\rho^{-1}(\emptyset)$, or $\mathbf{(ii)}$~$s^{\prime}{\in}{\cup}\rho^{{-}1}(B)$, or $\mathbf{(iii)}$~there exists $u{\in}B^{\prime}$ such that $s^{\prime}{\in}\rho^{{-}1}(\{u\})$, or $\mathbf{(iv)}$~there exists $u{\in}B^{\prime}$ such that $s^{\prime}{\in}\rho^{{-}1}(B
%
%
%
%
%
%
%
%
%
%
%
%
$\linebreak$
%
%
%
%
%
%
%
%
%
%
%
%
{\setminus}\{u\})$, or $\mathbf{(v)}$~$s^{\prime}{\in}B^{\prime}$.
We consider the following $5$~cases.
%
%
\\
\\
%
%
{\bf Case $\mathbf{(i)}$:}
Thus, $s^{\prime}{\in}A^{\prime}$, $s^{\prime}{\in}A$, $\rho^{\prime}(s^{\prime}){=}\emptyset$ and $\rho(s^{\prime}){=}\emptyset$.
Consequently, the least generated subframe of ${\mathcal F}_{A^{\prime},B^{\prime}}^{\rho^{\prime}}$ containing $s^{\prime}$ is isomorphic to the isolated point and the least generated subframe of ${\mathcal F}_{A,B}^{\rho}$ containing $s^{\prime}$ is isomorphic to the isolated point.
Hence, the former generated subframe is a bounded morphic image of the latter generated subframe.
%
%
\\
\\
%
%
{\bf Case $\mathbf{(ii)}$:}
Thus, $s^{\prime}{\in}A^{\prime}$, $s^{\prime}{\in}A$, $\rho^{\prime}(s^{\prime}){=}B^{\prime}$ and $\rho(s^{\prime}){=}B$.
Consequently, the least generated subframe of ${\mathcal F}_{A^{\prime},B^{\prime}}^{\rho^{\prime}}$ containing $s^{\prime}$ is of the form $(\{s^{\prime}\}{\cup}B^{\prime},(\{s^{\prime}\}{\times}B^{\prime}){\cup}
%
%
%
%
%
%
%
%
%
%
%
%
$\linebreak$
%
%
%
%
%
%
%
%
%
%
%
%
(B^{\prime}{\times}B^{\prime}))$ and the least generated subframe of ${\mathcal F}_{A,B}^{\rho}$ containing $s^{\prime}$ is of the form $(\{s^{\prime}\}{\cup}B,(\{s^{\prime}\}{\times}B){\cup}(B{\times}B))$.
Since $B^{\prime}{\subseteq}B$, then ${\parallel}B^{\prime}{\parallel}{\leq}{\parallel}B{\parallel}$.
Since $B^{\prime}{\not=}\emptyset$ and $B{\not=}\emptyset$, then by Lemma~\ref{lemma:inclusion:Bprime:B:for:generated:subframes:is:bmi}, the former generated subframe is a bounded morphic image of the latter generated subframe.
%
%
\\
\\
%
%
{\bf Case $\mathbf{(iii)}$:}
Hence, $s^{\prime}{\in}A^{\prime}$, $u{\in}B^{\prime}$, $s^{\prime}{\in}A$, $u{\in}B$, $\rho^{\prime}(s^{\prime}){=}\{u\}$ and $\rho(s^{\prime}){=}\{u\}$.
Thus, the least generated subframe of ${\mathcal F}_{A^{\prime},B^{\prime}}^{\rho^{\prime}}$ containing $s^{\prime}$ is of the form $(\{s^{\prime}\}{\cup}B^{\prime},\{(s^{\prime},
%
%
%
%
%
%
%
%
%
%
%
%
$\linebreak$
%
%
%
%
%
%
%
%
%
%
%
%
u)\}{\cup}(B^{\prime}{\times}B^{\prime}))$ and the least generated subframe of ${\mathcal F}_{A,B}^{\rho}$ containing $s^{\prime}$ is of the form $(\{s^{\prime}\}{\cup}B,\{(s^{\prime},u)\}{\cup}(B{\times}B))$.
Since $B^{\prime}{\subseteq}B$, then ${\parallel}B^{\prime}{\parallel}{\leq}{\parallel}B{\parallel}$.
Since $B^{\prime}{\not=}\emptyset$ and $B{\not=}\emptyset$, then by Lemma~\ref{lemma:inclusion:Bprime:B:for:generated:subframes:is:bmi:bis}, the former generated subframe is a bounded morphic image of the latter generated subframe.
%
%
\\
\\
%
%
{\bf Case $\mathbf{(iv)}$:}
Consequently, $s^{\prime}{\in}A^{\prime}$, $u{\in}B^{\prime}$, $s^{\prime}{\in}A$, $u{\in}B$, $\rho^{\prime}(s^{\prime}){=}B^{\prime}{\setminus}\{u\}$ and $\rho(s^{\prime}){=}B
%
%
%
%
%
%
%
%
%
%
%
%
$\linebreak$
%
%
%
%
%
%
%
%
%
%
%
%
{\setminus}\{u\}$.
Hence, the least generated subframe of ${\mathcal F}_{A^{\prime},B^{\prime}}^{\rho^{\prime}}$ containing $s^{\prime}$ is of the form $(\{s^{\prime}\}{\cup}B^{\prime},(\{s^{\prime}\}{\times}(B^{\prime}{\setminus}\{u\})){\cup}(B^{\prime}{\times}B^{\prime}))$ and the least generated subframe of ${\mathcal F}_{A,B}^{\rho}$ containing $s^{\prime}$ is of the form $(\{s^{\prime}\}{\cup}B,(\{s^{\prime}\}{\times}(B{\setminus}\{u\})){\cup}(B{\times}B))$.
Since $B^{\prime}{\subseteq}B$, then ${\parallel}B^{\prime}{\parallel}{\leq}{\parallel}B{\parallel}$.
Since $B^{\prime}{\not=}\emptyset$ and $B{\not=}\emptyset$, then by Lemma~\ref{lemma:inclusion:Bprime:B:for:generated:subframes:is:bmi:ter}, the former generated subframe is a bounded morphic image of the latter generated subframe.
%
%
\\
\\
%
%
{\bf Case $\mathbf{(v)}$:}
Thus, $s^{\prime}{\in}B^{\prime}$ and $s^{\prime}{\in}B$.
Consequently, the least generated subframe of ${\mathcal F}_{A^{\prime},B^{\prime}}^{\rho^{\prime}}$ containing $s^{\prime}$ is of the form $(B^{\prime},B^{\prime}{\times}B^{\prime})$ and the least generated subframe of ${\mathcal F}_{A,B}^{\rho}$ containing $s^{\prime}$ is of the form $(B,B{\times}B)$.
Since $B^{\prime}{\subseteq}B$, then ${\parallel}B^{\prime}{\parallel}{\leq}{\parallel}B{\parallel}$.
Since $B^{\prime}{\not=}\emptyset$ and $B{\not=}\emptyset$, then by Lemma~\ref{lemma:inclusion:Bprime:B:for:generated:subframes:is:bmi:quater}, the former generated subframe is a bounded morphic image of the latter generated subframe.
%
%
\\
\\
%
%
$\mathbf{(8)}$~Let us assume that some well-order on $A{\cup}B$ has been fixed.
Let $q^{\prime}{\in}\lsem1,q\rsem$ and $((s_{1},\ldots,s_{q^{\prime}{-}1}),(s^{\prime}_{1},\ldots,s^{\prime}_{q^{\prime}{-}1}))$ be a winning position of the Ehrenfeucht-Fra\"\i ss\'e game ${\mathcal G}_{q}({\mathcal F}_{A,B}^{\rho},{\mathcal F}_{A^{\prime},B^{\prime}}^{\rho^{\prime}})$.
Suppose the first player picked a member $s$ of ${\mathcal F}_{A,B}^{\rho}$ in its $q^{\prime}$th move.
Then
%
%
\begin{itemize}
%
%
\item if $s{=}s_{q^{\prime\prime}}$ for some $q^{\prime\prime}{\in}\lsem1,q^{\prime}{-}1\rsem$ and $s{\in}B$ then the second player picks $s^{\prime}_{q^{\prime\prime}}$ in $B^{\prime}$ in its $q^{\prime}$th move,
%
%
\item if $s{\not=}s_{q^{\prime\prime}}$ for every $q^{\prime\prime}{\in}\lsem1,q^{\prime}{-}1\rsem$, $s{\in}B$ and ${\parallel}\lbrack s\rbrack_{\approx}{\parallel}{\leq}q$ then the second player picks $s$ in $B^{\prime}$ in its $q^{\prime}$th move,
%
%
\item if $s{\not=}s_{q^{\prime\prime}}$ for every $q^{\prime\prime}{\in}\lsem1,q^{\prime}{-}1\rsem$, $s{\in}B$ and ${\parallel}\lbrack s\rbrack_{\approx}{\parallel}{>}q$ then the second player picks the least $s^{\prime}$ in $B^{\prime}$ such that
%
%
\begin{itemize}
%
%
\item $s{\approx}s^{\prime}$,
%
%
\item $s^{\prime}{\not=}s^{\prime}_{q^{\prime\prime}}$ for every $q^{\prime\prime}{\in}\lsem1,q^{\prime}{-}1\rsem$,
%
%
\item for all $q^{\prime\prime}{\in}\lsem1,q^{\prime}{-}1\rsem$, if $s_{q^{\prime\prime}}{\in}A$, $\rho(s_{q^{\prime\prime}}){\not=}\emptyset$ and $\rho(s_{q^{\prime\prime}}){\not=}B$ then
%
%
\begin{itemize}
%
%
\item $\rho(s_{q^{\prime\prime}}){=}\{s\}$ if and only if $\rho^{\prime}(s^{\prime}_{q^{\prime\prime}}){=}\{s^{\prime}\}$,
%
%
\item $\rho(s_{q^{\prime\prime}}){=}B{\setminus}\{s\}$ if and only if $\rho^{\prime}(s^{\prime}_{q^{\prime\prime}}){=}B^{\prime}{\setminus}\{s^{\prime}\}$,
%
%
\end{itemize}
%
%
\end{itemize}
%
%
in its $q^{\prime}$th move,
%
%
\item if $s{=}s_{q^{\prime\prime}}$ for some $q^{\prime\prime}{\in}\lsem1,q^{\prime}{-}1\rsem$ and $s{\in}A$ then the second player picks $s^{\prime}_{q^{\prime\prime}}$ in $A^{\prime}$ in its $q^{\prime}$th move,
%
%
\item if $s{\not=}s_{q^{\prime\prime}}$ for every $q^{\prime\prime}{\in}\lsem1,q^{\prime}{-}1\rsem$, $s{\in}A$ and either $\rho(s){=}\emptyset$, or $\rho(s){=}B$ then the second player picks $s$ in $A^{\prime}$ in its $q^{\prime}$th move,
%
%
\item if $s{\not=}s_{q^{\prime\prime}}$ for every $q^{\prime\prime}{\in}\lsem1,q^{\prime}{-}1\rsem$, $s{\in}A$ and either $\rho(s){=}\{t\}$, or $\rho(s){=}B{\setminus}\{t\}$ for some $t{\in}B$ such that ${\parallel}\lbrack t\rbrack_{\approx}{\parallel}{\leq}q$ then the second player picks $s$ in $A^{\prime}$ in its $q^{\prime}$th move,
%
%
\item if $s{\not=}s_{q^{\prime\prime}}$ for every $q^{\prime\prime}{\in}\lsem1,q^{\prime}{-}1\rsem$, $s{\in}A$ and $\rho(s){=}\{t\}$ for some $t{\in}B$ such that ${\parallel}\lbrack t\rbrack_{\approx}{\parallel}{>}q$ and $t{=}s_{q^{\prime\prime\prime}}$ for some $q^{\prime\prime\prime}{\in}\lsem1,q^{\prime}{-}1\rsem$ then the second player picks the least element in ${\rho^{\prime}}^{{-}1}(\{s^{\prime}_{q^{\prime\prime\prime}}\}){\setminus}\{s^{\prime}_{1},\ldots,s^{\prime}_{q^{\prime}{-}1}\}$,
%
%
\item if $s{\not=}s_{q^{\prime\prime}}$ for every $q^{\prime\prime}{\in}\lsem1,q^{\prime}{-}1\rsem$, $s{\in}A$ and $\rho(s){=}B{\setminus}\{t\}$ for some $t{\in}B$ such that ${\parallel}\lbrack t\rbrack_{\approx}{\parallel}{>}q$ and $t{=}s_{q^{\prime\prime\prime}}$ for some $q^{\prime\prime\prime}{\in}\lsem1,q^{\prime}{-}1\rsem$ then the second player picks the least element in ${\rho^{\prime}}^{{-}1}(B^{\prime}{\setminus}\{s^{\prime}_{q^{\prime\prime\prime}}\}){\setminus}\{s^{\prime}_{1},\ldots,s^{\prime}_{q^{\prime}{-}1}\}$,
%
%
\item if $s{\not=}s_{q^{\prime\prime}}$ for every $q^{\prime\prime}{\in}\lsem1,q^{\prime}{-}1\rsem$, $s{\in}A$ and $\rho(s){=}\{t\}$ for some $t{\in}B$ such that ${\parallel}\lbrack t\rbrack_{\approx}{\parallel}{>}q$ and $t{\not=}s_{q^{\prime\prime\prime}}$ for every $q^{\prime\prime\prime}{\in}\lsem1,q^{\prime}{-}1\rsem$ then $t^{\prime}$ being the least element in $B^{\prime}$ such that
%
%
\begin{itemize}
%
%
\item $t{\approx}t^{\prime}$,
%
%
\item $t^{\prime}{\not=}s^{\prime}_{q^{\prime\prime}}$ for every $q^{\prime\prime}{\in}\lsem1,q^{\prime}{-}1\rsem$,
%
%
\item for all $q^{\prime\prime}{\in}\lsem1,q^{\prime}{-}1\rsem$, if $s_{q^{\prime\prime}}{\in}A$, $\rho(s_{q^{\prime\prime}}){\not=}\emptyset$ and $\rho(s_{q^{\prime\prime}}){\not=}B$ then
%
%
\begin{itemize}
%
%
\item $\rho(s_{q^{\prime\prime}}){=}\{t\}$ if and only if $\rho^{\prime}(s^{\prime}_{q^{\prime\prime}}){=}\{t^{\prime}\}$,
%
%
\item $\rho(s_{q^{\prime\prime}}){=}B{\setminus}\{t\}$ if and only if $\rho^{\prime}(s^{\prime}_{q^{\prime\prime}}){=}B^{\prime}{\setminus}\{t^{\prime}\}$,
%
%
\end{itemize}
%
%
\end{itemize}
%
%
the second player picks the least element in ${\rho^{\prime}}^{{-}1}(\{t^{\prime}\}){\setminus}\{s^{\prime}_{1},\ldots,s^{\prime}_{q^{\prime}{-}1}\}$,
%
%
\item if $s{\not=}s_{q^{\prime\prime}}$ for every $q^{\prime\prime}{\in}\lsem1,q^{\prime}{-}1\rsem$, $s{\in}A$ and $\rho(s){=}B{\setminus}\{t\}$ for some $t{\in}B$ such that ${\parallel}\lbrack t\rbrack_{\approx}{\parallel}{>}q$ and $t{\not=}s_{q^{\prime\prime\prime}}$ for some $q^{\prime\prime\prime}{\in}\lsem1,q^{\prime}{-}1\rsem$ then $t^{\prime}$ being the least element in $B^{\prime}$ such that
%
%
\begin{itemize}
%
%
\item $t{\approx}t^{\prime}$,
%
%
\item $t^{\prime}{\not=}s^{\prime}_{q^{\prime\prime}}$ for every $q^{\prime\prime}{\in}\lsem1,q^{\prime}{-}1\rsem$,
%
%
\item for all $q^{\prime\prime}{\in}\lsem1,q^{\prime}{-}1\rsem$, if $s_{q^{\prime\prime}}{\in}A$, $\rho(s_{q^{\prime\prime}}){\not=}\emptyset$ and $\rho(s_{q^{\prime\prime}}){\not=}B$ then
%
%
\begin{itemize}
%
%
\item $\rho(s_{q^{\prime\prime}}){=}\{t\}$ if and only if $\rho^{\prime}(s^{\prime}_{q^{\prime\prime}}){=}\{t^{\prime}\}$,
%
%
\item $\rho(s_{q^{\prime\prime}}){=}B{\setminus}\{t\}$ if and only if $\rho^{\prime}(s^{\prime}_{q^{\prime\prime}}){=}B^{\prime}{\setminus}\{t^{\prime}\}$,
%
%
\end{itemize}
%
%
\end{itemize}
%
%
the second player picks the least element in ${\rho^{\prime}}^{{-}1}(B^{\prime}{\setminus}\{t^{\prime}\}){\setminus}\{s^{\prime}_{1},\ldots,s^{\prime}_{q^{\prime}{-}1}\}$.
%
%
\end{itemize}
%
%
Suppose the first player picked a member $s^{\prime}$ of ${\mathcal F}_{A^{\prime},B^{\prime}}^{\rho^{\prime}}$ in its $q^{\prime}$th move.
Then
%
%
\begin{itemize}
%
%
\item if $s^{\prime}{=}s^{\prime}_{q^{\prime\prime}}$ for some $q^{\prime\prime}{\in}\lsem1,q^{\prime}{-}1\rsem$ and $s^{\prime}{\in}B^{\prime}$ then the second player picks $s_{q^{\prime\prime}}$ in $B$ in its $q^{\prime}$th move,
%
%
\item if $s^{\prime}{\not=}s^{\prime}_{q^{\prime\prime}}$ for every $q^{\prime\prime}{\in}\lsem1,q^{\prime}{-}1\rsem$, $s^{\prime}{\in}B^{\prime}$ and ${\parallel}\lbrack s^{\prime}\rbrack_{\approx}{\parallel}{\leq}q$ then the second player picks $s^{\prime}$ in $B$ in its $q^{\prime}$th move,
%
%
\item if $s^{\prime}{\not=}s^{\prime}_{q^{\prime\prime}}$ for every $q^{\prime\prime}{\in}\lsem1,q^{\prime}{-}1\rsem$, $s^{\prime}{\in}B^{\prime}$ and ${\parallel}\lbrack s^{\prime}\rbrack_{\approx}{\parallel}{>}q$ then the second player picks the least $s$ in $B$ such that
%
%
\begin{itemize}
%
%
\item $s^{\prime}{\approx}s$,
%
%
\item $s{\not=}s_{q^{\prime\prime}}$ for every $q^{\prime\prime}{\in}\lsem1,q^{\prime}{-}1\rsem$,
%
%
\item for all $q^{\prime\prime}{\in}\lsem1,q^{\prime}{-}1\rsem$, if $s^{\prime}_{q^{\prime\prime}}{\in}A^{\prime}$, $\rho^{\prime}(s^{\prime}_{q^{\prime\prime}}){\not=}\emptyset$ and $\rho^{\prime}(s^{\prime}_{q^{\prime\prime}}){\not=}B^{\prime}$ then
%
%
\begin{itemize}
%
%
\item $\rho^{\prime}(s^{\prime}_{q^{\prime\prime}}){=}\{s^{\prime}\}$ if and only if $\rho(s_{q^{\prime\prime}}){=}\{s\}$,
%
%
\item $\rho^{\prime}(s^{\prime}_{q^{\prime\prime}}){=}B^{\prime}{\setminus}\{s^{\prime}\}$ if and only if $\rho(s_{q^{\prime\prime}}){=}B{\setminus}\{s\}$,
%
%
\end{itemize}
%
%
\end{itemize}
%
%
in its $q^{\prime}$th move,
%
%
\item if $s^{\prime}{=}s^{\prime}_{q^{\prime\prime}}$ for some $q^{\prime\prime}{\in}\lsem1,q^{\prime}{-}1\rsem$ and $s^{\prime}{\in}A^{\prime}$ then the second player picks $s_{q^{\prime\prime}}$ in $A$ in its $q^{\prime}$th move,
%
%
\item if $s^{\prime}{\not=}s^{\prime}_{q^{\prime\prime}}$ for every $q^{\prime\prime}{\in}\lsem1,q^{\prime}{-}1\rsem$, $s^{\prime}{\in}A^{\prime}$ and either $\rho^{\prime}(s^{\prime}){=}\emptyset$, or $\rho^{\prime}(s^{\prime}){=}B^{\prime}$ then the second player picks $s^{\prime}$ in $A$ in its $q^{\prime}$th move,
%
%
\item if $s^{\prime}{\not=}s^{\prime}_{q^{\prime\prime}}$ for every $q^{\prime\prime}{\in}\lsem1,q^{\prime}{-}1\rsem$, $s^{\prime}{\in}A^{\prime}$ and either $\rho^{\prime}(s^{\prime}){=}\{t^{\prime}\}$, or $\rho^{\prime}(s^{\prime}){=}B^{\prime}
%
%
%
%
%
%
%
%
%
%
%
%
$\linebreak$
%
%
%
%
%
%
%
%
%
%
%
%
{\setminus}\{t^{\prime}\}$ for some $t^{\prime}{\in}B^{\prime}$ such that ${\parallel}\lbrack t^{\prime}\rbrack_{\approx}{\parallel}{\leq}q$ then the second player picks $s^{\prime}$ in $A$ in its $q^{\prime}$th move,
%
%
\item if $s^{\prime}{\not=}s^{\prime}_{q^{\prime\prime}}$ for every $q^{\prime\prime}{\in}\lsem1,q^{\prime}{-}1\rsem$, $s^{\prime}{\in}A^{\prime}$ and $\rho(s^{\prime}){=}\{t^{\prime}\}$ for some $t^{\prime}{\in}B^{\prime}$ such that ${\parallel}\lbrack t^{\prime}\rbrack_{\approx}{\parallel}{>}q$ and $t^{\prime}{=}s^{\prime}_{q^{\prime\prime\prime}}$ for some $q^{\prime\prime\prime}{\in}\lsem1,q^{\prime}{-}1\rsem$ then the second player picks the least element in ${\rho}^{{-}1}(\{s_{q^{\prime\prime\prime}}\}){\setminus}\{s_{1},\ldots,s_{q^{\prime}{-}1}\}$,
%
%
\item if $s^{\prime}{\not=}s^{\prime}_{q^{\prime\prime}}$ for every $q^{\prime\prime}{\in}\lsem1,q^{\prime}{-}1\rsem$, $s^{\prime}{\in}A^{\prime}$ and $\rho^{\prime}(s^{\prime}){=}B^{\prime}{\setminus}\{t^{\prime}\}$ for some $t^{\prime}{\in}B^{\prime}$ such that ${\parallel}\lbrack t^{\prime}\rbrack_{\approx}{\parallel}{>}q$ and $t^{\prime}{=}s^{\prime}_{q^{\prime\prime\prime}}$ for some $q^{\prime\prime\prime}{\in}\lsem1,q^{\prime}{-}1\rsem$ then the second player picks the least element in ${\rho}^{{-}1}(B{\setminus}\{s_{q^{\prime\prime\prime}}\}){\setminus}\{s_{1},\ldots,s_{q^{\prime}{-}1}\}$,
%
%
\item if $s^{\prime}{\not=}s^{\prime}_{q^{\prime\prime}}$ for every $q^{\prime\prime}{\in}\lsem1,q^{\prime}{-}1\rsem$, $s^{\prime}{\in}A^{\prime}$ and $\rho^{\prime}(s^{\prime}){=}\{t^{\prime}\}$ for some $t^{\prime}{\in}B^{\prime}$ such that ${\parallel}\lbrack t^{\prime}\rbrack_{\approx}{\parallel}{>}q$ and $t^{\prime}{\not=}s^{\prime}_{q^{\prime\prime\prime}}$ for every $q^{\prime\prime\prime}{\in}\lsem1,q^{\prime}{-}1\rsem$ then $t$ being the least element in $B$ such that
%
%
\begin{itemize}
%
%
\item $t^{\prime}{\approx}t$,
%
%
\item $t{\not=}s_{q^{\prime\prime}}$ for every $q^{\prime\prime}{\in}\lsem1,q^{\prime}{-}1\rsem$,
%
%
\item for all $q^{\prime\prime}{\in}\lsem1,q^{\prime}{-}1\rsem$, if $s^{\prime}_{q^{\prime\prime}}{\in}A^{\prime}$, $\rho^{\prime}(s^{\prime}_{q^{\prime\prime}}){\not=}\emptyset$ and $\rho^{\prime}(s^{\prime}_{q^{\prime\prime}}){\not=}B^{\prime}$ then
%
%
\begin{itemize}
%
%
\item $\rho^{\prime}(s^{\prime}_{q^{\prime\prime}}){=}\{t^{\prime}\}$ if and only if $\rho(s_{q^{\prime\prime}}){=}\{t\}$,
%
%
\item $\rho^{\prime}(s^{\prime}_{q^{\prime\prime}}){=}B^{\prime}{\setminus}\{t^{\prime}\}$ if and only if $\rho(s_{q^{\prime\prime}}){=}B{\setminus}\{t\}$,
%
%
\end{itemize}
%
%
\end{itemize}
%
%
the second player picks the least element in ${\rho}^{{-}1}(\{t\}){\setminus}\{s_{1},\ldots,s_{q^{\prime}{-}1}\}$,
%
%
\item if $s^{\prime}{\not=}s^{\prime}_{q^{\prime\prime}}$ for every $q^{\prime\prime}{\in}\lsem1,q^{\prime}{-}1\rsem$, $s^{\prime}{\in}A^{\prime}$ and $\rho^{\prime}(s^{\prime}){=}B^{\prime}{\setminus}\{t^{\prime}\}$ for some $t^{\prime}{\in}B^{\prime}$ such that ${\parallel}\lbrack t^{\prime}\rbrack_{\approx}{\parallel}{>}q$ and $t^{\prime}{\not=}s^{\prime}_{q^{\prime\prime\prime}}$ for some $q^{\prime\prime\prime}{\in}\lsem1,q^{\prime}{-}1\rsem$ then $t$ being the least element in $B$ such that
%
%
\begin{itemize}
%
%
\item $t^{\prime}{\approx}t$,
%
%
\item $t{\not=}s_{q^{\prime\prime}}$ for every $q^{\prime\prime}{\in}\lsem1,q^{\prime}{-}1\rsem$,
%
%
\item for all $q^{\prime\prime}{\in}\lsem1,q^{\prime}{-}1\rsem$, if $s^{\prime}_{q^{\prime\prime}}{\in}A^{\prime}$, $\rho^{\prime}(s^{\prime}_{q^{\prime\prime}}){\not=}\emptyset$ and $\rho^{\prime}(s^{\prime}_{q^{\prime\prime}}){\not=}B^{\prime}$ then
%
%
\begin{itemize}
%
%
\item $\rho^{\prime}(s^{\prime}_{q^{\prime\prime}}){=}\{t^{\prime}\}$ if and only if $\rho(s_{q^{\prime\prime}}){=}\{t\}$,
%
%
\item $\rho^{\prime}(s^{\prime}_{q^{\prime\prime}}){=}B^{\prime}{\setminus}\{t^{\prime}\}$ if and only if $\rho(s_{q^{\prime\prime}}){=}B{\setminus}\{t\}$,
%
%
\end{itemize}
%
%
\end{itemize}
%
%
the second player picks the least element in ${\rho}^{{-}1}(B{\setminus}\{t\}){\setminus}\{s_{1},\ldots,s_{q^{\prime}{-}1}\}$.
%
%
\end{itemize}
%
%
Obviously, $((s_{1},\ldots,s_{q^{\prime}{-}1},s),(s^{\prime}_{1},\ldots,s^{\prime}_{q^{\prime}{-}1},s^{\prime}))$ is a winning position of the Eh\-renfeucht-Fra\"\i ss\'e game ${\mathcal G}_{q}({\mathcal F}_{A,B}^{\rho},{\mathcal F}_{A^{\prime},B^{\prime}}^{\rho^{\prime}})$.
%
%
\medskip
%
%
\end{proof}
%
%
%Let $q{\geq}3$.
Let $\{{\mathcal F}_{A_{i},B_{i}}^{\rho_{i}}:\ i{\in}I\}$ be a disjoint family of galaxies such that for all $i{\in}A$, either ${\mathcal F}_{A_{i},B_{i}}^{\rho_{i}}$ is a simple galaxy satisfying the conditions $(4)$ and $(5)$ of Lemma~\ref{lemma:the:alpha:reduction:is:always:dust:finite}, or ${\mathcal F}_{A_{i},B_{i}}^{\rho_{i}}$ is a non-simple galaxy.
$\{{\mathcal F}_{A_{i}^{\prime},B_{i}^{\prime}}^{\rho_{i}^{\prime}}:\ i{\in}I\}$ is a {\it $q$-$\gamma$-reduction of $\{{\mathcal F}_{A_{i},B_{i}}^{\rho_{i}}:\ i{\in}I\}$}\/ if for all $i{\in}I$, either ${\mathcal F}_{A_{i},B_{i}}^{\rho_{i}}$ is a simple galaxy satisfying the conditions $(4)$ and $(5)$ of Lemma~\ref{lemma:the:alpha:reduction:is:always:dust:finite} and ${\mathcal F}_{A_{i}^{\prime},B_{i}^{\prime}}^{\rho_{i}^{\prime}}$ is a $q$-$\gamma$-reduction of ${\mathcal F}_{A_{i},B_{i}}^{\rho_{i}}$, or ${\mathcal F}_{A_{i},B_{i}}^{\rho_{i}}$ is a non-simple galaxy and ${\mathcal F}_{A_{i}^{\prime},B_{i}^{\prime}}^{\rho_{i}^{\prime}}$ is equal to ${\mathcal F}_{A_{i},B_{i}}^{\rho_{i}}$.
%
%
\begin{proposition}\label{lemma:good:properties:of:gamma:reductions:indexed}
%
%
%Let $q{\geq}3$.
%Let $\{{\mathcal F}_{A_{i},B_{i}}^{\rho_{i}}:\ i{\in}I\}$ be a disjoint family of galaxies such that for all $i{\in}A$, either ${\mathcal F}_{A_{i},B_{i}}^{\rho_{i}}$ is a simple galaxy satisfying the conditions $(4)$ and $(5)$ of Lemma~\ref{lemma:the:alpha:reduction:is:always:dust:finite}, or ${\mathcal F}_{A_{i},B_{i}}^{\rho_{i}}$ is a non-simple galaxy.
%We shall say that $\{(W_{A_{i}^{\prime},B_{i}^{\prime}}^{\rho_{i}^{\prime}},R_{A_{i}^{\prime},B_{i}^{\prime}}^{\rho_{i}^{\prime}}):\ i{\in}I\}$ is {\it $q$-$\alpha$-reduced from $\{(W_{A_{i},B_{i}}^{\rho_{i}},R_{A_{i},B_{i}}^{\rho_{i}}):\ i{\in}I\}$}\/ if for all $i{\in}I$, either $(W_{A_{i}^{\prime},B_{i}^{\prime}}^{\rho_{i}^{\prime}},R_{A_{i}^{\prime},B_{i}^{\prime}}^{\rho_{i}^{\prime}})$ is equal to $(W_{A_{i},B_{i}}^{\rho_{i}},R_{A_{i},B_{i}}^{\rho_{i}})$, or $(W_{A_{i}^{\prime},B_{i}^{\prime}}^{\rho_{i}^{\prime}},R_{A_{i}^{\prime},B_{i}^{\prime}}^{\rho_{i}^{\prime}})$ is a $q$-$\alpha$-reduction of $(W_{A_{i},B_{i}}^{\rho_{i}},R_{A_{i},B_{i}}^{\rho_{i}})$.
Let $I^{ns}$ be the set of all $i{\in}I$ such that ${\mathcal F}_{A_{i},B_{i}}^{\rho_{i}}$ is non-simple.
If $\{{\mathcal F}_{A_{i}^{\prime},B_{i}^{\prime}}^{\rho_{i}^{\prime}}:\ i{\in}I\}$ is a $q$-$\gamma$-reduction of $\{{\mathcal F}_{A_{i},B_{i}}^{\rho_{i}}:\ i{\in}I\}$, $(W,R)$ is the union of $\{{\mathcal F}_{A_{i},B_{i}}^{\rho_{i}}:\ i{\in}I\}$ and $(W^{\prime},R^{\prime})$ is the union of $\{{\mathcal F}_{A_{i}^{\prime},B_{i}^{\prime}}^{\rho_{i}^{\prime}}:\ i{\in}I\}$ then
%
%
\begin{enumerate}
%
%
\item if $\{{\mathcal F}_{A_{i},B_{i}}^{\rho_{i}}:\ i{\in}I\}$ is dust-bounded then $\{{\mathcal F}_{A_{i}^{\prime},B_{i}^{\prime}}^{\rho_{i}^{\prime}}:\ i{\in}I\}$ is dust-bounded,
%
%
\item if $\{{\mathcal F}_{A_{i},B_{i}}^{\rho_{i}}:\ i{\in}I^{ns}\}$ is root-bounded and kernel-bounded then $\{{\mathcal F}_{A_{i}^{\prime},B_{i}^{\prime}}^{\rho_{i}^{\prime}}:\ i{\in}I\}$ is root-bounded and kernel-bounded,
%
%
\item for all $s^{\prime}{\in}W^{\prime}$, there exists $s{\in}W$ such that $(W^{\prime}_{s^{\prime}},R^{\prime}_{s^{\prime}})$ is a bounded morphic image of $(W_{s},R_{s})$,
%
%
\item the second player has a winning strategy in the Ehrenfeucht-Fra\"\i ss\'e game ${\mathcal G}_{q}((W,R),(W^{\prime},R^{\prime}))$.
%
%
\end{enumerate}
%
%
\end{proposition}
%
%
\begin{proof}
%
%
By Lemma~\ref{lemma:gamma:reduction:prebounded:dust:quasi:root:bounded:OK}.
%
%
\medskip
%
%
\end{proof}
%
%
%
%
\paragraph{$\delta$-reductions}
%
%
%In this section, we introduce a third type of reduction: the $\delta$-reduction of a disjoint indexed family of galaxies.
The main effect of a $\delta$-reduction on a disjoint indexed family of galaxies is to reduce its size.
%
%
\\
\\
%
%
%Let $Q,K{\geq}1$.
%Let $\{{\mathcal F}_{A_{i},B_{i}}^{\rho_{i}}:\ i{\in}I\}$ be a disjoint indexed family of galaxies such that for all $i{\in}I$, either ${\mathcal F}_{A_{i},B_{i}}^{\rho_{i}}$ is simple, ${\parallel}A_{i}{\parallel}{\leq}Q$ and ${\parallel}B_{i}{\parallel}{\leq}Q$, or ${\mathcal F}_{A_{i},B_{i}}^{\rho_{i}}$ is non-simple, ${\parallel}A_{i}{\parallel}{\leq}K$ and ${\parallel}B_{i}{\parallel}{\leq}K$.
%
%
Let $\{{\mathcal F}_{A_{i},B_{i}}^{\rho_{i}}:\ i{\in}I\}$ be a disjoint indexed family of galaxies.
%
%
\\
\\
%
%
Let $\simeq$ be the equivalence relation on $I$ defined by $i{\simeq}j$ if and only if ${\mathcal F}_{A_{i},B_{i}}^{\rho_{i}}$ is isomorphic with ${\mathcal F}_{A_{j},B_{j}}^{\rho_{j}}$.
For all $i{\in}I$, let $\lbrack i\rbrack_{\simeq}$ be the equivalence class of $i$ modulo $\simeq$.
%
%
\\
\\
%
%
For all $i{\in}I$, if ${\parallel}\lbrack i\rbrack_{\simeq}{\parallel}{\leq}q$ then let $X_{i}$ be $\lbrack i\rbrack_{\simeq}$ else let $X_{i}$ be an arbitrary subset of $\lbrack i\rbrack_{\simeq}$ of cardinality $q$.
Let $J$ be $\bigcup\{X_{i}:\ i{\in}I\}$.
Obviously, from each equivalence class modulo $\simeq$, $J$ contains at most $q$ elements.
Therefore, if $\simeq$ possesses $N$ equivalence classes then $J$ contains at most $q{\times}N$ elements.
%
%
\\
\\
%
%
$\{{\mathcal F}_{A_{j},B_{j}}^{\rho_{j}}:\ j{\in}J\}$ is a {\it $q$-$\delta$-reduction of $\{{\mathcal F}_{A_{i},B_{i}}^{\rho_{i}}:\ i{\in}I\}$.}
%Let $(W,R)$ be the union of $\{{\mathcal F}_{A_{i},B_{i}}^{\rho_{i}}:\ i{\in}I\}$ and $(W^{\prime},R^{\prime})$ be the union of $\{{\mathcal F}_{A_{j},B_{j}}^{\rho_{j}}:\ j{\in}J\}$.
%
%
%\\
%\\
%
%
%Let $h:\ I{\longrightarrow}J$ be a function such that for all $i{\in}I$, either $i{\in}J$ and $h(i){=}i$, or $i{\not\in}J$ and $h(i){\simeq}i$.
%Notice that the existence of such a function is a consequence of the definition of $J$.
%Notice also that $h$ is surjective.
%Moreover, for all $i{\in}I$, $h(i){\simeq}i$.
%For all $i,j{\in}I$, if $i{\simeq}j$ then let $g_{i,j}:\ {\mathcal F}_{A_{i},B_{i}}^{\rho_{i}}{\longrightarrow}{\mathcal F}_{A_{j},B_{j}}^{\rho_{j}}$ be an isomorphism.
%Notice that the existence of such an isomorphism is a consequence of the definition of $\simeq$.
%Let $f:\ W{\longrightarrow}W^{\prime}$ be the function such that for all $i{\in}I$ and for all $s{\in}A_{i}{\cup}B_{i}$, $f(s){=}g_{i,h(i)}(s)$.
%Notice that $f$ is surjective.
%
%
%\\
%\\
%
%
%Let us define a strategy $\pi$ of the second player in the Ehrenfeucht-Fra\"\i ss\'e game ${\mathcal G}_{q}((W,R),(W^{\prime},R^{\prime}))$.
%For all $q^{\prime}{\in}\lsem1,q\rsem$, we shall say that a member $i$ of $I$ is {\it left-fresh in a position $((s_{1},\ldots,s_{q^{\prime}{-}1}),(s^{\prime}_{1},\ldots,s^{\prime}_{q^{\prime}{-}1}))$ of the game}\/ if for all $q^{\prime\prime}{\in}\lsem1,q^{\prime}{-}1\rsem$, $s_{q^{\prime\prime}}$ is not a member of ${\mathcal F}_{A_{i},B_{i}}^{\rho_{i}}$.
%For all $q^{\prime}{\in}\lsem1,q\rsem$, we shall say that a member $j$ of $J$ is {\it right-fresh in a position $((s_{1},\ldots,s_{q^{\prime}{-}1}),(s^{\prime}_{1},\ldots,s^{\prime}_{q^{\prime}{-}1}))$ of the game}\/ if for all $q^{\prime\prime}{\in}\lsem1,q^{\prime}{-}1\rsem$, $s^{\prime}_{q^{\prime\prime}}$ is not a member of ${\mathcal F}_{A_{j},B_{j}}^{\rho_{j}}$.
%Let $q^{\prime}{\in}\lsem1,q\rsem$.
%In its $q^{\prime}$-th move, the second player will apply either $\mathbf{(\pi_{1})}$, or $\mathbf{(\pi_{2})}$, depending on the $q^{\prime}$-th move of the first player.
%
%
%\\
%\\
%
%
%$\mathbf{(\pi_{1})}$~If the first player picked in its $q^{\prime}$-th move a member $s$ of $(W,R)$ then there exists $i{\in}I$ such that $s$ is a member of ${\mathcal F}_{A_{i},B_{i}}^{\rho_{i}}$ and let the second player consider the following $2$~cases.
%
%
%\\
%\\
%
%
%{\bf Case where $i$ is left-fresh in the position $((s_{1},\ldots,s_{q^{\prime}{-}1}),(s^{\prime}_{1},\ldots,s^{\prime}_{i{-}1}))$:}
%As a consequence of the construction of $J$, there exists $j{\in}J$ such that $i{\simeq}j$ and $j$ is right-fresh in the position $((s_{1},\ldots,s_{q^{\prime}{-}1}),(s^{\prime}_{1},\ldots,s^{\prime}_{q^{\prime}{-}1}))$.
%Thus, let the second player pick in its $q^{\prime}$-th move the member $g_{i,j}(s)$ of $(W^{\prime},R^{\prime})$.
%
%
%\\
%\\
%
%
%{\bf Case where $i$ is not left-fresh in the position $((s_{1},\ldots,s_{q^{\prime}{-}1}),(s^{\prime}_{1},\ldots,
%
%
%
%
%
%
%
%
%
%
%
%
%$\linebreak$
%
%
%
%
%
%
%
%
%
%
%
%
%s^{\prime}_{q^{\prime}{-}1}))$:}
%As a consequence of the construction of $J$, there exists $q^{\prime\prime}{\in}\lsem1,q^{\prime}{-}1\rsem$ such that $s_{q^{\prime\prime}}$ is a member of ${\mathcal F}_{A_{i},B_{i}}^{\rho_{i}}$.
%Thus, let $j{\in}J$ be such that $i{\simeq}j$ and $s^{\prime}_{q^{\prime\prime}}$ is a member of ${\mathcal F}_{A_{j},B_{j}}^{\rho_{j}}$.
%Thus, let the second player pick in its $q^{\prime}$-th move the member $g_{i,j}(s)$ of $(W^{\prime},R^{\prime})$.
%
%
%\\
%\\
%
%
%$\mathbf{(\pi_{2})}$~If the first player picked in its $q^{\prime}$-th move a member $s^{\prime}$ of $(W^{\prime},R^{\prime})$ then there exists $j{\in}J$ such that $s^{\prime}$ is a member of ${\mathcal F}_{A_{j},B_{j}}^{\rho_{j}}$ and let the second player consider the following $2$~cases.
%
%
%\\
%\\
%
%
%{\bf Case where $j$ is right-fresh in the position $((s_{1},\ldots,s_{q^{\prime}{-}1}),(s^{\prime}_{1},\ldots,s^{\prime}_{q^{\prime}{-}1}))$:}
%As a consequence of the construction of $J$, there exists $i{\in}I$ such that $i{\simeq}j$ and $i$ is left-fresh in the position $((s_{1},\ldots,s_{q^{\prime}{-}1}),(s^{\prime}_{1},\ldots,s^{\prime}_{q^{\prime}{-}1}))$.
%Thus, let the second player pick in its $q^{\prime}$-th move the member $g_{i,j}^{{-}1}(s^{\prime})$ of $(W,R)$.
%
%
%\\
%\\
%
%
%{\bf Case where $j$ is not right-fresh in the position $((s_{1},\ldots,s_{q^{\prime}{-}1}),(s^{\prime}_{1},\ldots,
%
%
%
%
%
%
%
%
%
%
%
%
%$\linebreak$
%
%
%
%
%
%
%
%
%
%
%
%
%s^{\prime}_{q^{\prime}{-}1}))$:}
%As a consequence of the construction of $J$, there exists $q^{\prime\prime}{\in}\lsem1,q^{\prime}{-}1\rsem$ such that $s^{\prime}_{q^{\prime\prime}}$ is a member of ${\mathcal F}_{A_{j},B_{j}}^{\rho_{j}}$.
%Thus, let $i{\in}I$ be such that $i{\simeq}j$ and $s_{q^{\prime\prime}}$ is a member of ${\mathcal F}_{A_{i},B_{i}}^{\rho_{i}}$.
%Thus, let the second player pick in its $q^{\prime}$-th move the member $g_{i,j}^{{-}1}(s^{\prime})$ of $(W,R)$.
%
%
\begin{proposition}\label{lemma:good:properties:of:delta:reductions:indexed}
%
%
%Let $Q,K{\geq}1$.
%Let $\{{\mathcal F}_{A_{i},B_{i}}^{\rho_{i}}:\ i{\in}I\}$ be a disjoint indexed family of galaxies such that for all $i{\in}I$, either ${\mathcal F}_{A_{i},B_{i}}^{\rho_{i}}$ is simple, ${\parallel}A_{i}{\parallel}{\leq}Q$ and ${\parallel}B_{i}{\parallel}{\leq}Q$, or ${\mathcal F}_{A_{i},B_{i}}^{\rho_{i}}$ is non-simple, ${\parallel}A_{i}{\parallel}{\leq}K$ and ${\parallel}B_{i}{\parallel}{\leq}K$.
%Let $q{\geq}3$.
%Let $\{{\mathcal F}_{A_{i},B_{i}}^{\rho_{i}}:\ i{\in}I\}$ be a disjoint indexed family of galaxies.
%If $J$ is a $q$-$\delta$-reduction of $\{{\mathcal F}_{A_{i},B_{i}}^{\rho_{i}}:\ i{\in}I\}$, $(W,R)$ is the union of $\{{\mathcal F}_{A_{i},B_{i}}^{\rho_{i}}:\ i{\in}I\}$ and $(W^{\prime},R^{\prime})$ is the union of $\{{\mathcal F}_{A_{j},B_{j}}^{\rho_{j}}:\ j{\in}J\}$ then
%
%
If $(W,R)$ is the union of $\{{\mathcal F}_{A_{i},B_{i}}^{\rho_{i}}:\ i{\in}I\}$ and $(W^{\prime},R^{\prime})$ is the union of $\{{\mathcal F}_{A_{j},B_{j}}^{\rho_{j}}:\ j{\in}J\}$ then
%
%
\begin{enumerate}
%
%
\item If $\{{\mathcal F}_{A_{i},B_{i}}^{\rho_{i}}:\ i{\in}I\}$ is dust-boun\-ded (respectively root-boun\-ded, kernel-boun\-ded) then $\{{\mathcal F}_{A_{j},B_{j}}^{\rho_{j}}:\ j{\in}J\}$ is dust-boun\-ded (respectively root-boun\-ded, kernel-boun\-ded,
%
%
%\item ${\parallel}J{\parallel}{\leq}N_{Q,K}$,
%
%
\item for all $Q{\geq}1$, if for all $i{\in}I$, ${\parallel}A_{i}{\parallel}{\leq}Q$ and ${\parallel}B_{i}{\parallel}{\leq}Q$ then ${\parallel}J{\parallel}{\leq}q{\times}(Q{+}1)^{2}{\times}
%
%
%
%
%
%
%
%
%
%
%
%
$\linebreak$
%
%
%
%
%
%
%
%
%
%
%
%
2^{Q^{2}}$,
%
%
\item $(W^{\prime},R^{\prime})$ is a bounded morphic image of $(W,R)$,
%
%
\item the second player has a winning strategy in the Ehrenfeucht-Fra\"\i ss\'e game ${\mathcal G}_{q}((W,R),(W^{\prime},R^{\prime}))$.
%
%
\end{enumerate}
%
%
\end{proposition}
%
%
\begin{proof}
%
%
$\mathbf{(1)}$ and $\mathbf{(2)}$~By construction of $\{{\mathcal F}_{A_{j},B_{j}}^{\rho_{j}}:\ j{\in}J\}$.
%
%
\\
\\
%
%
$\mathbf{(3)}$~Let $h:\ I{\longrightarrow}J$ be a function such that for all $i{\in}I$, either $i{\in}J$ and $h(i){=}i$, or $i{\not\in}J$ and $h(i){\simeq}i$.
Notice that the existence of such a function is a consequence of the definition of $J$.
Notice also that $h$ is surjective.
Moreover, for all $i{\in}I$, $h(i){\simeq}i$.
For all $i{\in}I$, let $g_{i}$ be an isomorphism from ${\mathcal F}_{A_{i},B_{i}}^{\rho_{i}}$ to ${\mathcal F}_{A_{h(i)},B_{h(i)}}^{\rho_{h(i)}}$.
Let $f:\ W{\longrightarrow}W^{\prime}$ be the function such that for all $i{\in}I$ and for all $s{\in}A_{i}{\cup}B_{i}$, $f(s){=}g_{i}(s)$.
Notice that $f$ is a bounded morphism from $(W,R)$ to $(W^{\prime},R^{\prime})$.
Moreover, $f$ is surjective.
%
%
\\
\\
%
%
$\mathbf{(4)}$~For all $i{\in}I$ and for all $j{\in}J$, if $i{\simeq}j$ then let $f_{i,j}$ be an isomorphism from ${\mathcal F}_{A_{i},B_{i}}^{\rho_{i}}$ to ${\mathcal F}_{A_{j},B_{j}}^{\rho_{j}}$.
Let $q^{\prime}{\in}\lsem1,q\rsem$ and $((s_{1},\ldots,s_{q^{\prime}{-}1}),(s^{\prime}_{1},\ldots,s^{\prime}_{q^{\prime}{-}1}))$ be a winning position of the Ehrenfeucht-Fra\"\i ss\'e game ${\mathcal G}_{q}((W,R),(W^{\prime},R^{\prime}))$.
Suppose the first player picked a member $s$ of $(W,R)$ in its $q^{\prime}$th move.
Let $i{\in}I$ be such that $s$ is a member of ${\mathcal F}_{A_{i},B_{i}}^{\rho_{i}}$, $i_{1},\ldots,i_{q^{\prime}{-}1}{\in}I$ be such that $s_{1}$ is a member of ${\mathcal F}_{A_{i_{1}},B_{i_{1}}}^{\rho_{i_{1}}}$, $\ldots$, $s_{q^{\prime}{-}1}$ is a member of ${\mathcal F}_{A_{i_{q^{\prime}{-}1}},B_{i_{q^{\prime}{-}1}}}^{\rho_{i_{q^{\prime}{-}1}}}$ and $j_{1},\ldots,j_{q^{\prime}{-}1}{\in}J$ be such that $s^{\prime}_{1}$ is a member of ${\mathcal F}_{A_{j_{1}},B_{j_{1}}}^{\rho_{j_{1}}}$, $\ldots$, $s^{\prime}_{q^{\prime}{-}1}$ is a member of ${\mathcal F}_{A_{j_{q^{\prime}{-}1}},B_{j_{q^{\prime}{-}1}}}^{\rho_{j_{q^{\prime}{-}1}}}$.
Then
%
%
\begin{itemize}
%
%
\item if $i{=}i_{q^{\prime\prime}}$ for some $q^{\prime\prime}{\in}\lsem1,q^{\prime}{-}1\rsem$ then the second player picks $f_{i_{q^{\prime\prime}},j_{q^{\prime\prime}}}(s)$ in its $q^{\prime}$th move,
%
%
\item otherwise, $j$ being an element in $\lbrack i\rbrack_{\simeq}{\cap}(J{\setminus}\{j_{1},\ldots,j_{q^{\prime}{-}1}\})$, the second play\-er picks $f_{i,j}(s)$ in its $q^{\prime}$th move,
%
%
\end{itemize}
%
%
Suppose the first player picked a member $s^{\prime}$ of $(W^{\prime},R^{\prime})$ in its $q^{\prime}$th move.
Let $j{\in}J$ be such that $s^{\prime}$ is a member of ${\mathcal F}_{A_{j},B_{j}}^{\rho_{j}}$, $i_{1},\ldots,i_{q^{\prime}{-}1}{\in}I$ be such that $s_{1}$ is a member of ${\mathcal F}_{A_{i_{1}},B_{i_{1}}}^{\rho_{i_{1}}}$, $\ldots$, $s_{q^{\prime}{-}1}$ is a member of ${\mathcal F}_{A_{i_{q^{\prime}{-}1}},B_{i_{q^{\prime}{-}1}}}^{\rho_{i_{q^{\prime}{-}1}}}$ and $j_{1},\ldots,j_{q^{\prime}{-}1}{\in}J$ be such that $s^{\prime}_{1}$ is a member of ${\mathcal F}_{A_{j_{1}},B_{j_{1}}}^{\rho_{j_{1}}}$, $\ldots$, $s^{\prime}_{q^{\prime}{-}1}$ is a member of ${\mathcal F}_{A_{j_{q^{\prime}{-}1}},B_{j_{q^{\prime}{-}1}}}^{\rho_{j_{q^{\prime}{-}1}}}$.
Then
%
%
\begin{itemize}
%
%
\item if $j{=}j_{q^{\prime\prime}}$ for some $q^{\prime\prime}{\in}\lsem1,q^{\prime}{-}1\rsem$ then the second player picks $f^{{-}1}_{i_{q^{\prime\prime}},j_{q^{\prime\prime}}}(s^{\prime})$ in its $q^{\prime}$th move,
%
%
\item otherwise, $i$ being an element in $\lbrack j\rbrack_{\simeq}{\cap}(I{\setminus}\{i_{1},\ldots,i_{q^{\prime}{-}1}\})$, the second play\-er picks $f^{{-}1}_{i,j}(s)$ in its $q^{\prime}$th move,
%
%
\end{itemize}
%
%
Obviously, $((s_{1},\ldots,s_{q^{\prime}{-}1},s),(s^{\prime}_{1},\ldots,s^{\prime}_{q^{\prime}{-}1},s^{\prime}))$ is a winning position of the Eh\-renfeucht-Fra\"\i ss\'e game ${\mathcal G}_{q}((W,R),(W^{\prime},R^{\prime}))$.
%
%
\medskip
%
%
\end{proof}
%
%
%
%
\section{Propositional Modal Logic}
%
%
In this section, we introduce tools and techniques in Propositional Modal Logic that we will need later in our proofs about modal definability.
%
%
%
%
\paragraph{Modal syntax}
%
%
Let us consider a countably infinite set $\PVAR$ of {\it propositional variables}\/ (denoted $\mathbf{p}$, $\mathbf{q}$, $\ldots$).
%
%
\\
\\
%
%
The {\it modal formulas}\/ (denoted $\varphi$, $\psi$, $\ldots$) are inductively defined as follows:
%
%
\begin{itemize}
%
%
\item $\varphi,\psi::=\mathbf{p}\mid\bot\mid\neg\varphi\mid(\varphi\vee\psi)\mid\Box\varphi$.
%
%
\end{itemize}
%
%
We adopt the standard rules for omission of the parentheses.
%
%
\\
\\
%
%
We define the other Boolean constructs as usual.
The modal formula $\Diamond\varphi$ is obtained as the well-known abbreviation
%
%
\begin{itemize}
%
%
\item $\Diamond\varphi::=\neg\Box\neg\varphi$.
%
%
\end{itemize}
%
%
We will also use the abbreviations
%
%
\begin{itemize}
%
%
\item $\lbrack U\rbrack\varphi::=\varphi\wedge\Box\Box\varphi$,
%
%
\item $\langle U\rangle\varphi::=\neg\lbrack U\rbrack\neg\varphi$.
%
%
\end{itemize}
%
%
The {\it set of all propositional variables occurring in the modal formula $\varphi$}\/ (denoted $\vvar(\varphi)$) is defined as usual.
%
%
\\
\\
%
%
For all modal formulas $\varphi$, let $\size(\varphi)$ be the number of symbols occurring in $\varphi$.
%
%
%
%
\paragraph{Modal semantics}
%
%
Let $(W,R)$ be a frame.
%
%
\\
\\
%
%
A {\it valuation on $(W,R)$}\/ is a function $V:\ \PVAR{\longrightarrow}{\cal P}(W)$.
%
%
\\
\\
%
%
The {\it satisfiability of a modal formula $\varphi$ at $s{\in}W$ with respect to a valuation $V$ on $(W,R)$}\/ (denoted $(W,R),V,s{\models}\varphi$) is inductively defined as follows:
%
%
\begin{itemize}
%
%
\item $(W,R),V,s{\models}\mathbf{p}$ if and only if $s{\in}V(\mathbf{p})$,
%
%
\item $(W,R),V,s{\not\models}\bot$,
%
%
\item $(W,R),V,s{\models}\neg\varphi$ if and only if $(W,R),V,s{\not\models}\varphi$,
%
%
\item $(W,R),V,s{\models}\varphi\vee\psi$ if and only if $(W,R),V,s{\models}\varphi$ or $(W,R),V,s{\models}\psi$,
%
%
\item $(W,R),V,s{\models}\Box\varphi$ if and only if for all $t{\in}W$, if $s{R}t$ then $(W,R),V,t{\models}\varphi$.
%
%
\end{itemize}
%
%
As a result,
%
%
\begin{itemize}
%
%
\item $(W,R),V,s{\models}\Diamond\varphi$ if and only if there exists $t{\in}W$ such that $s{R}t$ and $(W,R),
%
%
%
%
%
%
%
%
%
%
%
%
$\linebreak$
%
%
%
%
%
%
%
%
%
%
%
%
V,t{\models}\varphi$.
%
%
\end{itemize}
%
%
Moreover, if $(W,R)$ is Euclidean then
%
%
\begin{itemize}
%
%
\item $(W,R),V,s{\models}\lbrack U\rbrack\varphi$ if and only if for all $t{\in}W_{s}$, $(W,R),V,t{\models}\varphi$,
%
%
\item $(W,R),V,s{\models}\langle U\rangle\varphi$ if and only if there exists $t{\in}W_{s}$ such that $(W,R),V,t{\models}
%
%
%
%
%
%
%
%
%
%
%
%
$\linebreak$
%
%
%
%
%
%
%
%
%
%
%
%
\varphi$.
%
%
\end{itemize}
%
%
A modal formula $\varphi$ is {\it true with respect to a valuation $V$ on $(W,R)$}\/ (denoted $(W,R),V{\models}\varphi$) if $\varphi$ is satisfied at all $s{\in}W$ with respect to $V$.
A set $\Sigma$ of modal formulas is {\it true with respect to a valuation $V$ on $(W,R)$}\/ (denoted $(W,R),V{\models}\Sigma$) if for all modal formulas $\varphi$ in $\Sigma$, $(W,R),V{\models}\varphi$.
A modal formula $\varphi$ is {\it valid in $(W,R)$}\/ (denoted $(W,R){\models}\varphi$) if $\varphi$ is true with respect to all valuations on $(W,R)$.
A set $\Sigma$ of modal formulas is {\it valid in $(W,R)$}\/ (denoted $(W,R){\models}\Sigma$) if for all modal formulas $\varphi$ in $\Sigma$, $(W,R){\models}\varphi$.
%For all classes ${\mathcal C}$ of frames and for all modal formulas $\varphi$, let ${\parallel}\varphi{\parallel}_{{\mathcal C}}$ be the class of all frames in ${\mathcal C}$ validating $\varphi$.
%
%
\begin{lemma}\label{lemma:validity:in:a:given:frame:is:in:coNP}
%
%
The following decision problem is in $\coNP$:
%
%
\begin{description}
%
%
\item[input:] a modal formula $\varphi$ and a finite frame $(W^{\prime\prime},R^{\prime\prime})$,
%
%
\item[output:] determine whether $(W^{\prime\prime},R^{\prime\prime}){\models}\varphi$.
%
%
\end{description}
%
%
\end{lemma}
%
%
\begin{proof}
%
%
See~\cite{Arnold:Crubille:1988,Clarke:et:al:1986}.
%
%
\medskip
%
%
\end{proof}
%
%
\begin{lemma}\label{lemma:about:upper:bound:on:the:kernel:1}
%
%
Let $\varphi$ be a modal formula.
If there exists $m{\in}\mathbf{N}^{+}$ such that ${\mathcal F}_{m}^{2}{\not\models}\varphi$ then there exists $m{\in}\lsem1,2^{{\parallel}\vvar(\varphi){\parallel}}\rsem$ such that ${\mathcal F}_{m}^{2}{\not\models}\varphi$.
%
%
\end{lemma}
%
%
\begin{lemma}\label{lemma:about:upper:bound:on:the:kernel:2}
%
%
Let $\varphi$ be a modal formula.
If there exists $n{\in}\mathbf{N}^{-}$ such that ${\mathcal F}_{2}^{n}{\not\models}\varphi$ then there exists $n{\in}\lsem{-}1,2^{{\parallel}\vvar(\varphi){\parallel}}\rsem$ such that ${\mathcal F}_{2}^{n}{\not\models}\varphi$.
%
%
\end{lemma}
%
%
A frame $(W^{\prime},R^{\prime})$ is {\it modally weaker than $(W,R)$}\/ (denoted $(W^{\prime},R^{\prime}){\preceq}(W,R)$) if for all modal formulas $\varphi$, if $(W^{\prime},R^{\prime}){\models}\varphi$ then $(W,R){\models}\varphi$.
%
%
\begin{lemma}[Generated Subframe Lemma]\label{Generated:Subframe:Lemma}
%
%
For all frames $(W^{\prime},R^{\prime})$, if $(W,
%
%
%
%
%
%
%
%
%
%
%
%
$\linebreak$
%
%
%
%
%
%
%
%
%
%
%
%
R)$ is a generated subframe of $(W^{\prime},R^{\prime})$ then $(W^{\prime},R^{\prime}){\preceq}(W,R)$.
%
%
\end{lemma}
%
%
\begin{proof}
%
%
See~\cite[Chapter~$3$]{Blackburn:deRijke:Venema:2001}.
%
%
\medskip
%
%
\end{proof}
%
%
\begin{lemma}[Bounded Morphism Lemma]\label{Bounded:Morphism:Lemma}
%
%
For all frames $(W^{\prime},R^{\prime})$, if $(W,
%
%
%
%
%
%
%
%
%
%
%
%
$\linebreak$
%
%
%
%
%
%
%
%
%
%
%
%
R)$ is a bounded morphic image of $(W^{\prime},R^{\prime})$ then $(W^{\prime},R^{\prime}){\preceq}(W,R)$.
%
%
\end{lemma}
%
%
\begin{proof}
%
%
See~\cite[Chapter~$3$]{Blackburn:deRijke:Venema:2001}.
%
%
\medskip
%
%
\end{proof}
%
%
\begin{lemma}\label{lemma:sigma:about:K2:galaxies:and:generated:by:s:subframes}
%
%
Let ${\mathcal F}_{A,B}^{\rho}$ be a headed galaxy.
For all modal formulas $\varphi$, if ${\mathcal F}_{A,B}^{\rho}{\not\models}\varphi$ then there exists $s{\in}A$ such that $(W_{s},R_{s}){\not\models}\varphi$ where $(W_{s},R_{s})$ is the least generated subframe of ${\mathcal F}_{A,B}^{\rho}$ containing $s$.
%
%
\end{lemma}
%
%
%\begin{lemma}\label{lemma:tau:about:L2:galaxies:and:generated:by:s:subframes}
%
%
%Let ${\mathcal F}_{A,B}^{\rho}$ be a galaxy in ${\mathcal L}_{2}$.
%For all modal formulas $\varphi$, if ${\mathcal F}_{A,B}^{\rho}{\not\models}\varphi$ then there exists $s{\in}A$ such that $(W_{s},R_{s}){\not\models}\varphi$ where $(W_{s},R_{s})$ is the least generated subframe of ${\mathcal F}_{A,B}^{\rho}$ containing $s$.
%
%
%\end{lemma}
%
%
\begin{lemma}\label{lemma:about:finite:subset:when:K5:shape:of:the:frame}
%
%
If $W{=}\{s\}{\cup}C$ and $R{=}(\{s\}{\times}B){\cup}(C{\times}C)$ for some sets $B,C$ and for some element $s$ such that $B{\not=}\emptyset$, $C{\setminus}B{\not=}\emptyset$, $B$ is finite, $B{\subseteq}C$ and $s{\not\in}C$ then for all modal formulas $\varphi$, if $(W,R){\not\models}\varphi$ then there exists a flower ${\mathcal F}_{m}^{n}$ such that $m{=}{\parallel}B{\parallel}$ and ${\mathcal F}_{m}^{n}{\not\models}\varphi$.
%
%
\end{lemma}
%
%
\begin{lemma}\label{lemma:about:finite:subset:when:K5:shape:of:the:frame:bis}
%
%
If $W{=}\{s\}{\cup}C$ and $R{=}(\{s\}{\times}B){\cup}(C{\times}C)$ for some sets $B,C$ and for some element $s$ such that $B{\not=}\emptyset$, $C{\setminus}B{\not=}\emptyset$, $C{\setminus}B$ is finite, $B{\subseteq}C$ and $s{\not\in}C$ then for all modal formulas $\varphi$, if $(W,R){\not\models}\varphi$ then there exists a flower ${\mathcal F}_{m}^{n}$ such that $n{=}{\parallel}C{\setminus}B{\parallel}$ and ${\mathcal F}_{m}^{n}{\not\models}\varphi$.
%
%
\end{lemma}
%
%
Let $q{\geq}3$.
%
%
\\
\\
%
%
Let $\{{\mathcal F}_{A_{i},B_{i}}^{\rho_{i}}:\ i{\in}I\}$ be a disjoint indexed family of galaxies.
%
%
\begin{lemma}\label{lemma:alpha:reduction:modal:logic}
%
%
%We shall say that $\{(W_{A_{i}^{\prime},B_{i}^{\prime}}^{\rho_{i}^{\prime}},R_{A_{i}^{\prime},B_{i}^{\prime}}^{\rho_{i}^{\prime}}):\ i{\in}I\}$ is {\it $q$-$\alpha$-reduced from $\{(W_{A_{i},B_{i}}^{\rho_{i}},R_{A_{i},B_{i}}^{\rho_{i}}):\ i{\in}I\}$}\/ if for all $i{\in}I$, either $(W_{A_{i}^{\prime},B_{i}^{\prime}}^{\rho_{i}^{\prime}},R_{A_{i}^{\prime},B_{i}^{\prime}}^{\rho_{i}^{\prime}})$ is equal to $(W_{A_{i},B_{i}}^{\rho_{i}},R_{A_{i},B_{i}}^{\rho_{i}})$, or $(W_{A_{i}^{\prime},B_{i}^{\prime}}^{\rho_{i}^{\prime}},R_{A_{i}^{\prime},B_{i}^{\prime}}^{\rho_{i}^{\prime}})$ is a $q$-$\alpha$-reduction of $(W_{A_{i},B_{i}}^{\rho_{i}},R_{A_{i},B_{i}}^{\rho_{i}})$.
If $\{{\mathcal F}_{A_{i}^{\prime},B_{i}^{\prime}}^{\rho_{i}^{\prime}}:\ i{\in}I\}$ is a $q$-$\alpha$-reduction of $\{{\mathcal F}_{A_{i},B_{i}}^{\rho_{i}}:\ i{\in}I\}$, $(W,R)$ is the union of $\{{\mathcal F}_{A_{i},B_{i}}^{\rho_{i}}:\ i{\in}I\}$ and $(W^{\prime},R^{\prime})$ is the union of $\{{\mathcal F}_{A_{i}^{\prime},B_{i}^{\prime}}^{\rho_{i}^{\prime}}:\ i{\in}I\}$ then $(W,R){\preceq}(W^{\prime},R^{\prime})$.
%
%
\end{lemma}
%
%
\begin{proof}
%
%
By Proposition~\ref{lemma:good:properties:of:alpha:reductions:indexed} and Lemma~\ref{Bounded:Morphism:Lemma}.
%
%
\medskip
%
%
\end{proof}
%
%
\begin{lemma}\label{lemma:gamma:reduction:modal:logic}
%
%
%Let $q{\geq}3$.
%Let $\{{\mathcal F}_{A_{i},B_{i}}^{\rho_{i}}:\ i{\in}I\}$ be a disjoint indexed family of galaxies.
%We shall say that $\{(W_{A_{i}^{\prime},B_{i}^{\prime}}^{\rho_{i}^{\prime}},R_{A_{i}^{\prime},B_{i}^{\prime}}^{\rho_{i}^{\prime}}):\ i{\in}I\}$ is {\it $q$-$\alpha$-reduced from $\{(W_{A_{i},B_{i}}^{\rho_{i}},R_{A_{i},B_{i}}^{\rho_{i}}):\ i{\in}I\}$}\/ if for all $i{\in}I$, either $(W_{A_{i}^{\prime},B_{i}^{\prime}}^{\rho_{i}^{\prime}},R_{A_{i}^{\prime},B_{i}^{\prime}}^{\rho_{i}^{\prime}})$ is equal to $(W_{A_{i},B_{i}}^{\rho_{i}},R_{A_{i},B_{i}}^{\rho_{i}})$, or $(W_{A_{i}^{\prime},B_{i}^{\prime}}^{\rho_{i}^{\prime}},R_{A_{i}^{\prime},B_{i}^{\prime}}^{\rho_{i}^{\prime}})$ is a $q$-$\alpha$-reduction of $(W_{A_{i},B_{i}}^{\rho_{i}},R_{A_{i},B_{i}}^{\rho_{i}})$.
If $\{{\mathcal F}_{A_{i}^{\prime},B_{i}^{\prime}}^{\rho_{i}^{\prime}}:\ i{\in}I\}$ is a $q$-$\gamma$-reduction of $\{{\mathcal F}_{A_{i},B_{i}}^{\rho_{i}}:\ i{\in}I\}$, $(W,R)$ is the union of $\{{\mathcal F}_{A_{i},B_{i}}^{\rho_{i}}:\ i{\in}I\}$ and $(W^{\prime},R^{\prime})$ is the union of $\{{\mathcal F}_{A_{i}^{\prime},B_{i}^{\prime}}^{\rho_{i}^{\prime}}:\ i{\in}I\}$ then
$(W,R){\preceq}(W^{\prime},R^{\prime})$.
%
%
\end{lemma}
%
%
\begin{proof}
%
%
By Proposition~\ref{lemma:good:properties:of:gamma:reductions:indexed} and Lemma~\ref{Bounded:Morphism:Lemma}.
%
%
\medskip
%
%
\end{proof}
%
%
\begin{lemma}\label{lemma:delta:reduction:modal:logic}
%
%
%Let $q{\geq}3$.
%Let $\{{\mathcal F}_{A_{i},B_{i}}^{\rho_{i}}:\ i{\in}I\}$ be a disjoint indexed family of galaxies.
If $\{{\mathcal F}_{A_{j},B_{j}}^{\rho_{j}}:\ j{\in}J\}$ is a $q$-$\delta$-reduction of $\{{\mathcal F}_{A_{i},B_{i}}^{\rho_{i}}:\ i{\in}I\}$, $(W,R)$ is the union of $\{{\mathcal F}_{A_{i},B_{i}}^{\rho_{i}}:\ i{\in}I\}$ and $(W^{\prime},R^{\prime})$ is the union of $\{{\mathcal F}_{A_{j},B_{j}}^{\rho_{j}}:\ j{\in}J\}$ then $(W,R){\preceq}(W^{\prime},R^{\prime})$.
%
%
\end{lemma}
%
%
\begin{proof}
%
%
By Proposition~\ref{lemma:good:properties:of:delta:reductions:indexed} and Lemma~\ref{Bounded:Morphism:Lemma}.
%
%
\medskip
%
%
\end{proof}
%
%
%
%
\paragraph{Jankov-Fine formulas}
%
%
Now, let us adapt to the isolated point and flowers, the standard concept of Jankov-Fine formulas.
%
%
\\
\\
%
%
Let $m{\in}\mathbf{N}^{+}{\cup}\{0\}$ and $n{\in}\mathbf{N}^{-}$ be such that either $m{\not=}0$, or $n{=}0$.
%
%
\\
\\
%
%
Let $\chi_{m}^{n}$ be the modal formula defined as follows:
%
%
\begin{itemize}
%
%
\item the negation of modal formula
%
%
\begin{itemize}
%
%
\item $\Box\bot$,
%
%
\end{itemize}
%
%
when $m{=}0$,
%
%
\item the negation of the conjunction of the modal formulas
%
%
\begin{itemize}
%
%
\item $p_{0}$,
%
%
\item $\lbrack U\rbrack\bigvee\{p_{k}:\ 0{\leq}k{\leq}m{+}n\}$,
%
%
\item $\lbrack U\rbrack(p_{k}\rightarrow\neg p_{l})$ for each $k,l{\in}\N$ such that $0{\leq}k{<}l{\leq}m{+}n$,
%
%
\item $\langle U\rangle p_{k}$ for each $k{\in}\N$ such that $0{\leq}k{\leq}m{+}n$,
%
%
\item $\lbrack U\rbrack(p_{0}\rightarrow\Box\neg p_{0})$,
%
%
\item $\lbrack U\rbrack(p_{0}\rightarrow\Diamond p_{k})$ for each $k{\in}\N$ such that $1{\leq}k{\leq}m$,
%
%
\item $\lbrack U\rbrack(p_{0}\rightarrow\Box\neg p_{k})$ for each $k{\in}\N$ such that $m{+}1{\leq}k{\leq}m{+}n$,
%
%
\item $\lbrack U\rbrack(p_{k}\rightarrow\Box\neg p_{0})$ for each $k{\in}\N$ such that $1{\leq}k{\leq}m{+}n$,
%
%
\item $\lbrack U\rbrack(p_{k}\rightarrow\Diamond p_{l})$ for each $k,l{\in}\N$ such that $1{\leq}k,l{\leq}m{+}n$,
%
%
\end{itemize}
%
%
when $m{\not=}0$ and $n{\in}\N$,
%
%
\item the negation of the conjunction of the modal formulas
%
%
\begin{itemize}
%
%
\item $p_{1}$,
%
%
\item $\lbrack U\rbrack\bigvee\{p_{k}:\ 1{\leq}k{\leq}m\}$,
%
%
\item $\lbrack U\rbrack(p_{k}\rightarrow\neg p_{l})$ for each $k,l{\in}\N$ such that $1{\leq}k{<}l{\leq}m$,
%
%
\item $\langle U\rangle p_{k}$ for each $k{\in}\N$ such that $1{\leq}k{\leq}m$,
%
%
\item $\lbrack U\rbrack(p_{k}\rightarrow\Diamond p_{l})$ for each $k,l{\in}\N$ such that $1{\leq}k,l{\leq}m$,
%
%
\end{itemize}
%
%
when $m{\not=}0$ and $n{=}{-}1$.
%
%
\end{itemize}
%
%
The modal formula $\chi_{m}^{n}$ is called {\it Jankov-Fine formula of the frame ${\mathcal F}_{m}^{n}$.}
%
%
\begin{lemma}\label{lemma:Jankov:Fine:ONE}
%
%
For all Euclidean frames $(W,R)$ and for all $s{\in}W$, the following conditions are equivalent:
%
%
\begin{itemize}
%
%
\item there exists a valuation $V$ on $(W,R)$ such that $(W,R),V,s{\not\models}\chi_{m}^{n}$,
%
%
\item there exists a surjective bounded morphism $f$ from $(W_{s},R_{s})$ to ${\mathcal F}_{m}^{n}$ such that $f(s)$ generates ${\mathcal F}_{m}^{n}$.
%
%
\end{itemize}
%
%
\end{lemma}
%
%
\begin{proof}
%
%
See~\cite[Chapter~$3$]{Blackburn:deRijke:Venema:2001}.
%
%
\medskip
%
%
\end{proof}
%
%
\begin{lemma}\label{lemma:Jankov:Fine:TWO}
%
%
${\mathcal F}_{m}^{n}$ does not validate its Jankov-Fine formula.
%
%
\end{lemma}
%
%
\begin{proof}
%
%
By Lemma~\ref{lemma:Jankov:Fine:ONE}.
%
%
\medskip
%
%
\end{proof}
%
%
%
%
\paragraph{Normal modal logics}
%
%
A {\it normal modal logic}\/ is a set $\L$ of formulas such that
%
%
\begin{itemize}
%
%
\item $\L$ is closed under the rule of uniform substitution,
%
%
\item $\L$ contains all propositional tautologies,
%
%
\item for all modal formulas $\varphi,\psi$, $\Box(\varphi\rightarrow\psi)\rightarrow(\Box\varphi\rightarrow\Box\psi){\in}\L$,
%
%
\item for all modal formulas $\varphi,\psi$, if $\varphi{\in}\L$ and $\varphi\rightarrow\psi{\in}\L$ then $\psi{\in}\L$,
%
%
\item for all modal formulas $\varphi$, if $\varphi{\in}\L$ then $\Box\varphi{\in}\L$.
%
%
\end{itemize}
%
%
From now on, {\bf when we write ``modal logic'', we mean ``normal modal logic''.}
%
%
\\
\\
%
%
A modal logic $\L$ is {\it consistent}\/ if $\bot{\not\in}\L$.
%
%
\\
\\
%
%
A consistent modal logic $\L$ is {\it Euclidean}\/ if for all modal formulas $\varphi$, $\Diamond\varphi\rightarrow\Box\Diamond\varphi{\in}\L$~---~i.e. $\L$ is an extension of the modal logic $\K5$.
%
%
\\
\\
%
%
The extensions of the modal logic $\K5$ have been studied by Nagle~\cite{Nagle:1981} and Nagle and Thomason~\cite{Nagle:Thomason:1985}.
They possess remarkable properties: they have the poly-size model property, they are finitely axiomatizable, etc.
%
%
\\
\\
%
%
For all modal formulas $\varphi$, let $\K5{\oplus}\varphi$ be the least extension of $\K5$ containing $\varphi$.
%
%
\\
\\
%
%
For all Euclidean modal logics $\L$, let $\Fr(\L)$ be the class of all frames $(W,R)$ such that $(W,R){\models}\L$.
%
%
\\
\\
%
%
For all Euclidean modal logics $\L$, let
%
%
\begin{itemize}
%
%
\item $\mathtt{S}_{\L}{=}\{(m,n):\ m{\in}\mathbf{N}^{+},\ n{\in}\mathbf{N}^{-}$ and ${\mathcal F}_{m}^{n}{\models}\L\}$.
%
%
\end{itemize}
%
%
\begin{lemma}\label{lemma:S:L:is:closed:subset:of:N:plus:times:N:minus}
%
%
For all Euclidean modal logics $\L$, $\mathtt{S}_{\L}$ is a closed subset of $\mathbf{N}^{+}{\times}\mathbf{N}^{-}$.
%
%
\end{lemma}
%
%
\begin{proof}
%
%
By Lemmas~\ref{lemma:about:bounded:morphisms:in:the:situation:of F:m:n:frames} and~\ref{Bounded:Morphism:Lemma}.
%
%
\medskip
%
%
\end{proof}
%
%
\begin{lemma}\label{lemma:alpha}
%
%
Let ${\mathcal F}^{\rho}_{A,B}$ be a galaxy in ${\mathcal K}_{2}$.
If $\{2\}{\times}\mathbf{N}^{-}{\subseteq}\mathtt{S}_{\L}$ then ${\mathcal F}^{\rho}_{A,B}{\models}\L$.
%
%
\end{lemma}
%
%
\begin{proof}
%
%
For the sake of the contradiction, suppose $\{2\}{\times}\mathbf{N}^{-}{\subseteq}\mathtt{S}_{\L}$ and ${\mathcal F}^{\rho}_{A,B}{\not\models}\L$.
Hence, there exists a modal formula $\varphi{\in}\L$ such that ${\mathcal F}^{\rho}_{A,B}{\not\models}\varphi$.
Since ${\mathcal F}^{\rho}_{A,B}$ is a galaxy in ${\mathcal K}_{2}$, then ${\mathcal F}^{\rho}_{A,B}$ is headed and by Lemma~\ref{lemma:sigma:about:K2:galaxies:and:generated:by:s:subframes}, there exists $s{\in}A$ such that $(W_{s},R_{s}){\not\models}\varphi$ where $(W_{s},R_{s})$ is the least generated subframe of ${\mathcal F}_{A,B}^{\rho}$ containing $s$.
Moreover, obviously, $W_{s}{=}\{s\}{\cup}C^{\prime}$ and $R{=}(\{s\}{\times}B^{\prime}){\cup}(C^{\prime}{\times}C^{\prime})$ for some sets $B^{\prime},C^{\prime}$ such that ${\parallel}B^{\prime}{\parallel}{=}2$, ${\parallel}C^{\prime}{\setminus}B^{\prime}{\parallel}{\geq}2$, $B^{\prime}{\subseteq}C^{\prime}$ and $s{\not\in}C^{\prime}$.
Thus, by Lemma~\ref{lemma:about:finite:subset:when:K5:shape:of:the:frame}, there exists a flower ${\mathcal F}_{m}^{n}$ such that $m{=}2$ and ${\mathcal F}_{m}^{n}{\not\models}\varphi$.
Since $\{2\}{\times}\mathbf{N}^{-}{\subseteq}\mathtt{S}_{\L}$, then $(m,n){\in}\mathtt{S}_{\L}$.
Consequently, ${\mathcal F}_{m}^{n}{\models}\L$.
Since $\varphi{\in}\L$, then ${\mathcal F}_{m}^{n}{\models}\varphi$: a contradiction.
%
%
\medskip
%
%
\end{proof}
%
%
\begin{lemma}\label{lemma:beta}
%
%
Let ${\mathcal F}^{\rho}_{A,B}$ be a galaxy in ${\mathcal L}_{2}$.
If $\mathbf{N}^{+}{\times}\{2\}{\subseteq}\mathtt{S}_{\L}$ then ${\mathcal F}^{\rho}_{A,B}{\models}\L$.
%
%
\end{lemma}
%
%
\begin{proof}
%
%
For the sake of the contradiction, suppose $\mathbf{N}^{+}{\times}\{2\}{\subseteq}\mathtt{S}_{\L}$ and ${\mathcal F}^{\rho}_{A,B}{\not\models}\L$.
Hence, there exists a modal formula $\varphi{\in}\L$ such that ${\mathcal F}^{\rho}_{A,B}{\not\models}\varphi$.
Since ${\mathcal F}^{\rho}_{A,B}$ is a galaxy in ${\mathcal L}_{2}$, then ${\mathcal F}^{\rho}_{A,B}$ is headed and by Lemma~\ref{lemma:sigma:about:K2:galaxies:and:generated:by:s:subframes}, there exists $s{\in}A$ such that $(W_{s},R_{s}){\not\models}\varphi$ where $(W_{s},R_{s})$ is the least generated subframe of ${\mathcal F}_{A,B}^{\rho}$ containing $s$.
Moreover, obviously, $W_{s}{=}\{s\}{\cup}C^{\prime}$ and $R{=}(\{s\}{\times}B^{\prime}){\cup}(C^{\prime}{\times}C^{\prime})$ for some sets $B^{\prime},C^{\prime}$ such that ${\parallel}B^{\prime}{\parallel}{\geq}2$, ${\parallel}C^{\prime}{\setminus}B^{\prime}{\parallel}{=}2$, $B^{\prime}{\subseteq}C^{\prime}$ and $s{\not\in}C^{\prime}$.
Thus, by Lemma~\ref{lemma:about:finite:subset:when:K5:shape:of:the:frame:bis}, there exists a flower ${\mathcal F}_{m}^{n}$ such that $n{=}2$ and ${\mathcal F}_{m}^{n}{\not\models}\varphi$.
Since $\mathbf{N}^{+}{\times}\{2\}{\subseteq}\mathtt{S}_{\L}$, then $(m,n){\in}\mathtt{S}_{\L}$.
Consequently, ${\mathcal F}_{m}^{n}{\models}\L$.
Since $\varphi{\in}\L$, then ${\mathcal F}_{m}^{n}{\models}\varphi$: a contradiction.
%
%
\medskip
%
%
\end{proof}
%
%
%
%
\section{Model Theory of Classical First-Order Logic}\label{section:first:order:syntax:and:semantics}
%
%
In this section, we introduce tools and techniques in Model Theory of Classical First-Order Logic that we will need later in our proofs about modal definability.
%
%
%
%
\paragraph{First-order syntax}
%
%
Let us consider a countably infinite set $\IVAR$ of {\it individual variables}\/ (denoted $\mathbf{x}$, $\mathbf{y}$, $\ldots$).
%
%
\\
\\
%
%
The {\it first-order formulas}\/ (denoted $A$, $B$, $\ldots$) are inductively defined as follows:
%
%
\begin{itemize}
%
%
\item $A,B::=\mathbf{R}(\mathbf{x},\mathbf{y})\mid\mathbf{x}{=}\mathbf{y}\mid\neg A\mid(A\vee B)\mid\forall\mathbf{x}A$.
%
%
\end{itemize}
%
%
We adopt the standard rules for omission of the parentheses.
%
%
\\
\\
%
%
We define the other Boolean constructs as usual.
The first-order formula $\exists\mathbf{x}A$ is obtained as the well-known abbreviation
%
%
\begin{itemize}
%
%
\item $\exists\mathbf{x}A::=\neg\forall\mathbf{x}\neg A$.
%
%
\end{itemize}
%
%
We will also use the abbreviations
%
%
\begin{itemize}
%
%
\item $\exists^{{=}1}\mathbf{x}A(\mathbf{x})::=\exists\mathbf{x}(A(\mathbf{x})\wedge\forall\mathbf{y}(A(\mathbf{y})\rightarrow\mathbf{y}{=}\mathbf{x}))$,
%
%
\item $\exists^{{=}2}\mathbf{x}A(\mathbf{x})::=\exists\mathbf{x_{1}}\exists\mathbf{x_{2}}(A(\mathbf{x_{1}})\wedge A(\mathbf{x_{2}})\wedge\mathbf{x_{1}}{\not=}\mathbf{x_{2}}\wedge\forall\mathbf{y}(A(\mathbf{y})\rightarrow\mathbf{y}{=}\mathbf{x_{1}}\vee\mathbf{y}{=}\mathbf{x_{2}}))$,
%
%
\item $\{\mathbf{x_{1}},\mathbf{x_{2}}\}{\cap}\{\mathbf{y_{1}},\mathbf{y_{2}}\}{=}\emptyset::=\mathbf{x_{1}}{\not=}\mathbf{y_{1}}\wedge\mathbf{x_{1}}{\not=}\mathbf{y_{2}}\wedge\mathbf{x_{2}}{\not=}\mathbf{y_{1}}\wedge\mathbf{x_{2}}{\not=}\mathbf{y_{2}}$,
%
%
\item $\mathbf{R}(\mathbf{z}){=}\{\mathbf{x_{i}},\mathbf{y_{j}}\}::=\forall\mathbf{t}(\mathbf{R}(\mathbf{z},\mathbf{t})\leftrightarrow\mathbf{t}{=}\mathbf{x_{i}}\vee\mathbf{t}{=}\mathbf{y_{j}})$,
%
%
\item $\mathbf{R}(\mathbf{z}){=}\{\mathbf{x_{i}},\mathbf{y_{j}}\}^{\mathbf{c}}::=\forall\mathbf{t}(\mathbf{R}(\mathbf{z},\mathbf{t})\leftrightarrow\mathbf{R}(\mathbf{t},\mathbf{t})\wedge\mathbf{t}{\not=}\mathbf{x_{i}}\wedge\mathbf{t}{\not=}\mathbf{y_{j}})$,
%
%
\item $\{\mathbf{x_{1}},\mathbf{x_{2}}\}{=}\{\mathbf{y_{1}},\mathbf{y_{2}}\}::=(\mathbf{x_{1}}{=}\mathbf{y_{1}}\wedge\mathbf{x_{2}}{=}\mathbf{y_{2}})\vee(\mathbf{x_{1}}{=}\mathbf{y_{2}}\wedge\mathbf{x_{2}}{=}\mathbf{y_{1}})$.
%
%
\end{itemize}
%
%
The {\it set of all free individual variables occurring in the first-order formula $A$}\/ (denoted $\fiv(A)$) is defined as usual.
The first-order formula $A$ is a {\it sentence}\/ if $\fiv(A){=}\emptyset$.
%
%
\\
\\
%
%
For all first-order formulas $A$, let $\size(A)$ be the number of symbols occurring in $A$.
%
%
\\
\\
%
%
The {\it quantifier depth of the first-order formula $A$}\/ (denoted $\qd(A)$) is defined as usual.
For all first-order formulas $A$, let $\qdd(A){=}\max\{\qd(A),3\}$.
%
%
\\
\\
%
%
Obviously, for all first-order formulas $A$, $\qd(A){\leq}\qdd(A)$ and $\qdd(A){\geq}3$.
%
%
\\
\\
%
%
For all integers $k$ such that $k{\geq}4$ and for all first-order formulas $A$, let $\Pi_{k}^{A}{=}\{(m,
%
%
%
%
%
%
%
%
%
%
%
%
$\linebreak$
%
%
%
%
%
%
%
%
%
%
%
%
n):\ m{\in}\mathbf{N}^{+}$ and $n{\in}\mathbf{N}^{-}$ are such that if $m{\geq}2$ and $n{\geq}2$ then $m{+}n{\leq}k$, if $m{=}1$ then $n{\leq}\qdd(A)$ and if either $n{=}{-}1$, or $n{=}0$, or $n{=}1$ then $m{\leq}\qdd(A)\}$.
%
%
\begin{lemma}\label{lemma:ll:finitely:many:PI:PI:pairs}
%
%
For all integers $k$ such that $k{\geq}4$ and for all first-order formulas $A$, $\Pi_{k}^{A}{\subseteq}\lsem1,\max\{k,\qdd(A)\}\rsem{\times}\lsem{-}1,\max\{k,\qdd(A)\}\rsem$.
%
%
\end{lemma}
%
%
%
%
\paragraph{First-order semantics}
%
%
Let $(W,R)$ be a frame.
%
%
\\
\\
%
%
An {\it assignment on $(W,R)$}\/ is a function $g:\ \IVAR{\longrightarrow}W$.
The {\it update of an assignment $g$ on $(W,R)$ with respect to $s{\in}W$ and an individual variable $\mathbf{x}$}\/ is the assignment $g_{s}^{\mathbf{x}}$ on $(W,R)$ such that for all individual variables $\mathbf{y}$,
%
%
\begin{itemize}
%
%
\item $g_{s}^{\mathbf{x}}(\mathbf{y}){=}s$ when $\mathbf{y}{=}\mathbf{x}$,
%
%
\item $g_{s}^{\mathbf{x}}(\mathbf{y}){=}g(\mathbf{y})$ when $\mathbf{y}{\not=}\mathbf{x}$.
%
%
\end{itemize}
%
%
The {\it satisfiability of a first-order formula $A$ with respect to an assignment $g$ on $(W,R)$}\/ (denoted $(W,R),g\models A$) is inductively defined as follows:
%
%
\begin{itemize}
%
%
\item $(W,R),g{\models}\mathbf{R}(\mathbf{x},\mathbf{y})$ if and only if $g(\mathbf{x}){R}g(\mathbf{y})$,
%
%
\item $(W,R),g{\models}\mathbf{x}{=}\mathbf{y}$ if and only if $g(\mathbf{x}){=}g(\mathbf{y})$,
%
%
%\item $(W,R),g{\not\models}\bot$,
%
%
\item $(W,R),g{\models}\neg A$ if and only if $(W,R),g{\not\models}A$,
%
%
\item $(W,R),g{\models}A\vee B$ if and only if $(W,R),g{\models}A$ or $(W,R),g{\models}B$,
%
%
\item $(W,R),g{\models}\forall\mathbf{x}A$ if and only if for all $s{\in}W$, $(W,R),g_{s}^{\mathbf{x}}{\models}A$.
%
%
\end{itemize}
%
%
As a result,
%
%
\begin{itemize}
%
%
\item $(W,R),g{\models}\exists\mathbf{x}A$ if and only if there exists $s{\in}W$ such that $(W,R),g_{s}^{\mathbf{x}}{\models}A$.
%
%
\end{itemize}
%
%
Moreover,
%
%
\begin{itemize}
%
%
\item $(W,R),g{\models}\exists^{{=}1}\mathbf{x}A(\mathbf{x})$ if and only if there exists exactly one $s{\in}W$ such that $(W,R),g_{s}^{\mathbf{x}}{\models}A$,
%
%
\item $(W,R),g{\models}\exists^{{=}2}\mathbf{x}A(\mathbf{x})$ if and only if there exists exactly two $s{\in}W$ such that $(W,R),g_{s}^{\mathbf{x}}{\models}A$,
%
%
\item $(W,R),g{\models}\{\mathbf{x_{1}},\mathbf{x_{2}}\}{\cap}\{\mathbf{y_{1}},\mathbf{y_{2}}\}{=}\emptyset$ if and only if $\{g(\mathbf{x_{1}}),g(\mathbf{x_{2}})\}{\cap}\{g(\mathbf{y_{1}}),
%
%
%
%
%
%
%
%
%
%
%
%
$\linebreak$
%
%
%
%
%
%
%
%
%
%
%
%
g(\mathbf{y_{2}})\}{=}\emptyset$,
%
%
\item $(W,R),g{\models}\mathbf{R}(\mathbf{z}){=}\{\mathbf{x_{i}},\mathbf{y_{j}}\}$ if and only if $R(g(\mathbf{z})){=}\{g(\mathbf{x_{i}}),g(\mathbf{y_{j}})\}$,
%
%
\item $(W,R),g{\models}\mathbf{R}(\mathbf{z}){=}\{\mathbf{x_{i}},\mathbf{y_{j}}\}^{\mathbf{c}}$ if and only if $R(g(\mathbf{z})){=}\{t{\in}W:\ t{R}t\}{\setminus}\{g(\mathbf{x_{i}}),
%
%
%
%
%
%
%
%
%
%
%
%
$\linebreak$
%
%
%
%
%
%
%
%
%
%
%
%
g(\mathbf{y_{j}})\}$
%\mathbf{R}(\mathbf{z}){=}\{\mathbf{x_{i}},\mathbf{y_{j}}\}^{\mathbf{c}}
%
%
\item $(W,R),g{\models}\{\mathbf{x_{1}},\mathbf{x_{2}}\}{=}\{\mathbf{y_{1}},\mathbf{y_{2}}\}$ if and only if $\{g(\mathbf{x_{1}}),g(\mathbf{x_{2}})\}{=}\{g(\mathbf{y_{1}}),
%
%
%
%
%
%
%
%
%
%
%
%
$\linebreak$
%
%
%
%
%
%
%
%
%
%
%
%
g(\mathbf{y_{2}})\}$.
%
%
\end{itemize}
%
%
From now on, for all first-order formulas $A(\mathbf{x_{1}},\ldots,\mathbf{x_{m}})$ with free individual variables $\mathbf{x_{1}},\ldots,\mathbf{x_{m}}$ and for all $s_{1},\ldots,s_{m}\in W$, {\bf when we write ``$(W,R){\models}A(\mathbf{x_{1}},
%
%
%
%
%
%
%
%
%
%
%
%
$\linebreak$
%
%
%
%
%
%
%
%
%
%
%
%
\ldots,\mathbf{x_{m}})\ \lbrack s_{1},\ldots,s_{m}\rbrack$'', we shall mean ``for all assignments $g$ on $(W,R)$, if $g(\mathbf{x_{1}}){=}s_{1}$, $\ldots$, $g(\mathbf{x_{m}}){=}s_{m}$ then $(W,R),g{\models}A(\mathbf{x_{1}},\ldots,\mathbf{x_{m}})$''.}
%
%
\\
\\
%
%
A first-order formula $A$ is {\it valid in $(W,R)$}\/ (denoted $(W,R){\models}A$) if $A$ is satisfied with respect to all assignments on $(W,R)$.
A set $\Sigma$ of first-order formulas is {\it valid in $(W,R)$}\/ (denoted $(W,R){\models}\Sigma$) if for all first-order formulas $A$ in $\Sigma$, $(W,R){\models}A$.
%For all classes ${\mathcal C}$ of frames and for all sentences $A$, let ${\parallel}A{\parallel}_{{\mathcal C}}$ be the class of all frames in ${\mathcal C}$ validating $A$.
%
%
\begin{lemma}\label{lemma:validity:in:a:given:frame:is:in:coNP:fol}
%
%
The following decision problem is in $\PSPACE$:
%
%
\begin{description}
%
%
\item[input:] a sentence $A$ and a finite frame $(W^{\prime\prime},R^{\prime\prime})$,
%
%
\item[output:] determine whether $(W^{\prime\prime},R^{\prime\prime}){\models}A$.
%
%
\end{description}
%
%
\end{lemma}
%
%
\begin{proof}
%
%
See~\cite{Vardi:1982}.
%
%
\medskip
%
%
\end{proof}
%
%
%Let $T$ be a set of sentences.
The {\it theory of a set $T$ of sentences}\/ is the set of all sentences satisfied on every frame satisfying $T$'s sentences.
A set of sentences is a {\it theory}\/ if it is the theory of a set of sentences.
%
%
\\
\\
%
%
A theory $T$ is {\it hereditarily undecidable}\/ if for all theories $T^{\prime}$, if $T{\supseteq}T^{\prime}$ then $T^{\prime}$ is undecidable.
%
%
\\
\\
%
%
For all classes ${\mathcal C}$ of frames, let $\Th({\mathcal C})$ be the set (called {\it theory of ${\mathcal C}$}) of all sentences which are valid in every ${\mathcal C}$'s frame.
%
%
\begin{lemma}\label{lemma:gamma}
%
%
If $\{2\}{\times}\mathbf{N}^{-}{\subseteq}\mathtt{S}_{\L}$ then $\Th(\Fr(\L))$ is contained in the first-order
%
%
%
%
%
%
%
%
%
%
%
%
\linebreak
%
%
%
%
%
%
%
%
%
%
%
%
theory of ${\mathcal K}_{2}$.
%
%
\end{lemma}
%
%
\begin{proof}
%
%
For the sake of the contradiction, suppose $\{2\}{\times}\mathbf{N}^{-}{\subseteq}\mathtt{S}_{\L}$ and $\Th(\Fr(\L))$ is not contained in the first-order theory of ${\mathcal K}_{2}$.
Hence, there exists a sentence $A$ such that $A$ is valid in every $\Fr(\L)$'s frame and $A$ is not valid in some ${\mathcal K}_{2}$'s galaxy.
Thus, let ${\mathcal F}_{A,B}^{\rho}$ be a ${\mathcal K}_{2}$'s galaxy such that ${\mathcal F}_{A,B}^{\rho}{\not\models}A$.
Since $\{2\}{\times}\mathbf{N}^{-}{\subseteq}\mathtt{S}_{\L}$, then by Lemma~\ref{lemma:alpha}, ${\mathcal F}_{A,B}^{\rho}{\models}\L$.
Since $A$ is valid in every $\Fr(\L)$'s frame, then ${\mathcal F}_{A,B}^{\rho}{\models}A$: a contradiction.
%
%
\medskip
%
%
\end{proof}
%
%
\begin{lemma}\label{lemma:delta}
%
%
If $\mathbf{N}^{+}{\times}\{2\}{\subseteq}\mathtt{S}_{\L}$ then $\Th(\Fr(\L))$ is contained in the first-order
%
%
%
%
%
%
%
%
%
%
%
%
\linebreak
%
%
%
%
%
%
%
%
%
%
%
%
theory of ${\mathcal L}_{2}$.
%
%
\end{lemma}
%
%
\begin{proof}
%
%
For the sake of the contradiction, suppose $\mathbf{N}^{+}{\times}\{2\}{\subseteq}\mathtt{S}_{\L}$ and $\Th(\Fr(\L))$ is not contained in the first-order theory of ${\mathcal L}_{2}$.
Hence, there exists a sentence $A$ such that $A$ is valid in every $\Fr(\L)$'s frame and $A$ is not valid in some ${\mathcal L}_{2}$'s galaxy.
Thus, let ${\mathcal F}_{A,B}^{\rho}$ be a ${\mathcal L}_{2}$'s galaxy such that ${\mathcal F}_{A,B}^{\rho}{\not\models}A$.
Since $\mathbf{N}^{+}{\times}\{2\}{\subseteq}\mathtt{S}_{\L}$, then by Lemma~\ref{lemma:beta}, ${\mathcal F}_{A,B}^{\rho}{\models}\L$.
Since $A$ is valid in every $\Fr(\L)$'s frame, then ${\mathcal F}_{A,B}^{\rho}{\models}A$: a contradiction.
%
%
\medskip
%
%
\end{proof}
%
%
For all $q{\in}\N$, a frame $(W^{\prime},R^{\prime})$ is {\it elementary $q$-equivalent to $(W,R)$}\/ (denoted $(W^{\prime},R^{\prime}){\equiv_{q}}(W,R)$) if for all sentences $A$, if $\qd(A){\leq}q$ then $(W^{\prime},R^{\prime}){\models}A$ if and only if $(W,R){\models}A$.
%
%
\begin{lemma}\label{ehrenfeucht:theorem}
%
%
For all $q{\in}\N$ and for all frames $(W^{\prime},R^{\prime})$, if the second player has a winning strategy in the Ehrenfeucht-Fra\"\i ss\'e game ${\mathcal G}_{q}((W,R),(W^{\prime},R^{\prime}))$ then $(W^{\prime},R^{\prime}){\equiv_{q}}(W,R)$.
%
%
\end{lemma}
%
%
\begin{proof}
%
%
See~\cite[Theorem~$2.2.8$ at Page~$18$]{Ebbinghaus:Flum:1999}.
%
%
\medskip
%
%
\end{proof}
%
%
Let $q{\geq}3$.
%
%
\\
\\
%
%
Let $\{{\mathcal F}_{A_{i},B_{i}}^{\rho_{i}}:\ i{\in}I\}$ be a disjoint indexed family of galaxies.
%
%
\begin{lemma}\label{lemma:alpha:reduction:fol}
%
%
%We shall say that $\{(W_{A_{i}^{\prime},B_{i}^{\prime}}^{\rho_{i}^{\prime}},R_{A_{i}^{\prime},B_{i}^{\prime}}^{\rho_{i}^{\prime}}):\ i{\in}I\}$ is {\it $q$-$\alpha$-reduced from $\{(W_{A_{i},B_{i}}^{\rho_{i}},R_{A_{i},B_{i}}^{\rho_{i}}):\ i{\in}I\}$}\/ if for all $i{\in}I$, either $(W_{A_{i}^{\prime},B_{i}^{\prime}}^{\rho_{i}^{\prime}},R_{A_{i}^{\prime},B_{i}^{\prime}}^{\rho_{i}^{\prime}})$ is equal to $(W_{A_{i},B_{i}}^{\rho_{i}},R_{A_{i},B_{i}}^{\rho_{i}})$, or $(W_{A_{i}^{\prime},B_{i}^{\prime}}^{\rho_{i}^{\prime}},R_{A_{i}^{\prime},B_{i}^{\prime}}^{\rho_{i}^{\prime}})$ is a $q$-$\alpha$-reduction of $(W_{A_{i},B_{i}}^{\rho_{i}},R_{A_{i},B_{i}}^{\rho_{i}})$.
If $\{{\mathcal F}_{A_{i}^{\prime},B_{i}^{\prime}}^{\rho_{i}^{\prime}}:\ i{\in}I\}$ is a $q$-$\alpha$-reduction of $\{{\mathcal F}_{A_{i},B_{i}}^{\rho_{i}}:\ i{\in}I\}$, $(W,R)$ is the union of $\{{\mathcal F}_{A_{i},B_{i}}^{\rho_{i}}:\ i{\in}I\}$ and $(W^{\prime},R^{\prime})$ is the union of $\{{\mathcal F}_{A_{i}^{\prime},B_{i}^{\prime}}^{\rho_{i}^{\prime}}:\ i{\in}I\}$ then $(W,R){\equiv_{q}}(W^{\prime},R^{\prime})$.
%
%
\end{lemma}
%
%
\begin{proof}
%
%
By Proposition~\ref{lemma:good:properties:of:alpha:reductions:indexed} and Lemma~\ref{ehrenfeucht:theorem}.
%
%
\medskip
%
%
\end{proof}
%
%
%\begin{lemma}\label{lemma:beta:reduction:fol}
%
%
%Let $\{(W_{A_{i},B_{i}}^{\rho_{i}},R_{A_{i},B_{i}}^{\rho_{i}}):\ i{\in}I\}$ be a disjoint indexed family of simple galaxies, $q{\in}\N$ and $\{(W_{A_{i}^{\prime},B_{i}}^{\rho_{i}^{\prime}},R_{A_{i}^{\prime},B_{i}}^{\rho_{i}^{\prime}}):\ i{\in}I\}$ be $q$-$\beta$-reduced from $\{(W_{A_{i},B_{i}}^{\rho_{i}},R_{A_{i},B_{i}}^{\rho_{i}}):\ i{\in}I\}$.
%If $(W,R)$ is the union of $\{(W_{A_{i},B_{i}}^{\rho_{i}},R_{A_{i},B_{i}}^{\rho_{i}}):\ i{\in}I\}$ and $(W^{\prime},R^{\prime})$ is the union of $\{(W_{A_{i}^{\prime},B_{i}}^{\rho_{i}^{\prime}},R_{A_{i}^{\prime},B_{i}}^{\rho_{i}^{\prime}}):\ i{\in}I\}$ then $(W,R){\equiv_{q}}(W^{\prime},R^{\prime})$.
%
%
%\end{lemma}
%
%
\begin{lemma}\label{lemma:gamma:reduction:fol}
%
%
%Let $q{\geq}3$.
%Let $\{{\mathcal F}_{A_{i},B_{i}}^{\rho_{i}}:\ i{\in}I\}$ be a disjoint family of galaxies.
%We shall say that $\{(W_{A_{i}^{\prime},B_{i}^{\prime}}^{\rho_{i}^{\prime}},R_{A_{i}^{\prime},B_{i}^{\prime}}^{\rho_{i}^{\prime}}):\ i{\in}I\}$ is {\it $q$-$\alpha$-reduced from $\{(W_{A_{i},B_{i}}^{\rho_{i}},R_{A_{i},B_{i}}^{\rho_{i}}):\ i{\in}I\}$}\/ if for all $i{\in}I$, either $(W_{A_{i}^{\prime},B_{i}^{\prime}}^{\rho_{i}^{\prime}},R_{A_{i}^{\prime},B_{i}^{\prime}}^{\rho_{i}^{\prime}})$ is equal to $(W_{A_{i},B_{i}}^{\rho_{i}},R_{A_{i},B_{i}}^{\rho_{i}})$, or $(W_{A_{i}^{\prime},B_{i}^{\prime}}^{\rho_{i}^{\prime}},R_{A_{i}^{\prime},B_{i}^{\prime}}^{\rho_{i}^{\prime}})$ is a $q$-$\alpha$-reduction of $(W_{A_{i},B_{i}}^{\rho_{i}},R_{A_{i},B_{i}}^{\rho_{i}})$.
If $\{{\mathcal F}_{A_{i}^{\prime},B_{i}^{\prime}}^{\rho_{i}^{\prime}}:\ i{\in}I\}$ is a $q$-$\gamma$-reduction of $\{{\mathcal F}_{A_{i},B_{i}}^{\rho_{i}}:\ i{\in}I\}$, $(W,R)$ is the union of $\{{\mathcal F}_{A_{i},B_{i}}^{\rho_{i}}:\ i{\in}I\}$ and $(W^{\prime},R^{\prime})$ is the union of $\{{\mathcal F}_{A_{i}^{\prime},B_{i}^{\prime}}^{\rho_{i}^{\prime}}:\ i{\in}I\}$ then $(W,R){\equiv_{q}}(W^{\prime},R^{\prime})$.
%
%
\end{lemma}
%
%
\begin{proof}
%
%
By Proposition~\ref{lemma:good:properties:of:gamma:reductions:indexed} and Lemma~\ref{ehrenfeucht:theorem}.
%
%
\medskip
%
%
\end{proof}
%
%
\begin{lemma}\label{lemma:delta:reduction:fol}
%
%
%Let $q{\geq}3$.
%Let $\{{\mathcal F}_{A_{i},B_{i}}^{\rho_{i}}:\ i{\in}I\}$ be a disjoint indexed family of galaxies.
If $\{{\mathcal F}_{A_{j},B_{j}}^{\rho_{j}}:\ j{\in}J\}$ is a $q$-$\delta$-reduction of $\{{\mathcal F}_{A_{i},B_{i}}^{\rho_{i}}:\ i{\in}I\}$, $(W,R)$ is the union of $\{{\mathcal F}_{A_{i},B_{i}}^{\rho_{i}}:\ i{\in}I\}$ and $(W^{\prime},R^{\prime})$ is the union of $\{{\mathcal F}_{A_{j},B_{j}}^{\rho_{j}}:\ j{\in}J\}$ then $(W,R){\equiv_{q}}(W^{\prime},R^{\prime})$.
%
%
\end{lemma}
%
%
\begin{proof}
%
%
By Proposition~\ref{lemma:good:properties:of:delta:reductions:indexed} and Lemma~\ref{ehrenfeucht:theorem}.
%
%
\medskip
%
%
\end{proof}
%
%
%
%
\paragraph{Rooted translation}
%
%
The {\it rooted translation of a first-order formula $A$ with respect to an individual variable $\mathbf{x}$ not occurring in $A$}\/ (denoted $\tau(\mathbf{x},A)$) is inductively defined as follows:
%
%
\begin{itemize}
%
%
\item $\tau(\mathbf{x},\mathbf{R}(\mathbf{y},\mathbf{z}))$ is $\mathbf{R}(\mathbf{y},\mathbf{z})$,
%
%
\item $\tau(\mathbf{x},\mathbf{y}=\mathbf{z})$ is $\mathbf{y}=\mathbf{z}$,
%
%
\item $\tau(\mathbf{x},\neg A)$ is $\neg\tau(\mathbf{x},A)$,
%
%
\item $\tau(\mathbf{x},A\vee B)$ is $\tau(\mathbf{x},A)\vee\tau(\mathbf{x},B)$,
%
%
%\item $\tau(\mathbf{x},\forall\mathbf{y}A)$ is $\forall\mathbf{y}(\mathbf{x}=\mathbf{y}\vee\mathbf{R}(\mathbf{x},\mathbf{y})\vee\exists\mathbf{z}(\mathbf{R}(\mathbf{x},\mathbf{z})\wedge\mathbf{R}(\mathbf{z},\mathbf{y}))\rightarrow\tau(\mathbf{x},A))$.
%
%
\item $\tau(\mathbf{x},\forall\mathbf{y}A)$ is $\forall\mathbf{y}(\mathbf{x}=\mathbf{y}\vee\exists\mathbf{z}(\mathbf{R}(\mathbf{x},\mathbf{z})\wedge\mathbf{R}(\mathbf{z},\mathbf{y}))\rightarrow\tau(\mathbf{x},A))$.
%
%
\end{itemize}
%
%
%\begin{lemma}\label{lemma:about:the:rooted:translation:preliminary:statement}
%
%
%Let $(W,R)$ be an Euclidean frame and $g$ be an assignment on $(W,R)$.
%For all sentences $A$ and for all individual variables $\mathbf{x}$ not occurring in $A$, the following conditions are equivalent:
%
%
%\begin{itemize}
%
%
%\item $(W,R),g{\models}\tau(\mathbf{x},A)$,
%
%
%\item $(W_{g(x)},R_{g(x)}){\models}A$.
%
%
%\end{itemize}
%
%
%\end{lemma}
%
%
\begin{lemma}\label{lemma:about:the:rooted:translation}
%
%
Let $(W,R)$ be an Euclidean frame.
For all sentences $A$ and for all individual variables $\mathbf{x}$ not occurring in $A$, the following conditions are equivalent:
%
%
\begin{itemize}
%
%
\item $(W,R){\models}\forall\mathbf{x}\tau(\mathbf{x},A)$,
%
%
\item for all $s{\in}W$, $(W_{s},R_{s}){\models}A$.
%
%
\end{itemize}
%
%
\end{lemma}
%
%
%
%
\section{Definability problems}\label{section:definability:problems}
%
%
In this section, we introduce $3$~definability problems: first-order definability problem, modal definability problem and correspondence problem.
We remind what is known concerning their computability.
%
%
%
%
%\paragraph{First-order definability problem}
\\
\\
%
%
%Let ${\mathcal C}$ be a class of frames.
A modal formula $\varphi$ is {\it first-order definable in a class ${\mathcal C}$ of frames}\/ if there exists a sentence $A$ such that for all frames ${\mathcal F}$ in ${\mathcal C}$, ${\mathcal F}{\models}\varphi$ if and only if ${\mathcal F}{\models}A$.
As is well-known, in the class of all frames, $\Box\Diamond p\rightarrow\Diamond\Box p$ is not first-order definable~\cite[Chapter~$3$]{Blackburn:deRijke:Venema:2001}.
The {\it first-order definability problem in a class ${\mathcal C}$ of frames}\/ is the following decision problem: determine whether a given modal formula is first-order definable in ${\mathcal C}$.
%
%
\\
\\
%
%
%For all classes ${\mathcal C}$ of frames, let $\fod({\mathcal C})$ be the set of all first-order definable modal formulas in ${\mathcal C}$.
%
%
%
%
%\paragraph{Modal definability problem}
%\\
%\\
%
%
%Let ${\mathcal C}$ be a class of frames.
A sentence $A$ is {\it modally definable in a class ${\mathcal C}$ of frames}\/ if there exists a modal formula $\varphi$ such that for all frames ${\mathcal F}$ in ${\mathcal C}$, ${\mathcal F}{\models}\varphi$ if and only if ${\mathcal F}{\models}A$.
As is well-known, in the class of all frames, $\exists\mathbf{x}\mathbf{R}(\mathbf{x},\mathbf{x})$ is not modally definable~\cite[Chapter~$3$]{Blackburn:deRijke:Venema:2001}.
The {\it modal definability problem in a class ${\mathcal C}$ of frames}\/ is the following decision problem: determine whether a given sentence is modally definable in ${\mathcal C}$.
%
%
%\end{description}
%
%
%For all classes ${\mathcal C}$ of frames, let $\mdef({\mathcal C})$ be the set of all modally definable sentences in ${\mathcal C}$.
%
%
%
%
%\paragraph{Correspondence problem}
\\
\\
%
%
%Let ${\mathcal C}$ be a class of frames.
A modal formula $\varphi$ and a sentence $A$ {\it correspond in a class ${\mathcal C}$ of frames}\/ if for all frames ${\mathcal F}$ in ${\mathcal C}$, ${\mathcal F}{\models}\varphi$ if and only if ${\mathcal F}{\models}A$.
As is well-known, in the class of all frames, $\Diamond p\rightarrow\Box\Diamond p$ corresponds with $\forall\mathbf{x}\forall\mathbf{y}\forall\mathbf{z}(\mathbf{R}(\mathbf{x},\mathbf{y})\wedge\mathbf{R}(\mathbf{x},\mathbf{z})\rightarrow\mathbf{R}(\mathbf{y},\mathbf{z})\wedge\mathbf{R}(\mathbf{z},\mathbf{y}))$~\cite[Chapter~$3$]{Blackburn:deRijke:Venema:2001}.
The {\it correspondence problem in a class ${\mathcal C}$ of frames}\/ is the following decision problem: determine whether a given modal formula and a given sentence correspond in ${\mathcal C}$.
%
%
%
%
%\paragraph{Computability}
\\
\\
%
%
%The question of the correspondence between modal formulas and sentences is concomitant with the creation of the relational semantics of modal logic.
%Kripke~\cite{Kripke:1963} already observed that some sentences are modally definable: transitivity vs $\Box p\rightarrow\Box\Box p$, symmetry vs $p\rightarrow\Box\Diamond p$, etc.
%Less than $20$ years have elapsed between Kripke's observation and the development of Correspondence Theory culminating in the publication of the book ``Modal Logic and Classical Logic''~\cite{vanBenthem:1983}: in 1975, Sahlqvist~\cite{Sahlqvist:1975} isolated a large set of modal formulas which guarantee completeness with respect to first-order definable classes of frames, whereas van Benthem~\cite{vanBenthem:1975} and Goldblatt~\cite{Goldblatt:1975} independently noticed that McKinsey formula $\Box\Diamond p\rightarrow\Diamond\Box p$ has no first-order correspondent.
%
%
%\\
%\\
%
%
%Since the first-order conditions corresponding to Sahlqvist formulas are effectively computable~\cite{Sahlqvist:1975}, then it is natural to ask whether Sahlqvist fragment contains all modal formulas possessing first-order correspondents.
%Since the first-order conditions corresponding to Sahlqvist formulas are effectively computable~\cite{Sahlqvist:1975}, then it is natural to ask whether Sahlqvist fragment modally defines all first-order correspondents of modal formulas.
%This question has received a negative answer, the conjunction $(\Box\Diamond p\rightarrow\Diamond\Box p)\wedge(\Box p\rightarrow\Box\Box p)$ possessing a first-order correspondent not modally definable by a Sahlqvist formula~\cite[Chapter~$3$]{Blackburn:deRijke:Venema:2001}.
%However, notice that this statement should be carefully considered, seeing that the ``Sahlqvist formulas'' considered by Blackburn {\it et al.}\/ constitute a strict subset of the formulas considered by Sahlqvist.
%Indeed, we do not know whether the first-order correspondent of $(\Box\Diamond p\rightarrow\Diamond\Box p)\wedge(\Box p\rightarrow\Box\Box p)$ is modally definable by a formula considered by Sahlqvist.
%
%
%\\
%\\
%
%
%In any case, owing to the significance of Correspondence Theory, it is natural to ask whether the first-order definability problem, the modal definability problem and the correspondence problem are decidable.
Owing to the significance of Correspondence Theory~\cite{vanBenthem:1983,vanBenthem:1984}, it is natural to ask whether the first-order definability problem, the modal definability problem and the correspondence problem are decidable.
%
%
%\\
%\\
%
%
The answers to these questions have been firstly obtained by Chagrova in her doctoral thesis~\cite{Chagrova:1989} and then further developed in~\cite{Chagrov:Chagrova:1995,Chagrov:Chagrova:2006,Chagrov:Chagrova:2007,Chagrova:1991}: the first-order definability problem, the modal definability problem and the correspondence problem are undecidable.
%
%
%\begin{itemize}
%
%
%\item there is no algorithm which when presented with an arbitrary modal formula as input, will terminate after finitely many steps returning either ``yes'' if there is a first-order correspondent, or ``no'' if there is not,
%
%
%\item there is no algorithm which when presented with an arbitrary sentence as input, will terminate after finitely many steps returning either ``yes'' if there is a modal correspondent, or ``no'' if there is not,
%
%
%\item there is no algorithm which when presented with an arbitrary modal formula and an arbitrary sentence as input, will terminate after finitely many steps returning either ``yes'' if they correspond, or ``no'' if they do not.
%
%
%\end{itemize}
%
%
\\
\\
%
%
Chagrova's results concern the class of all frames.
They have been obtained by reductions from accessibility problems in Minsky machines~\cite{Chagrov:Chagrova:2006}.
%As a result, it is not obvious how similar results can be obtained with respect to restricted classes of frames (the class of all transitive frames, the class of all symmetric frames, etc).
%However, notice that Chagrov and Chagrova~\cite{Chagrov:Chagrova:1995} have also considered the first-order definability problem with respect to the class of all irreflexive transitive frames.
Although Chagrov and Chagrova~\cite{Chagrov:Chagrova:1995} have also considered the first-order definability problem with respect to the class of all irreflexive transitive frames, it is not obvious how similar results can be obtained with respect to restricted classes of frames such as the class of all transitive frames, the class of all symmetric frames, etc.
%
%
\\
\\
%
%
With respect to the class of frames determined by $\S5$, it follows from $\S5$~Reduction Theorem in~\cite[Page~$51$]{Hughes:Cresswell:1968} and Lemma~$9.7$ in~\cite[Page~$99$]{vanBenthem:1983} that every modal formula is first-order definable.
Moreover, as proved in~\cite{Balbiani:Tinchev:2005,Balbiani:Tinchev:2006}, modal definability with respect to the class of frames determined by $\S5$ is $\PSPACE$-complete.
%
%
%\\
%\\
%
%
With respect to the classes of frames determined by $\KD45$ and $\K45$, similar results hold as well~\cite{Georgiev:2017a,Georgiev:2017b}.
%
%
\\
\\
%
%
Therefore, a first question naturally comes: is there an extension of $\K5$ such that the class of frames it determines gives rise both to a trivial first-order definability problem and an undecidable modal definability problem?
This question has been positively answered: with respect to the class of frames determined by $\K5$, first-order definability is trivial, whereas modal definability is undecidable~\cite{Balbiani:Georgiev:Tinchev:2018}.
%
%
\\
\\
%
%
The truth is that the classes of frames determined by all extensions of $\K5$ give rise to a trivial first-order definability problem.
Therefore, a second question naturally comes: characterize those extensions of $\K5$ such that the classes of frames they determine give rise to an undecidable modal definability problem.
Our answer will be given in Theorem~\ref{proposition:finale}: for all Euclidean modal logics $\L$, the problem of deciding the modal definability of sentences with respect to $\Fr(\L)$ is undecidable when $\mathtt{S}_{\L}{\setminus}((\{1\}{\times}\mathbf{N}^{-}){\cup}(\mathbf{N}^{+}{\times}\{{-}1,0,1\}))$ is infinite and decidable in $\EXPSPACE$ otherwise.
%
%
%
%
%\paragraph{Characterizing modal definability}
\\
\\
%
%
The following result joins up the modal definability character of a sentence and the behaviour of this sentence with respect to generated subframes and bounded morphic images.
%It will be used in Section~\ref{definability:results:section} where we prove that the problem of deciding the modal definability of sentences is decidable with respect to classes of frames determined by Euclidean modal logics $\L$ such that $\mathtt{S}_{\L}{\setminus}((\{1\}{\times}\mathbf{N}^{-}){\cup}0,1\}))$ is finite.
We will need it later in our proofs about modal definability.
%
%
\begin{proposition}\label{proposition:characterizing:modal:definability}
%
%
Let $\L$ be an Euclidean modal logic.
For all sentences $A$, $A$ is modally definable in $\Fr(\L)$ if and only if
%
%
\begin{itemize}
%
%
\item for all Euclidean frames $(W,R)$, if $(W,R){\models}\L$ then $(W,R){\models}A$ if and only if for all $s{\in}W$, $(W_{s},R_{s}){\models}A$,
%
%
\item for all flowers ${\mathcal F}_{m}^{n},{\mathcal F}_{m^{\prime}}^{n^{\prime}}$, if ${\mathcal F}_{m}^{n}{\models}\L$, ${\mathcal F}_{m}^{n}{\models}A$ and ${\mathcal F}_{m^{\prime}}^{n^{\prime}}$ is a bounded morphic image of ${\mathcal F}_{m}^{n}$ then ${\mathcal F}_{m^{\prime}}^{n^{\prime}}{\models}A$.
%
%
\end{itemize}
%
%
\end{proposition}
%
%
\begin{proof}
%
%
Let $A$ be a sentence.
%
%
\\
\\
%
%
$(\Rightarrow)$
By Lemmas~\ref{Generated:Subframe:Lemma} and~\ref{Bounded:Morphism:Lemma}.
%
%
\\
\\
%
%
$(\Leftarrow)$
Suppose
%
%
\begin{itemize}
%
%
\item for all Euclidean frames $(W,R)$, if $(W,R){\models}\L$ then $(W,R){\models}A$ if and only if for all $s{\in}W$, $(W_{s},R_{s}){\models}A$,
%
%
\item for all flowers ${\mathcal F}_{m}^{n},{\mathcal F}_{m^{\prime}}^{n^{\prime}}$, if ${\mathcal F}_{m}^{n}{\models}\L$, ${\mathcal F}_{m}^{n}{\models}A$ and ${\mathcal F}_{m^{\prime}}^{n^{\prime}}$ is a bounded morphic image of ${\mathcal F}_{m}^{n}$ then ${\mathcal F}_{m^{\prime}}^{n^{\prime}}{\models}A$.
%
%
\end{itemize}
%
%
We consider the following $2$~cases.
%
%
\\
\\
%
%
{\bf Case where $\Fr(\L){\models}A$:}
Hence, the modal formula $\top$ and the sentence $A$ correspond in $\Fr(\L)$.
Thus, $A$ is modally definable in $\Fr(\L)$.
%
%
\\
\\
%
%
{\bf Case where $\Fr(\L){\not\models}A$:}
Consequently, let $(W^{\prime},R^{\prime})$ be an Euclidean frame such that $(W^{\prime},R^{\prime}){\models}\L$ and $(W^{\prime},R^{\prime}){\not\models}A$.
Since for all Euclidean frames $(W,R)$, if $(W,R){\models}\L$ then $(W,R){\models}A$ if and only if for all $s{\in}W$, $(W_{s},R_{s}){\models}A$, then let $s^{\prime}{\in}W^{\prime}$ be such that $(W^{\prime}_{s^{\prime}},R^{\prime}_{s^{\prime}}){\not\models}A$.
Since $(W^{\prime},R^{\prime}){\models}\L$, then by Lemma~\ref{Generated:Subframe:Lemma}, $(W^{\prime}_{s^{\prime}},R^{\prime}_{s^{\prime}}){\models}\L$.
Moreover, by Lemma~\ref{lemma:generated:subframe:from:a:single:point:has:specific:shapes}, either $R^{\prime}_{s^{\prime}}{=}W^{\prime}_{s^{\prime}}{\times}W^{\prime}_{s^{\prime}}$, or $W^{\prime}_{s^{\prime}}{=}\{s^{\prime}\}{\cup}C^{\prime}$ and $R^{\prime}_{s^{\prime}}{=}(\{s^{\prime}\}{\times}B^{\prime}){\cup}(C^{\prime}{\times}C^{\prime})$ for some sets $B^{\prime},C^{\prime}$ such that either $B^{\prime}{\not=}\emptyset$, or $C^{\prime}{=}\emptyset$, $B^{\prime}{\subseteq}C^{\prime}$ and $s^{\prime}{\not\in}C^{\prime}$.
Let ${\mathcal C}$ be the class of all rooted Euclidean frames $(W,R)$ such that $(W,R){\models}\L$ and $(W,R){\not\models}A$.
Since $(W^{\prime}_{s^{\prime}},R^{\prime}_{s^{\prime}}){\not\models}A$ and $(W^{\prime}_{s^{\prime}},R^{\prime}_{s^{\prime}}){\models}\L$, then $(W^{\prime}_{s^{\prime}},R^{\prime}_{s^{\prime}})$ is in ${\mathcal C}$.
Let $q{=}\qdd(A)$.
%Notice that $\qd(A){\leq}q$ and $q{\geq}3$.
Let
%
%
\begin{itemize}
%
%
\item $\mathtt{S}_{\L}^{A}{=}\{(m,n):\ m{\in}\lsem0,q\rsem$, $n{\in}\lsem{-}1,q\rsem$, either $m{\not=}0$, or $n{=}0$, ${\mathcal F}_{m}^{n}{\models}\L$ and ${\mathcal F}_{m}^{n}{\not\models}A\}$.
%
%
\end{itemize}
%
%
Obviously, $\mathtt{S}_{\L}^{A}{\subseteq}\mathtt{S}_{\L}{\cup}\{(0,0\}$.
%
%
\\
\\
%
%
We claim that $\mathtt{S}_{\L}^{A}$ is non-empty.
We consider the following $2$~cases.
%
%
\\
\\
%
%
{\bf Case where $R^{\prime}_{s^{\prime}}{=}W^{\prime}_{s^{\prime}}{\times}W^{\prime}_{s^{\prime}}$:}
Let $(W^{\prime\prime},R^{\prime\prime})$ be an Euclidean frame obtained from $(W^{\prime}_{s^{\prime}},R^{\prime}_{s^{\prime}})$ as follows:
%
%
\begin{itemize}
%
%
\item if ${\parallel}W^{\prime}_{s^{\prime}}{\parallel}{\leq}q$ then let $W^{\prime\prime}$ be $W^{\prime}_{s^{\prime}}$ else let $W^{\prime\prime}$ be an arbitrary subset of $W^{\prime}_{s^{\prime}}$ of cardinality $q$,
%
%
\item let $R^{\prime\prime}$ be $W^{\prime\prime}{\times}W^{\prime\prime}$.
%
%
\end{itemize}
%
%
Obviously, either ${\parallel}W^{\prime}_{s^{\prime}}{\parallel}{\leq}q$ and ${\parallel}W^{\prime\prime}{\parallel}{=}{\parallel}W^{\prime}_{s^{\prime}}{\parallel}$, or ${\parallel}W^{\prime}_{s^{\prime}}{\parallel}{>}q$ and ${\parallel}W^{\prime\prime}{\parallel}{=}q$.
Hen\-ce, by Lemma~\ref{lemma:about:partitions:q:greater:smaller}, $(W^{\prime\prime},R^{\prime\prime})$ is a bounded morphic image of $(W^{\prime}_{s^{\prime}},R^{\prime}_{s^{\prime}})$.
Moreover, by Lemma~\ref{lemma:about:partitions:q:greater:smaller:games}, the second player has a winning strategy in the Ehrenfeucht-Fra\"\i ss\'e game ${\mathcal G}_{q}((W^{\prime}_{s^{\prime}},R^{\prime}_{s^{\prime}}),(W^{\prime\prime},R^{\prime\prime}))$.
Thus, by Lemmas~\ref{Bounded:Morphism:Lemma} and~\ref{ehrenfeucht:theorem}, $(W^{\prime}_{s^{\prime}},R^{\prime}_{s^{\prime}})
%
%
%
%
%
%
%
%
%
%
%
%
$\linebreak$
%
%
%
%
%
%
%
%
%
%
%
%
{\preceq}(W^{\prime\prime},R^{\prime\prime})$ and $(W^{\prime\prime},R^{\prime\prime}){\equiv_{q}}(W^{\prime}_{s^{\prime}},R^{\prime}_{s^{\prime}})$.
Since $(W^{\prime}_{s^{\prime}},R^{\prime}_{s^{\prime}}){\not\models}A$ and $(W^{\prime}_{s^{\prime}},R^{\prime}_{s^{\prime}}){\models}\L$, then $(W^{\prime\prime},R^{\prime\prime}){\not\models}A$ and $(W^{\prime\prime},R^{\prime\prime}){\models}\L$.
Let $m{=}{\parallel}W^{\prime\prime}{\parallel}$ and $n{=}{-}1$.
Obviously, $m{\in}\lsem0,q\rsem$, $n{\in}\lsem{-}1,q\rsem$ and either $m{\not=}0$, or $n{=}0$.
Moreover, $(W^{\prime\prime},R^{\prime\prime})$ is isomorphic with ${\mathcal F}_{m}^{n}$.
Since $(W^{\prime\prime},R^{\prime\prime}){\not\models}A$ and $(W^{\prime\prime},R^{\prime\prime}){\models}\L$, then ${\mathcal F}_{m}^{n}{\not\models}A$ and ${\mathcal F}_{m}^{n}{\models}\L$.
Since $m{\in}\lsem0,q\rsem$, $n{\in}\lsem{-}1,q\rsem$ and either $m{\not=}0$, or $n{=}0$, then $(m,n)$ is in $\mathtt{S}_{\L}^{A}$.
%Consequently, $\mathtt{S}_{\L}^{A}$ is non-empty.
%
%
\\
\\
%
%
{\bf Case where $W^{\prime}_{s^{\prime}}{=}\{s^{\prime}\}{\cup}C^{\prime}$ and $R^{\prime}_{s^{\prime}}{=}(\{s^{\prime}\}{\times}B^{\prime}){\cup}(C^{\prime}{\times}C^{\prime})$ for some sets $B^{\prime},
%
%
%
%
%
%
%
%
%
%
%
%
$\linebreak$
%
%
%
%
%
%
%
%
%
%
%
%
C^{\prime}$ such that either $B^{\prime}{\not=}\emptyset$, or $C^{\prime}{=}\emptyset$, $B^{\prime}{\subseteq}C^{\prime}$ and $s^{\prime}{\not\in}C^{\prime}$:}
Let $(W^{\prime\prime},R^{\prime\prime})$ be an Euclidean frame obtained from $(W^{\prime}_{s^{\prime}},R^{\prime}_{s^{\prime}})$ as follows:
%
%
\begin{itemize}
%
%
\item if ${\parallel}B^{\prime}{\parallel}{\leq}q$ then let $B^{\prime\prime}$ be $B^{\prime}$ else let $B^{\prime\prime}$ be an arbitrary subset of $B^{\prime}$ of cardinality $q$,
%
%
\item if ${\parallel}C^{\prime}{\setminus}B^{\prime}{\parallel}{\leq}q$ then let $D^{\prime\prime}$ be $C^{\prime}{\setminus}B^{\prime}$ else let $D^{\prime\prime}$ be an arbitrary subset of $C^{\prime}{\setminus}B^{\prime}$ of cardinality $q$,
%
%
\item let $C^{\prime\prime}$ be $B^{\prime\prime}{\cup}D^{\prime\prime}$,
%
%
\item let $W^{\prime\prime}$ be $\{s^{\prime}\}{\cup}C^{\prime\prime}$,
%
%
\item let $R^{\prime\prime}$ be $(\{s^{\prime}\}{\times}B^{\prime\prime}){\cup}(C^{\prime\prime}{\times}C^{\prime\prime})$.
%
%
\end{itemize}
%
%
Obviously, either ${\parallel}B^{\prime}{\parallel}{\leq}q$ and ${\parallel}B^{\prime}{\parallel}{=}{\parallel}B^{\prime\prime}{\parallel}$, or ${\parallel}B^{\prime}{\parallel}{>}q$ and ${\parallel}B^{\prime\prime}{\parallel}{=}q$ and either ${\parallel}C^{\prime}{\setminus}B^{\prime}{\parallel}{\leq}q$ and ${\parallel}C^{\prime}{\setminus}B^{\prime}{\parallel}{=}{\parallel}C^{\prime\prime}{\setminus}B^{\prime\prime}{\parallel}$, or ${\parallel}C^{\prime}{\setminus}B^{\prime}{\parallel}{>}q$ and ${\parallel}C^{\prime\prime}{\setminus}B^{\prime\prime}{\parallel}{=}q$.
Consequently, by Lemma~\ref{lemma:about:frames:C:B:smaller:greater}, $(W^{\prime\prime},R^{\prime\prime})$ is a bounded morphic image of $(W^{\prime}_{s^{\prime}},R^{\prime}_{s^{\prime}})$.
Moreover, by Lemma~\ref{lemma:about:frames:C:B:smaller:greater:games}, the second player has a winning strategy in the Ehrenfeucht-Fra\"\i ss\'e game ${\mathcal G}_{q}((W^{\prime}_{s^{\prime}},R^{\prime}_{s^{\prime}}),(W^{\prime\prime},R^{\prime\prime}))$.
Hence, by Lemmas~\ref{Bounded:Morphism:Lemma}
%
%
%
%
%
%
%
%
%
%
%
%
\linebreak
%
%
%
%
%
%
%
%
%
%
%
%
and~\ref{ehrenfeucht:theorem}, $(W^{\prime}_{s^{\prime}},R^{\prime}_{s^{\prime}}){\preceq}(W^{\prime\prime},R^{\prime\prime})$ and $(W^{\prime\prime},R^{\prime\prime}){\equiv_{q}}(W^{\prime}_{s^{\prime}},R^{\prime}_{s^{\prime}})$.
Since $(W^{\prime}_{s^{\prime}},R^{\prime}_{s^{\prime}}){\not\models}A$ and $(W^{\prime}_{s^{\prime}},R^{\prime}_{s^{\prime}}){\models}\L$, then $(W^{\prime\prime},R^{\prime\prime}){\not\models}A$ and $(W^{\prime\prime},R^{\prime\prime}){\models}\L$.
Let $m{=}{\parallel}B^{\prime\prime}{\parallel}$ and $n{=}{\parallel}D^{\prime\prime}{\parallel}$.
Obviously, $m{\in}\lsem0,q\rsem$, $n{\in}\lsem{-}1,q\rsem$ and either $m{\not=}0$, or $n{=}0$.
Moreover, $(W^{\prime\prime},R^{\prime\prime})$ is isomorphic with ${\mathcal F}_{m}^{n}$.
Since $(W^{\prime\prime},R^{\prime\prime}){\not\models}A$ and $(W^{\prime\prime},R^{\prime\prime}){\models}\L$, then ${\mathcal F}_{m}^{n}{\not\models}A$ and ${\mathcal F}_{m}^{n}{\models}\L$.
Since $m{\in}\lsem0,q\rsem$, $n{\in}\lsem{-}1,q\rsem$ and either $m{\not=}0$, or $n{=}0$, then $(m,n)$ is in $\mathtt{S}_{\L}^{A}$.
%Consequently, $\mathtt{S}_{\L}^{A}$ is non-empty.
%
%
\\
\\
%
%
Thus, $\mathtt{S}_{\L}^{A}$ is non-empty.
%
%
\\
\\
%
%
We claim that $\mathtt{S}_{\L}^{A}$ is finite.
Obviously, $\mathtt{S}_{\L}^{A}{\subseteq}\lsem0,q\rsem{\times}\lsem{-}1,q\rsem$.
Consequently, $\mathtt{S}_{\L}^{A}$ is finite.
%
%
\\
\\
%
%
Let
%
%
\begin{itemize}
%
%
\item $\varphi_{\L}^{A}{=}\bigwedge\{\chi_{m}^{n}:\ (m,n){\in}\mathtt{S}_{\L}^{A}\}$.
%
%
\end{itemize}
%
%
%\\
%\\
%
%
We claim that the modal formula $\varphi_{\L}^{A}$ and the sentence $A$ correspond in $\Fr(\L)$.
If not, let $(W^{\prime},R^{\prime})$ be an Euclidean frame such that $(W^{\prime},R^{\prime}){\models}\L$ and either $(W^{\prime},R^{\prime}){\models}\varphi_{\L}^{A}$ and $(W^{\prime},R^{\prime}){\not\models}A$, or $(W^{\prime},R^{\prime}){\not\models}\varphi_{\L}^{A}$ and $(W^{\prime},R^{\prime}){\models}A$.
In the former case, since for all Euclidean frames $(W,R)$, if $(W,R){\models}\L$ then $(W,R){\models}A$ if and only if for all $s{\in}W$, $(W_{s},R_{s}){\models}A$, then let $s^{\prime}{\in}W^{\prime}$ be such that $(W^{\prime}_{s^{\prime}},R^{\prime}_{s^{\prime}}){\not\models}A$.
Since $(W^{\prime},R^{\prime}){\models}\L$ and $(W^{\prime},R^{\prime}){\models}\varphi_{\L}^{A}$, then by Lemma~\ref{Generated:Subframe:Lemma}, $(W^{\prime}_{s^{\prime}},R^{\prime}_{s^{\prime}}){\models}\L$ and $(W^{\prime}_{s^{\prime}},R^{\prime}_{s^{\prime}}){\models}\varphi_{\L}^{A}$.
Since $(W^{\prime}_{s^{\prime}},R^{\prime}_{s^{\prime}}){\not\models}A$, then let $(W^{\prime\prime},R^{\prime\prime})$ be the Euclidean frame and $(m,n)$ be the couple of integers obtained from $(W^{\prime}_{s^{\prime}},R^{\prime}_{s^{\prime}})$ as above in the argument showing that $\mathtt{S}_{\L}^{A}$ is non-empty.
Remind from this argument that $(W^{\prime}_{s^{\prime}},R^{\prime}_{s^{\prime}}){\preceq}(W^{\prime\prime},R^{\prime\prime})$ and $(W^{\prime\prime},R^{\prime\prime}){\equiv_{q}}(W^{\prime}_{s^{\prime}},R^{\prime}_{s^{\prime}})$.
Moreover, $m{\in}\lsem0,q\rsem$, $n{\in}\lsem{-}1,q\rsem$ and either $m{\not=}0$, or $n{=}0$.
In other respect, $(W^{\prime\prime},R^{\prime\prime})$ is isomorphic with ${\mathcal F}_{m}^{n}$.
Since $(W^{\prime}_{s^{\prime}},R^{\prime}_{s^{\prime}}){\not\models}A$, $(W^{\prime}_{s^{\prime}},R^{\prime}_{s^{\prime}}){\models}\L$ and $(W^{\prime}_{s^{\prime}},R^{\prime}_{s^{\prime}}){\models}\varphi_{\L}^{A}$, then ${\mathcal F}_{m}^{n}{\not\models}A$, ${\mathcal F}_{m}^{n}{\models}\L$ and ${\mathcal F}_{m}^{n}{\models}\varphi_{\L}^{A}$.
%Hence, there exists $(m,n){\in}\mathtt{S}_{\L}^{A}$ such that $(W^{\prime\prime},R^{\prime\prime})$ and ${\mathcal F}_{m}^{n}$ are isomorphic.
Since $m{\in}\lsem0,q\rsem$, $n{\in}\lsem{-}1,q\rsem$ and either $m{\not=}0$, or $n{=}0$, then $(m,n){\in}\mathtt{S}_{\L}^{A}$.
Since ${\mathcal F}_{m}^{n}{\models}\varphi_{\L}^{A}$, then ${\mathcal F}_{m}^{n}{\models}\chi_{m}^{n}$: a contradiction with Lemma~\ref{lemma:Jankov:Fine:TWO}.
In the latter case, there exists a valuation $V^{\prime}$ on $(W^{\prime},R^{\prime})$ and there exists $s^{\prime}{\in}W^{\prime}$ such that $(W^{\prime},R^{\prime}),V^{\prime},s^{\prime}{\not\models}\varphi_{\L}^{A}$.
Moreover, since for all Euclidean frames $(W,R)$, if $(W,R){\models}\L$ then $(W,R){\models}A$ if and only if for all $s{\in}W$, $(W_{s},R_{s}){\models}A$, then $(W^{\prime}_{s^{\prime}},R^{\prime}_{s^{\prime}}){\models}A$.
Hence, there exists $(m,n){\in}\mathtt{S}_{\L}^{A}$ such that $(W^{\prime},R^{\prime}),V^{\prime},s^{\prime}{\not\models}\chi_{m}^{n}$.
Thus, by Lemma~\ref{lemma:Jankov:Fine:ONE}, there exists a surjective bounded morphism $f$ from $(W^{\prime}_{s^{\prime}},R^{\prime}_{s^{\prime}})$ to ${\mathcal F}_{m}^{n}$ such that $f(s^{\prime})$ generates ${\mathcal F}_{m}^{n}$.
We consider the following $3$~cases.
%
%
\\
\\
%
%
{\bf Case where $m{=}0$:}
Since either $m{\not=}0$, or $n{=}0$, then $n{=}0$.
Since $(m,n){\in}\mathtt{S}_{\L}^{A}$ and $m{=}0$, then ${\mathcal F}_{0}^{0}{\not\models}A$.
Moreover, since $(W^{\prime},R^{\prime}),V^{\prime},s^{\prime}{\not\models}\chi_{m}^{n}$ and $m{=}0$, then $(W^{\prime},R^{\prime}),V^{\prime},s^{\prime}{\models}\Box\bot$.
Consequently, $(W^{\prime}_{s^{\prime}},R^{\prime}_{s^{\prime}})$ is isomorphic with ${\mathcal F}_{0}^{0}$.
Since $(W^{\prime}_{s^{\prime}},R^{\prime}_{s^{\prime}}){\models}A$, then ${\mathcal F}_{0}^{0}{\models}A$: a contradiction.
%
%
\\
\\
%
%
{\bf Case where $m{\not=}0$ and $n{\in}\N$:}
Since $f$ is a surjective bounded morphism from $(W^{\prime}_{s^{\prime}},R^{\prime}_{s^{\prime}})$ to ${\mathcal F}_{m}^{n}$ such that $f(s^{\prime})$ generates ${\mathcal F}_{m}^{n}$, then by Lemma~\ref{lemma:about:what:happens:when:it:exists:bmi:flower:integer}, $W^{\prime}_{s^{\prime}}{=}\{s^{\prime}\}{\cup}C^{\prime}$ and $R^{\prime}_{s^{\prime}}{=}(\{s^{\prime}\}{\times}B^{\prime}){\cup}(C^{\prime}{\times}C^{\prime})$ for some sets $B^{\prime},C^{\prime}$ such that $B^{\prime}{\subseteq}
%
%
%
%
%
%
%
%
%
%
%
%
$\linebreak$
%
%
%
%
%
%
%
%
%
%
%
%
C^{\prime}$, $s^{\prime}{\not\in}C^{\prime}$, ${\parallel}B^{\prime}{\parallel}{\geq}m$ and ${\parallel}C^{\prime}{\setminus}B^{\prime}{\parallel}{\geq}n$.
Since $(m,n){\in}\mathtt{S}_{\L}^{A}$, then ${\mathcal F}_{m}^{n}{\not\models}A$, $m{\leq}q$ and $n{\leq}q$.
Let $m^{\prime}{=}\min\{{\parallel}B^{\prime}{\parallel},q\}$ and $n^{\prime}{=}\min\{{\parallel}C^{\prime}{\setminus}B^{\prime}{\parallel},q\}$.
Obviously, either ${\parallel}B^{\prime}{\parallel}{\leq}q$ and ${\parallel}B^{\prime}{\parallel}{=}m^{\prime}$, or ${\parallel}B^{\prime}{\parallel}{>}q$ and $m^{\prime}{=}q$ and either ${\parallel}C^{\prime}{\setminus}B^{\prime}{\parallel}{\leq}q$ and ${\parallel}C^{\prime}{\setminus}B^{\prime}{\parallel}
%
%
%
%
%
%
%
%
%
%
%
%
$\linebreak$
%
%
%
%
%
%
%
%
%
%
%
%
{=}n^{\prime}$, or ${\parallel}C^{\prime}{\setminus}B^{\prime}{\parallel}{>}q$ and $n^{\prime}{=}q$.
Hence, by Lemma~\ref{lemma:about:frames:C:B:smaller:greater:games}, the second player has a winning strategy in the Ehren\-feucht-Fra\"\i ss\'e game ${\mathcal G}_{q}((W_{s^{\prime}}^{\prime},R_{s^{\prime}}^{\prime}),{\mathcal F}_{m^{\prime}}^{n^{\prime}})$.
Thus, by Lemma~\ref{ehrenfeucht:theorem}, $(W_{s^{\prime}}^{\prime},R_{s^{\prime}}^{\prime}){\equiv_{q}}{\mathcal F}_{m^{\prime}}^{n^{\prime}}$.
Moreover, since ${\parallel}B^{\prime}{\parallel}{\geq}m$, ${\parallel}C^{\prime}{\setminus}B^{\prime}{\parallel}{\geq}n$, $m{\leq}q$ and $n{\leq}q$, then $m{\leq}m^{\prime}$ and $n{\leq}n^{\prime}$.
Consequently, $(m,n){\ll}(m^{\prime},n^{\prime})$.
Hence, by Lemma~\ref{lemma:about:bounded:morphisms:in:the:situation:of F:m:n:frames}, ${\mathcal F}_{m}^{n}$ is a bounded morphic image of ${\mathcal F}_{m^{\prime}}^{n^{\prime}}$.
%Thus, ${\mathcal F}_{m^{\prime}}^{n^{\prime}}{\preceq}{\mathcal F}_{m}^{n}$.
Since $(W^{\prime}_{s^{\prime}},R^{\prime}_{s^{\prime}}){\models}A$ and $(W^{\prime}_{s^{\prime}},R^{\prime}_{s^{\prime}}){\equiv_{q}}{\mathcal F}_{m^{\prime}}^{n^{\prime}}$, then ${\mathcal F}_{m^{\prime}}^{n^{\prime}}{\models}A$.
Since ${\mathcal F}_{m}^{n}$ is a bounded morphic image of ${\mathcal F}_{m^{\prime}}^{n^{\prime}}$, then ${\mathcal F}_{m}^{n}{\models}A$: a contradiction.
%
%
\\
\\
%
%
{\bf Case where $m{\not=}0$ and $n{=}{-}1$:}
Since $f$ is a surjective bounded morphism from $(W^{\prime}_{s^{\prime}},R^{\prime}_{s^{\prime}})$ to ${\mathcal F}_{m}^{n}$ such that $f(s^{\prime})$ generates ${\mathcal F}_{m}^{n}$, then by Lemma~\ref{lemma:about:what:happens:when:it:exists:bmi:flower:minus:one}, either $R^{\prime}_{s^{\prime}}{=}W^{\prime}_{s^{\prime}}{\times}W^{\prime}_{s^{\prime}}$ and ${\parallel}W^{\prime}_{s^{\prime}}{\parallel}{\geq}m$, or $W^{\prime}_{s^{\prime}}{=}\{s^{\prime}\}{\cup}C^{\prime}$ and $R^{\prime}_{s^{\prime}}{=}(\{s^{\prime}\}{\times}B^{\prime}){\cup}(C^{\prime}{\times}C^{\prime})$ for some sets $B^{\prime},C^{\prime}$ such that $B^{\prime}{\subseteq}C^{\prime}$, $s^{\prime}{\not\in}C^{\prime}$ and ${\parallel}B^{\prime}{\parallel}{\geq}m$.
In the former case, since $(m,n){\in}\mathtt{S}_{\L}^{A}$ and $n{=}{-}1$, then ${\mathcal F}_{m}^{{-}1}{\not\models}A$ and $m{\leq}q$.
Let $m^{\prime}{=}\min\{{\parallel}W^{\prime}_{s^{\prime}}{\parallel},q\}$.
Obviously, either ${\parallel}W^{\prime}_{s^{\prime}}{\parallel}{\leq}q$ and ${\parallel}W^{\prime}_{s^{\prime}}{\parallel}{=}m^{\prime}$, or ${\parallel}W^{\prime}_{s^{\prime}}{\parallel}{>}q$ and $m^{\prime}{=}q$.
Thus, by Lemma~\ref{lemma:about:partitions:q:greater:smaller:games}, the second player has a winning strategy in the Ehren\-feucht-Fra\"\i ss\'e game ${\mathcal G}_{q}((W^{\prime}_{s^{\prime}},R^{\prime}_{s^{\prime}}),{\mathcal F}_{m^{\prime}}^{{-}1})$.
Consequently, by Lemma~\ref{ehrenfeucht:theorem}, $(W^{\prime}_{s^{\prime}},R^{\prime}_{s^{\prime}}){\equiv_{q}}{\mathcal F}_{m^{\prime}}^{{-}1}$.
%$(W^{\prime}_{s^{\prime}},R^{\prime}_{s^{\prime}}){\equiv_{q}}{\mathcal F}_{m^{\prime}}^{{-}1}$.
Moreover, since ${\parallel}W^{\prime}_{s^{\prime}}{\parallel}{\geq}m$ and $m{\leq}q$, then $m{\leq}m^{\prime}$.
Hence, $(m,{-}1){\ll}(m^{\prime},{-}1)$.
Thus, by Lemma~\ref{lemma:about:bounded:morphisms:in:the:situation:of F:m:n:frames}, ${\mathcal F}_{m}^{{-}1}$ is a bounded morphic image of ${\mathcal F}_{m^{\prime}}^{{-}1}$.
Since $(W^{\prime}_{s^{\prime}},R^{\prime}_{s^{\prime}})
%
%
%
%
%
%
%
%
%
%
%
%
$\linebreak$
%
%
%
%
%
%
%
%
%
%
%
%
{\models}A$ and $(W^{\prime}_{s^{\prime}},R^{\prime}_{s^{\prime}}){\equiv_{q}}{\mathcal F}_{m^{\prime}}^{{-}1}$, then ${\mathcal F}_{m^{\prime}}^{{-}1}{\models}A$.
Since ${\mathcal F}_{m}^{{-}1}$ is a bounded morphic image of ${\mathcal F}_{m^{\prime}}^{{-}1}$, then ${\mathcal F}_{m}^{{-}1}{\models}A$: a contradiction.
In the latter case, since $(m,n){\in}\mathtt{S}_{\L}^{A}$ and $n{=}{-}1$, then ${\mathcal F}_{m}^{{-}1}{\not\models}A$ and $m{\leq}q$.
Let $m^{\prime}{=}\min\{{\parallel}C^{\prime}{\parallel},q\}$.
Obviously, either ${\parallel}C^{\prime}{\parallel}{\leq}q$ and ${\parallel}C^{\prime}{\parallel}{=}m^{\prime}$, or ${\parallel}C^{\prime}{\parallel}{>}q$ and $m^{\prime}{=}q$.
Consequently, by Lemma~\ref{lemma:about:frames:C:B:smaller:greater:games}, the second player has a winning strategy in the Ehren\-feucht-Fra\"\i ss\'e game ${\mathcal G}_{q}((C^{\prime},C^{\prime}
%
%
%
%
%
%
%
%
%
%
%
%
$\linebreak$
%
%
%
%
%
%
%
%
%
%
%
%
{\times}C^{\prime}),{\mathcal F}_{m^{\prime}}^{{-}1})$.
Hence, by Lemma~\ref{ehrenfeucht:theorem}, $(C^{\prime},C^{\prime}{\times}C^{\prime}){\equiv_{q}}{\mathcal F}_{m^{\prime}}^{{-}1}$.
%$(W^{\prime}_{s^{\prime}},R^{\prime}_{s^{\prime}}){\equiv_{q}}{\mathcal F}_{m^{\prime}}^{{-}1}$.
Moreover, since $B^{\prime}{\subseteq}C^{\prime}$, ${\parallel}B^{\prime}{\parallel}{\geq}m$ and $m{\leq}q$, then $m{\leq}m^{\prime}$.
Thus, $(m,{-}1){\ll}(m^{\prime},{-}1)$.
Consequently, by Lemma~\ref{lemma:about:bounded:morphisms:in:the:situation:of F:m:n:frames}, ${\mathcal F}_{m}^{{-}1}$ is a bounded morphic image of ${\mathcal F}_{m^{\prime}}^{{-}1}$.
Obviously, $(C^{\prime},C^{\prime}{\times}C^{\prime})$ is a bounded morphic image of $(W^{\prime}_{s^{\prime}},R^{\prime}_{s^{\prime}})$.
Since $(W^{\prime}_{s^{\prime}},R^{\prime}_{s^{\prime}}){\models}A$ and $(C^{\prime},C^{\prime}{\times}C^{\prime}){\equiv_{q}}
%
%
%
%
%
%
%
%
%
%
%
%
$\linebreak$
%
%
%
%
%
%
%
%
%
%
%
%
{\mathcal F}_{m^{\prime}}^{{-}1}$, then ${\mathcal F}_{m^{\prime}}^{{-}1}{\models}A$.
Since ${\mathcal F}_{m}^{{-}1}$ is a bounded morphic image of ${\mathcal F}_{m^{\prime}}^{{-}1}$, then ${\mathcal F}_{m}^{{-}1}{\models}A$: a contradiction.
%
%
\\
\\
%
%
Hence, the modal formula $\varphi_{\L}^{A}$ and the sentence $A$ correspond in $\Fr(\L)$.
%
%
\\
\\
%
%
Thus, $A$ is modally definable in $\Fr(\L)$.
%
%
\medskip
%
%
\end{proof}
%
%
%
%
\section{Relative interpretations}\label{section:about:relative:interpretations}
%
%
In this section, we adapt to our purpose a theorem about relative interpretations.
See~\cite[Page~$272$]{Ershov:1980} for a more general presentation.
%As far as we know, the concept of relative interpretation can only be seen in the book of Ershov about {\it Decision Problems and Constructive Models}\/~\cite{Ershov:1980}.
We will use it later for the proofs that the theories of the classes of frames determined by some Euclidean modal logics are undecidable.
%The importance of the concept of relative interpretation lies in the fact that zzzzz~\cite[Theorem $zzzzz$ at Page~$zzzzz$]{Ershov:1980}.
%
%
\\
\\
%
%
Let ${\mathcal C}$ and ${\mathcal C}^{\prime}$ be classes of frames.
%
%
\\
\\
%
%
We say that ${\mathcal C}$ is {\it relatively elementary definable in ${\mathcal C}^{\prime}$}\/ if there exists first-order formulas $\mathbf{U}(\mathbf{x_{1}},\mathbf{x_{2}})$, $\mathbf{E}(\mathbf{x_{1}},\mathbf{x_{2}},\mathbf{x_{3}},\mathbf{x_{4}})$ and $\mathbf{A}(\mathbf{\mathbf{x_{1}},\mathbf{x_{2}},x_{3}},\mathbf{x_{4}})$ such that for all frames $(W,R)$ in ${\mathcal C}$, there exists a frame $(W^{\prime},R^{\prime})$ in ${\mathcal C}^{\prime}$ such that
%
%
\begin{itemize}
%
%
\item $X$ being the set of all couples $(s_{1},s_{2})$ such that $s_{1},s_{2}{\in}W^{\prime}$ are such that $(W^{\prime},R^{\prime}){\models}\mathbf{U}(\mathbf{x_{1}},\mathbf{x_{2}})\ \lbrack s_{1},s_{2}\rbrack$,
%
%
\item $\eta$ being the binary relation on $X$ such that for all $(s_{1},s_{2}),(t_{1},t_{2})$ in $X$, $(s_{1},s_{2})\eta(t_{1},t_{2})$ if and only if $(W^{\prime},R^{\prime}){\models}\mathbf{E}(\mathbf{x_{1}},\mathbf{x_{2}},\mathbf{y_{1}},\mathbf{y_{2}})\ \lbrack s_{1},s_{2},t_{1},t_{2}\rbrack$,
%
%
\item $S$ being the binary relation on $X$ such that for all $(s_{1},s_{2}),(t_{1},t_{2})$ in $X$, $(s_{1},s_{2})S(t_{1},t_{2})$ if and only if $(W^{\prime},R^{\prime}){\models}\mathbf{A}(\mathbf{x_{1}},\mathbf{x_{2}},\mathbf{y_{1}},\mathbf{y_{2}})\ \lbrack s_{1},s_{2},t_{1},t_{2}\rbrack$,
%
%
\end{itemize}
%
%
$X$ is non-empty, $\eta$ is a congruence on $(X,S)$ and $(X,S)/\eta$ is isomorphic to $(W,R)$.
%
%
\begin{proposition}[Relative elementary definability]\label{Syntactic:Restriction:Lemma:hereditarily:undecidable}
%
%
%Let ${\mathcal C}$ and ${\mathcal C}^{\prime}$ be classes of frames.
If ${\mathcal C}$ is relatively
%
%
%
%
%
%
%
%
%
%
%
%
\linebreak
%
%
%
%
%
%
%
%
%
%
%
%
elementary definable in ${\mathcal C}^{\prime}$ and $\Th({\mathcal C})$ is hereditarily undecidable then $\Th({\mathcal C}^{\prime})$ is hereditarily undecidable.
%
%
\end{proposition}
%
%
\begin{proof}
%
%
See~\cite[Page~$272$]{Ershov:1980}.
%
%
\medskip
%
%
\end{proof}
%
%
\begin{lemma}\label{lemma:relatively:elementary:definable:K2}
%
%
The class of all irreflexive symmetric frames with at least $2$~elements is relatively elementary definable in ${\mathcal K}_{2}$.
%
%
\end{lemma}
%
%
\begin{proof}
%
%
Let $(W,R)$ be an arbitrary irreflexive symmetric frame with at least $2$ elements.
Let $E$ be $\{\{s_{1},s_{2}\}:\ s_{1},s_{2}{\in}W$ and $s_{1}{R}s_{2}\}$.
Notice that for all $s_{1},s_{2}{\in}W$, if $\{s_{1},s_{2}\}{\in}E$ then $s_{1}{\not=}s_{2}$.
Without loss of generality, we can assume that $E$ and $W{\times}\{(3,4)\}$ are disjoint.
Let $A$ be $E{\cup}(W{\times}\{3,4\})$.
Let $B$ be $W{\times}\{1,2\}$.
Let ${\mathcal F}^{\rho}_{A,B}$ be the galaxy such that
%
%
\begin{itemize}
%
%
\item for all $\{s_{1},s_{2}\}{\in}E$, $\rho(\{s_{1},s_{2}\}){=}\{(s_{1},1),(s_{2},1)\}$,
%
%
\item for all $s{\in}W$, $\rho((s,3)){=}\{(s,1),(s,2)\}$,
%
%
\item for all $s{\in}W$, $\rho((s,4)){=}\{(s,1),(s,2)\}$.
%
%
\end{itemize}
%
%
%Since $\{2\}{\times}\mathbf{N}^{-}{\subseteq}\mathtt{S}_{\L}$, then ${\mathcal F}^{\rho}_{A,B}{\models}\L$.
Obviously, ${\mathcal F}^{\rho}_{A,B}$ is in ${\mathcal K}_{2}$.
%
%
\\
\\
%
%
Let $X$ be $\{(a_{1},a_{2}):\ a_{1},a_{2}{\in}W^{\rho}_{A,B}$ and ${\mathcal F}^{\rho}_{A,B}{\models}\mathbf{U}(\mathbf{x_{1}},\mathbf{x_{2}})\ \lbrack a_{1},a_{2}\rbrack\}$ where
%
%
\begin{itemize}
%
%
\item $\mathbf{U}(\mathbf{x_{1}},\mathbf{x_{2}})::=\mathbf{x_{1}}{\not=}\mathbf{x_{2}}\wedge\mathbf{R}(\mathbf{x_{1}},\mathbf{x_{2}})\wedge\mathbf{R}(\mathbf{x_{2}},\mathbf{x_{1}})\wedge\exists^{{=}2}\mathbf{y}(\mathbf{R}(\mathbf{y}){=}\{\mathbf{x_{1}},\mathbf{x_{2}}\}\wedge\neg\exists\mathbf{z}\mathbf{R}(\mathbf{z},\mathbf{y}))$.
%
%
\end{itemize}
%
%
Obviously, for all $s{\in}W$, $((s,1),(s,2)),((s,2),(s,1)){\in}X$.
Hence, $X$ is non-empty.
%
%
\\
\\
%
%
Obviously, for all $(a_{1},a_{2}){\in}X$, there exists $s{\in}W$ such that $\{a_{1},a_{2}\}{=}\{(s,1),(s,
%
%
%
%
%
%
%
%
%
%
%
%
$\linebreak$
%
%
%
%
%
%
%
%
%
%
%
%
2)\}$.
%
%
\\
\\
%
%
Let $\bowtie$ be the binary relation on $X$ such that for all $(a_{1},a_{2}),(b_{1},b_{2}){\in}X$, $(a_{1},a_{2}){\bowtie}
%
%
%
%
%
%
%
%
%
%
%
%
$\linebreak$
%
%
%
%
%
%
%
%
%
%
%
%
(b_{1},b_{2})$ if and only if ${\mathcal F}^{\rho}_{A,B}{\models}\mathbf{E}(\mathbf{x_{1}},\mathbf{x_{2}},\mathbf{y_{1}},\mathbf{y_{2}})\ \lbrack a_{1},a_{2},b_{1},b_{2}\rbrack$ where
%
%
\begin{itemize}
%
%
\item $\mathbf{E}(\mathbf{x_{1}},\mathbf{x_{2}},\mathbf{y_{1}},\mathbf{y_{2}})::=\mathbf{U}(\mathbf{x_{1}},\mathbf{x_{2}})\wedge\mathbf{U}(\mathbf{y_{1}},\mathbf{y_{2}})\wedge\{\mathbf{x_{1}},\mathbf{x_{2}}\}{=}\{\mathbf{y_{1}},\mathbf{y_{2}}\}$.
%
%
\end{itemize}
%
%
Obviously, $\bowtie$ is an equivalence relation on $X$.
For all $(a_{1},a_{2}){\in}X$, the equivalence class of $(a_{1},a_{2})$ modulo $\bowtie$ will be denoted $\lbrack(a_{1},a_{2})\rbrack$.
%
%
\\
\\
%
%
Since for all $(a_{1},a_{2}){\in}X$, there exists $s{\in}W$ such that $\{a_{1},a_{2}\}{=}\{(s,1),(s,2)\}$, then for all $(a_{1},a_{2}),(b_{1},b_{2}){\in}X$, $(a_{1},a_{2}){\bowtie}(b_{1},b_{2})$ if and only if there exists $s{\in}W$ such that $\{a_{1},a_{2}\}{=}\{(s,1),(s,2)\}$ and $\{b_{1},b_{2}\}{=}\{(s,1),(s,2)\}$.
%
%
\\
\\
%
%
Let $S$ be the binary relation on $X$ such that for all $(a_{1},a_{2}),(b_{1},b_{2}){\in}X$, $(a_{1},a_{2}){S}
%
%
%
%
%
%
%
%
%
%
%
%
$\linebreak$
%
%
%
%
%
%
%
%
%
%
%
%
(b_{1},b_{2})$ if and only if ${\mathcal F}^{\rho}_{A,B}{\models}\mathbf{A}(\mathbf{x_{1}},\mathbf{x_{2}},\mathbf{y_{1}},\mathbf{y_{2}})\ \lbrack a_{1},a_{2},b_{1},b_{2}\rbrack$ where
%
%
\begin{itemize}
%
%
\item $\mathbf{A}(\mathbf{x_{1}},\mathbf{x_{2}},\mathbf{y_{1}},\mathbf{y_{2}})::=\mathbf{U}(\mathbf{x_{1}},\mathbf{x_{2}})\wedge\mathbf{U}(\mathbf{y_{1}},\mathbf{y_{2}})\wedge\{\mathbf{x_{1}},\mathbf{x_{2}}\}{\cap}\{\mathbf{y_{1}},\mathbf{y_{2}}\}{=}\emptyset\wedge
%
%
%
%
%
%
%
%
%
%
%
%
$\linebreak$
%
%
%
%
%
%
%
%
%
%
%
%
\bigvee\{\exists^{{=}1}\mathbf{z}(\mathbf{R}(\mathbf{z}){=}\{\mathbf{x_{i}},\mathbf{y_{j}}\}\wedge\neg\exists\mathbf{t}\mathbf{R}(\mathbf{t},\mathbf{z})):\ 1{\leq}i,j{\leq}2\}$.
%
%
\end{itemize}
%
%
We claim that for all $(a_{1},a_{2}),(b_{1},b_{2}),(c_{1},c_{2}),(d_{1},d_{2}){\in}X$, if $(a_{1},a_{2}){\bowtie}(c_{1},c_{2})$ and $(b_{1},b_{2}){\bowtie}(d_{1},d_{2})$ then $(a_{1},a_{2}){S}(b_{1},b_{2})$ if and only if $(c_{1},c_{2}){S}(d_{1},d_{2})$.
If not, let $(a_{1},a_{2}),(b_{1},b_{2}),(c_{1},c_{2}),(d_{1},d_{2}){\in}X$ be such that $(a_{1},a_{2}){\bowtie}(c_{1},c_{2})$, $(b_{1},b_{2}){\bowtie}(d_{1},
%
%
%
%
%
%
%
%
%
%
%
%
$\linebreak$
%
%
%
%
%
%
%
%
%
%
%
%
d_{2})$ and either $(a_{1},a_{2}){S}(b_{1},b_{2})$ and not $(c_{1},c_{2}){S}(d_{1},d_{2})$, or not $(a_{1},a_{2}){S}(b_{1},b_{2})$ and $(c_{1},c_{2}){S}(d_{1},d_{2})$.
Without loss of generality, suppose $(a_{1},a_{2}){S}(b_{1},b_{2})$ and not $(c_{1},c_{2}){S}(d_{1},d_{2})$.
Thus, ${\mathcal F}^{\rho}_{A,B}{\models}\mathbf{A}(\mathbf{x_{1}},\mathbf{x_{2}},\mathbf{y_{1}},\mathbf{y_{2}})\ \lbrack a_{1},a_{2},b_{1},b_{2}\rbrack$ and ${\mathcal F}^{\rho}_{A,B}{\not\models}
%
%
%
%
%
%
%
%
%
%
%
%
$\linebreak$
%
%
%
%
%
%
%
%
%
%
%
%
\mathbf{A}(\mathbf{x_{1}},\mathbf{x_{2}},\mathbf{y_{1}},\mathbf{y_{2}})\ \lbrack c_{1},c_{2},d_{1},d_{2}\rbrack$.
Consequently,
%
%
\begin{description}
%
%
\item[$(C_{1})$] $\{a_{1},a_{2}\}{\cap}\{b_{1},b_{2}\}{=}\emptyset$,
%
%
\item[$(C_{2})$] there exists $i,j{\in}\{1,2\}$ such that for exactly $1$~$e$ in ${\mathcal F}_{A,B}^{\rho}$, $R_{A,B}^{\rho}(e){=}\{a_{i},b_{j}\}$ and ${R_{A,B}^{\rho}}^{{-}1}(e){=}\emptyset$,
%
%
\item[$(C_{3})$] at least one of the following conditions holds:
%
%
\begin{description}
%
%
\item[$(C_{3}^{a})$] $\{c_{1},c_{2}\}{\cap}\{d_{1},d_{2}\}{\not=}\emptyset$,
%
%
\item[$(C_{3}^{b})$] for all $k,l{\in}\{1,2\}$, either for all $f$ in ${\mathcal F}_{A,B}^{\rho}$, if $R_{A,B}^{\rho}(f){=}\{c_{k},d_{l}\}$ then ${R_{A,B}^{\rho}}^{{-}1}(f){\not=}\emptyset$, or there exists $f^{\prime},f^{\prime\prime}$ in ${\mathcal F}_{A,B}^{\rho}$ such that $f^{\prime}{\not=}f^{\prime\prime}$,
%
%
%
%
%
%
%
%
%
%
%
%
\linebreak$
%
%
%
%
%
%
%
%
%
%
%
%
R_{A,B}^{\rho}(f^{\prime}){=}\{c_{k},d_{l}\}$, $R_{A,B}^{\rho}(f^{\prime\prime}){=}\{c_{k},d_{l}\}$, ${R_{A,B}^{\rho}}^{{-}1}(f^{\prime}){=}\emptyset$ and
%
%
%
%
%
%
%
%
%
%
%
%
\linebreak$
%
%
%
%
%
%
%
%
%
%
%
%
{R_{A,B}^{\rho}}^{{-}1}(f^{\prime\prime}){=}\emptyset$.
%
%
\end{description}
%
%
\end{description}
%
%
%Since $(a_{1},a_{2}),(b_{1},b_{2}),(c_{1},c_{2}),(d_{1},d_{2}){\in}X$, ${\mathcal F}^{\rho}_{A,B}{\models}\mathbf{U}(\mathbf{x_{1}},\mathbf{x_{2}})\ \lbrack a_{1},a_{2}\rbrack$, ${\mathcal F}^{\rho}_{A,B}{\models}\mathbf{U}(\mathbf{x_{1}},\mathbf{x_{2}})\ \lbrack b_{1},b_{2}\rbrack$, ${\mathcal F}^{\rho}_{A,B}{\models}\mathbf{U}(\mathbf{x_{1}},\mathbf{x_{2}})\ \lbrack c_{1},c_{2}\rbrack$ and ${\mathcal F}^{\rho}_{A,B}{\models}\mathbf{U}(\mathbf{x_{1}},\mathbf{x_{2}})\ \lbrack d_{1},d_{2}\rbrack$.
Let $s,t,u,v{\in}W$ be such that $\{a_{1},a_{2}\}{=}\{(s,1),(s,2)\}$, $\{b_{1},b_{2}\}{=}\{(t,1),(t,2)\}$, $\{c_{1},c_{2}\}{=}\{(u,1),(u,2)\}$ and $\{d_{1},d_{2}\}{=}\{(v,1),(v,2)\}$.
Since $(a_{1},a_{2}){\bowtie}(c_{1},c_{2})$ and $(b_{1},b_{2}){\bowtie}(d_{1},d_{2})$, then ${\mathcal F}^{\rho}_{A,B}{\models}\mathbf{E}(\mathbf{x_{1}},\mathbf{x_{2}},\mathbf{y_{1}},\mathbf{y_{2}})\ \lbrack a_{1},a_{2},c_{1},c_{2}\rbrack$ and ${\mathcal F}^{\rho}_{A,B}{\models}\mathbf{E}(\mathbf{x_{1}},
%
%
%
%
%
%
%
%
%
%
%
%
$\linebreak$
%
%
%
%
%
%
%
%
%
%
%
%
\mathbf{x_{2}},\mathbf{y_{1}},\mathbf{y_{2}})\ \lbrack b_{1},b_{2},d_{1},d_{2}\rbrack$.
Hence, %${\mathcal F}^{\rho}_{A,B}{\models}\mathbf{U}(\mathbf{x_{1}},\mathbf{x_{2}})\ \lbrack a_{1},a_{2}\rbrack$, ${\mathcal F}^{\rho}_{A,B}{\models}\mathbf{U}(\mathbf{x_{1}},\mathbf{x_{2}})\ \lbrack b_{1},b_{2}\rbrack$, ${\mathcal F}^{\rho}_{A,B}{\models}\mathbf{U}(\mathbf{x_{1}},\mathbf{x_{2}})\ \lbrack c_{1},c_{2}\rbrack$ and ${\mathcal F}^{\rho}_{A,B}{\models}\mathbf{U}(\mathbf{x_{1}},\mathbf{x_{2}})\ \lbrack d_{1},d_{2}\rbrack$.
%Moreover,
$\{a_{1},a_{2}\}{=}\{c_{1},c_{2}\}$ and $\{b_{1},b_{2}\}{=}\{d_{1},d_{2}\}$.
Since $\{a_{1},a_{2}\}{=}\{(s,1),(s,2)\}$, $\{b_{1},b_{2}\}{=}\{(t,1),(t,2)\}$, $\{c_{1},c_{2}\}{=}\{(u,1),(u,2)\}$ and $\{d_{1},
%
%
%
%
%
%
%
%
%
%
%
%
$\linebreak$
%
%
%
%
%
%
%
%
%
%
%
%
d_{2}\}{=}\{(v,1),(v,2)\}$, then $s{=}u$ and $t{=}v$.
Moreover, by $(C_{1})$, $s{\not=}t$.
In other respect, by $(C_{2})$, there exists $k^{\prime},l^{\prime}{\in}\{1,2\}$ such that $a_{i}{=}(s,k^{\prime})$, $b_{j}{=}(t,l^{\prime})$ and for exactly $1$~$e$ in ${\mathcal F}_{A,B}^{\rho}$, $R_{A,B}^{\rho}(e){=}\{(s,k^{\prime}),(t,l^{\prime})\}$ and ${R_{A,B}^{\rho}}^{{-}1}(e){=}\emptyset$.
Since $\{c_{1},c_{2}\}{=}\{(u,1),(u,2)\}$, $s{=}u$, $\{d_{1},d_{2}\}{=}\{(v,1),(v,2)\}$, $t{=}v$ and $s{\not=}t$, then $\{c_{1},
%
%
%
%
%
%
%
%
%
%
%
%
$\linebreak$
%
%
%
%
%
%
%
%
%
%
%
%
c_{2}\}{\cap}\{d_{1},d_{2}\}{=}\emptyset$.
Moreover, there exists $k,l{\in}\{1,2\}$ such that $c_{k}{=}(s,k^{\prime})$ and $d_{l}{=}(t,l^{\prime})$.
Hence, $(C_{3}^{a})$ does not hold.
Thus, by $(C_{3})$, $(C_{3}^{b})$ holds and either for all $f$ in ${\mathcal F}_{A,B}^{\rho}$, if $R_{A,B}^{\rho}(f){=}\{(s,k^{\prime}),(t,l^{\prime})\}$ then ${R_{A,B}^{\rho}}^{{-}1}(f){\not=}\emptyset$, or there exists $f^{\prime},f^{\prime\prime}$ in ${\mathcal F}_{A,B}^{\rho}$ such that $f^{\prime}{\not=}f^{\prime\prime}$, $R_{A,B}^{\rho}(f^{\prime}){=}\{(s,k^{\prime}),(t,l^{\prime})\}$, $R_{A,B}^{\rho}(f^{\prime\prime}){=}\{(s,k^{\prime}),(t,l^{\prime})\}$, ${R_{A,B}^{\rho}}^{{-}1}(f^{\prime}){=}\emptyset$ and ${R_{A,B}^{\rho}}^{{-}1}(f^{\prime\prime}){=}\emptyset$.
In the former case, since $R_{A,B}^{\rho}(e){=}\{(s,k^{\prime}),
%
%
%
%
%
%
%
%
%
%
%
%
$\linebreak$
%
%
%
%
%
%
%
%
%
%
%
%
(t,l^{\prime})\}$, then ${R_{A,B}^{\rho}}^{{-}1}(e){\not=}\emptyset$: a contradiction.
In the latter case, since $a_{i}{=}(s,k^{\prime})$, $b_{j}{=}(t,l^{\prime})$, then there exists $f^{\prime},f^{\prime\prime}$ in ${\mathcal F}_{A,B}^{\rho}$ such that $f^{\prime}{\not=}f^{\prime\prime}$, $R_{A,B}^{\rho}(f^{\prime}){=}\{a_{i},b_{j}\}$, $R_{A,B}^{\rho}(f^{\prime\prime}){=}\{a_{i},b_{j}\}$, ${R_{A,B}^{\rho}}^{{-}1}(f^{\prime}){=}\emptyset$ and ${R_{A,B}^{\rho}}^{{-}1}(f^{\prime\prime}){=}\emptyset$: a contradiction.
Consequently, for all $(a_{1},a_{2}),(b_{1},b_{2}),(c_{1},c_{2}),(d_{1},d_{2}){\in}X$, if $(a_{1},a_{2}){\bowtie}(c_{1},c_{2})$ and $(b_{1},b_{2}){\bowtie}(d_{1},d_{2})$ then $(a_{1},a_{2}){S}(b_{1},b_{2})$ if and only if $(c_{1},c_{2}){S}(d_{1},d_{2})$.
%
%
\\
\\
%
%
Let $(W^{\prime},R^{\prime})$ be the quotient of $(X,S)$ modulo $\bowtie$~---~i.e.
%
%
\begin{itemize}
%
%
\item $W^{\prime}$ is the quotient set of $X$ modulo $\bowtie$,
%
%
\item $R^{\prime}$ is the binary relation on $W^{\prime}$ such that for all $(a_{1},a_{2}),(b_{1},b_{2}){\in}X$, $\lbrack(a_{1},a_{2})\rbrack{R^{\prime}}\lbrack(b_{1},b_{2})\rbrack$ if and only if $(a_{1},a_{2}){S}(b_{1},b_{2})$.
%
%
\end{itemize}
%
%
Since for all $(a_{1},a_{2}){\in}X$, there exists $s{\in}W$ such that $\{a_{1},a_{2}\}{=}\{(s,1),(s,2)\}$ and for all $(a_{1},a_{2}),(b_{1},b_{2}){\in}X$, $(a_{1},a_{2}){\bowtie}(b_{1},b_{2})$ if and only if there exists $s{\in}W$ such that $\{a_{1},a_{2}\}{=}\{(s,1),(s,2)\}$ and $\{b_{1},b_{2}\}{=}\{(s,1),(s,2)\}$, then let $\kappa:\ W^{\prime}{\longrightarrow}W$ be the function such that for all $(a_{1},a_{2}){\in}X$, $\kappa(\lbrack(a_{1},a_{2})\rbrack)$ is the unique $s{\in}W$ such that $\{a_{1},a_{2}\}{=}\{(s,1),(s,2)\}$.
Obviously, $\kappa$ is bijective.
%for all $s{\in}W$, $((s,1),(s,2)),((s,2),(s,1)){\in}X$ and
%
\\
\\
%
%
We claim that for all $(a_{1},a_{2}),(b_{1},b_{2}){\in}X$, $\lbrack(a_{1},a_{2})\rbrack{R^{\prime}}\lbrack(b_{1},b_{2})\rbrack$ if and only if $\kappa(\lbrack(a_{1},a_{2})\rbrack){R}\kappa(\lbrack(b_{1},b_{2})\rbrack)$.
Let $(a_{1},a_{2}),(b_{1},b_{2}){\in}X$.
Let $s$ be the unique element in $W$ such that $\{a_{1},a_{2}\}{=}\{(s,1),(s,2)\}$ and $t$ be the unique element in $W$ such that $\{b_{1},b_{2}\}{=}\{(t,1),(t,2)\}$.
Obviously, the following conditions are equivalent:
%
%
\begin{itemize}
%
%
\item $\lbrack(a_{1},a_{2})\rbrack{R^{\prime}}\lbrack(b_{1},b_{2})\rbrack$,
%
%
\item $(a_{1},a_{2}){S}(b_{1},b_{2})$,
%
%
\item ${\mathcal F}^{\rho}_{A,B}{\models}\mathbf{A}(\mathbf{x_{1}},\mathbf{x_{2}},\mathbf{y_{1}},\mathbf{y_{2}})\ \lbrack a_{1},a_{2},b_{1},b_{2}\rbrack$,
%
%
\item $s{\not=}t$ and there exists $i,j{\in}\{1,2\}$ such that for exactly $1$~$e$ in ${\mathcal F}_{A,B}^{\rho}$, $R_{A,B}^{\rho}(e){=}
%
%
%
%
%
%
%
%
%
%
%
%
$\linebreak$
%
%
%
%
%
%
%
%
%
%
%
%
\{a_{i},b_{j}\}$ and ${R_{A,B}^{\rho}}^{{-}1}(e){=}\emptyset$,
%
%
\item $\kappa(\lbrack(a_{1},a_{2})\rbrack){R}\kappa(\lbrack(b_{1},b_{2})\rbrack)$.
%
%
\end{itemize}
%
%
Hence, for all $(a_{1},a_{2}),(b_{1},b_{2}){\in}X$, $\lbrack(a_{1},a_{2})\rbrack{R^{\prime}}\lbrack(b_{1},b_{2})\rbrack$ if and only if $\kappa(\lbrack(a_{1},a_{2})\rbrack){R}
%
%
%
%
%
%
%
%
%
%
%
%
$\linebreak$
%
%
%
%
%
%
%
%
%
%
%
%
\kappa(\lbrack(b_{1},b_{2})\rbrack)$.
%
%
\\
\\
%
%
All in all, we have proved that, $(W^{\prime},R^{\prime})$ is isomorphic with $(W,R)$.
%
%
\\
\\
%
%
Thus, we have shown that the class of all irreflexive symmetric frames with at least $2$ elements is relatively elementary definable in ${\mathcal K}_{2}$.
%
%
\medskip
%
%
\end{proof}
%
%
\begin{lemma}\label{lemma:relatively:elementary:definable:L2}
%
%
The class of all irreflexive symmetric frames with at least $2$~elements is relatively elementary definable in ${\mathcal L}_{2}$.
%
%
\end{lemma}
%
%
\begin{proof}
%
%
Let $(W,R)$ be an arbitrary irreflexive symmetric frame with at least $2$ elements.
Let $E$ be $\{\{s_{1},s_{2}\}:\ s_{1},s_{2}{\in}W$ and $s_{1}{R}s_{2}\}$.
Notice that for all $s_{1},s_{2}{\in}W$, if $\{s_{1},s_{2}\}{\in}E$ then $s_{1}{\not=}s_{2}$.
Without loss of generality, we can assume that $E$ and $W{\times}\{(3,4)\}$ are disjoint.
Let $A$ be $E{\cup}(W{\times}\{3,4\})$.
Let $B$ be $W{\times}\{1,2\}$.
Let ${\mathcal F}^{\rho}_{A,B}$ be the galaxy such that
%
%
\begin{itemize}
%
%
\item for all $\{s_{1},s_{2}\}{\in}E$, $\rho(\{s_{1},s_{2}\}){=}B{\setminus}\{(s_{1},1),(s_{2},1)\}$,
%
%
\item for all $s{\in}W$, $\rho((s,3)){=}B{\setminus}\{(s,1),(s,2)\}$,
%
%
\item for all $s{\in}W$, $\rho((s,4)){=}B{\setminus}\{(s,1),(s,2)\}$.
%
%
\end{itemize}
%
%
%Since $\mathbf{N}^{+}{\times}\{2\}{\subseteq}\mathtt{S}_{\L}$, then ${\mathcal F}^{\rho}_{A,B}{\models}\L$.
Obviously, ${\mathcal F}^{\rho}_{A,B}$ is in ${\mathcal L}_{2}$.
%
%
\\
\\
%
%
Let $X$ be $\{(a_{1},a_{2}):\ a_{1},a_{2}{\in}W^{\rho}_{A,B}$ and ${\mathcal F}^{\rho}_{A,B}{\models}\mathbf{U}(\mathbf{x_{1}},\mathbf{x_{2}})\ \lbrack a_{1},a_{2}\rbrack\}$ where
%
%
\begin{itemize}
%
%
\item $\mathbf{U}(\mathbf{x_{1}},\mathbf{x_{2}})::=\mathbf{x_{1}}{\not=}\mathbf{x_{2}}\wedge\mathbf{R}(\mathbf{x_{1}},\mathbf{x_{2}})\wedge\mathbf{R}(\mathbf{x_{2}},\mathbf{x_{1}})\wedge\exists^{{=}2}\mathbf{y}(\mathbf{R}(\mathbf{y}){=}\{\mathbf{x_{1}},\mathbf{x_{2}}\}^{\mathbf{c}}\wedge\neg\exists\mathbf{z}\mathbf{R}(\mathbf{z},\mathbf{y}))$.
%
%
\end{itemize}
%
%
Obviously, for all $s{\in}W$, $((s,1),(s,2)),((s,2),(s,1)){\in}X$.
Hence, $X$ is non-empty.
%
%
\\
\\
%
%
Obviously, for all $(a_{1},a_{2}){\in}X$, there exists $s{\in}W$ such that $\{a_{1},a_{2}\}{=}\{(s,1),(s,
%
%
%
%
%
%
%
%
%
%
%
%
$\linebreak$
%
%
%
%
%
%
%
%
%
%
%
%
2)\}$.
%
%
\\
\\
%
%
Let $\bowtie$ be the binary relation on $X$ such that for all $(a_{1},a_{2}),(b_{1},b_{2}){\in}X$, $(a_{1},a_{2}){\bowtie}
%
%
%
%
%
%
%
%
%
%
%
%
$\linebreak$
%
%
%
%
%
%
%
%
%
%
%
%
(b_{1},b_{2})$ if and only if ${\mathcal F}^{\rho}_{A,B}{\models}\mathbf{E}(\mathbf{x_{1}},\mathbf{x_{2}},\mathbf{y_{1}},\mathbf{y_{2}})\ \lbrack a_{1},a_{2},b_{1},b_{2}\rbrack$ where
%
%
\begin{itemize}
%
%
\item $\mathbf{E}(\mathbf{x_{1}},\mathbf{x_{2}},\mathbf{y_{1}},\mathbf{y_{2}})::=\mathbf{U}(\mathbf{x_{1}},\mathbf{x_{2}})\wedge\mathbf{U}(\mathbf{y_{1}},\mathbf{y_{2}})\wedge\{\mathbf{x_{1}},\mathbf{x_{2}}\}{=}\{\mathbf{y_{1}},\mathbf{y_{2}}\}$.
%
%
\end{itemize}
%
%
Obviously, $\bowtie$ is an equivalence relation on $X$.
For all $(a_{1},a_{2}){\in}X$, the equivalence class of $(a_{1},a_{2})$ modulo $\bowtie$ will be denoted $\lbrack(a_{1},a_{2})\rbrack$.
%
%
\\
\\
%
%
Since for all $(a_{1},a_{2}){\in}X$, there exists $s{\in}W$ such that $\{a_{1},a_{2}\}{=}\{(s,1),(s,2)\}$, then for all $(a_{1},a_{2}),(b_{1},b_{2}){\in}X$, $(a_{1},a_{2}){\bowtie}(b_{1},b_{2})$ if and only if there exists $s{\in}W$ such that $\{a_{1},a_{2}\}{=}\{(s,1),(s,2)\}$ and $\{b_{1},b_{2}\}{=}\{(s,1),(s,2)\}$.
%
%
\\
\\
%
%
Let $S$ be the binary relation on $X$ such that for all $(a_{1},a_{2}),(b_{1},b_{2}){\in}X$, $(a_{1},a_{2}){S}
%
%
%
%
%
%
%
%
%
%
%
%
$\linebreak$
%
%
%
%
%
%
%
%
%
%
%
%
(b_{1},b_{2})$ if and only if ${\mathcal F}^{\rho}_{A,B}{\models}\mathbf{A}(\mathbf{x_{1}},\mathbf{x_{2}},\mathbf{y_{1}},\mathbf{y_{2}})\ \lbrack a_{1},a_{2},b_{1},b_{2}\rbrack$ where
%
%
\begin{itemize}
%
%
\item $\mathbf{A}(\mathbf{x_{1}},\mathbf{x_{2}},\mathbf{y_{1}},\mathbf{y_{2}})::=\mathbf{U}(\mathbf{x_{1}},\mathbf{x_{2}})\wedge\mathbf{U}(\mathbf{y_{1}},\mathbf{y_{2}})\wedge\{\mathbf{x_{1}},\mathbf{x_{2}}\}{\cap}\{\mathbf{y_{1}},\mathbf{y_{2}}\}{=}\emptyset\wedge
%
%
%
%
%
%
%
%
%
%
%
%
$\linebreak$
%
%
%
%
%
%
%
%
%
%
%
%
\bigvee\{\exists^{{=}1}\mathbf{z}(\mathbf{R}(\mathbf{z}){=}\{\mathbf{x_{i}},\mathbf{y_{j}}\}^{\mathbf{c}}\wedge\neg\exists\mathbf{t}\mathbf{R}(\mathbf{t},\mathbf{z})):\ 1{\leq}i,j{\leq}2\}$.
%
%
\end{itemize}
%
%
We claim that for all $(a_{1},a_{2}),(b_{1},b_{2}),(c_{1},c_{2}),(d_{1},d_{2}){\in}X$, if $(a_{1},a_{2}){\bowtie}(c_{1},c_{2})$ and $(b_{1},b_{2}){\bowtie}(d_{1},d_{2})$ then $(a_{1},a_{2}){S}(b_{1},b_{2})$ if and only if $(c_{1},c_{2}){S}(d_{1},d_{2})$.
If not, let $(a_{1},a_{2}),(b_{1},b_{2}),(c_{1},c_{2}),(d_{1},d_{2}){\in}X$ be such that $(a_{1},a_{2}){\bowtie}(c_{1},c_{2})$, $(b_{1},b_{2}){\bowtie}(d_{1},
%
%
%
%
%
%
%
%
%
%
%
%
$\linebreak$
%
%
%
%
%
%
%
%
%
%
%
%
d_{2})$ and either $(a_{1},a_{2}){S}(b_{1},b_{2})$ and not $(c_{1},c_{2}){S}(d_{1},d_{2})$, or not $(a_{1},a_{2}){S}(b_{1},b_{2})$ and $(c_{1},c_{2}){S}(d_{1},d_{2})$.
Without loss of generality, suppose $(a_{1},a_{2}){S}(b_{1},b_{2})$ and not $(c_{1},c_{2}){S}(d_{1},d_{2})$.
Thus, ${\mathcal F}^{\rho}_{A,B}{\models}\mathbf{A}(\mathbf{x_{1}},\mathbf{x_{2}},\mathbf{y_{1}},\mathbf{y_{2}})\ \lbrack a_{1},a_{2},b_{1},b_{2}\rbrack$ and ${\mathcal F}^{\rho}_{A,B}{\not\models}
%
%
%
%
%
%
%
%
%
%
%
%
$\linebreak$
%
%
%
%
%
%
%
%
%
%
%
%
\mathbf{A}(\mathbf{x_{1}},\mathbf{x_{2}},\mathbf{y_{1}},\mathbf{y_{2}})\ \lbrack c_{1},c_{2},d_{1},d_{2}\rbrack$.
Consequently,
%
%
\begin{description}
%
%
\item[$(C_{1})$] $\{a_{1},a_{2}\}{\cap}\{b_{1},b_{2}\}{=}\emptyset$,
%
%
\item[$(C_{2})$] there exists $i,j{\in}\{1,2\}$ such that for exactly $1$~$e$ in ${\mathcal F}_{A,B}^{\rho}$, $R_{A,B}^{\rho}(e){=}B{\setminus}\{a_{i},
%
%
%
%
%
%
%
%
%
%
%
%
$\linebreak$
%
%
%
%
%
%
%
%
%
%
%
%
b_{j}\}$ and ${R_{A,B}^{\rho}}^{{-}1}(e){=}\emptyset$,
%
%
\item[$(C_{3})$] at least one of the following conditions holds:
%
%
\begin{description}
%
%
\item[$(C_{3}^{a})$] $\{c_{1},c_{2}\}{\cap}\{d_{1},d_{2}\}{\not=}\emptyset$,
%
%
\item[$(C_{3}^{b})$] for all $k,l{\in}\{1,2\}$, either for all $f$ in ${\mathcal F}_{A,B}^{\rho}$, if $R_{A,B}^{\rho}(f){=}B{\setminus}\{c_{k},d_{l}\}$ then ${R_{A,B}^{\rho}}^{{-}1}(f){\not=}\emptyset$, or there exists $f^{\prime},f^{\prime\prime}$ in ${\mathcal F}_{A,B}^{\rho}$ such that $f^{\prime}{\not=}f^{\prime\prime}$, $R_{A,B}^{\rho}(f^{\prime}){=}B{\setminus}\{c_{k},d_{l}\}$, $R_{A,B}^{\rho}(f^{\prime\prime}){=}B{\setminus}\{c_{k},d_{l}\}$, ${R_{A,B}^{\rho}}^{{-}1}(f^{\prime}){=}\emptyset$ and
%
%
%
%
%
%
%
%
%
%
%
%
\linebreak$
%
%
%
%
%
%
%
%
%
%
%
%
{R_{A,B}^{\rho}}^{{-}1}(f^{\prime\prime}){=}\emptyset$.
%
%
\end{description}
%
%
\end{description}
%
%
%Since $(a_{1},a_{2}),(b_{1},b_{2}),(c_{1},c_{2}),(d_{1},d_{2}){\in}X$, ${\mathcal F}^{\rho}_{A,B}{\models}\mathbf{U}(\mathbf{x_{1}},\mathbf{x_{2}})\ \lbrack a_{1},a_{2}\rbrack$, ${\mathcal F}^{\rho}_{A,B}{\models}\mathbf{U}(\mathbf{x_{1}},\mathbf{x_{2}})\ \lbrack b_{1},b_{2}\rbrack$, ${\mathcal F}^{\rho}_{A,B}{\models}\mathbf{U}(\mathbf{x_{1}},\mathbf{x_{2}})\ \lbrack c_{1},c_{2}\rbrack$ and ${\mathcal F}^{\rho}_{A,B}{\models}\mathbf{U}(\mathbf{x_{1}},\mathbf{x_{2}})\ \lbrack d_{1},d_{2}\rbrack$.
Let $s,t,u,v{\in}W$ be such that $\{a_{1},a_{2}\}{=}\{(s,1),(s,2)\}$, $\{b_{1},b_{2}\}{=}\{(t,1),(t,2)\}$, $\{c_{1},c_{2}\}{=}\{(u,1),(u,2)\}$ and $\{d_{1},d_{2}\}{=}\{(v,1),(v,2)\}$.
Since $(a_{1},a_{2}){\bowtie}(c_{1},c_{2})$ and $(b_{1},b_{2}){\bowtie}(d_{1},d_{2})$, then ${\mathcal F}^{\rho}_{A,B}{\models}\mathbf{E}(\mathbf{x_{1}},\mathbf{x_{2}},\mathbf{y_{1}},\mathbf{y_{2}})\ \lbrack a_{1},a_{2},c_{1},c_{2}\rbrack$ and ${\mathcal F}^{\rho}_{A,B}{\models}\mathbf{E}(\mathbf{x_{1}},
%
%
%
%
%
%
%
%
%
%
%
%
$\linebreak$
%
%
%
%
%
%
%
%
%
%
%
%
\mathbf{x_{2}},\mathbf{y_{1}},\mathbf{y_{2}})\ \lbrack b_{1},b_{2},d_{1},d_{2}\rbrack$.
Hence, %${\mathcal F}^{\rho}_{A,B}{\models}\mathbf{U}(\mathbf{x_{1}},\mathbf{x_{2}})\ \lbrack a_{1},a_{2}\rbrack$, ${\mathcal F}^{\rho}_{A,B}{\models}\mathbf{U}(\mathbf{x_{1}},\mathbf{x_{2}})\ \lbrack b_{1},b_{2}\rbrack$, ${\mathcal F}^{\rho}_{A,B}{\models}\mathbf{U}(\mathbf{x_{1}},\mathbf{x_{2}})\ \lbrack c_{1},c_{2}\rbrack$ and ${\mathcal F}^{\rho}_{A,B}{\models}\mathbf{U}(\mathbf{x_{1}},\mathbf{x_{2}})\ \lbrack d_{1},d_{2}\rbrack$.
%Moreover,
$\{a_{1},a_{2}\}{=}\{c_{1},c_{2}\}$ and $\{b_{1},b_{2}\}{=}\{d_{1},d_{2}\}$.
Since $\{a_{1},a_{2}\}{=}\{(s,1),(s,2)\}$, $\{b_{1},b_{2}\}{=}\{(t,1),(t,2)\}$, $\{c_{1},c_{2}\}{=}\{(u,1),(u,2)\}$ and $\{d_{1},
%
%
%
%
%
%
%
%
%
%
%
%
$\linebreak$
%
%
%
%
%
%
%
%
%
%
%
%
d_{2}\}{=}\{(v,1),(v,2)\}$, then $s{=}u$ and $t{=}v$.
Moreover, by $(C_{1})$, $s{\not=}t$.
In other respect, by $(C_{2})$, there exists $k^{\prime},l^{\prime}{\in}\{1,2\}$ such that $a_{i}{=}(s,k^{\prime})$, $b_{j}{=}(t,l^{\prime})$ and for exactly $1$~$e$ in ${\mathcal F}_{A,B}^{\rho}$, $R_{A,B}^{\rho}(e){=}B{\setminus}\{(s,k^{\prime}),(t,l^{\prime})\}$ and ${R_{A,B}^{\rho}}^{{-}1}(e){=}\emptyset$.
Since $\{c_{1},c_{2}\}{=}\{(u,1),(u,2)\}$, $s{=}u$, $\{d_{1},d_{2}\}{=}\{(v,1),(v,2)\}$, $t{=}v$ and $s{\not=}t$, then $\{c_{1},
%
%
%
%
%
%
%
%
%
%
%
%
$\linebreak$
%
%
%
%
%
%
%
%
%
%
%
%
c_{2}\}{\cap}\{d_{1},d_{2}\}{=}\emptyset$.
Moreover, there exists $k,l{\in}\{1,2\}$ such that $c_{k}{=}(s,k^{\prime})$ and $d_{l}{=}(t,l^{\prime})$.
Thus $(C_{3}^{a})$ does not hold.
Consequently, by $(C_{3})$, $(C_{3}^{b})$ holds and either for all $f$ in ${\mathcal F}_{A,B}^{\rho}$, if $R_{A,B}^{\rho}(f){=}B{\setminus}\{(s,k^{\prime}),(t,l^{\prime})\}$ then ${R_{A,B}^{\rho}}^{{-}1}(f){\not=}\emptyset$, or there exists $f^{\prime},f^{\prime\prime}$ in ${\mathcal F}_{A,B}^{\rho}$ such that $f^{\prime}{\not=}f^{\prime\prime}$, $R_{A,B}^{\rho}(f^{\prime}){=}B{\setminus}\{(s,k^{\prime}),(t,l^{\prime})\}$, $R_{A,B}^{\rho}(f^{\prime\prime}){=}
%
%
%
%
%
%
%
%
%
%
%
%
$\linebreak$
%
%
%
%
%
%
%
%
%
%
%
%
B{\setminus}\{(s,k^{\prime}),(t,l^{\prime})\}$, ${R_{A,B}^{\rho}}^{{-}1}(f^{\prime}){=}\emptyset$ and ${R_{A,B}^{\rho}}^{{-}1}(f^{\prime\prime}){=}\emptyset$.
In the former case, since $R_{A,B}^{\rho}(e){=}B{\setminus}\{(s,k^{\prime}),(t,l^{\prime})\}$, then ${R_{A,B}^{\rho}}^{{-}1}(e){\not=}\emptyset$: a contradiction.
In the latter case, since $a_{i}{=}(s,k^{\prime})$, $b_{j}{=}(t,l^{\prime})$, then there exists $f^{\prime},f^{\prime\prime}$ in ${\mathcal F}_{A,B}^{\rho}$ such that $f^{\prime}{\not=}f^{\prime\prime}$, $R_{A,B}^{\rho}(f^{\prime}){=}B{\setminus}\{a_{i},b_{j}\}$, $R_{A,B}^{\rho}(f^{\prime\prime}){=}B{\setminus}\{a_{i},b_{j}\}$, ${R_{A,B}^{\rho}}^{{-}1}(f^{\prime}){=}\emptyset$ and ${R_{A,B}^{\rho}}^{{-}1}(f^{\prime\prime}){=}
%
%
%
%
%
%
%
%
%
%
%
%
$\linebreak$
%
%
%
%
%
%
%
%
%
%
%
%
\emptyset$: a contradiction.
Hence, for all $(a_{1},a_{2}),(b_{1},b_{2}),(c_{1},c_{2}),(d_{1},d_{2}){\in}X$, if $(a_{1},a_{2})
%
%
%
%
%
%
%
%
%
%
%
%
$\linebreak$
%
%
%
%
%
%
%
%
%
%
%
%
{\bowtie}(c_{1},c_{2})$ and $(b_{1},b_{2}){\bowtie}(d_{1},d_{2})$ then $(a_{1},a_{2}){S}(b_{1},b_{2})$ if and only if $(c_{1},c_{2}){S}(d_{1},
%
%
%
%
%
%
%
%
%
%
%
%
$\linebreak$
%
%
%
%
%
%
%
%
%
%
%
%
d_{2})$.
%
%
\\
\\
%
%
Let $(W^{\prime},R^{\prime})$ be the quotient of $(X,S)$ modulo $\bowtie$~---~i.e.
%
%
\begin{itemize}
%
%
\item $W^{\prime}$ is the quotient set of $X$ modulo $\bowtie$,
%
%
\item $R^{\prime}$ is the binary relation on $W^{\prime}$ such that for all $(a_{1},a_{2}),(b_{1},b_{2}){\in}X$, $\lbrack(a_{1},a_{2})\rbrack{R^{\prime}}\lbrack(b_{1},b_{2})\rbrack$ if and only if $(a_{1},a_{2}){S}(b_{1},b_{2})$.
%
%
\end{itemize}
%
%
Since for all $(a_{1},a_{2}){\in}X$, there exists $s{\in}W$ such that $\{a_{1},a_{2}\}{=}\{(s,1),(s,2)\}$ and for all $(a_{1},a_{2}),(b_{1},b_{2}){\in}X$, $(a_{1},a_{2}){\bowtie}(b_{1},b_{2})$ if and only if there exists $s{\in}W$ such that $\{a_{1},a_{2}\}{=}\{(s,1),(s,2)\}$ and $\{b_{1},b_{2}\}{=}\{(s,1),(s,2)\}$, then let $\kappa:\ W^{\prime}{\longrightarrow}W$ be the function such that for all $(a_{1},a_{2}){\in}X$, $\kappa(\lbrack(a_{1},a_{2})\rbrack)$ is the unique $s{\in}W$ such that $\{a_{1},a_{2}\}{=}\{(s,1),(s,2)\}$.
Obviously, $\kappa$ is bijective.
%for all $s{\in}W$, $((s,1),(s,2)),((s,2),(s,1)){\in}X$ and
%
\\
\\
%
%
We claim that for all $(a_{1},a_{2}),(b_{1},b_{2}){\in}X$, $\lbrack(a_{1},a_{2})\rbrack{R^{\prime}}\lbrack(b_{1},b_{2})\rbrack$ if and only if $\kappa(\lbrack(a_{1},a_{2})\rbrack){R}\kappa(\lbrack(b_{1},b_{2})\rbrack)$.
Let $(a_{1},a_{2}),(b_{1},b_{2}){\in}X$.
Let $s$ be the unique element in $W$ such that $\{a_{1},a_{2}\}{=}\{(s,1),(s,2)\}$ and $t$ be the unique element in $W$ such that $\{b_{1},b_{2}\}{=}\{(t,1),(t,2)\}$.
Obviously, the following conditions are equivalent:
%
%
\begin{itemize}
%
%
\item $\lbrack(a_{1},a_{2})\rbrack{R^{\prime}}\lbrack(b_{1},b_{2})\rbrack$,
%
%
\item $(a_{1},a_{2}){S}(b_{1},b_{2})$,
%
%
\item ${\mathcal F}^{\rho}_{A,B}{\models}\mathbf{A}(\mathbf{x_{1}},\mathbf{x_{2}},\mathbf{y_{1}},\mathbf{y_{2}})\ \lbrack a_{1},a_{2},b_{1},b_{2}\rbrack$,
%
%
\item $s{\not=}t$ and there exists $i,j{\in}\{1,2\}$ such that for exactly $1$~$e$ in ${\mathcal F}_{A,B}^{\rho}$, $R_{A,B}^{\rho}(e){=}
%
%
%
%
%
%
%
%
%
%
%
%
$\linebreak$
%
%
%
%
%
%
%
%
%
%
%
%
B{\setminus}\{a_{i},b_{j}\}$ and ${R_{A,B}^{\rho}}^{{-}1}(e){=}\emptyset$,
%
%
\item $\kappa(\lbrack(a_{1},a_{2})\rbrack){R}\kappa(\lbrack(b_{1},b_{2})\rbrack)$.
%
%
\end{itemize}
%
%
Thus, for all $(a_{1},a_{2}),(b_{1},b_{2}){\in}X$, $\lbrack(a_{1},a_{2})\rbrack{R^{\prime}}\lbrack(b_{1},b_{2})\rbrack$ if and only if $\kappa(\lbrack(a_{1},a_{2})\rbrack){R}
%
%
%
%
%
%
%
%
%
%
%
%
$\linebreak$
%
%
%
%
%
%
%
%
%
%
%
%
\kappa(\lbrack(b_{1},b_{2})\rbrack)$.
%
%
\\
\\
%
%
All in all, we have proved that, $(W^{\prime},R^{\prime})$ is isomorphic with $(W,R)$.
%
%
\\
\\
%
%
Consequently, we have shown that the class of all irreflexive symmetric frames with at least $2$ elements is relatively elementary definable in ${\mathcal L}_{2}$.
%
%
\medskip
%
%
\end{proof}
%
%
\begin{lemma}\label{lemma:first:case:2:N:moins}
%
%
Let $\L$ be an Euclidean modal logic.
If $\{2\}{\times}\mathbf{N}^{-}{\subseteq}\mathtt{S}_{\L}$ then $\Th(\Fr(\L))$ is undecidable.
%
%
\end{lemma}
%
%
\begin{proof}
%
%
Suppose $\{2\}{\times}\mathbf{N}^{-}{\subseteq}\mathtt{S}_{\L}$.
Since the first-order theory of the class of all irreflexive symmetric frames with at least $2$ elements is hereditarily undecidable~\cite[Page~$273$]{Ershov:1980}, then by Proposition~\ref{Syntactic:Restriction:Lemma:hereditarily:undecidable} and Lemma~\ref{lemma:relatively:elementary:definable:K2}, the first-order theory of ${\mathcal K}_{2}$ is hereditarily undecidable.
Since $\{2\}{\times}\mathbf{N}^{-}{\subseteq}\mathtt{S}_{\L}$, then by Lemma~\ref{lemma:gamma}, $\Th(\Fr(\L))$ is contained in the first-order theory of ${\mathcal K}_{2}$.
Since the first-order theory of ${\mathcal K}_{2}$ is hereditarily undecidable, then $\Th(\Fr(\L))$ is undecidable.
%
%
\medskip
%
%
\end{proof}
%
%
\begin{lemma}\label{lemma:second:case:N:plus:1}
%
%
Let $\L$ be an Euclidean modal logic.
If $\mathbf{N}^{+}{\times}\{2\}{\subseteq}\mathtt{S}_{\L}$ then $\Th(\Fr(\L))$ is undecidable.
%
%
\end{lemma}
%
%
\begin{proof}
%
%
Suppose $\mathbf{N}^{+}{\times}\{2\}{\subseteq}\mathtt{S}_{\L}$.
Since the first-order theory of the class of all irreflexive symmetric frames with at least $2$ elements is hereditarily undecidable~\cite[Page~$273$]{Ershov:1980}, then by Proposition~\ref{Syntactic:Restriction:Lemma:hereditarily:undecidable} and Lemma~\ref{lemma:relatively:elementary:definable:L2}, the first-order theory of ${\mathcal L}_{2}$ is hereditarily undecidable.
Since $\mathbf{N}^{+}{\times}\{2\}{\subseteq}\mathtt{S}_{\L}$, then by Lemma~\ref{lemma:delta}, $\Th(\Fr(\L))$ is contained in the first-order theory of ${\mathcal L}_{2}$.
Since the first-order theory of ${\mathcal L}_{2}$ is hereditarily undecidable, then $\Th(\Fr(\L))$ is undecidable.
%
%
\medskip
%
%
\end{proof}
%
%
%
%
\section{Relativizations}\label{section:about:relativization}
%
%
In this section, we adapt to our purpose a theorem about relativizations.
See~\cite[Theorem $5.1.1$ at Page~$203$]{Hodges:1993} for a more general presentation.
As far as we know, the first systematic treatment of the concept of relativization can be seen in the book of Tarski, Mostowski and Robinson about {\it Undecidable Theories}\/~\cite{Tarski:Mostowski:Robinson:1953}.
We will use it later for the proof that the modal definability problem with respect to classes of frames determined by some Euclidean modal logics is undecidable.
%
%
\\
\\
%
%
The {\it relativization of a first-order formula $C$ with respect to a first-order formula $A$ and an individual variable $\mathbf{x}$}\/ (denoted $(C)_{\mathbf{x}}^{A}$) is inductively defined as follows:
%
%
\begin{itemize}
%
%
\item $(\mathbf{R}(\mathbf{y},\mathbf{z}))_{\mathbf{x}}^{A}$ is $\mathbf{R}(\mathbf{y},\mathbf{z})$,
%
%
\item $(\mathbf{y}=\mathbf{z})_{\mathbf{x}}^{A}$ is $\mathbf{y}=\mathbf{z}$,
%
%
\item $(\neg C)_{\mathbf{x}}^{A}$ is $\neg(C)_{\mathbf{x}}^{A}$,
%
%
\item $(C\vee D)_{\mathbf{x}}^{A}$ is $(C)_{\mathbf{x}}^{A}\vee(D)_{\mathbf{x}}^{A}$,
%
%
\item $(\forall\mathbf{y}C)_{\mathbf{x}}^{A}$ is $\forall\mathbf{y}(A\lbrack\mathbf{x}/\mathbf{y}\rbrack\rightarrow(C)_{\mathbf{x}}^{A})$.
%
%
\end{itemize}
%
%
In the above definition, $A\lbrack\mathbf{x}/\mathbf{y}\rbrack$ denotes the first-order formula obtained from the first-order formula $A$ by replacing every free occurrence of the individual variable $\mathbf{x}$ in $A$ by the individual variable $\mathbf{y}$.
%
%
\\
\\
%
%
From now on, {\bf when we write $(C)_{\mathbf{x}}^{A}$, we shall assume that the sets of individual variables occurring in $A$ and $C$ are disjoint.}
%
%
\\
\\
%
%
The reader may easily verify by induction on the first-order formula $C$ that $fiv((C)_{\mathbf{x}}^{A})\subseteq(fiv(A)\setminus\{\mathbf{x}\}){\cup}fiv(C)$.
Hence, if $C$ is a sentence then $fiv((C)_{\mathbf{x}}^{A})\subseteq fiv(A)\setminus\{\mathbf{x}\}$.
%
%
\\
\\
%
%
Let $(W,R),(W^{\prime},R^{\prime})$ be frames.
%
%
\\
\\
%
%
$(W^{\prime},R^{\prime})$ is a {\it relativized reduct of $(W,R)$}\/ if there exists a first-order formula $A(\mathbf{x_{1}},\ldots,\mathbf{x_{m}},\mathbf{y})$ and there exists $s_{1},\ldots,s_{m}{\in}W$ such that $(W^{\prime},R^{\prime})$ is the restriction of $(W,R)$ to the set of all $t{\in}W$ such that $(W,R)\models A(\mathbf{x_{1}},\ldots,\mathbf{x_{m}},\mathbf{y})\ \lbrack s_{1},\ldots,
%
%
%
%
%
%
%
%
%
%
%
%
$\linebreak$
%
%
%
%
%
%
%
%
%
%
%
%
s_{m},t\rbrack$.
In this case, $(W^{\prime},R^{\prime})$ is the {\it relativized reduct of $(W,R)$ with respect to the first-order formula $A(\mathbf{x_{1}},\ldots,\mathbf{x_{m}},\mathbf{y})$ and $s_{1},\ldots,s_{m}{\in}W$.}
Obviously, $(W,R)$ possesses a relativized reduct with respect to a first-order formula $A(\mathbf{x_{1}},\ldots,\mathbf{x_{m}},\mathbf{y})$ and $s_{1},\ldots,s_{m}{\in}W$ if and only if $(W,R)\models\exists\mathbf{y}A(\mathbf{x_{1}},\ldots,\mathbf{x_{m}},\mathbf{y})\ \lbrack s_{1},\ldots,s_{m}\rbrack$.
%
%
\\
\\
As proved in~\cite[Theorem $5.1.1$ at Page~$203$]{Hodges:1993}, for all first-order formulas $A(\mathbf{x_{1}},
%
%
%
%
%
%
%
%
%
%
%
%
$\linebreak$
%
%
%
%
%
%
%
%
%
%
%
%
\ldots,\mathbf{x_{m}},\mathbf{y})$ and for all $s_{1},\ldots,s_{m}{\in}W$, if $(W^{\prime},R^{\prime})$ is the relativized reduct of $(W,R)$ with respect to $A(\mathbf{x_{1}},\ldots,\mathbf{x_{m}},\mathbf{y})$ and $s_{1},\ldots,s_{m}$ then for all first-order formulas $C(\mathbf{z_{1}},\ldots,\mathbf{z_{n}})$ and for all $u^{\prime}_{1},\ldots,u^{\prime}_{n}{\in}W^{\prime}$, the following conditions are equivalent:
%
%
\begin{itemize}
%
%
\item $(W,R)\models(C(\mathbf{z_{1}},\ldots,\mathbf{z_{n}}))_{\mathbf{y}}^{A(\mathbf{x_{1}},\ldots,\mathbf{x_{m}},\mathbf{y})}\ \lbrack s_{1},\ldots,s_{m},u^{\prime}_{1},\ldots,u^{\prime}_{n}\rbrack$,
%
%
\item $(W^{\prime},R^{\prime})\models C(\mathbf{z_{1}},\ldots,\mathbf{z_{n}})\ \lbrack u^{\prime}_{1},\ldots,u^{\prime}_{n}\rbrack$.
%
%
\end{itemize}
%\begin{proposition}[Relativization Theorem]\label{Syntactic:Restriction:Lemma}
%
%
%Let $A(\mathbf{x_{1}},\ldots,\mathbf{x_{m}},\mathbf{y})$ be a first-order formula and $s_{1},\ldots,s_{m}{\in}W$.
%If $(W^{\prime},R^{\prime})$ is the relativized reduct of $(W,R)$ with respect to $A(\mathbf{x_{1}},\ldots,\mathbf{x_{m}},\mathbf{y})$ and $s_{1},\ldots,s_{m}$ then for all first-order formulas $C(\mathbf{z_{1}},\ldots,\mathbf{z_{n}})$ and for all $u^{\prime}_{1},\ldots,u^{\prime}_{n}{\in}W^{\prime}$, the following conditions are equivalent:
%
%
%\begin{itemize}
%
%
%\item $(W,R)\models(C(\mathbf{z_{1}},\ldots,\mathbf{z_{n}}))_{\mathbf{y}}^{A(\mathbf{x_{1}},\ldots,\mathbf{x_{m}},\mathbf{y})}\ \lbrack s_{1},\ldots,s_{m},u^{\prime}_{1},\ldots,u^{\prime}_{n}\rbrack$,
%
%
%\item $(W^{\prime},R^{\prime})\models C(\mathbf{z_{1}},\ldots,\mathbf{z_{n}})\ \lbrack u^{\prime}_{1},\ldots,u^{\prime}_{n}\rbrack$.
%
%
%\end{itemize}
%
%
%\end{proposition}
%
%
The importance of the concept of relativization lies in the fact that it is at the heart of the concept of stability~\cite{Balbiani:Tinchev:2017,Balbiani:Tinchev:2019}.
%
%
\\
\\
%
%
A class ${\cal C}$ of frames is {\it stable}\/ if there exists a first-order formula $A(\mathbf{x_{1}},\ldots,\mathbf{x_{m}},\mathbf{y})$ and there exists a sentence $B$ such that
%
%
\begin{itemize}
%
%
\item for all frames $(W^{\prime},R^{\prime})$ in ${\cal C}$ and for all $s^{\prime}_{1},\ldots,s^{\prime}_{m}{\in}W^{\prime}$, the relativized reduct of $(W^{\prime},R^{\prime})$ with respect to $A(\mathbf{x_{1}},\ldots,\mathbf{x_{m}},\mathbf{y})$ and $s^{\prime}_{1},\ldots,s^{\prime}_{m}$ is in ${\cal C}$,
%
\item for all frames $(W_{0},R_{0})$ in ${\cal C}$, there exists frames $(W^{\prime},R^{\prime}),(W^{\prime\prime},R^{\prime\prime})$ in ${\cal C}$ and there exists $s^{\prime}_{1},\ldots,s^{\prime}_{m}{\in}W^{\prime}$ such that $(W_{0},R_{0})$ is the relativized reduct of $(W^{\prime},R^{\prime})$ with respect to $A(\mathbf{x_{1}},\ldots,\mathbf{x_{m}},\mathbf{y})$ and $s^{\prime}_{1},\ldots,s^{\prime}_{m}$, $(W^{\prime},R^{\prime})\models B$, $(W^{\prime\prime},R^{\prime\prime})\not\models B$ and $(W^{\prime},R^{\prime})\preceq(W^{\prime\prime},R^{\prime\prime})$.
%
%
\end{itemize}
%
%
In this case, $(A(\mathbf{x_{1}},\ldots,\mathbf{x_{m}},\mathbf{y}),B)$ is called {\it witness of the stability of ${\cal C}$.}
%
%
\\
\\
%
%
The importance of the concept of stability lies in its capacity to connect the problem of deciding the validity of sentences to the problem of deciding the modal definability of sentences.
%
%
\begin{proposition}\label{Validity:And:Definability}
%
%
Let ${\mathcal C}$ be a class of frames.
If ${\cal C}$ is stable then the problem of deciding the validity of sentences in ${\cal C}$ is reducible to the problem of deciding the modal definability of sentences with respect to ${\cal C}$.
%
%
\end{proposition}
%
%
\begin{proof}
%
%
See~\cite{Balbiani:Tinchev:2017,Balbiani:Tinchev:2019}.
%
%
\medskip
%
%
\end{proof}
%
%
By Proposition~\ref{Validity:And:Definability}, Balbiani and Tinchev~\cite{Balbiani:Tinchev:2017,Balbiani:Tinchev:2019} have proved that the problem of deciding the modal definability of sentences is undecidable with respect to classes of frames such as the class of all transitive frames, the class of all symmetric frames, etc.
%
%
\begin{lemma}\label{lemma:L:modal:logic:Fr:L:is:stable}
%
%
Let $\L$ be an Euclidean modal logic.
Then, $\Fr(\L)$ is stable.
%
%
\end{lemma}
%
%
\begin{proof}
%
%
We consider the following $2$~cases.
%
%
\\
\\
%
%
{\bf Case where $\Diamond\top{\in}\L$:}
Let
%
%
\begin{itemize}
%
%
\item $A(\mathbf{x_{1}},\mathbf{x_{2}},\mathbf{y})$ be the first-order formula $\forall\mathbf{x}(\mathbf{R}(\mathbf{x},\mathbf{x_{1}})\vee\mathbf{R}(\mathbf{x_{1}},\mathbf{x})\leftrightarrow\mathbf{x}=\mathbf{x_{1}})\wedge\forall\mathbf{x}(\mathbf{R}(\mathbf{x},\mathbf{x_{2}})\vee\mathbf{R}(\mathbf{x_{2}},\mathbf{x})\leftrightarrow\mathbf{x}=\mathbf{x_{2}})\rightarrow\mathbf{y}\not=\mathbf{x_{1}}\wedge\mathbf{y}\not=\mathbf{x_{2}}$,
%
%
\item $B$ be the sentence $\exists\mathbf{x_{1}}\exists\mathbf{x_{2}}(\forall\mathbf{x}(\mathbf{R}(\mathbf{x},\mathbf{x_{1}})\vee\mathbf{R}(\mathbf{x_{1}},\mathbf{x})\leftrightarrow\mathbf{x}=\mathbf{x_{1}})\wedge\forall\mathbf{x}(\mathbf{R}(\mathbf{x},
%
%
%
%
%
%
%
%
%
%
%
%
$\linebreak$
%
%
%
%
%
%
%
%
%
%
%
%
\mathbf{x_{2}})\vee\mathbf{R}(\mathbf{x_{2}},\mathbf{x})\leftrightarrow\mathbf{x}=\mathbf{x_{2}})\wedge\mathbf{x_{1}}\not=\mathbf{x_{2}})$.
%
%
\end{itemize}
%
%
Obviously, for all Euclidean frames $(W,R)$ in $\Fr(\L)$, for all $s_{1},s_{2}{\in}W$ and for all Euclidean frames $(W^{\prime},R^{\prime})$, if $(W^{\prime},R^{\prime})$ is the relativized reduct of $(W,R)$ with respect to $A(\mathbf{x_{1}},\mathbf{x_{2}},\mathbf{y})$ and $s_{1},s_{2}$ then $(W^{\prime},R^{\prime})$ is in $\Fr(\L)$.
Now, let $(W_{0},R_{0})$ be an Euclidean frame in $\Fr(\L)$.
Let $(W,R)$ be the Euclidean frame in $\Fr(\L)$ such that
%
%
\begin{itemize}
%
%
\item $W{=}W_{0}{\cup}\{s_{1},s_{2}\}$,
%
%
\item $R{=}R_{0}{\cup}\{(s_{1},s_{1}),(s_{2},s_{2})\}$,
%
%
\end{itemize}
%
%
$s_{1}$ and $s_{2}$ being $2$~new distinct elements.
Let $(W^{\prime},R^{\prime})$ be the Euclidean frame in $\Fr(\L)$ such that
%
%
\begin{itemize}
%
%
\item $W^{\prime}{=}\{s^{\prime}\}$,
%
%
\item $R^{\prime}{=}\{(s^{\prime},s^{\prime})\}$.
%
%
\end{itemize}
%
%
Obviously, $(W_{0},R_{0})$ is the relativized reduct of $(W,R)$ with respect to $A(\mathbf{x_{1}},\mathbf{x_{2}},
%
%
%
%
%
%
%
%
%
%
%
%
$\linebreak$
%
%
%
%
%
%
%
%
%
%
%
%
\mathbf{y})$ and $s_{1},s_{2}$, $(W,R)\models B$, $(W^{\prime},R^{\prime})\not\models B$ and $(W,R)\preceq(W^{\prime},R^{\prime})$.
Hence, $(A(\mathbf{x_{1}},\mathbf{x_{2}},\mathbf{y}),B)$ is a witness of the stability of $\Fr(\L)$.
%
%
\\
\\
%
%
{\bf Case where $\Diamond\top{\not\in}\L$:}
Let
%
%
\begin{itemize}
%
%
\item $A(\mathbf{x_{1}},\mathbf{x_{2}},\mathbf{y})$ be the first-order formula $\forall\mathbf{x}\neg\mathbf{R}(\mathbf{x_{1}},\mathbf{x})\wedge\forall\mathbf{x}\neg\mathbf{R}(\mathbf{x_{2}},\mathbf{x})\rightarrow\mathbf{y}\not=\mathbf{x_{1}}\wedge\mathbf{y}\not=\mathbf{x_{2}}$,
%
%
\item $B$ be the sentence $\exists\mathbf{x_{1}}\exists\mathbf{x_{2}}(\forall\mathbf{x}\neg\mathbf{R}(\mathbf{x_{1}},\mathbf{x})\wedge\forall\mathbf{x}\neg\mathbf{R}(\mathbf{x_{2}},\mathbf{x})\wedge\mathbf{x_{1}}\not=\mathbf{x_{2}})$.
%
%
\end{itemize}
%
%
Obviously, for all Euclidean frames $(W,R)$ in $\Fr(\L)$, for all $s_{1},s_{2}{\in}W$ and for all Euclidean frames $(W^{\prime},R^{\prime})$, if $(W^{\prime},R^{\prime})$ is the relativized reduct of $(W,R)$ with respect to $A(\mathbf{x_{1}},\mathbf{x_{2}},\mathbf{y})$ and $s_{1},s_{2}$ then $(W^{\prime},R^{\prime})$ is in $\Fr(\L)$.
Now, let $(W_{0},R_{0})$ be an Euclidean frame in $\Fr(\L)$.
Let $(W,R)$ be the Euclidean frame in $\Fr(\L)$ such that
%
%
\begin{itemize}
%
%
\item $W{=}W_{0}{\cup}\{s_{1},s_{2}\}$,
%
%
\item $R{=}R_{0}$,
%
%
\end{itemize}
%
%
$s_{1}$ and $s_{2}$ being $2$~new distinct elements.
Let $(W^{\prime},R^{\prime})$ be the Euclidean frame in $\Fr(\L)$ such that
%
%
\begin{itemize}
%
%
\item $W^{\prime}{=}\{s^{\prime}\}$,
%
%
\item $R^{\prime}{=}\emptyset$.
%
%
\end{itemize}
%
%
Obviously, $(W_{0},R_{0})$ is the relativized reduct of $(W,R)$ with respect to $A(\mathbf{x_{1}},\mathbf{x_{2}},
%
%
%
%
%
%
%
%
%
%
%
%
$\linebreak$
%
%
%
%
%
%
%
%
%
%
%
%
\mathbf{y})$ and $s_{1},s_{2}$, $(W,R)\models B$, $(W^{\prime},R^{\prime})\not\models B$ and $(W,R)\preceq(W^{\prime},R^{\prime})$.
Thus, $(A(\mathbf{x_{1}},\mathbf{x_{2}},\mathbf{y}),B)$ is a witness of the stability of $\Fr(\L)$.
%
%
\medskip
%
%
\end{proof}
%
%
\begin{lemma}\label{Lemma:53:bis}
%
%
Let $\L$ be an Euclidean modal logic.
If $\Th(\Fr(\L))$ is undecidable then the problem of deciding the modal definability of sentences with respect to $\Fr(\L)$ is undecidable.
%
%
\end{lemma}
%
%
\begin{proof}
%
%
By Proposition~\ref{Validity:And:Definability} and Lemma~\ref{lemma:L:modal:logic:Fr:L:is:stable}.
%
%
\medskip
%
%
\end{proof}
%
%
%Lemma~\ref{lemma:L:modal:logic:Fr:L:is:stable} will be used in Section~\ref{definability:results:section} where we prove that the problem of deciding the modal definability of sentences is undecidable with respect to classes of frames determined by Euclidean modal logics $\L$ such that $\mathtt{S}_{\L}{\setminus}((\{1\}{\times}\mathbf{N}^{-}){\cup}0,1\}))$ is infinite.
%
%
%
%
%\section{Definability results}\label{definability:results:section}
%
%
%
%
\section{Computability of modal definability}\label{section:about:computability:of:modal:definability}
%
%
The following results show that for all Euclidean modal logics $\L$, the computability of $\Th(\Fr(\L))$ only depends on the cardinality of $\mathtt{S}_{\L}{\setminus}((\{1\}{\times}\mathbf{N}^{-}){\cup}(\mathbf{N}^{+}{\times}\{{-}1,
%
%
%
%
%
%
%
%
%
%
%
%
$\linebreak$
%
%
%
%
%
%
%
%
%
%
%
%
0,1\}))$.
%
%
\begin{lemma}\label{decidability:if:setminus:is:finite}
%
%
%Let $\L$ be an Euclidean modal logic.
%If $\mathtt{S}_{\L}{\setminus}((\{1\}{\times}\mathbf{N}^{-}){\cup}(\mathbf{N}^{+}{\times}\{{-}1,
%
%
%
%
%
%
%
%
%
%
%
%
%$\linebreak$
%
%
%
%
%
%
%
%
%
%
%
%
%0,1\}))$ is finite then $\Th(\Fr(\L))$ is decidable.
Let $\L$ be an Euclidean modal logic such that $\mathtt{S}_{\L}{\setminus}((\{1\}{\times}
%
%
%
%
%
%
%
%
%
%
%
%
$\linebreak$
%
%
%
%
%
%
%
%
%
%
%
%
\mathbf{N}^{-}){\cup}(\mathbf{N}^{+}{\times}\{{-}1,0,1\}))$ is finite.
Let $\varphi$ be a modal formula such that $\L{=}\K5{\oplus}\varphi$.
Let $k$ be the least integer such that $k{\geq}4$ and $(\ast)$~for all $m{\in}\mathbf{N}^{+}$ and for all $n{\in}\mathbf{N}^{-}$, if $(m,n){\in}\mathtt{S}_{\L}{\setminus}((\{1\}{\times}\mathbf{N}^{-}){\cup}(\mathbf{N}^{+}{\times}\{{-}1,0,1\}))$ then $m{+}n{\leq}k$.
For all sentences $A$, the following conditions are equivalent:
%
%
\begin{itemize}
%
%
\item $A{\in}\Th(\Fr(\L))$,
%
%
\item for all finite Euclidean frames $(W,R)$, if ${\parallel}W{\parallel}{\leq}2{\times}q{\times}((Q{\times}K){+}1)^{2}{\times}
%
%
%
%
%
%
%
%
%
%
%
%
$\linebreak$
%
%
%
%
%
%
%
%
%
%
%
%
2^{(Q{\times}K)^{2}}{\times}Q{\times}K$ and $(W,R)$ validates $\varphi$ then $(W,R)$ validates $A$,
%
%
\end{itemize}
%
%
$q$ denoting $\qdd(A)$, $Q$ denoting $2{\times}q{\times}(q{\times}(q{+}1)^{2}+1)$ and $K$ denoting $2^{k}$.
%
%
\end{lemma}
%
%
\begin{proof}
%
%
Since~$(\ast)$, then for all $m{\in}\mathbf{N}^{+}$ and for all $n{\in}\mathbf{N}^{-}$, if ${\mathcal F}_{m}^{n}{\models}\L$ then either $m{=}1$, or $n{=}{-}1$, or $n{=}0$, or $n{=}1$, or $m{\not=}1$, $n{\not=}{-}1$, $n{\not=}0$, $n{\not=}1$ and $m{+}n{\leq}k$.
Thus, by Lemma~\ref{lemma:about:flowers:simple:galaxies}, for all $m{\in}\mathbf{N}^{+}$ and for all $n{\in}\mathbf{N}^{-}$, if ${\mathcal F}_{m}^{n}{\models}\L$ then either ${\mathcal F}_{m}^{n}$ is simple, or ${\mathcal F}_{m}^{n}$ is non-simple and $m{+}n{\leq}k$.
%
%
\\
\\
%
%
We claim that for all galaxies ${\mathcal F}_{A,B}^{\rho}$, if ${\mathcal F}_{A,B}^{\rho}{\models}\L$ then either ${\mathcal F}_{A,B}^{\rho}$ is simple, or ${\mathcal F}_{A,B}^{\rho}$ is non-simple and ${\parallel}B{\parallel}{\leq}k$.
If not, let ${\mathcal F}_{A,B}^{\rho}$ be a galaxy such that ${\mathcal F}_{A,B}^{\rho}{\models}\L$, ${\mathcal F}_{A,B}^{\rho}$ is non-simple and ${\parallel}B{\parallel}{>}k$.
Consequently, there exists $s{\in}A$ such that ${\parallel}\rho(s){\parallel}{>}1$ and ${\parallel}B{\setminus}\rho(s){\parallel}{>}1$.
Let $(W_{s},R_{s})$ be the least generated subframe of ${\mathcal F}_{A,B}^{\rho}$ containing $s$.
Obviously, $W_{s}{=}\{s\}{\cup}B$ and $R_{s}{=}(\{s\}{\times}\rho(s)){\cup}(B{\times}B)$.
Moreover, since ${\mathcal F}_{A,B}^{\rho}{\models}\L$, then by Lemma~\ref{Generated:Subframe:Lemma}, $(W_{s},R_{s}){\models}\L$.
Let $m^{\prime}{=}\min\{{\parallel}\rho(s){\parallel},k\}$ and $n^{\prime}{=}\min\{{\parallel}B{\setminus}\rho(s){\parallel},k\}$.
Obviously, $m^{\prime}{>}1$, $n^{\prime}{>}1$ and $m^{\prime}{+}n^{\prime}{>}k$.
Moreover, since ${\parallel}B{\parallel}{>}k$, ${\parallel}\rho(s){\parallel}{>}1$ and ${\parallel}B{\setminus}\rho(s){\parallel}{>}1$, then by Lemma~\ref{lemma:about:k:s:new:element:bmi:of:W:R}, ${\mathcal F}_{m^{\prime}}^{n^{\prime}}$ is a bounded morphic image of $(W_{s},R_{s})$.
Since $(W_{s},R_{s}){\models}\L$, then by Lemma~\ref{Bounded:Morphism:Lemma}, ${\mathcal F}_{m^{\prime}}^{n^{\prime}}{\models}\L$.
Since $m^{\prime}{>}1$ and $n^{\prime}{>}1$, then by Lemma~\ref{lemma:about:flowers:simple:galaxies}, ${\mathcal F}_{m^{\prime}}^{n^{\prime}}$ is non-simple.
Since for all $m{\in}\mathbf{N}^{+}$ and for all $n{\in}\mathbf{N}^{-}$, if ${\mathcal F}_{m}^{n}{\models}\L$ then either ${\mathcal F}_{m}^{n}$ is simple, or ${\mathcal F}_{m}^{n}$ is non-simple and $m{+}n{\leq}k$, then either ${\mathcal F}_{m^{\prime}}^{n^{\prime}}{\not\models}\L$, or $m^{\prime}{+}n^{\prime}{\leq}k$: a contradiction.
Hence, for all galaxies ${\mathcal F}_{A,B}^{\rho}$, if ${\mathcal F}_{A,B}^{\rho}{\models}\L$ then either ${\mathcal F}_{A,B}^{\rho}$ is simple, or ${\mathcal F}_{A,B}^{\rho}$ is non-simple and ${\parallel}B{\parallel}{\leq}k$.
%
%
\\
\\
%
%
Thus, by Lemmas~\ref{lemma:about:generated:subframe:and:disjoint:unions} and~\ref{Generated:Subframe:Lemma},
%
%
\begin{itemize}
%
%
\item for all disjoint indexed families $\{{\mathcal F}_{A_{i},B_{i}}^{\rho_{i}}:\ i{\in}I\}$ of galaxies, if the union of $\{{\mathcal F}_{A_{i},B_{i}}^{\rho_{i}}:\ i{\in}I\}$ validates $\L$ then for all $i{\in}I$, either ${\mathcal F}_{A_{i},B_{i}}^{\rho_{i}}$ is simple, or ${\mathcal F}_{A_{i},B_{i}}^{\rho_{i}}$ is non-simple and ${\parallel}B_{i}{\parallel}{\leq}k$.
%
%
\end{itemize}
%
%
Now, let $A$ be an arbitrary sentence.
Suppose $A{\not\in}\Th(\Fr(\L))$.
Consequently, there exists an Euclidean frame $(W,R)$ such that $(W,R)$ validates $\L$ and $(W,R)$ does not validate $A$.
Hence, by Lemma~\ref{lemma:every:frame:disjoint:union:galaxies}, there exists a disjoint indexed family $\{{\mathcal F}_{A_{i},B_{i}}^{\rho_{i}}:\ i{\in}I\}$ of galaxies such that the union $(W,R)$ of $\{{\mathcal F}_{A_{i},B_{i}}^{\rho_{i}}:\ i{\in}I\}$ validates $\L$ and $(W,R)$ does not validate $A$.
Let $q{=}\qdd(A)$.
%Notice that $\qd(A){\leq}q$ and $q{\geq}3$.
%
%
\\
\\
%
%
Let $\{{\mathcal F}_{A_{i}^{\prime},B_{i}^{\prime}}^{\rho_{i}^{\prime}}:\ i{\in}I\}$ be a $q$-$\alpha$-reduction of $\{{\mathcal F}_{A_{i},B_{i}}^{\rho_{i}}:\ i{\in}I\}$.
Since $(W,R)$ validates $\L$, then for all $i{\in}I$, either ${\mathcal F}_{A_{i},B_{i}}^{\rho_{i}}$ is simple, or ${\mathcal F}_{A_{i},B_{i}}^{\rho_{i}}$ is non-simple and ${\parallel}B_{i}{\parallel}{\leq}k$.
Let $I^{ns}$ be the set of all $i{\in}I$ such that ${\mathcal F}_{A_{i},B_{i}}^{\rho_{i}}$ is non-simple.
Thus, $\{{\mathcal F}_{A_{i},B_{i}}^{\rho_{i}}:\ i{\in}I^{ns}\}$ is kernel-bounded.
Consequently, by Lemma~\ref{lemma:the:alpha:reduction:is:always:dust:finite}, for all $i{\in}I$,
%
%
\begin{itemize}
%
%
\item $A_{i}^{\prime}{\subseteq}A_{i}$ and if ${\parallel}A_{i}^{\prime}{\parallel}{\leq}2$ then $A_{i}^{\prime}{=}A_{i}$,
%
%
\item $B_{i}^{\prime}{=}B_{i}$,
%
%
\item for all $s^{\prime}{\in}A_{i}^{\prime}$, $\rho_{i}^{\prime}(s^{\prime}){=}\rho_{i}(s^{\prime})$,
%
%
\item ${\parallel}{\rho_{i}^{\prime}}^{{-}1}(\emptyset){\parallel}{\leq}q$ and ${\parallel}{\rho_{i}^{\prime}}^{{-}1}(B_{i}^{\prime}){\parallel}{\leq}q$,
%
%
\item for all $t^{\prime}{\in}B_{i}^{\prime}$, ${\parallel}{\rho_{i}^{\prime}}^{{-}1}(\{t^{\prime}\}){\parallel}{\leq}q$ and for all $t^{\prime}{\in}B_{i}^{\prime}$, ${\parallel}{\rho_{i}^{\prime}}^{{-}1}(B_{i}^{\prime}{\setminus}\{t^{\prime}\}){\parallel}{\leq}q$,
%
%
\item either ${\mathcal F}_{A_{i}^{\prime},B_{i}^{\prime}}^{\rho_{i}^{\prime}}$ is simple, or ${\mathcal F}_{A_{i}^{\prime},B_{i}^{\prime}}^{\rho_{i}^{\prime}}$ is non-simple, ${\parallel}B_{i}^{\prime}{\parallel}{\leq}k$ and ${\parallel}A_{i}^{\prime}{\parallel}{\leq}q{\times}2^{k}$.
%
%
%\item ${\mathcal F}_{A_{i}^{\prime},B_{i}^{\prime}}^{\rho_{i}^{\prime}}$ is a bounded morphic image of ${\mathcal F}_{A_{i},B_{i}}^{\rho_{i}}$,
%
%
%\item the second player has a winning strategy in the Ehrenfeucht-Fra\"\i ss\'e game ${\mathcal G}_{q}({\mathcal F}_{A_{i},B_{i}}^{\rho_{i}},{\mathcal F}_{A_{i}^{\prime},B_{i}^{\prime}}^{\rho_{i}^{\prime}})$.
%
%
\end{itemize}
%
%
%For all $i{\in}I$, if $(W_{A_{i}^{\prime},B_{i}^{\prime}}^{\rho_{i}^{\prime}},R_{A_{i}^{\prime},B_{i}^{\prime}}^{\rho_{i}^{\prime}})$ is simple then let $(W_{A_{i}^{\prime\prime},B_{i}^{\prime\prime}}^{\rho_{i}^{\prime\prime}},R_{A_{i}^{\prime\prime},B_{i}^{\prime\prime}}^{\rho_{i}^{\prime\prime}})$ be a $q$-$\beta$-reduction of $(W_{A_{i}^{\prime},B_{i}^{\prime}}^{\rho_{i}^{\prime}},R_{A_{i}^{\prime},B_{i}^{\prime}}^{\rho_{i}^{\prime}})$, whereas if $(W_{A_{i}^{\prime},B_{i}^{\prime}}^{\rho_{i}^{\prime}},R_{A_{i}^{\prime},B_{i}^{\prime}}^{\rho_{i}^{\prime}})$ is non-simple, ${\parallel}\{s^{\prime}{\in}A_{i}^{\prime}:\ \rho_{i}^{\prime}(s^{\prime}){\not=}\emptyset$ and $\rho_{i}^{\prime}(s^{\prime}){\not=}B_{i}^{\prime}\}{\parallel}{\leq}q{\times}(2^{k}{-}2)$ and ${\parallel}B_{i}^{\prime}{\parallel}{\leq}k$ then let $(W_{A_{i}^{\prime\prime},B_{i}^{\prime\prime}}^{\rho_{i}^{\prime\prime}},R_{A_{i}^{\prime\prime},B_{i}^{\prime\prime}}^{\rho_{i}^{\prime\prime}})$ be $(W_{A_{i}^{\prime},B_{i}^{\prime}}^{\rho_{i}^{\prime}},R_{A_{i}^{\prime},B_{i}^{\prime}}^{\rho_{i}^{\prime}})$.
%Consequently, by Lemma~\ref{lemma:beta:reduction:quasi:root:bounded:dust:OK} and Conditions~$(C_{1})$, $(C_{2})$ and~$(C_{3})$, for all $i{\in}I$,
%
%
%\begin{description}
%
%
%\item[$(C_{4})$] ${\parallel}\{s^{\prime\prime}{\in}A_{i}^{\prime\prime}:\ \rho_{i}^{\prime\prime}(s^{\prime\prime}){=}\emptyset\}{\parallel}{\leq}q$,
%
%
%\item[$(C_{5})$] ${\parallel}\{s^{\prime\prime}{\in}A_{i}^{\prime\prime}:\ \rho_{i}^{\prime\prime}(s^{\prime\prime}){=}B_{i}^{\prime\prime}\}{\parallel}{\leq}q$,
%
%
%\item[$(C_{6})$] either $(W_{A_{i}^{\prime\prime},B_{i}^{\prime\prime}}^{\rho_{i}^{\prime\prime}},R_{A_{i}^{\prime\prime},B_{i}^{\prime\prime}}^{\rho_{i}^{\prime\prime}})$ is simple and for all $s^{\prime\prime}{\in}B_{i}^{\prime\prime}$, ${\parallel}{\rho_{i}^{\prime\prime}}^{{-}1}(s^{\prime\prime}){\parallel}{\leq}q$, or $(W_{A_{i}^{\prime\prime},B_{i}^{\prime\prime}}^{\rho_{i}^{\prime\prime}},R_{A_{i}^{\prime\prime},B_{i}^{\prime\prime}}^{\rho_{i}^{\prime\prime}})$ is non-simple, ${\parallel}\{s^{\prime\prime}{\in}A_{i}^{\prime\prime}:\ \rho_{i}^{\prime\prime}(s^{\prime\prime}){\not=}\emptyset$ and $\rho_{i}^{\prime\prime}(s^{\prime\prime}){\not=}B_{i}^{\prime\prime}\}{\parallel}{\leq}q{\times}(2^{k}{-}2)$ and ${\parallel}B_{i}^{\prime\prime}{\parallel}{\leq}k$.
%
%
%\end{description}
%
%
Moreover, by Proposition~\ref{lemma:good:properties:of:alpha:reductions:indexed}, $\{{\mathcal F}_{A_{i}^{\prime},B_{i}^{\prime}}^{\rho_{i}^{\prime}}:\ i{\in}I\}$ is dust-bounded and $\{{\mathcal F}_{A_{i}^{\prime},B_{i}^{\prime}}^{\rho_{i}^{\prime}}:\ i{\in}I^{ns}\}$ is root-bounded and kernel-bounded.
%$(W,R)$ being the union of $\{{\mathcal F}_{A_{i},B_{i}}^{\rho_{i}}:\ i{\in}I\}$ and
In other respect, $(W^{\prime},R^{\prime})$ being the union of $\{{\mathcal F}_{A_{i}^{\prime},B_{i}^{\prime}}^{\rho_{i}^{\prime}}:\ i{\in}I\}$, $(W^{\prime},R^{\prime})$ is a bounded morphic image of $(W,R)$ and the second player has a winning strategy in the Ehrenfeucht-Fra\"\i ss\'e game ${\mathcal G}_{q}((W,R),(W^{\prime},R^{\prime}))$.
%Thus,
As a result, by Lemmas~\ref{lemma:alpha:reduction:modal:logic} and~\ref{lemma:alpha:reduction:fol}, $(W,R){\preceq}(W^{\prime},R^{\prime})$ and $(W,R){\equiv_{q}}(W^{\prime},R^{\prime})$.
%
%
\\
\\
%
%
Let $\{{\mathcal F}_{A_{i}^{\prime\prime},B_{i}^{\prime\prime}}^{\rho_{i}^{\prime\prime}}:\ i{\in}I\}$ be a $q$-$\gamma$-reduction of $\{{\mathcal F}_{A_{i}^{\prime},B_{i}^{\prime}}^{\rho_{i}^{\prime}}:\ i{\in}I\}$.
%For all $i{\in}I$, if $(W_{A_{i}^{\prime},B_{i}^{\prime}}^{\rho_{i}^{\prime}},R_{A_{i}^{\prime},B_{i}^{\prime}}^{\rho_{i}^{\prime}})$ is non-simple then let $(W_{A_{i}^{\prime\prime},B_{i}^{\prime\prime}}^{\rho_{i}^{\prime\prime}},R_{A_{i}^{\prime\prime},B_{i}^{\prime\prime}}^{\rho_{i}^{\prime\prime}})$ be $(W_{A_{i}^{\prime},B_{i}^{\prime}}^{\rho_{i}^{\prime}},R_{A_{i}^{\prime},B_{i}^{\prime}}^{\rho_{i}^{\prime}})$, whereas if $(W_{A_{i}^{\prime},B_{i}^{\prime}}^{\rho_{i}^{\prime}},R_{A_{i}^{\prime},B_{i}^{\prime}}^{\rho_{i}^{\prime}})$ is simple then let $(W_{A_{i}^{\prime\prime},B_{i}^{\prime\prime}}^{\rho_{i}^{\prime\prime}},R_{A_{i}^{\prime\prime},B_{i}^{\prime\prime}}^{\rho_{i}^{\prime\prime}})$ be a $q$-$\gamma$-reduction of $(W_{A_{i}^{\prime},B_{i}^{\prime}}^{\rho_{i}^{\prime}},R_{A_{i}^{\prime},B_{i}^{\prime}}^{\rho_{i}^{\prime}})$.
Hence, by Lemma~\ref{lemma:gamma:reduction:prebounded:dust:quasi:root:bounded:OK}, for all $i{\in}I$,
%
%
\begin{itemize}
%
%
\item $A_{i}^{\prime\prime}{\subseteq}A_{i}^{\prime}$,
%
%
\item $B_{i}^{\prime\prime}{\subseteq}B_{i}^{\prime}$ and if ${\parallel}B_{i}^{\prime\prime}{\parallel}{\leq}2$ then $B_{i}^{\prime\prime}{=}B_{i}^{\prime}$,
%
%
\item for all $s^{\prime\prime}{\in}A_{i}^{\prime\prime}$, $\rho_{i}^{\prime\prime}(s^{\prime\prime}){\subseteq}\rho_{i}^{\prime}(s^{\prime\prime})$,
%
%
\item ${\parallel}{\rho_{i}^{\prime\prime}}^{{-}1}(\emptyset){\parallel}{\leq}q$ and ${\parallel}{\rho_{i}^{\prime\prime}}^{{-}1}(B_{i}^{\prime\prime}){\parallel}{\leq}q$,
%
%
\item for all $t^{\prime\prime}{\in}B_{i}^{\prime\prime}$, ${\parallel}{\rho_{i}^{\prime\prime}}^{{-}1}(\{t^{\prime\prime}\}){\parallel}{\leq}q$ and for all $t^{\prime\prime}{\in}B_{i}^{\prime\prime}$, ${\parallel}{\rho_{i}^{\prime\prime}}^{{-}1}(B_{i}^{\prime\prime}{\setminus}\{t^{\prime\prime}\}){\parallel}{\leq}q$,
%
%
\item either ${\mathcal F}_{A_{i}^{\prime\prime},B_{i}^{\prime\prime}}^{\rho_{i}^{\prime\prime}}$ is simple, ${\parallel}B_{i}^{\prime\prime}{\parallel}{\leq}q{\times}(q{+}1)^{2}$ and ${\parallel}A_{i}^{\prime\prime}{\parallel}{\leq}2{\times}q{\times}(q{\times}(q{+}1)^{2}+1)$, or ${\mathcal F}_{A_{i}^{\prime\prime},B_{i}^{\prime\prime}}^{\rho_{i}^{\prime\prime}}$ is non-simple, ${\parallel}B_{i}^{\prime\prime}{\parallel}{\leq}k$ and ${\parallel}A_{i}^{\prime\prime}{\parallel}{\leq}q{\times}2^{k}$.
%
%
%\item ${\mathcal F}_{A_{i}^{\prime\prime},B_{i}^{\prime\prime}}^{\rho_{i}^{\prime\prime}}$ is a bounded morphic image of ${\mathcal F}_{A_{i}^{\prime},B_{i}^{\prime}}^{\rho_{i}^{\prime}}$,
%
%
%\item the second player has a winning strategy in the Ehrenfeucht-Fra\"\i ss\'e game ${\mathcal G}_{q}({\mathcal F}_{A_{i}^{\prime},B_{i}^{\prime}}^{\rho_{i}^{\prime}},{\mathcal F}_{A_{i}^{\prime\prime},B_{i}^{\prime\prime}}^{\rho_{i}^{\prime\prime}})$.
%
%
\end{itemize}
%
%
As a result, $\{{\mathcal F}_{A_{i}^{\prime\prime},B_{i}^{\prime\prime}}^{\rho_{i}^{\prime\prime}}:\ i{\in}I\}$ is dust-bounded, root-bounded and kernel-bounded, seeing that for all $i{\in}I$, ${\parallel}A_{i}^{\prime\prime}{\parallel}{\leq}Q{\times}K$ and ${\parallel}B_{i}^{\prime\prime}{\parallel}{\leq}Q{\times}K$ where $Q{=}2{\times}q{\times}(q{\times}(q{+}
%
%
%
%
%
%
%
%
%
%
%
%
$\linebreak$
%
%
%
%
%
%
%
%
%
%
%
%
1)^{2}+1)$ and $K{=}2^{k}$.
Moreover, $(W^{\prime\prime},R^{\prime\prime})$ being the union of $\{{\mathcal F}_{A_{i}^{\prime\prime},B_{i}^{\prime\prime}}^{\rho_{i}^{\prime\prime}}:\ i{\in}I\}$, by Lemmas~\ref{lemma:gamma:reduction:modal:logic} and~\ref{lemma:gamma:reduction:fol}, $(W^{\prime},R^{\prime}){\preceq}(W^{\prime\prime},R^{\prime\prime})$ and $(W^{\prime},R^{\prime}){\equiv_{q}}(W^{\prime\prime},R^{\prime\prime})$.
%Thus, for all $i{\in}I$,
%
%
%\begin{itemize}
%
%
%\item ${\parallel}W_{A_{i}^{\prime\prime},B_{i}^{\prime\prime}}^{\rho_{i}^{\prime\prime}}{\parallel}{\leq}\max\{2{\times}q{\times}(q{\times}(q{+}1)^{2}+1){+}q{\times}(q{+}1)^{2},q{\times}2^{k}{+}k\}$.
%
%
%\end{itemize}
%
%
%Consequently, $k^{\prime}$ denoting $\max\{2{\times}q{\times}(q{\times}(q{+}1)^{2}+1){+}q{\times}(q{+}1)^{2},q{\times}2^{k}{+}k\}$, the disjoint indexed family $\{(W_{A_{i}^{\prime\prime},B_{i}^{\prime\prime}}^{\rho_{i}^{\prime\prime}},R_{A_{i}^{\prime\prime},B_{i}^{\prime\prime}}^{\rho_{i}^{\prime\prime}}):\ i{\in}I\}$ of galaxies is $k^{\prime}$-prebounded.
%
%
\\
\\
%
%
Let $\{{\mathcal F}_{A_{j}^{\prime\prime},B_{j}^{\prime\prime}}^{\rho_{j}^{\prime\prime}}:\ j{\in}J\}$ be a $q$-$\delta$-reduction of $\{{\mathcal F}_{A_{i}^{\prime\prime},B_{i}^{\prime\prime}}^{\rho_{i}^{\prime\prime}}:\ i{\in}I\}$.
Since $\{{\mathcal F}_{A_{i}^{\prime\prime},B_{i}^{\prime\prime}}^{\rho_{i}^{\prime\prime}}:\ i{\in}I\}$ is dust-bounded, root-bounded and kernel-bounded, then by Proposition~\ref{lemma:good:properties:of:delta:reductions:indexed}, $\{{\mathcal F}_{A_{j}^{\prime\prime},B_{j}^{\prime\prime}}^{\rho_{j}^{\prime\prime}}:\ j{\in}J\}$ is dust-bounded, root-bounded and kernel-bounded.
Moreover, since for all $i{\in}I$, ${\parallel}A_{i}^{\prime\prime}{\parallel}{\leq}Q{\times}K$ and ${\parallel}B_{i}^{\prime\prime}{\parallel}{\leq}Q{\times}K$, then by Proposition~\ref{lemma:good:properties:of:delta:reductions:indexed}, ${\parallel}J{\parallel}{\leq}q{\times}(Q{\times}K{+}1)^{2}{\times}2^{(Q{\times}K)^{2}}$.
%Since the disjoint indexed family $\{(W_{A_{i}^{\prime\prime},B_{i}^{\prime\prime}}^{\rho_{i}^{\prime\prime}},R_{A_{i}^{\prime\prime},B_{i}^{\prime\prime}}^{\rho_{i}^{\prime\prime}}):\ i{\in}I\}$ of galaxies is $k^{\prime}$-prebounded, by Lemma~\ref{lemma:good:properties:of:delta:reductions:indexed}, ${\parallel}J{\parallel}{\leq}q{\times}N_{k^{\prime}}$.
In other respect, $(W^{\prime\prime\prime},R^{\prime\prime\prime})$ being the union of $\{{\mathcal F}_{A_{j}^{\prime\prime},B_{j}^{\prime\prime}}^{\rho_{j}^{\prime\prime}}:\ j{\in}J\}$, $(W^{\prime\prime\prime},R^{\prime\prime\prime})$ is a bounded morphic image of $(W^{\prime\prime},R^{\prime\prime})$ and the second player has a winning strategy in the Ehrenfeucht-Fra\"\i ss\'e game ${\mathcal G}_{q}((W^{\prime\prime},R^{\prime\prime}),(W^{\prime\prime\prime},R^{\prime\prime\prime}))$.
Hence, by Lemmas~\ref{lemma:delta:reduction:modal:logic} and~\ref{lemma:delta:reduction:fol}, $(W^{\prime\prime},R^{\prime\prime}){\preceq}(W^{\prime\prime\prime},R^{\prime\prime\prime})$ and $(W^{\prime\prime},R^{\prime\prime}){\equiv_{q}}(W^{\prime\prime\prime},R^{\prime\prime\prime})$.
%
%
\\
\\
%
%
Since $(W,R)$ validates $\L$, $(W,R)$ does not validate $A$, $(W,R){\preceq}(W^{\prime},R^{\prime})$, $(W,R)
%
%
%
%
%
%
%
%
%
%
%
%
$\linebreak$
%
%
%
%
%
%
%
%
%
%
%
%
{\equiv_{q}}(W^{\prime},R^{\prime})$, $(W^{\prime},R^{\prime}){\preceq}(W^{\prime\prime},R^{\prime\prime})$ and $(W^{\prime},R^{\prime}){\equiv_{q}}(W^{\prime\prime},R^{\prime\prime})$, then $(W^{\prime\prime\prime},R^{\prime\prime\prime})$ validates $\L$ and $(W^{\prime\prime\prime},R^{\prime\prime\prime})$ does not validate $A$.
In other respect, ${\parallel}W^{\prime\prime\prime}{\parallel}{\leq}2{\times}q{\times}(Q{\times}
%
%
%
%
%
%
%
%
%
%
%
%
$\linebreak$
%
%
%
%
%
%
%
%
%
%
%
%
K{+}1)^{2}{\times}2^{(Q{\times}K)^{2}}{\times}Q{\times}K$.
%
%
\\
\\
%
%
All in all, we have proved that for all sentences $A$, $q$ denoting $\qdd(A)$, $Q$ denoting $2{\times}q{\times}(q{\times}(q{+}1)^{2}+1)$ and $K$ denoting $2^{k}$, the following conditions are equivalent:
%
%
\begin{itemize}
%
%
\item $A{\in}\Th(\Fr(\L))$,
%
%
\item for all finite Euclidean frames $(W,R)$, if ${\parallel}W{\parallel}{\leq}2{\times}q{\times}((Q{\times}K){+}1)^{2}{\times}
%
%
%
%
%
%
%
%
%
%
%
%
$\linebreak$
%
%
%
%
%
%
%
%
%
%
%
%
2^{(Q{\times}K)^{2}}{\times}Q{\times}K$ and $(W,R)$ validates $\varphi$ then $(W,R)$ validates $A$.
%
%
\end{itemize}
%
%
\medskip
%
%
\end{proof}
%
%
\begin{lemma}\label{lemma:infinite:undecidable:case}
%
%
Let $\L$ be an Euclidean modal logic.
If $\mathtt{S}_{\L}{\setminus}((\{1\}{\times}\mathbf{N}^{-}){\cup}(\mathbf{N}^{+}{\times}\{{-}1,
%
%
%
%
%
%
%
%
%
%
%
%
$\linebreak$
%
%
%
%
%
%
%
%
%
%
%
%
0,1\}))$ is infinite then $\Th(\Fr(\L))$ is undecidable.
%
%
\end{lemma}
%
%
\begin{proof}
%
%
Suppose $\mathtt{S}_{\L}{\setminus}((\{1\}{\times}\mathbf{N}^{-}){\cup}(\mathbf{N}^{+}{\times}\{{-}1,0,1\}))$ is infinite.
Hence, by Lemmas~\ref{lemma:if:S:L:setminus:some:pairs:is:infinite:then:two:consequences} and~\ref{lemma:S:L:is:closed:subset:of:N:plus:times:N:minus}, either $\{2\}{\times}\mathbf{N}^{-}{\subseteq}\mathtt{S}_{\L}$, or $\mathbf{N}^{+}{\times}\{2\}{\subseteq}\mathtt{S}_{\L}$.
Thus, by Lemmas~\ref{lemma:first:case:2:N:moins} and~\ref{lemma:second:case:N:plus:1}, $\Th(\Fr(\L))$ is undecidable.
%
%
\medskip
%
%
\end{proof}
%
%
\begin{proposition}\label{proposition:beta:beta:nsc:for:decidability:theory}
%
%
For all Euclidean modal logics $\L$, $\Th(\Fr(\L))$ is decidable if and only if $\mathtt{S}_{\L}{\setminus}((\{1\}{\times}\mathbf{N}^{-}){\cup}(\mathbf{N}^{+}{\times}\{{-}1,0,1\}))$ is finite.
In that case, $\Th(\Fr(\L))$ is in $\EXPSPACE$.
%
%
\end{proposition}
%
%
\begin{proof}
%
%
By Lemmas~\ref{lemma:validity:in:a:given:frame:is:in:coNP}, \ref{lemma:validity:in:a:given:frame:is:in:coNP:fol}, \ref{decidability:if:setminus:is:finite} and~\ref{lemma:infinite:undecidable:case}.
%
%
\medskip
%
%
\end{proof}
%
%
%\begin{proposition}\label{proposition:beta:when:L:is:finite:and:varphi:moreover,K:equal:2k:sentence}
%
%
%
%
%\end{proposition}
%
%
%\begin{proposition}\label{proposition:alpha:alpha:when:decidable:theory:in:PSPACE}
%
%
%For all Euclidean modal logics $\L$, if $\mathtt{S}_{\L}{\setminus}((\{1\}{\times}\mathbf{N}^{-}){\cup}(\mathbf{N}^{+}{\times}
%
%
%
%
%
%
%
%
%
%
%
%
%$\linebreak$
%
%
%
%
%
%
%
%
%
%
%
%
%\{{-}1,0,1\}))$ is finite then $\Th(\Fr(\L))$ is in $\EXPSPACE$.
%
%
%\end{proposition}
%
%
Remind that for all Euclidean modal logics $\L$, the problem of deciding the first-order definability of modal formulas with respect to $\Fr(\L)$ is trivial~\cite{Balbiani:Georgiev:Tinchev:2018}.
The following results show that for all Euclidean modal logics $\L$, the computability of the problem of deciding the modal definability of sentences with respect to $\Fr(\L)$ only depends on the cardinality of $\mathtt{S}_{\L}{\setminus}((\{1\}{\times}\mathbf{N}^{-}){\cup}(\mathbf{N}^{+}{\times}\{{-}1,0,1\}))$.
%
%
\begin{lemma}\label{proposition:gamma:gamma:L:finite:sl:setminus:ensemble:TFCAE:sentence:A}
%
%
Let $\L$ be an Euclidean modal logic such that $\mathtt{S}_{\L}{\setminus}((\{1\}{\times}\mathbf{N}^{-}){\cup}(\mathbf{N}^{+}{\times}
%
%
%
%
%
%
%
%
%
%
%
%
$\linebreak$
%
%
%
%
%
%
%
%
%
%
%
%
\{{-}1,0,1\}))$ is finite.
Let $\varphi$ be a modal formula such that $\L{=}\K5{\oplus}\varphi$.
Let $k$ be the least integer such that $k{\geq}4$ and $(\ast)$~for all $m{\in}\mathbf{N}^{+}$ and for all $n{\in}\mathbf{N}^{-}$, if $(m,n){\in}\mathtt{S}_{\L}{\setminus}((\{1\}{\times}\mathbf{N}^{-}){\cup}(\mathbf{N}^{+}{\times}\{{-}1,0,1\}))$ then $m{+}n{\leq}k$.
For all sentences $A$ and for all individual variables $\mathbf{x}$ not occurring in $A$, the following conditions are equivalent:
%
%
\begin{itemize}
%
%
\item $A$ is modally definable in $\Fr(\L)$,
%
%
\item $A{\leftrightarrow}\forall\mathbf{x}\tau(\mathbf{x},A)$ is in $\Th(\Fr(\L))$ and for all $m,m^{\prime}{\in}\mathbf{N}^{+}$ and for all $n,n^{\prime}{\in}\mathbf{N}^{-}$, if ${\mathcal F}_{m}^{n}{\models}\varphi$, $(m,n){\in}\Pi_{k}^{A}$, $m^{\prime}{\leq}m$, $n^{\prime}{\leq}n$ and ${\mathcal F}_{m}^{n}{\models}A$ then ${\mathcal F}_{m^{\prime}}^{n^{\prime}}{\models}A$.
%
%
\end{itemize}
%
%
\end{lemma}
%
%
\begin{proof}
%
%
Let $A$ be a sentence and $\mathbf{x}$ be an individual variable not occurring in $A$.
%
%
\\
\\
%
%
$\mathbf{(\Rightarrow)}$~Suppose $A$ is modally definable in $\Fr(\L)$.
Hence, by Proposition~\ref{proposition:characterizing:modal:definability}, $\mathbf{(i)}$~for all Euclidean frames $(W,R)$, if $(W,R){\models}\L$ then $(W,R){\models}A$ if and only if for all $s{\in}W$, $(W_{s},R_{s}){\models}A$ and $\mathbf{(ii)}$~for all flowers ${\mathcal F}_{m^{\prime\prime}}^{n^{\prime\prime}},{\mathcal F}_{m^{\prime\prime\prime}}^{n^{\prime\prime\prime}}$, if ${\mathcal F}_{m^{\prime\prime}}^{n^{\prime\prime}}{\models}\L$, ${\mathcal F}_{m^{\prime\prime}}^{n^{\prime\prime}}{\models}A$ and ${\mathcal F}_{m^{\prime\prime\prime}}^{n^{\prime\prime\prime}}$ is a bounded morphic image of ${\mathcal F}_{m^{\prime\prime}}^{n^{\prime\prime}}$ then ${\mathcal F}_{m^{\prime\prime\prime}}^{n^{\prime\prime\prime}}{\models}A$.
Thus, by Lemma~\ref{lemma:about:the:rooted:translation}, for all Euclidean frames $(W,R)$, if $(W,R){\models}\L$ then $(W,R){\models}A$ if and only if $(W,R){\models}\forall\mathbf{x}\tau(\mathbf{x},A)$.
Consequently, $A{\leftrightarrow}\forall\mathbf{x}\tau(\mathbf{x},A)$ is in $\Th(\Fr(\L))$.
For the sake of the contradiction, suppose there exists $m,m^{\prime}{\in}\mathbf{N}^{+}$ and there exists $n,n^{\prime}{\in}\mathbf{N}^{-}$ such that ${\mathcal F}_{m}^{n}{\models}\varphi$, $(m,n){\in}\Pi_{k}^{A}$, $m^{\prime}{\leq}m$, $n^{\prime}{\leq}n$, ${\mathcal F}_{m}^{n}{\models}A$ and ${\mathcal F}_{m^{\prime}}^{n^{\prime}}{\not\models}A$.
Hence, ${\mathcal F}_{m}^{n}{\models}\L$.
Moreover, by Lemma~\ref{lemma:about:bounded:morphisms:in:the:situation:of F:m:n:frames}, ${\mathcal F}_{m^{\prime}}^{n^{\prime}}$ is a bounded morphic image of ${\mathcal F}_{m}^{n}$.
Since ${\mathcal F}_{m}^{n}{\models}A$, then by $\mathbf{(ii)}$, ${\mathcal F}_{m^{\prime}}^{n^{\prime}}{\models}A$: a contradiction.
%
%
\\
\\
%
%
$\mathbf{(\Leftarrow)}$~Suppose $\mathbf{(iii)}$~$A{\leftrightarrow}\forall\mathbf{x}\tau(\mathbf{x},A)$ is in $\Th(\Fr(\L))$ and $\mathbf{(iv)}$~for all $m,m^{\prime}{\in}\mathbf{N}^{+}$ and for all $n,n^{\prime}{\in}\mathbf{N}^{-}$, if ${\mathcal F}_{m}^{n}{\models}\varphi$, $(m,n){\in}\Pi_{k}^{A}$, $m^{\prime}{\leq}m$, $n^{\prime}{\leq}n$ and ${\mathcal F}_{m}^{n}{\models}A$ then ${\mathcal F}_{m^{\prime}}^{n^{\prime}}{\models}A$.
For the sake of the contradiction, suppose $A$ is not modally definable in $\Fr(\L)$.
Since $\mathbf{(iii)}$, then for all Euclidean frames $(W,R)$, if $(W,R){\models}\L$ then $(W,R){\models}A$ if and only if $(W,R){\models}\forall\mathbf{x}\tau(\mathbf{x},A)$.
Thus, by Lemma~\ref{lemma:about:the:rooted:translation}, for all Euclidean frames $(W,R)$, if $(W,R){\models}\L$ then $(W,R){\models}A$ if and only if for all $s{\in}W$, $(W_{s},R_{s}){\models}A$.
Since $A$ is not modally definable in $\Fr(\L)$, then by Proposition~\ref{proposition:characterizing:modal:definability}, there exists flowers ${\mathcal F}_{m^{\prime\prime}}^{n^{\prime\prime}},{\mathcal F}_{m^{\prime\prime\prime}}^{n^{\prime\prime\prime}}$ such that ${\mathcal F}_{m^{\prime\prime}}^{n^{\prime\prime}}{\models}\L$, ${\mathcal F}_{m^{\prime\prime}}^{n^{\prime\prime}}{\models}A$, ${\mathcal F}_{m^{\prime\prime\prime}}^{n^{\prime\prime\prime}}$ is a bounded morphic image of ${\mathcal F}_{m^{\prime\prime}}^{n^{\prime\prime}}$ and ${\mathcal F}_{m^{\prime\prime\prime}}^{n^{\prime\prime\prime}}{\not\models}A$.
Consequently, ${\mathcal F}_{m^{\prime\prime}}^{n^{\prime\prime}}{\models}\varphi$.
Moreover, by Lemma~\ref{lemma:about:bounded:morphisms:in:the:situation:of F:m:n:frames}, $m^{\prime\prime\prime}{\leq}m^{\prime\prime}$ and $n^{\prime\prime\prime}{\leq}n^{\prime\prime}$.
Let ${\mathcal S}{=}\{(m,n):\ m{\in}\mathbf{N}^{+},\ n{\in}\mathbf{N}^{-},\ {\mathcal F}_{m}^{n}{\models}\varphi,\ {\mathcal F}_{m}^{n}{\models}A$ and there exists $m^{\prime}{\in}\mathbf{N}^{+}$ and $n^{\prime}{\in}\mathbf{N}^{-}$ such that $m^{\prime}{\leq}m$, $n^{\prime}{\leq}n$ and ${\mathcal F}_{m^{\prime}}^{n^{\prime}}{\not\models}A\}$.
Since ${\mathcal F}_{m^{\prime\prime}}^{n^{\prime\prime}}{\models}A$, ${\mathcal F}_{m^{\prime\prime\prime}}^{n^{\prime\prime\prime}}{\not\models}A$, ${\mathcal F}_{m^{\prime\prime}}^{n^{\prime\prime}}{\models}\varphi$, $m^{\prime\prime\prime}{\leq}m^{\prime\prime}$ and $n^{\prime\prime\prime}{\leq}n^{\prime\prime}$, then $(m^{\prime\prime},n^{\prime\prime}){\in}{\mathcal S}$.
Let $(m,n)$ be a $\ll$-minimal element in ${\mathcal S}$.
Hence, ${\mathcal F}_{m}^{n}{\models}\varphi$, ${\mathcal F}_{m}^{n}{\models}A$ and there exists $m^{\prime}{\in}\mathbf{N}^{+}$ and $n^{\prime}{\in}\mathbf{N}^{-}$ such that $m^{\prime}{\leq}m$, $n^{\prime}{\leq}n$ and ${\mathcal F}_{m^{\prime}}^{n^{\prime}}{\not\models}A$.
Thus, $(m,n){\in}\mathtt{S}_{\L}$.
Moreover,
%\mathtt{S}_{\L}$ such that ${\mathcal F}_{m}^{n}{\models}A$ and there exists $m^{\prime}{\in}\mathbf{N}^{+}$ and $n^{\prime}{\in}\mathbf{N}^{-}$ such that $m^{\prime}{\leq}m$, $n^{\prime}{\leq}n$ and ${\mathcal F}_{m^{\prime}}^{n^{\prime}}{\not\models}A$.
by $\mathbf{(iv)}$, $(m,n){\not\in}\Pi_{k}^{A}$.
Thus, either $m{\geq}2$, $n{\geq}2$ and $m{+}n{>}k$, or $m{=}1$ and $n{>}\qdd(A)$, or either $n{=}{-}1$, or $n{=}0$, or $n{=}1$ and $m{>}\qdd(A)$.
In the first case, since $(m,n){\in}\mathtt{S}_{\L}$, then $(m,n){\in}\mathtt{S}_{\L}{\setminus}((\{1\}{\times}\mathbf{N}^{-}){\cup}(\mathbf{N}^{+}{\times}\{{-}1,0,1\}))$.
Consequently, by~$(\ast)$, $m{+}n{\leq}k$: a contradiction.
In the second case, since ${\mathcal F}_{m}^{n}{\models}A$ and ${\mathcal F}_{m^{\prime}}^{n^{\prime}}{\not\models}A$, then ${\mathcal F}_{m}^{n}{\not\equiv_{\qdd(A)}}{\mathcal F}_{m^{\prime}}^{n^{\prime}}$.
Since $m^{\prime}{\leq}m$, $n^{\prime}{\leq}n$, $m{=}1$ and $n{>}\qdd(A)$, then $(m,n{-}1){\in}{\mathcal S}$: a contradiction with the $\ll$-minimality of $(m,n)$ in ${\mathcal S}$.
In the third case, since ${\mathcal F}_{m}^{n}{\models}A$ and ${\mathcal F}_{m^{\prime}}^{n^{\prime}}{\not\models}A$, then ${\mathcal F}_{m}^{n}{\not\equiv_{\qdd(A)}}{\mathcal F}_{m^{\prime}}^{n^{\prime}}$.
Since $m^{\prime}{\leq}m$, $n^{\prime}{\leq}n$, either $n{=}{-}1$, or $n{=}0$, or $n{=}1$ and $m{>}\qdd(A)$, then $(m{-}1,n){\in}{\mathcal S}$: a contradiction with the $\ll$-minimality of $(m,n)$ in ${\mathcal S}$.
% DELTA DELTA DELTA DELTA DELTA
%
%
\medskip
%
%
\end{proof}
%
%
%\begin{lemma}\label{proposition:decidability:of:modal:definability:first:situation}
%
%
%Let $\L$ be an Euclidean modal logic.
%If $\mathtt{S}_{\L}{\setminus}((\{1\}{\times}\mathbf{N}^{-}){\cup}(\mathbf{N}^{+}{\times}\{{-}1,
%
%
%
%
%
%
%
%
%
%
%
%
%$\linebreak$
%
%
%
%
%
%
%
%
%
%
%
%
%0,1\}))$ is finite then the problem of deciding the modal definability of sentences with respect to $\Fr(\L)$ is decidable.
%
%
%\end{lemma}
%
%
\begin{lemma}\label{proposition:undecidability:of:modal:definability}
%
%
Let $\L$ be an Euclidean modal logic.
If $\mathtt{S}_{\L}{\setminus}((\{1\}{\times}\mathbf{N}^{-}){\cup}(\mathbf{N}^{+}{\times}\{{-}1,
%
%
%
%
%
%
%
%
%
%
%
%
$\linebreak$
%
%
%
%
%
%
%
%
%
%
%
%
0,1\}))$ is infinite then the problem of deciding the modal definability of sentences with respect to $\Fr(\L)$ is undecidable.
%
%
\end{lemma}
%
%
\begin{proof}
%
%
By Lemmas~\ref{Lemma:53:bis} and~\ref{lemma:infinite:undecidable:case}.
%
%
\medskip
%
%
\end{proof}
%
%
\begin{theorem}\label{proposition:finale}
%
%
For all Euclidean modal logics $\L$, the problem of deciding the modal definability of sentences with respect to $\Fr(\L)$ is decidable if and only if $\mathtt{S}_{\L}{\setminus}((\{1\}{\times}\mathbf{N}^{-}){\cup}(\mathbf{N}^{+}{\times}\{{-}1,0,1\}))$ is finite.
In that case,
%
%
%\end{theorem}
%
%
%\begin{theorem}\label{proposition:finale:complexity}
%
%
%For all Euclidean modal logics $\L$, if $\mathtt{S}_{\L}{\setminus}((\{1\}{\times}\mathbf{N}^{-}){\cup}(\mathbf{N}^{+}{\times}\{{-}1,0,
%
%
%
%
%
%
%
%
%
%
%
%
%$\linebreak$
%
%
%
%
%
%
%
%
%
%
%
%
%1\}))$ is finite then
the problem of deciding the modal definability of sentences with respect to $\Fr(\L)$ is in $\EXPSPACE$.
%
%
\end{theorem}
%
%
\begin{proof}
%
%
By Lemmas~\ref{lemma:validity:in:a:given:frame:is:in:coNP}, \ref{lemma:ll:finitely:many:PI:PI:pairs}, \ref{lemma:validity:in:a:given:frame:is:in:coNP:fol}, \ref{proposition:gamma:gamma:L:finite:sl:setminus:ensemble:TFCAE:sentence:A} and~\ref{proposition:undecidability:of:modal:definability} and Proposition~\ref{proposition:beta:beta:nsc:for:decidability:theory}.
%
%
\medskip
%
%
\end{proof}
%
%
%\begin{proposition}\label{proposition:finale:pspace:definab}
%
%
%For all modal logics $\L$, if $\mathtt{S}_{\L}{\setminus}((\{1\}{\times}\mathbf{N}^{-}){\cup}(\mathbf{N}^{+}{\times}\{{-}1,0,1\}))$ is finite then the problem of deciding the modal definability of sentences with respect to $\Fr(\L)$ is $\PSPACE$-complete.
%
%
%\end{proposition}
%
%
%\begin{lemma}\label{proposition:decidability:of:correspondence:problem:first:situation}
%
%
%Let $\L$ be an Euclidean modal logic.
%If $\mathtt{S}_{\L}{\setminus}((\{1\}{\times}\mathbf{N}^{-}){\cup}(\mathbf{N}^{+}{\times}\{{-}1,
%
%
%
%
%
%
%
%
%
%
%
%
%$\linebreak$
%
%
%
%
%
%
%
%
%
%
%
%
%0,1\}))$ is finite then the problem of deciding the correspondence of modal formulas and sentences with respect to $\Fr(\L)$ is decidable.
%
%
%\end{lemma}
%
%
%\begin{lemma}\label{proposition:undecidability:of:correspondence:problem}
%
%
%Let $\L$ be an Euclidean modal logic.
%If $\mathtt{S}_{\L}{\setminus}((\{1\}{\times}\mathbf{N}^{-}){\cup}(\mathbf{N}^{+}{\times}\{{-}1,
%
%
%
%
%
%
%
%
%
%
%
%
%$\linebreak$
%
%
%
%
%
%
%
%
%
%
%
%
%0,1\}))$ is infinite then the problem of deciding the correspondence of modal formulas and sentences with respect to $\Fr(\L)$ is undecidable.
%
%
%\end{lemma}
%
%
\begin{theorem}\label{proposition:finale:correspondence:problem}
%
%
For all Euclidean modal logics $\L$, the problem of deciding the correspondence of modal formulas and sentences with respect to $\Fr(\L)$ is decidable if and only if $\mathtt{S}_{\L}{\setminus}((\{1\}{\times}\mathbf{N}^{-}){\cup}(\mathbf{N}^{+}{\times}\{{-}1,0,1\}))$ is finite.
In that case,
%
%
%\end{theorem}
%
%
%\begin{theorem}\label{proposition:finale:correspondence:problem:complexity}
%
%
%For all Euclidean modal logics $\L$, if $\mathtt{S}_{\L}{\setminus}((\{1\}{\times}\mathbf{N}^{-}){\cup}(\mathbf{N}^{+}{\times}\{{-}1,0,
%
%
%
%
%
%
%
%
%
%
%
%
%$\linebreak$
%
%
%
%
%
%
%
%
%
%
%
%
%1\}))$ is finite then
the problem of deciding the correspondence of modal formulas and sentences with respect to $\Fr(\L)$ is in $\EXPSPACE$.
%
%
\end{theorem}
%
%
\begin{proof}
%
%
Let $\L$ be an Euclidean modal logics.
%
%
\\
\\
%
%
Suppose $\mathtt{S}_{\L}{\setminus}((\{1\}{\times}\mathbf{N}^{-}){\cup}(\mathbf{N}^{+}{\times}\{{-}1,0,1\}))$ is infinite.
Hence, by Proposition~\ref{proposition:beta:beta:nsc:for:decidability:theory}, $\Th(\Fr(\L))$ is undecidable.
Since for all sentences $A$, the modal formula $\top$ and $A$ correspond in $\Fr(\L)$ if and only if $A$ is in $\Th(\Fr(\L))$, then the problem of deciding the correspondence of modal formulas and sentences with respect to $\Fr(\L)$ is undecidable.
%
%
\\
\\
%
%
Suppose Suppose $\mathtt{S}_{\L}{\setminus}((\{1\}{\times}\mathbf{N}^{-}){\cup}(\mathbf{N}^{+}{\times}\{{-}1,0,1\}))$ is finite.
Thus, by Proposition~\ref{proposition:beta:beta:nsc:for:decidability:theory}, $\Th(\Fr(\L))$ is in $\EXPSPACE$.
According to Balbiani {\it et al.}\/~\cite[Lem\-ma~$18$]{Balbiani:Georgiev:Tinchev:2018}, given an arbitrary modal formula $\varphi$, one can easily construct a sentence $A(\varphi)$ such that $\size(A(\varphi)){=}{\mathcal O}(\size(\varphi))$ and $\varphi$ and $A(\varphi)$ correspond in the class of all Euclidean frames.
%%%%%%%
%%%%%%%A VERIFIER : $\size(A(\varphi)){=}{\mathcal O}(\size(\varphi))$ !!!!!!!
%%%%%%%
Thus, for all modal formulas $\varphi$ and for all sentences $A$, $\varphi$ and $A$ correspond in $\Fr(\L)$ if and only if $A\leftrightarrow A(\varphi)$ is in $\Th(\Fr(\L))$.
Since $\Th(\Fr(\L))$ is in $\EXPSPACE$, then the problem of deciding the correspondence of modal formulas and sentences with respect to $\Fr(\L)$ is in $\EXPSPACE$.
% ICI ICI ICI ICI ICI
%
%
\medskip
%
%
\end{proof}
%
%
Now, it is time to assert and prove the following additional results.
%
%
\begin{corollary}\label{corollary:modal:definability:complexity}
%
%
The following decision problem is in $\NEXPTIME$:
%
%
\begin{description}
%
%
\item[input:] a modal formula $\varphi$,
%
%
\item[output:] determine whether the problem of deciding the modal definability of sentences with respect to $\Fr(\K5{\oplus}\varphi)$ is decidable.
%
%
\end{description}
%
%
\end{corollary}
%
%
\begin{proof}
%
%
Let $\varphi$ be a modal formula.
Obviously, $l$ denoting $2^{{\parallel}\vvar(\varphi){\parallel}}$, the following conditions are equivalent:
%
%
\begin{enumerate}
%
%
\item the modal definability of sentences with respect to $\Fr(\K5{\oplus}\varphi)$ is decidable,
%
%
\item $\mathtt{S}_{\K5{\oplus}\varphi}{\setminus}((\{1\}{\times}\mathbf{N}^{-}){\cup}(\mathbf{N}^{+}{\times}\{{-}1,0,1\}))$ is finite,
%
%
\item $\{2\}{\times}\mathbf{N}^{-}{\not\subseteq}\mathtt{S}_{\K5{\oplus}\varphi}$ and $\mathbf{N}^{+}{\times}\{2\}{\not\subseteq}\mathtt{S}_{\K5{\oplus}\varphi}$,
%
%
\item there exists $m{\in}\mathbf{N}^{+}$ such that ${\mathcal F}_{m}^{2}{\not\models}\varphi$ and there exists $n{\in}\mathbf{N}^{-}$ such that ${\mathcal F}_{2}^{n}{\not\models}\varphi$,
%
%
\item there exists $m{\in}\lsem1,l\rsem$ such that ${\mathcal F}_{m}^{2}{\not\models}\varphi$ and there exists $n{\in}\lsem{-}1,l\rsem$ such that ${\mathcal F}_{2}^{n}{\not\models}\varphi$,
%
%
\item ${\mathcal F}_{l}^{2}{\not\models}\varphi$ and either ${\mathcal F}_{2}^{{-}1}{\not\models}\varphi$, or ${\mathcal F}_{2}^{0}{\not\models}\varphi$, or ${\mathcal F}_{2}^{l}{\not\models}\varphi$.
%
%
\end{enumerate}
%
%
Indeed, the equivalence between~$\mathbf{(1)}$ and~$\mathbf{(2)}$ is a consequence of Theorem~\ref{proposition:finale}, the equivalence between~$\mathbf{(2)}$ and~$\mathbf{(3)}$ is a consequence of Lemmas~\ref{lemma:if:S:L:setminus:some:pairs:is:infinite:then:two:consequences} and~\ref{lemma:S:L:is:closed:subset:of:N:plus:times:N:minus}, the equivalence between~$\mathbf{(4)}$ and~$\mathbf{(5)}$ is a consequence of Lemmas~\ref{lemma:about:upper:bound:on:the:kernel:1} and~\ref{lemma:about:upper:bound:on:the:kernel:2} and the equivalence between~$\mathbf{(5)}$ and~$\mathbf{(6)}$ is a consequence of Lemmas~\ref{lemma:about:bounded:morphisms:in:the:situation:of F:m:n:frames} and~\ref{Bounded:Morphism:Lemma}.
Hence, by Lemma~\ref{lemma:validity:in:a:given:frame:is:in:coNP}, the decision problem considered in Corollary~\ref{corollary:modal:definability:complexity} is in $\NEXPTIME$.
%
%
\medskip
%
%
\end{proof}
%
%
\begin{corollary}\label{corollary:correspondence:complexity}
%
%
The following decision problem is in $\NEXPTIME$:
%
%
\begin{description}
%
%
\item[input:] a modal formula $\varphi$,
%
%
\item[output:] determine whether the problem of deciding the correspondence of modal formulas and sentences with respect to $\Fr(\K5{\oplus}\varphi)$ is decidable.
%
%
\end{description}
%
%
\end{corollary}
%
%
\begin{proof}
%
%
Let $\varphi$ be a modal formula.
Obviously, $l$ denoting $2^{{\parallel}\vvar(\varphi){\parallel}}$, the following conditions are equivalent:
%
%
\begin{enumerate}
%
%
\item the correspondence of modal formulas and sentences with respect to $\Fr(\K5
%
%
%
%
%
%
%
%
%
%
%
%
$\linebreak$
%
%
%
%
%
%
%
%
%
%
%
%
{\oplus}\varphi)$ is decidable,
%
%
\item $\mathtt{S}_{\K5{\oplus}\varphi}{\setminus}((\{1\}{\times}\mathbf{N}^{-}){\cup}(\mathbf{N}^{+}{\times}\{{-}1,0,1\}))$ is finite,
%
%
\item $\{2\}{\times}\mathbf{N}^{-}{\not\subseteq}\mathtt{S}_{\K5{\oplus}\varphi}$ and $\mathbf{N}^{+}{\times}\{2\}{\not\subseteq}\mathtt{S}_{\K5{\oplus}\varphi}$,
%
%
\item there exists $m{\in}\mathbf{N}^{+}$ such that ${\mathcal F}_{m}^{2}{\not\models}\varphi$ and there exists $n{\in}\mathbf{N}^{-}$ such that ${\mathcal F}_{2}^{n}{\not\models}\varphi$,
%
%
\item there exists $m{\in}\lsem1,l\rsem$ such that ${\mathcal F}_{m}^{2}{\not\models}\varphi$ and there exists $n{\in}\lsem{-}1,l\rsem$ such that ${\mathcal F}_{2}^{n}{\not\models}\varphi$,
%
%
\item ${\mathcal F}_{l}^{2}{\not\models}\varphi$ and either ${\mathcal F}_{2}^{{-}1}{\not\models}\varphi$, or ${\mathcal F}_{2}^{0}{\not\models}\varphi$, or ${\mathcal F}_{2}^{l}{\not\models}\varphi$.
%
%
\end{enumerate}
%
%
Indeed, the equivalence between~$\mathbf{(1)}$ and~$\mathbf{(2)}$ is a consequence of Theorem~\ref{proposition:finale:correspondence:problem}, the equivalence between~$\mathbf{(2)}$ and~$\mathbf{(3)}$ is a consequence of Lemmas~\ref{lemma:if:S:L:setminus:some:pairs:is:infinite:then:two:consequences} and~\ref{lemma:S:L:is:closed:subset:of:N:plus:times:N:minus}, the equivalence between~$\mathbf{(4)}$ and~$\mathbf{(5)}$ is a consequence of Lemmas~\ref{lemma:about:upper:bound:on:the:kernel:1} and~\ref{lemma:about:upper:bound:on:the:kernel:2} and the equivalence between~$\mathbf{(5)}$ and~$\mathbf{(6)}$ is a consequence of Lemmas~\ref{lemma:about:bounded:morphisms:in:the:situation:of F:m:n:frames} and~\ref{Bounded:Morphism:Lemma}.
Hence, by Lemma~\ref{lemma:validity:in:a:given:frame:is:in:coNP}, the decision problem considered in Corollary~\ref{corollary:correspondence:complexity} is in $\NEXPTIME$.
%
%
\medskip
%
%
\end{proof}
%
%
%\begin{proposition}\label{proposition:finale:correspondence:problem:pspace}
%
%
%For all modal logics $\L$, if $\mathtt{S}_{\L}{\setminus}((\{1\}{\times}\mathbf{N}^{-}){\cup}(\mathbf{N}^{+}{\times}\{{-}1,0,1\}))$ is finite then the problem of deciding the correspondence of modal formulas and sentences with respect to $\Fr(\L)$ is $\PSPACE$-complete.
%
%
%\end{proposition}
%
%
%\begin{proof}
%
%
%zzzzz
%
%
%\medskip
%
%
%\end{proof}
%
%
%\begin{lemma}\label{lemma:first:order:definability:finite:frames}
%
%
%Let $\L$ be a modal logic.
%For all modal formulas $\varphi$ a,d for all sentences $A$,
%
%
%\begin{enumerate}
%
%
%\item if $\varphi$ is first-order definable in $\Fr_{\fin}(\L)$ then $\varphi$ is first-order definable in $\Fr(\L)$,
%
%
%\item if $A$ is modally definable in $\Fr_{\fin}(\L)$ then $A$ is modally definable in $\Fr(\L)$,
%
%
%\item if $\varphi$ and $A$ correspond in $\Fr_{\fin}(\L)$ then $\varphi$ and $A$ correspond in $\Fr(\L)$.
%
%
%\end{enumerate}
%
%
%\end{lemma}
%
%
%\begin{proof}
%
%
%zzzzz
%
%
%\medskip
%
%
%\end{proof}
%
%
%
%
%\section{Computational complexity}\label{section:computational:complexity}
%
%
%zzzzz
%
%
%
%
\section{Conclusion}\label{section:conclusion}
%
%
This paper was about the computability of the modal definability problem in classes of frames determined by Euclidean modal logics.
We have characterized those Euclidean modal logics such that the classes of frames they determine give rise to an undecidable modal definability problem.
Much remains to be done.
%
%
\\
\\
%
%
A first obvious question is whether the problem of deciding the modal definability of sentences considered in Theorem~\ref{proposition:finale} and the problem of deciding the correspondence of modal formulas and sentences considered in Theorem~\ref{proposition:finale:correspondence:problem}~---~when they are decidable~---~are $\EXPSPACE$-hard.
%
%
\\
\\
%
%
A second obvious question is whether there exists other classes of frames for which modal definability is decidable.
Is the modal definability problem in classes of frames determined by extensions of $\S4.3$ decidable?
What about the first-order definability problem in classes of frames determined by extensions of $\S4.3$?
Another question is whether there exists other modal logics for which the first-order definability problem is trivial in the classes of frames they determine.
%
%
\\
\\
%
%
It is also worth to consider restrictions or extensions of the ordinary language of modal logic.
For example, one can consider either the implication fragment based on $\rightarrow$ and $\Box$, or the positive fragment based on $\vee$, $\wedge$, $\Box$ and $\Diamond$, or the tense extension, or the extension with the universal modality.
For such restrictions or extensions of the ordinary language of modal logic, what is the computability of the first-order definability problem and the modal definability problem?
%
%
\\
\\
%
%
And in the end, there is the open question of the existence of modal logics determining classes of frames in which the modal definability problem is decidable and the first-order definability problem is undecidable.
We conjecture that such modal logics do not exist.
%
%
%
%
\section*{Funding}
%
%
The preparation of this paper has been supported by the {\it Rila Programme} (project $\mathbf{KP{-}06{-}RILA{/}4}$ of the Bulgarian Ministry of Education and Science and project $\mathbf{48055QG}$ of the French Ministry for Europe and Foreign Affairs and the French Ministry of Higher Education and Research).
%
%
%
%
\section*{Acknowledgement}
%
%
Special acknowledgement is heartily granted to our colleagues of the Faculty of Mathematics and Informatics (Sofia, Bulgaria) and the Toulouse Institute of Computer Science Research (Toulouse, France) for many stimulating discussions about the subject of this paper.
%We also make a point of strongly thanking the referees for their feedback: their useful suggestions have been essential for improving the correctness and the readability of a preliminary version of this paper.
%
%
%
%
\bibliographystyle{named}
%
%
\begin{thebibliography}{}
%
%
\bibitem{Arnold:Crubille:1988}
{\sc Arnold,} A., and P. {\sc Crubill\'e,}
`A linear algorithm to solve fixed-point equations on transition systems',
\emph{Information Processing Letters}\/ {\bf 29} (1988) 57--66.
%
%
%\bibitem{Baader:1998}
%{\sc Baader,} F.,
%`On the complexity of Boolean unification',
%\emph{Information Processing Letters}\/ {\bf 67} (1998) 215--220.
%
%
%\bibitem{Baader:Ghilardi:2011}
%{\sc Baader,} F., and S. {\sc Ghilardi,}
%`Unification in modal and description logics',
%\emph{Logic Journal of the IGPL}\/ {\bf 19} (2011) 705--730.
%
%
%\bibitem{Baader:Morawska:2009}
%{\sc Baader,} F., and B. {\sc Morawska,}
%`Unification in the description logic $\mathcal{EL}$',
%In: \emph{Rewriting Techniques and Applications,}\/
%Springer (2009) 350--364.
%
%
%\bibitem{Baader:Narendran:2001}
%{\sc Baader,} F., and P. {\sc Narendran,}
%`Unification of concept terms in description logics',
%\emph{Journal of Symbolic Computation}\/ {\bf 31} (2001) 277--305.
%
%
%\bibitem{Baader:Snyder:2001}
%{\sc Baader,} F., and W. {\sc Snyder,}
%`Unification theory',
%In: \emph{Handbook of Automated Reasoning,}\/
%Elsevier (2001) 439--526.
%
%
%\bibitem{Babenyshev:Rybakov:2010}
%{\sc Babenyshev,} S., and V. {\sc Rybakov,}
%`Linear temporal logic $\LTL$: basis for admissible rules',
%\emph{Journal of and Computation}\/ {\bf 21} (2010) 157--177.
%
%
%\bibitem{Babenyshev:Rybakov:2011}
%{\sc Babenyshev,} S., and V. {\sc Rybakov,}
%`Unification in linear temporal logic $\LTL$',
%\emph{Annals of Pure and Applied Logic}\/ {\bf 162} (2011) 991--1000.
%
%
%\bibitem{Babenyshev:Rybakov:Schmidt:Tishkovsky:2010}
%{\sc Babenyshev,} S., V. {\sc Rybakov,} R. {\sc Schmidt,} and D. {\sc Tishkovsky,}
%`A tableau method for checking rule admissibility in $\S4$',
%\emph{Electronic Notes in Theoretical Computer Science}\/ {\bf 262} (2010) 17--32.
%
%
%\bibitem{Balbiani:2019}
%{\sc Balbiani,} P.,
%`Remarks about the unification type of several non-symmetric non-transitive modal logics',
%\emph{Logic Journal of the IGPL}\/ {\bf 27} (2019) 639--658.
%
%
%\bibitem{Balbiani:Dunchev:Tinchev:zzzzz}
%{\sc Balbiani,} P., T. {\sc Dunchev,} and T. {\sc Tinchev,}
%`Modal definability over a class of structures with two equivalence relations',
%zzzzz
%
%
%\bibitem{Balbiani:Gencer:2017a}
%{\sc Balbiani,} P., and \c{C}. {\sc Gencer,}
%`$\KD$ is nullary',
%\emph{Journal of Applied Non-Classical Logics}\/ {\bf 27} (2017) 196--205.
%
%
%\bibitem{Balbiani:Gencer:2017b}
%{\sc Balbiani,} P., and \c{C}. {\sc Gencer,}
%`Unification in epistemic logics',
%\emph{Journal of Applied Non-Classical Logics}\/ {\bf 27} (2017) 91--105.
%
%
%\bibitem{Balbiani:Gencer:to:appear}
%{\sc Balbiani,} P., and \c{C}. {\sc Gencer,}
%`About the unification type of modal logics between $\KB$ and $\KTB$',
%\emph{Studia Logica}\/ {\bf 108} (2020) 941--966.
%
%
\bibitem{Balbiani:Georgiev:Tinchev:2018}
{\sc Balbiani,} P., D. {\sc Georgiev,} and T. {\sc Tinchev,}
`Modal correspondence theory in the class of all Euclidean frames',
\emph{Journal of Logic and Computation}\/ {\bf 28} (2018) 119--131.
%
%
%\bibitem{Balbiani:Kikot:2012}
%{\sc Balbiani,} P., and S. {\sc Kikot,}
%`Sahlqvist theorems for precontact logics',
%In: \emph{Advances in Modal Logic. Vol. 9,}\/
%College Publications (2012) 55-70.
%
%
%\bibitem{Balbiani:Mojtahedi:submitted}
%{\sc Balbiani,} P., and M. {\sc Mojtahedi,}
%`Unification in the implication fragment of Classical Propositional Logic',
%\emph{Logic Journal of the IGPL}\/ (to appear).
%
%
%\bibitem{Balbiani:Shapirovsky:Shehtman:2006}
%{\sc Balbiani,} P., I. {\sc Shapirovsky,} and V. {\sc Shehtman,}
%`Every world can see a Sahlqvist world',
%In: \emph{Advances in Modal Logic,}\/
%College Publications (2006) 69--85.
%
%
\bibitem{Balbiani:Tinchev:2005}
{\sc Balbiani,} P., and T. {\sc Tinchev,}
`Decidability and complexity of definability over the class of all partitions',
In: \emph{Proceedings of the 5th Panhellenic Logic Symposium,}\/
University of Athens (2005) 26--33.
%
%
\bibitem{Balbiani:Tinchev:2006}
{\sc Balbiani,} P., and T. {\sc Tinchev,}
`Definability over the class of all partitions',
\emph{Journal of Logic and Computation}\/ {\bf 16} (2006) 541--557.
%
%
%\bibitem{Balbiani:Tinchev:2014}
%{\sc Balbiani,} P., and T. {\sc Tinchev,}
%`Definability and canonicity for Boolean logic with a binary relation',
%\emph{Fundamenta Informatic\ae}\/ {\bf 129} (2014) 301--327.
%
%
%\bibitem{Balbiani:Tinchev:2016}
%{\sc Balbiani,} P., and T. {\sc Tinchev,}
%`Unification in modal logic $\Alt_{1}$',
%In: \emph{Advances in Modal Logic,}\/
%College Publications (2016) 117--134.
%
%
\bibitem{Balbiani:Tinchev:2017}
{\sc Balbiani,} P., and T. {\sc Tinchev,}
`Undecidable problems for modal definability',
\emph{Journal of Logic and Computation}\/ {\bf 27} (2017) 901--920.
%
%
%\bibitem{Balbiani:Tinchev:2018}
%{\sc Balbiani,} P., and T. {\sc Tinchev,}
%`Elementary unification in modal logic $\KD45$',
%\emph{Journal of Applied Logics}\/ {\bf 5} (2018) 301--317.
%
%
\bibitem{Balbiani:Tinchev:2019}
{\sc Balbiani,} P., and T. {\sc Tinchev,}
`Decidable and undecidable problems for first-order definability and modal definability',
In: \emph{Logic, Language, and Computation,}\/
Springer (2021) 214--236.
%
%
%\bibitem{Balbiani:Tinchev:Vakarelov:2007}
%{\sc Balbiani,} P., T. {\sc Tinchev,} and D. {\sc Vakarelov,}
%`Modal logics for region-based theories of space',
%\emph{Fundamenta Informatic\ae}\/ {\bf 81} (2007) 29--82.
%
%
%\bibitem{Baltag:Moss:2004}
%{\sc Baltag,} A., and L. {\sc Moss.}
%`Logics for epistemic programs',
%\emph{Synthese}\/ {\bf 139} (2004) 165--224.
%
%
\bibitem{vanBenthem:1975}
{\sc Van Benthem,} J.,
`A note on modal formulas and relational properties',
\emph{Journal of Symbolic Logic}\/ {\bf 40} (1975) 85--88.
%
%
\bibitem{vanBenthem:1983}
{\sc Van Benthem,} J.,
`Modal Logic and Classical Logic',
Bibliopolis (1983).
%
%
\bibitem{vanBenthem:1984}
{\sc Van Benthem,} J.,
`Correspondence theory',
In: \emph{Handbook of Philosophical Logic: Volume II: Extensions of Classical Logic,}\/
Reidel (1984) 167--247.
%
%
%\bibitem{Bezhanishvili:Marx:2003}
%{\sc Bezhanishvili,} N., M. {\sc Marx,}
%`All proper normal extensions of $\S5$-square have the polynomial size model property',
%\emph{Studia Logica}\/ {\bf 73} (2003) 367--382.
%
%
%\bibitem{Biere:et:al:2009}
%{\sc Biere,} A., M. {\sc Heule,} H. {\sc van Maaren,} and T. {\sc Walsh,}
%\emph{Handbook of Satisfiability,}\/
%IOS Press (2009).
%
%
\bibitem{Blackburn:deRijke:Venema:2001}
{\sc Blackburn,} P., M. {\sc de Rijke,} and Y. {\sc Venema,}
\emph{Modal Logic,}\/
Cambridge University Press (2001).
%
%
%\bibitem{Bolander:Andersen:2011}
%{\sc Bolander,} T., and M. {\sc Andersen,}
%`Epistemic planning for single- and multi-agent systems',
%\emph{Journal of Applied Non-Classical Logics}\/ {\bf 21} (2011) 9--34.
%
%
%\bibitem{Celani:2005}
%{\sc Celani,} S.,
%`Modal Tarski algebras',
%\emph{Reports on Mathematical Logic}\/ {\bf 39} (2005) 113--126.
%
%
%\bibitem{Chagrov:1992}
%{\sc Chagrov,} A.,
%`Decidable modal logic with undecidable admissibility problem',
%\emph{Algebra and Logic}\/ {\bf 31} (1992) 53--61.
%
%
\bibitem{Chagrov:Chagrova:1995}
{\sc Chagrov,} A., and L. {\sc Chagrova,}
`Algorithmic problems concerning first-order definability of modal formulas on the class of all finite frames',
\emph{Studia Logica}\/ {\bf 55} (1995) 421--448.
%
%
\bibitem{Chagrov:Chagrova:2006}
{\sc Chagrov,} A., and L. {\sc Chagrova,}
`The truth about algorithmic problems in correspondence theory',
In: \emph{Advances in Modal Logic. Vol. 6,}\/
College Publications (2006) 121--138.
%
%
\bibitem{Chagrov:Chagrova:2007}
{\sc Chagrov,} A., and L. {\sc Chagrova,}
`Demise of the algorithmic agenda in the correspondence theory~?',
\emph{Logical Investigations}\/ {\bf 13} (2007) 224--248.
%
%
%\bibitem{Chagrov:Zakharyaschev:1997}
%{\sc Chagrov,} A., and M. {\sc Zakharyaschev,}
%\emph{Modal Logic,}\/
%Oxford University Press (1997).
%
%
\bibitem{Chagrova:1989}
{\sc Chagrova,} L.,
\emph{On the Problem of Definability of Propositional Formulas of Intuitionistic Logic by Formulas of Classical First-Order Logic,}\/
Doctoral thesis of the University of Kalinin (1989).
%
%
\bibitem{Chagrova:1991}
{\sc Chagrova,} L.,
`An undecidable problem in correspondence theory',
\emph{The Journal of Symbolic Logic}\/ {\bf 56} (1991) 1261--1272.
%
%
%\bibitem{Chandra:et:al:1981}
%{\sc Chandra,} A., D. {\sc Kozen,} and L. {\sc Stockmeyer,}
%`Alternation',
%\emph{Journal of the Association for Computing Machinery}\/ {\bf 28} (1981) 114--133.
%
%
%\bibitem{Chellas:1980}
%{\sc Chellas,} B.,
%\emph{Modal Logic. An Introduction,}\/
%Cambridge University Press (1980).
%
%
%\bibitem{Church:Quine:1952}
%{\sc Church,} A., and W. {\sc Quine,}
%`Some theorems on definability and decidability',
%\emph{The Journal of Symbolic Logic}\/ {\bf 17} (1952) 179--187.
%
%
%\bibitem{Cintula:Metcalfe:2010}
%{\sc Cintula,} P., and G. {\sc Metcalfe,}
%`Admissible rules in the implication-negation fragment of intuitionistic logic',
%\emph{Annals of Pure and Applied Logic}\/ {\bf 162} (2010) 162--171.
%
%
\bibitem{Clarke:et:al:1986}
{\sc Clarke,} E., E. {\sc Emerson,} and A. {\sc Sistla,}
`Automatic verification of finite-state concurrent systems using temporal logic specifications',
\emph{ACM Transactions on Programming Languages and Systems}\/ {\bf 8} (1986) 244--263.
%
%
%\bibitem{Conradie:Goranko:2008}
%{\sc Conradie,} W., and V. {\sc Goranko,}
%`Algorithmic correspondence and completeness in modal logic. IV. Semantic extensions of $\SQEMA$',
%\emph{Journal of Applied Non-Classical Logics}\/ {\bf 18} (2008) 175--211.
%
%
%\bibitem{Conradie:Goranko:Vakarelov:2006}
%{\sc Conradie,} W., V. {\sc Goranko,} and D. {\sc Vakarelov,}
%`Algorithmic correspondence and completeness in modal logic. II. Polyadic and hybrid extensions of the algorithm $\SQEMA$',
%\emph{Journal of Logic and Computation}\/ {\bf 16} (2006) 579--612.
%
%
%\bibitem{Conradie:Goranko:Vakarelov:2009}
%{\sc Conradie,} W., V. {\sc Goranko,} and D. {\sc Vakarelov,}
%`Algorithmic correspondence and completeness in modal logic. III. Extensions of the algorithm $\SQEMA$ with substitutions',
%\emph{Fundamenta Informatic\ae}\/ {\bf 92} (2009) 307--343.
%
%
%\bibitem{Conradie:Goranko:Vakarelov:2010}
%{\sc Conradie,} W., V. {\sc Goranko,} and D. {\sc Vakarelov,}
%`Algorithmic correspondence and completeness in modal logic. V. Recursive extensions of $\SQEMA$',
%\emph{Journal of Applied Logic}\/ {\bf 8} (2010) 319--333.
%
%
%\bibitem{Demri:1998}
%{\sc Demri,} S.,
%`A class of decidable information logics',
%\emph{Theoretical Computer Science}\/ {\bf 195} (1998) 33--60.
%
%
%\bibitem{Diagne:1989}
%{\sc Diagne,} S.,
%\emph{Boole,}\/
%Belin (1989).
%
%
%\bibitem{VanDitmarsch:2014}
%{\sc Van Ditmarsch,} H.,
%`Dynamics of lying',
%\emph{Synthese}\/ {\bf 191} (2014) 745--777.
%
%
%\bibitem{VanDitmarsch:et:al:2014}
%{\sc Van Ditmarsch,} H., J. {\sc Fan,} W. {\sc van der Hoek,} and P. {\sc Iliev,}
%`Some exponential lower bounds on formula-size in modal logic',
%In: \emph{Advances in Modal Logic,}\/
%College Publications (2014) 139--157.
%
%
%\bibitem{VanDitmarsch:et:al:2008}
%{\sc Van Ditmarsch,} H., W. {\sc van der Hoek,} and B. {\sc Kooi,}
%\emph{Dynamic Epistemic Logic,}\/
%Springer (2008).
%
%
%\bibitem{Dunn:Hardegree:2001}
%{\sc Dunn,} M., and G. {\sc Hardegree,}
%\emph{Algebraic Methods in Philosophical Logic,}\/
%Oxford University Press (2001).
%
%
%\bibitem{Dzik:2003}
%{\sc Dzik,} W.,
%`Unitary unification of $\S5$ modal logics and its extensions',
%\emph{Bulletin of the Section of Logic}\/ {\bf 32} (2003) 19--26.
%
%
%\bibitem{Dzik:2007}
%{\sc Dzik,} W.,
%\emph{Unification Types in Logic,}\/
%Wydawnicto Uniwersytetu Slaskiego (2007).
%
%
%\bibitem{Dzik:2011}
%{\sc Dzik,} W.,
%`Remarks on projective unifiers',
%\emph{Bulletin of the Section of Logic}\/ {\bf 40} (2011) 37--46.
%
%
%\bibitem{Dzik:Wojtylak:2012}
%{\sc Dzik,} W., and P. {\sc Wojtylak,}
%`Projective unification in modal logic',
%\emph{Logic Journal of the IGPL}\/ {\bf 20} (2012) 121--153.
%
%
%\bibitem{Dzik:Wojtylak:2013}
%{\sc Dzik,} W., and P. {\sc Wojtylak,}
%`Modal consequence relations extending $\S4.3$: an application of projective unification',
%\emph{Notre Dame Journal of Formal Logic}\/ {\bf 57} (2013) 523--549.
%
%
\bibitem{Ebbinghaus:Flum:1999}
{\sc Ebbinghaus,} H.-D., and J. {\sc Flum,}
\emph{Finite Model Theory,}\/
Springer (1999).
%
%
%\bibitem{Emerson:1990}
%{\sc Emerson,} E.,
%`Temporal and modal logic',
%In: \emph{Handbook of Theoretical Computer Science, Volume B, Formal Models and Semantics,}\/
%Elsevier (1990) 995--1072.
%
%
%\bibitem{Endriss:et:al:2016}
%{\sc Endriss,} U., U. {\sc Grandi,} R. {\sc de Haan,} and J. {\sc Lang,}
%`Succinctness of languages for judgment aggregation',
%In: \emph{Proceedings of the Fifteenth International Conference on the Principles of Knowledge Representation and Reasoning,}\/
%AAAI Press (2016) 176--185.
%
%
\bibitem{Ershov:1980}
{\sc Ershov,} Y.,
\emph{Decidability Problems and Constructive Models,}\/
Nauka (1980).
%
%
%\bibitem{Fagin:et:al:1992}
%{\sc Fagin,} R., J. {\sc Halpern,} and M. {\sc Vardi,}
%`What is an inference rule?',
%\emph{Journal of Symbolic Logic}\/ {\bf 57} (1992) 1018--1045.
%
%
%\bibitem{Fernandez:Gil:2012}
%{\sc Fern\'andez Gil,} O.,
%\emph{Hybrid Unification in the Description Logic $\mathcal{EL}$,}\/
%Master Thesis of Technische Universit\"at Dresden (2012).
%
%
%\bibitem{French:Humberstone:2015}
%{\sc French,} R., and L. {\sc Humberstone,}
%`An observation concerning Porte's rule in modal logics',
%\emph{Bulletin of the Section of Logic}\/ {\bf 44} (2015) 25--31.
%
%
%\bibitem{Gabbay:Kurucz:Wolter:Zakharyaschev:2003}
%{\sc Gabbay,} D., A. {\sc Kurucz,} F. {\sc Wolter,} and M. {\sc Zakharyaschev,}
%\emph{Many-Dimensional Modal Logics: Theory and Applications,}\/
%Elsevier Science (2003).
%
%
%\bibitem{Gabbay:Shehtman:1998}
%{\sc Gabbay,} D., V. {\sc Shehtman,}
%`Products of modal logics, part 1',
%\emph{Logic Journal of the IGPL}\/ {\bf 6} (1998) 73--146.
%
%
%\bibitem{Gencer:deJongh:2009}
%{\sc Gencer,} \c{C}., and D. {\sc de Jongh,}
%`Unifiability in extensions of $\K4$',
%\emph{Logic Journal of the IGPL}\/ {\bf 17} (2009) 159--172.
%
%
%\bibitem{Ghilardi:1999}
%{\sc Ghilardi,} S.,
%`Unification in intuitionistic logic',
%\emph{Journal of Symbolic Logic}\/ {\bf 64} (1999) 859--880.
%
%
%\bibitem{Ghilardi:2000}
%{\sc Ghilardi,} S.,
%`Best solving modal equations',
%\emph{Annals of Pure and Applied Logic}\/ {\bf 102} (2000) 183--198.
%
%
%\bibitem{Ghilardi:Sacchetti:2004}
%{\sc Ghilardi,} S., and L. {\sc Sacchetti,}
%`Filtering unification and most general unifiers in modal logic',
%\emph{Journal of Symbolic Logic}\/ {\bf 69} (2004) 879--906.
%
%
\bibitem{Georgiev:2017a}
{\sc Georgiev,} D., 
`Definability in the class of all $\KD45$-frames --- computability and complexity',
\emph{Journal of Applied Non-Classical Logics}\/ {\bf 27} (2017) 1--26.
%
%
\bibitem{Georgiev:2017b}
{\sc Georgiev,} D., 
\emph{Algorithmic Methods for Non-Classical Logics,}\/
Doctoral thesis of Sofia University St. Kliment Ohridski (2017).
%
%
%\bibitem{Georgiev:Tinchev:2008}
%{\sc Georgiev,} G., and T. {\sc Tinchev,}
%`Second-order logic on equivalence relations',
%\emph{Journal of Applied Non-Classical Logics}\/ {\bf 18} (2008) 229--246.
%
%
%\bibitem{Georgieva:2021}
%{\sc Georgieva,} Y., 
%\emph{Modal Definability: Two Commuting Equivalence Relations,}\/
%Master thesis of Sofia University St. Kliment Ohridski (2021).
%
%
\bibitem{Goldblatt:1975}
{\sc Goldblatt,} R.,
`First-order definability in modal logic',
\emph{Journal of Symbolic Logic}\/ {\bf 40} (1975) 35--40.
%
%
%\bibitem{Goldblatt:1992}
%{\sc Goldblatt,} R.,
%\emph{Logics of Time and Computation (Second Edition),}\/
%CSLI Publications (1992).
%
%
%\bibitem{Goldblatt:Thomason:1975}
%{\sc Goldblatt,} R., and S. {\sc Thomason,}
%`Axiomatic classes in propositional modal logic',
%In: \emph{Algebra and Logic,}\/
%Springer (1975) 163--173.
%
%
%\bibitem{Goranko:Passy:1992}
%{\sc Goranko,} V., and S. {\sc Passy,}
%`Using the universal modality: gains and questions',
%\emph{Journal of Logic and Computation}\/ {\bf 2} (1992) 5--30.
%
%
%\bibitem{Goranko:Vakarelov:2006}
%{\sc Goranko,} V., and D. {\sc Vakarelov,}
%`Elementary canonical formul\ae: extending Sahlqvist's theorem',
%\emph{Annals of Pure and Applied Logic}\/ {\bf 141} (2006) 180--217.
%
%
%\bibitem{Halpern:Rego:2007}
%Halpern, J., R\^ego, L.:
%{\it Characterizing the $NP$-$PSPACE$ gap in the satisfiability problem for modal logic.}
%Journal of Logic and Computation {\bf 17} (2007) 795--806.
%
%
%\bibitem{Hemaspaandra:1996}
%{\sc Hemaspaandra,} E.,
%`The price of universality',
%\emph{Notre Dame Journal of Formal Logic}\/ {\bf 37} (1996) 174--203.
%
%
\bibitem{Hodges:1993}
{\sc Hodges,} W.,
\emph{Model Theory,}\/
Cambridge University Press (1993).
%
%
%\bibitem{Hodkinson:Reynolds:2007}
%{\sc Hodkinson,} I., and M. {\sc Reynolds,}
%`Temporal logic',
%In: \emph{Handbook of Modal Logic,}\/
%Elsevier (2007) 655--720.
%
%
\bibitem{Hughes:Cresswell:1968}
{\sc Hughes,} G., and M. {\sc Cresswell,}
\emph{An Introduction to Modal Logic,}\/
Methuen (1968).
%
%
%\bibitem{Iemhoff:2001}
%{\sc Iemhoff,} R.,
%`On the admissible rules of intuitionistic propositional logic',
%\emph{Journal of Symbolic Computation}\/ {\bf 66} (2001) 281--294.
%
%
%\bibitem{Iemhoff:2016}
%{\sc Iemhoff,} R.,
%`A syntactic approach to unification in transitive reflexive modal logics',
%\emph{Notre Dame Journal of Formal Logic}\/ {\bf 57} (2016) 233--247.
%
%
%\bibitem{Jansana:1994}
%{\sc Jansana,} R.,
%`Some logics related to von Wright's logic of place',
%\emph{Notre Dame Journal of Formal Logic}\/ {\bf 35} (1994) 88--98.
%
%
%\bibitem{Jerabek:2007}
%{\sc Je\u{r}\'abek,} E.,
%`Complexity of admissible rules',
%\emph{Archive for Mathematical Logic}\/ {\bf 46} (2007) 73--92.
%
%
%\bibitem{Jerabek:2013}
%{\sc Je\u{r}\'abek,} E.,
%`Logics with directed unification',
%In: \emph{Algebra and Coalgebra meet Proof Theory,}\/
%Workshop at Utrecht University (2013).
%
%
%\bibitem{Jerabek:2015}
%{\sc Je\u{r}\'abek,} E.,
%`Blending margins: the modal logic $\K$ has nullary unification type',
%\emph{Journal of Logic and Computation}\/ {\bf 25} (2015) 1231--1240.
%
%
%\bibitem{Kalmar:1937}
%Kalm\'ar, L.:
%{\it Zur\"uckf\"uhrung des Entscheidungsproblems auf den Fall von Formeln mit einer einzigen, bin\"aren, Funktionsvariablen.}
%Compositio Mathematica {\bf 4} (1937) 137--144.
%
%
%\bibitem{Kikot:Zolin:2013}
%{\sc Kikot,} S., and E. {\sc Zolin,}
%`Modal definability of first-order formulas with free variables and query answering',
%\emph{Journal of Applied Logic}\/ {\bf 11} (2013) 190--216.
%
%
%\bibitem{Kontchakov:et:al:2014}
%{\sc Kontchakov,} R., I. {\sc Pratt-Hartmann,} and M. {\sc Zakharyaschev,}
%`Spatial reasoning with $RCC8$ and connectedness constraints in Euclidean spaces',
%\emph{Artificial Intelligence}\/ {\bf 217} (2014) 43--75.
%
%
%\bibitem{Kost:2018}
%{\sc Kost,} S.,
%`Projective unification in transitive modal logics',
%\emph{Logic Journal of the IGPL}\/ {\bf 26} (2018) 548--566.
%
%
%\bibitem{Kracht:1999}
%{\sc Kracht,} M.,
%\emph{Tools and Techniques in Modal Logic,}\/
%Elsevier (1999).
%
%
%\bibitem{Kracht:Wolter:1991}
%{\sc Kracht,} M., F. {\sc Wolter,}
%`Properties of independently axiomatizable bimodal logics',
%\emph{Journal of Symbolic Logic}\/ {\bf 56} (1991) 1469--1485.
%
%
\bibitem{Kripke:1963}
{\sc Kripke,} S.,
`Semantical analysis of modal logic. I. Normal modal propositional calculi',
\emph{Zeitschrift f\"ur mathematische Logik und Grundlagen der Mathematik}\/ {\bf 9} (1963) 67--96.
%
%
%\bibitem{Kurucz:2007}
%{\sc Kurucz,} A.,
%`Combining modal logics',
%In: \emph{Handbook of Modal Logic,}\/
%Elsevier (2007) 869--924.
%
%
%\bibitem{Lavrov:1963}
%Lavrov, I.:
%{\it Effective inseparability of the set of identically true formulas and the set of formulas with finite counterexamples for certain elementary theories.}
%Algebra i Logika {\bf 2} (1963) 5--18.
%
%
%\bibitem{Lowenheim:1908}
%{\sc L\"owenheim,} L.,
%`\"Uber das Aufl\"osungsproblem im logischen Klassenkalk\"ul',
%\emph{Sitzungsberichte der Berliner mathematischen Gesellschaft}\/ {\bf 7} (1908) 89--94.
%
%
%\bibitem{Lutz:2006}
%{\sc Lutz,} C.,
%`Complexity and succinctness of Public Announcement Logic',
%In: \emph{Proceedings of the Fifth International Joint Conference on Autonomous Agents and Multiagent Systems,}\/
%ACM (2006) 137--143.
%
%
%\bibitem{Lutz:Walther:Wolter:2007}
%{\sc Lutz,} C., D. {\sc Walther,} and F. {\sc Wolter,}
%`Conservative extensions in expressive description logics',
%In: \emph{Proceedings of the Twentieth International Joint Conference on Artificial Intelligence,}\/
%AAAI Press (2007) 453--458.
%
%
%\bibitem{Martin:Nipkow:1988}
%{\sc Martin,} U., T. {\sc Nipkow,}
%`Unification in Boolean rings',
%\emph{Journal of Automated Reasoning}\/ {\bf 4} (1988) 381--396.
%
%
%\bibitem{Martin:Nipkow:1989}
%{\sc Martin,} U., T. {\sc Nipkow,}
%`Boolean unification~---~the story so far',
%\emph{Journal of Symbolic Computation}\/ {\bf 7} (1989) 275--293.
%
%
%\bibitem{Miyazaki:2004}
%{\sc Miyazaki,} Y.,
%`Normal modal logics containing $\KTB$ with some finiteness conditions',
%In: \emph{Advances in Modal Logic,}\/
%College Publications (2004) 171--190.
%
%
\bibitem{Nagle:1981}
{\sc Nagle,} M.,
`The decidability of normal $\K5$ logics',
\emph{Journal of Symbolic Logic}\/ {\bf 46} (1981) 319--328.
%
%
\bibitem{Nagle:Thomason:1985}
{\sc Nagle,} M., S. {\sc Thomason,}
`The extensions of the modal logic $\K5$',
\emph{Journal of Symbolic Logic} {\bf 50} (1985) 102--109.
%
%
%\bibitem{Papadimitriou:1994}
%Papadimitriou, C.:
%{\it Computational Complexity.}
%Addison-Wesley (1994).
%
%
%\bibitem{Porte:1981}
%{\sc Porte,} J.,
%`The deducibilities of $\S5$',
%\emph{Journal of Philosophical Logic}\/ {\bf 10} (1981) 409--422.
%
%
%\bibitem{Reynolds:Zakharyaschev:2001}
%{\sc Reynolds,} M., M. {\sc Zakharyaschev,}
%`On the products of linear modal logics',
%\emph{Journal of Logic and Compution}\/ {\bf 11} (2001) 909--931.
%
%
%\bibitem{deRijke:1992}
%{\sc De Rijke,} M.,
%`The modal logic of inequality',
%\emph{The Journal of Symbolic Logic}\/ {\bf 57} (1992) 566--584.
%
%
%\bibitem{Rogers:1952}
%Rogers, H.:
%{\it Some Results on Definability and Decidability in Elementary Theories.}
%Doctoral thesis of the University of Princeton (1952).
%
%
%\bibitem{Rogers:1956}
%Rogers, H.:
%{\it Certain logical reduction and decision problems.}
%Annals of Mathematics {\bf 64} (1956) 264--284.
%
%
%\bibitem{Rogers:1967}
%Rogers, H.:
%{\it Theory of Recursive Functions and Effective Computability.}
%McGraw-Hill (1967).
%
%
%\bibitem{Rudeanu:1974}
%{\sc Rudeanu,} S.,
%\emph{Boolean Functions and Equations,}\/
%Elsevier (1974).
%
%
%\bibitem{Rybakov:1984}
%{\sc Rybakov,} V.,
%`A criterion for admissibility of rules in the model system $\S4$ and the intuitionistic logic.
%\emph{Algebra and Logic}\/ {\bf 23} (1984) 369--384.
%
%
%\bibitem{Rybakov:1985}
%{\sc Rybakov,} V.,
%`Bases of admissible rules of the logics $\S4$ and $\Int$',
%\emph{Algebra and Logic}\/ {\bf 24} (1985) 55--68.
%
%
%\bibitem{Rybakov:1997}
%{\sc Rybakov,} V.,
%\emph{Admissibility of Logical Inference Rules,}\/
%Elsevier (1997).
%
%
%\bibitem{Rybakov:2001}
%{\sc Rybakov,} V.,
%`Construction of an explicit basis for rules admissible in modal system $\S4$',
%\emph{Mathematical Logic Quarterly}\/ {\bf 47} (2001) 441--446.
%
%
%\bibitem{Rybakov:2008}
%{\sc Rybakov,} V.,
%`Linear temporal logic with until and next, logical consecutions',
%\emph{Annals of Pure and Applied Logic}\/ {\bf 155} (2008) 32--45.
%
%
%\bibitem{Rybakov:Terziler:Gencer:1999}
%{\sc Rybakov,} V., M. {\sc Terziler,} and \c{C}. {\sc Gencer,}
%`An essay on unification and inference rules for modal logics',
%\emph{Bulletin of the Section of Logic}\/ {\bf 28} (1999) 145--157.
%
%
\bibitem{Sahlqvist:1975}
{\sc Sahlqvist,} H.,
`Completeness and correspondence in the first and second order semantics for modal logic',
In: \emph{Proceedings of the Third Scandinavian Logic Symposium,}\/
North-Holland (1975) 110--143.
%
%
%\bibitem{Segerberg:1986}
%{\sc Segerberg,} K.,
%`Modal logics with functional alternative relations',
%\emph{Notre Dame Journal of Formal Logic}\/ {\bf 27} (1986) 504--522.
%
%
%\bibitem{Shapirovsky:2006}
%{\sc Shapirovsky,} I.,
%`Downward-directed transitive frames with universal relations',
%In: \emph{Advances in Modal Logic,}\/
%College Publications (2006) 413--428.
%
%
%\bibitem{Shapirovsky:Shehtman:2016}
%{\sc Shapirovsky,} I., and V. {\sc Shehtman,}
%`Local tabularity without transitivi\-ty',
%In: \emph{Advances in Modal Logic,}\/
%College Publications (2016) 520--534.
%
%
%\bibitem{Shoenfield:1967}
%Shoenfield, J.:
%{\it Mathematical Logic.}
%Addison-Wesley (1967).
%
%
%\bibitem{Stockmeyer:1974}
%Stockmeyer, L.:
%{\it The Complexity of Decision Problems in Automata Theory.}
%Doctoral thesis of the Massachusetts Institute of Technology (1974).
%
%
%\bibitem{Stockmeyer:1977}
%Stockmeyer, L.:
%{\it The polynomial-time hierarchy.}
%Theoretical Computer Science {\bf 3} (1977) 1--22.
%
%
%\bibitem{Taitslin:1962}
%Ta\v\i tslin, M.:
%{\it Effective inseparability of the sets of identically true and finitely refutable formulas of elementary lattice theory.}
%Algebra i Logika {\bf 1} (1962) 24--38.
%
%
%\bibitem{Tarski:1949}
%Tarski, A.:
%{\it Undecidability of the theories of lattices and projective geometries.}
%The Journal of Symbolic Logic {\bf 14} (1949) 77--78.
%
%
\bibitem{Tarski:Mostowski:Robinson:1953}
{\sc Tarski,} A., A. {\sc Mostowski,} R. {\sc Robinson,}
\emph{Undecidable Theories,}\/
North-Holland (1953).
%
%
\bibitem{Vardi:1982}
{\sc Vardi,} M.,
`The complexity of relational query languages',
In: {\it Proceedings of the 14th ACM Symposium on Theory of Computing,}\/
ACM (1982) 137--146.
%
%
%\bibitem{Vaught:1960}
%Vaught, R.:
%{\it Sentences true in all constructive models.}
%The Journal of Symbolic Logic {\bf 25} (1960) 39--53.
%
%
%\bibitem{Wolter:1998}
%{\sc Wolter,} F.,
%`Fusions of modal logics revisited',
%In: \emph{Advances in Modal Logic,}\/
%CSLI Publications (1998) 361--379.
%
%
%\bibitem{Wolter:Zakharyaschev:2008}
%{\sc Wolter,} F., and M. {\sc Zakharyaschev,}
%`Undecidability of the unification and admissibility problems for modal and description logics',
%\emph{ACM Transactions on Computational Logic}\/ {\bf 9} (2008) 25:1--25:20.
%
%
\end{thebibliography}
%
%
%
%
\end{document}
%
%
%
%
